\documentclass[10pt,landscape,a4paper,twocolumn]{article}
\usepackage[utf8]{inputenc}
\usepackage[english]{babel}
\usepackage{listings}
\usepackage{xcolor}
\usepackage{geometry}
\usepackage{fancyhdr}
\usepackage{hyperref}
\usepackage{tocloft}

\geometry{top=1cm, bottom=1.5cm, left=1cm, right=1cm, includehead, headheight=15pt, headsep=0.3cm}

\definecolor{codegreen}{rgb}{0,0.6,0}
\definecolor{codegray}{rgb}{0.5,0.5,0.5}
\definecolor{codepurple}{rgb}{0.58,0,0.82}
\definecolor{backcolour}{rgb}{0.96,0.96,0.96}
\definecolor{linkcolor}{rgb}{0,0.3,0.6}
\definecolor{keywordcolor}{rgb}{0.8,0.1,0.4}

\hypersetup{
    colorlinks=true,
    linkcolor=linkcolor,
    filecolor=magenta,      
    urlcolor=blue,
    pdftitle={ECPC Reference Document},
    pdfauthor={Tarek Accepted 3shan Zh2t}
}

\lstdefinestyle{mystyle}{
    backgroundcolor=\color{backcolour},   
    commentstyle=\itshape\color{codegreen},
    keywordstyle=\bfseries\color{keywordcolor},
    numberstyle=\tiny\color{codegray},
    stringstyle=\color{codepurple},
    basicstyle=\ttfamily\footnotesize,
    breakatwhitespace=false,         
    breaklines=true,                 
    captionpos=b,                    
    keepspaces=true,                 
    numbers=left,                    
    numbersep=5pt,                  
    showspaces=false,                
    showstringspaces=false,
    showtabs=false,                  
    tabsize=4,
    frame=single,
    rulecolor=\color{codegray}
}
\lstset{style=mystyle}

\pagestyle{fancy}
\fancyhf{}
\fancyhead[L]{\textbf{Tarek Accepted 3shan Zh2t} - Tanta University}
\fancyhead[C]{ECPC Team Reference Document}
\fancyhead[R]{\textbf{Page \thepage}}
\renewcommand{\headrulewidth}{0.4pt}

\fancypagestyle{plain}{
    \fancyhf{}
    \fancyhead[L]{\textbf{Tarek Accepted 3shan Zh2t} - Tanta University}
    \fancyhead[C]{ECPC Team Reference Document}
    \fancyhead[R]{\textbf{Page \thepage}}
    \renewcommand{\headrulewidth}{0.4pt}
}

\setlength{\cftsecnumwidth}{2.5em}
\setlength{\cftsubsecnumwidth}{3em}
\setcounter{tocdepth}{3}

\begin{document}
\tableofcontents
\newpage
\section{Template}
\subsection{Main}
\begin{lstlisting}[language=C++]
#include <bits/stdc++.h>
#include <ext/pb_ds/assoc_container.hpp>
#include <ext/pb_ds/tree_policy.hpp>

using namespace std;
using namespace __gnu_pbds;

#define El_Hefnawys                   \
    ios_base::sync_with_stdio(false); \
    cin.tie(nullptr);                 \
    cout.tie(nullptr)

typedef long long ll;
typedef unsigned long long ull;

#define int ll

typedef vector<vector<int>> GRAPH;

const int MOD = 1e9 + 7;
const int OO = 1e15;
const double PI = acos(-1.0);
const double EPS = 1e-9;

int dx4[4] = {0, 1, 0, -1};
int dy4[4] = {1, 0, -1, 0};

int dx8[8] = {0, 0, 1, -1, 1, 1, -1, -1};
int dy8[8] = {1, -1, 0, 0, 1, -1, 1, -1};

void solve()
{
}

signed main()
{
    El_Hefnawys;
    int t{1};
    // cin >> t;
    while (t--)
    {
        solve();
    }
    return 0;
}

\end{lstlisting}
\section{Data Structures}
\subsection{Centroid decomposition}
\subsubsection{Centroid decomposition}
\begin{lstlisting}[language=C++]
/**
 *  THE ULTIMATE CP TEMPLATE BLUEPRINT: CENTROID DECOMPOSITION
 *
 * 1.  WHAT PROBLEM DOES THIS SOLVE?
 * - Keywords: "Paths of length K", "Distance between all pairs", "Update node color, find closest red node".
 * - Classic Scenarios: You need to answer heavy path queries on a massive tree where standard DFS
 *   would Time Limit Exceed (TLE).
 * - The Magic: We find the "center of gravity" (Centroid) of the tree, make it a root, and "remove" it.
 *   This shatters the graph into a forest. We repeat this recursively. The resulting "Centroid Tree"
 *   connects these centers and mathematically guarantees a maximum depth of O(log N).
 *   You can safely brute-force climbing from any node to the root of the Centroid Tree in just ~17 steps!
 *
 * 2.  HOW TO USE IT (THE BLACK BOX)
 * - Initialization: Pass the size of the tree (1-indexed) and the adjacency list.
 *       CentroidDecomposition cd(N, adj);
 *
 * - Navigation: The `cd.centroid_parent[u]` array tells you the parent of node `u` in the new tree.
 * - Distance: Use `cd.get_dist(u, v)` to find the distance between nodes in the ORIGINAL tree.
 *
 * - Complexity:
 *       Time: O(N log N) to build the tree and the LCA tables.
 *       Space: O(N log N) memory for binary lifting tables.
 *
 * 3.  HOW TO ADAPT IT (THE DIALS & SWITCHES)
 * - How to solve a problem with it:
 *   Instead of answering queries globally, answer them for "all paths passing through the current centroid".
 *   For point updates (e.g., painting a node red), start at the updated node and climb up its
 *   `centroid_parent` path, updating the answer at each ancestor using `get_dist(u, ancestor)`.
 */

#include <bits/stdc++.h>
using namespace std;

struct CentroidDecomposition
{
    int n;
    vector<vector<int>> adj;

    // Variables for the Centroid Tree
    vector<int> sz;
    vector<bool> removed;
    vector<int> centroid_parent;

    // Variables for O(1) distance queries in the original tree (Binary Lifting)
    const int LOG = 20;
    vector<vector<int>> up;
    vector<int> depth;

    CentroidDecomposition(int nodes, const vector<vector<int>> &graph)
    {
        n = nodes;
        adj = graph;

        sz.assign(n + 1, 0);
        removed.assign(n + 1, false);
        centroid_parent.assign(n + 1, -1);

        up.assign(n + 1, vector<int>(LOG, 0));
        depth.assign(n + 1, 0);

        // 1. Precompute LCA for fast distance queries
        lca_dfs(1, 1);

        // 2. Build the Centroid Tree
        decompose(1, -1);
    }

    // Standard Binary Lifting DFS to compute depths and ancestors
    void lca_dfs(int u, int p)
    {
        up[u][0] = p;
        for (int i = 1; i < LOG; i++)
        {
            up[u][i] = up[up[u][i - 1]][i - 1];
        }
        for (int v : adj[u])
        {
            if (v != p)
            {
                depth[v] = depth[u] + 1;
                lca_dfs(v, u);
            }
        }
    }

    int get_lca(int u, int v)
    {
        if (depth[u] < depth[v])
            swap(u, v);
        int diff = depth[u] - depth[v];
        for (int i = 0; i < LOG; i++)
        {
            if ((diff >> i) & 1)
                u = up[u][i];
        }
        if (u == v)
            return u;
        for (int i = LOG - 1; i >= 0; i--)
        {
            if (up[u][i] != up[v][i])
            {
                u = up[u][i];
                v = up[v][i];
            }
        }
        return up[u][0];
    }

    // Get the shortest path distance between u and v in the ORIGINAL tree
    int get_dist(int u, int v)
    {
        return depth[u] + depth[v] - 2 * depth[get_lca(u, v)];
    }

    // Your exact dfs function, slightly modified to ignore "removed" nodes!
    void get_sz(int u, int p)
    {
        sz[u] = 1;
        for (int v : adj[u])
        {
            if (v == p || removed[v])
                continue; // Ignore dead ends
            get_sz(v, u);
            sz[u] += sz[v];
        }
    }

    // Your exact centroid function, slightly modified to ignore "removed" nodes!
    int get_centroid(int u, int p, int total_nodes)
    {
        for (int v : adj[u])
        {
            if (v == p || removed[v])
                continue; // Ignore dead ends

            if (sz[v] > total_nodes / 2)
            {
                return get_centroid(v, u, total_nodes);
            }
        }
        return u;
    }

    // The master function that shatters the graph and builds the new tree
    void decompose(int u, int p)
    {
        // 1. Calculate the subtree sizes of the current shattered component
        get_sz(u, -1);
        int total_nodes = sz[u];

        // 2. Find the exact center of gravity (Centroid) for this component
        int c = get_centroid(u, -1, total_nodes);

        // 3. Link it to the Centroid Tree!
        centroid_parent[c] = p;

        // 4. "Delete" this centroid to shatter the tree into smaller forests
        removed[c] = true;

        // 5. Recursively find the centroids of the remaining pieces
        for (int v : adj[c])
        {
            if (!removed[v])
            {
                decompose(v, c);
            }
        }
    }
};

void solve()
{
    int n, q;
    // Read N nodes and Q queries
    if (!(cin >> n >> q))
        return;

    vector<vector<int>> adj(n + 1);
    for (int i = 0; i < n - 1; i++)
    {
        int u, v;
        cin >> u >> v;
        adj[u].push_back(v);
        adj[v].push_back(u);
    }

    // Build the magical O(log N) depth tree
    CentroidDecomposition cd(n, adj);

    // EXAMPLE QUERY: "Update node U's color, then find closest node with that color"
    // To solve this, you would start at U and literally climb the centroid tree!
    /*
    int u = 5;
    int curr = u;
    while (curr != -1) {
        int distance_in_original_tree = cd.get_dist(u, curr);
        // update_answer(curr, distance_in_original_tree);
        curr = cd.centroid_parent[curr];
    }
    */
}

int main()
{
    // Fast I/O
    ios_base::sync_with_stdio(false);
    cin.tie(NULL);
    solve();
    return 0;
}
\end{lstlisting}
\subsection{DSU}
\subsubsection{DSU(Hassan)}
\begin{lstlisting}[language=C++]
class DSU {
public:
    vector<int>parent,sizes,mn,mx;
    DSU(int n) {
        parent.resize(n+1);
        sizes.resize(n+1);
        mx.resize(n+1);
        mn.resize(n+1);
        for (int i = 1; i <= n; i++)
            parent[i] = mn[i] = mx[i] = i,sizes[i] = 1  ;
    }

    int findRoot(int u ) {
        return parent[u] == u ? u : parent[u] = findRoot(parent[u]);
    }

    bool merge(int u, int v) {
        int root_u = findRoot(u), root_v = findRoot(v);
        if (root_u == root_v)
            return false;

        if (sizes[root_u] > sizes[root_v])
            swap(root_u, root_v);

        mx[root_v] = max(mx[root_v], mx[root_u]);
        mn[root_v] = min(mn[root_v], mn[root_u]);
        parent[root_u] = root_v;
        sizes[root_v] += sizes[root_u];
        return true;
    }
};

void Hassan() {
    int n,q;cin >>n>>q;
    DSU d(n);
    while (q--) {
        string op;int u,v;cin >> op ;
        if (op == "get") {
            cin >>v;
            int r = d.findRoot(v);
            cout << d.mn[r] <<" " << d.mx[r] <<" "<< d.sizes[r] <<endl;
        }else {
            cin >>u>>v;
            d.merge(u,v);
        }
    }
}

\end{lstlisting}
\subsubsection{DSU}
\begin{lstlisting}[language=C++]
class DSU
{
private:
    vector<int> parent, size;
    int forests{};

    void link(int x, int y)
    {
        if (size[x] > size[y])
            swap(x, y);
        parent[x] = y;
        size[y] += size[x];
    }

public:
    DSU(int n)
    {
        parent = vector<int>(n), size = vector<int>(n);
        forests = n;
        for (int i = 0; i < n; i++)
            parent[i] = i, size[i] = 1;
    }

    int find_set(int x)
    {
        if (x == parent[x])
            return x;
        return parent[x] = find_set(parent[x]);
    }

    int union_sets(int x, int y)
    {
        x = find_set(x), y = find_set(y);
        if (x != y)
        {
            link(x, y);
            forests--;
        }
        return x != y;
    }

    bool same_set(int x, int y)
    {
        x = find_set(x), y = find_set(y);
        return x == y;
    }

    int get_size(int x)
    {
        return size[find_set(x)];
    }

    int get_forests() const
    {
        return forests;
    }
};

\end{lstlisting}
\subsubsection{DSUOnTree}
\begin{lstlisting}[language=C++]
/**
 *  THE ULTIMATE CP TEMPLATE BLUEPRINT: DSU ON TREE (SACK / SMALL TO LARGE)
 *
 * 1.  WHAT PROBLEM DOES THIS SOLVE?
 * - Keywords: "Subtree queries", "Offline queries", "Distinct colors in subtree", "Frequencies at depth".
 * - Classic Scenarios: You have a tree and Q offline queries asking about properties within 
 *   the subtree of node U (e.g., "Can we form a palindrome using nodes at depth H in U's subtree?").
 *   A naive DFS for each query gives TLE O(N^2).
 * - The Magic: We use "Small to Large" merging. We find the "heavy" child (largest subtree). 
 *   We process light children and ERASE their data, then process the heavy child and KEEP its data. 
 *   Finally, we add light children back. Since a node is a "light child" at most O(log N) times, 
 *   the complexity drops to O(N log N).
 *
 * 2.  HOW TO USE IT (THE BLACK BOX)
 * - Initialization: Pass the number of nodes, the string of characters, and number of queries.
 *       DSUOnTree dsu_tree(N, s, Q);
 *
 * - Updates: Add edges to build the tree, and add queries offline.
 *       dsu_tree.addEdge(u, v);
 *       dsu_tree.addQuery(node, depth, query_index);
 *
 * - Execution: Call solve(root) to process everything and get the answers.
 *       vector<string> answers = dsu_tree.process(1);
 *
 * - Complexity:
 *       Time: O(N log N) processing time.
 *       Space: O(N) memory.
 *
 * 3.  HOW TO ADAPT IT (THE DIALS & SWITCHES)
 * - 0-indexed or 1-indexed?
 *   The template naturally handles 1-indexed nodes. If your string is 0-indexed, it automatically 
 *   shifts it to 1-indexed internally for safety.
 * - Changing the problem condition:
 *   Modify the `modifyState()` function to count frequencies, update maximums, or manage a `std::set`.
 *   Then modify the "Answer queries" section in `dfsDSU()` to check your specific condition.
 */

#include <bits/stdc++.h>
using namespace std;

struct DSUOnTree
{
private:
    int n;
    string s;
    vector<vector<int>> adj;
    vector<vector<pair<int, int>>> queries; // {depth, query_index}
    
    vector<int> sz;
    vector<int> lvl;
    vector<int> mask; // Bitmask to track odd/even frequencies of characters
    vector<bool> big; // Marks the heavy child
    vector<string> ans;

    // Step 1: Precalculate subtree sizes and depths
    void dfsSize(int node, int parent, int level)
    {
        lvl[node] = level;
        sz[node] = 1;
        for (int child : adj[node])
        {
            if (child != parent)
            {
                dfsSize(child, node, level + 1);
                sz[node] += sz[child];
            }
        }
    }

    // Step 2: Add or remove a node's contribution from our global state
    void modifyState(int node, int parent, int x)
    {
        // ? THE MAGIC: We toggle the i-th bit to flip the parity of the character's frequency
        mask[lvl[node]] ^= (1 << (s[node] - 'a'));

        for (int child : adj[node])
        {
            // ? CRITICAL: Do NOT traverse the heavy child if it's already marked!
            if (child != parent && !big[child])
            {
                modifyState(child, node, x);
            }
        }
    }

    // Step 3: The core DSU on Tree logic
    void dfsDSU(int node, int parent, bool keep)
    {
        int maxSize = -1, bigChild = -1;
        
        // Find the heavy child
        for (int child : adj[node])
        {
            if (child != parent && sz[child] > maxSize)
            {
                maxSize = sz[child];
                bigChild = child;
            }
        }

        // Process all light children first, and ERASE their data (keep = 0)
        for (int child : adj[node])
        {
            if (child != parent && child != bigChild)
            {
                dfsDSU(child, node, 0);
            }
        }

        // Process the heavy child, and KEEP its data (keep = 1)
        if (bigChild != -1)
        {
            dfsDSU(bigChild, node, 1);
            big[bigChild] = true; // Mark as heavy so `modifyState` ignores it
        }

        // Add the current node and its light children back to the state
        modifyState(node, parent, 1);

        // --- ANSWER QUERIES FOR THIS SUBTREE ---
        for (auto& q : queries[node])
        {
            int target_depth = q.first;
            int query_id = q.second;

            // ? THE MAGIC: If at most 1 bit is set, the characters can form a palindrome.
            if (__builtin_popcount(mask[target_depth]) <= 1)
            {
                ans[query_id] = "Yes";
            }
            else
            {
                ans[query_id] = "No";
            }
        }

        // Cleanup: Unmark the heavy child
        if (bigChild != -1)
        {
            big[bigChild] = false;
        }

        // ? CRITICAL: If this whole subtree was a light child, erase it completely!
        if (keep == 0)
        {
            modifyState(node, parent, -1);
        }
    }

public:
    DSUOnTree(int n, string input_s, int q_count)
    {
        this->n = n;
        // Pad the string so it becomes 1-indexed to match the nodes
        this->s = " " + input_s; 
        
        adj.resize(n + 1);
        queries.resize(n + 1);
        sz.assign(n + 1, 0);
        lvl.assign(n + 1, 0);
        mask.assign(n + 1, 0);
        big.assign(n + 1, false);
        ans.resize(q_count);
    }

    void addEdge(int u, int v)
    {
        adj[u].push_back(v);
        adj[v].push_back(u);
    }

    void addQuery(int node, int depth, int query_index)
    {
        queries[node].push_back({depth, query_index});
    }

    vector<string> process(int root = 1)
    {
        dfsSize(root, 0, 1);
        dfsDSU(root, 0, 0);
        return ans;
    }
};

void solve()
{
    int n, q;
    // Example Input: N nodes, Q queries
    if (!(cin >> n >> q)) return;

    string s;
    cin >> s;

    DSUOnTree dsu_tree(n, s, q);

    for (int i = 1; i < n; i++)
    {
        int u, v;
        cin >> u >> v;
        dsu_tree.addEdge(u, v);
    }

    for (int i = 0; i < q; i++)
    {
        int v, h;
        cin >> v >> h;
        // 0-indexed queries
        dsu_tree.addQuery(v, h, i); 
    }

    // Process all queries offline
    vector<string> results = dsu_tree.process(1);

    for (const string& res : results)
    {
        cout << res << "\n";
    }
}

int main()
{
    ios_base::sync_with_stdio(false);
    cin.tie(NULL);
    // solve();
    return 0;
}

\end{lstlisting}
\subsubsection{DSU move erase}
\begin{lstlisting}[language=C++]
struct DSU
{
    vector<int> _t;

    DSU(int N) : _t(2 * N, -1)
    {
        iota(_t.begin(), _t.begin() + N, N);
    }

    int find(int x)
    {
        return _t[x] < 0 ? x : _t[x] = find(_t[x]);
    }

    bool join(int x, int y)
    {
        auto fx = find(x), fy = find(y);
        if (is_del(x) || is_del(y) || fx == fy)
            return false;
        if (_t[fx] > _t[fy])
            swap(fx, fy);
        _t[fx] += _t[fy];
        _t[fy] = fx;
        return true;
    }

    bool move(int x, int y)
    {
        auto fx = find(x), fy = find(y);
        if (is_del(x) || is_del(y) || fx == fy)
            return false;
        _t[fx]++;
        _t[fy]--;
        _t[x] = fy;
        return true;
    }

    bool erase(int x)
    {
        auto fx = find(x);
        if (is_del(x))
            return false;
        _t[fx]++;
        _t[x] = -1;
        return true;
    }

    bool is_del(int x)
    {
        return _t[x] == -1;
    }

    int sz(int x)
    {
        return -_t[find(x)];
    }
};

\end{lstlisting}
\subsubsection{DSU with rollback}
\begin{lstlisting}[language=C++]
class DSU
{
private:
    vector<int> parent, size;
    stack<int> changes;
    int forests{};

    void link(int x, int y)
    {
        if (size[x] > size[y])
            swap(x, y);
        changes.push(x);
        parent[x] = y;
        size[y] += size[x];
    }

public:
    DSU(int n)
    {
        parent = vector<int>(n), size = vector<int>(n), changes = stack<int>();
        forests = n;
        for (int i = 0; i < n; i++)
            parent[i] = i, size[i] = 1;
    }

    int find_set(int x)
    {
        while (x != parent[x])
            x = parent[x];
        return x;
    }

    int union_sets(int x, int y)
    {
        x = find_set(x), y = find_set(y);
        if (x != y)
        {
            link(x, y);
            forests--;
        }
        return x != y;
    }

    void rollback()
    {
        if (changes.empty())
            return;
        int x = changes.top();
        changes.pop();
        size[parent[x]] -= size[x];
        parent[x] = x;
        forests++;
    }

    bool same_set(int x, int y)
    {
        x = find_set(x), y = find_set(y);
        return x == y;
    }

    int get_size(int x)
    {
        return size[find_set(x)];
    }

    int get_forests() const
    {
        return forests;
    }
};

\end{lstlisting}
\subsubsection{KRT}
\begin{lstlisting}[language=C++]
/**
 *  THE ULTIMATE CP TEMPLATE BLUEPRINT: KRUSKAL RECONSTRUCTION TREE (KRT)
 *
 * 1.  WHAT PROBLEM DOES THIS SOLVE?
 * - Keywords: "Minimax path", "Reachability under weight limit", "Max edge on path".
 * - Classic Scenarios:
 *   A) You are asked to find a path between U and V that MINIMIZES the MAXIMUM edge weight along the way.
 *   B) You are at city U. You can only use roads with a weight <= W. How many cities can you reach?
 * - The Magic: During Kruskal's MST algorithm, instead of just merging two components, we create a
 *   BRAND NEW artificial node to represent the edge! We make the roots of the two components children
 *   of this new artificial node. The result is a tree where the "leaves" are the original cities, and
 *   the "internal nodes" are the edges! The weight of the ancestors strictly increases as you go up.
 *
 * 2.  HOW TO USE IT (THE BLACK BOX)
 * - Initialization: Pass the number of nodes (1-indexed) and a vector of Edges.
 *       vector<Edge> edges = {{1, 2, 10}, {2, 3, 20}, {1, 3, 5}};
 *       KRT krt(3, edges);
 *
 * - Minimax Path Query:
 *       // Find the maximum weight edge on the optimal path between U and V
 *       long long max_edge = krt.getMaxEdgeOnPath(U, V);
 *
 * - Reachability Query:
 *       // From U, how many cities can I reach using ONLY edges with weight <= W?
 *       int cities = krt.getReachableLeaves(U, W);
 *
 * - Complexity:
 *       Time: O(E log E) to sort edges and build. O(log N) per query!
 *       Space: O(N log N) memory for Binary Lifting tables.
 *
 * 3.  HOW TO ADAPT IT (THE DIALS & SWITCHES)
 * - Max Spanning Tree instead of Min Spanning Tree?
 *   If you want a path that MAXIMIZES the MINIMUM edge, change the `operator<` in the `Edge` struct
 *   to `return w > other.w;`. Then, in `getReachableLeaves`, change the jump condition to
 *   `val[p] >= min_w`.
 * - Disconnected Graphs?
 *   The template natively supports disconnected graphs! `getMaxEdgeOnPath` returns `-1` if U and V
 *   are completely disconnected.
 */

#include <bits/stdc++.h>
using namespace std;

struct Edge
{
    int u, v;
    long long w;
    // Sort ascending for standard Min-MST KRT
    bool operator<(const Edge &other) const
    {
        return w < other.w;
    }
};

struct KRT
{
private:
    int n, krt_nodes;
    const int LOG = 21; // Safe for up to 10^6 nodes

    // DSU parent array (crucial that we don't use size/rank here, we MUST link to the new node)
    vector<int> dsu_parent;

    // KRT structure
    vector<long long> val;    // The weight of the edge this node represents
    vector<vector<int>> adj;  // Tree edges
    vector<int> depth;        // Depth for LCA
    vector<vector<int>> up;   // Binary Lifting table
    vector<int> leaves_count; // How many original cities are in this subtree?

    int find(int i)
    {
        if (i == dsu_parent[i])
            return i;
        // Path compression is perfectly safe here because we only use DSU to find the current root
        return dsu_parent[i] = find(dsu_parent[i]);
    }

    // Standard Binary Lifting DFS
    void dfs(int u, int p)
    {
        up[u][0] = p;
        for (int i = 1; i < LOG; i++)
        {
            up[u][i] = up[up[u][i - 1]][i - 1];
        }

        // Base case: If it has no children, it's an original city (Leaf)
        if (adj[u].empty())
        {
            leaves_count[u] = 1;
        }
        else
        {
            leaves_count[u] = 0;
        }

        for (int v : adj[u])
        {
            if (v != p)
            {
                depth[v] = depth[u] + 1;
                dfs(v, u);
                // Bubble up the leaf count
                leaves_count[u] += leaves_count[v];
            }
        }
    }

public:
    KRT(int n, vector<Edge> &edges)
    {
        this->n = n;
        // The KRT will have exactly N original leaves, and at most N-1 artificial edge-nodes
        int max_nodes = 2 * n + 5;

        dsu_parent.resize(max_nodes);
        iota(dsu_parent.begin(), dsu_parent.end(), 0);

        val.assign(max_nodes, 0);
        adj.resize(max_nodes);
        leaves_count.assign(max_nodes, 0);
        depth.assign(max_nodes, 0);
        up.assign(max_nodes, vector<int>(LOG, 0));

        // 1. Sort the edges to process in Kruskal's order
        sort(edges.begin(), edges.end());

        krt_nodes = n; // Nodes 1 to N are the original cities

        // 2. Build the Tree
        for (auto &e : edges)
        {
            int ru = find(e.u);
            int rv = find(e.v);

            if (ru != rv)
            {
                krt_nodes++; // Create a new artificial node for this edge
                val[krt_nodes] = e.w;

                // Connect the new node to the roots of the two components
                adj[krt_nodes].push_back(ru);
                adj[krt_nodes].push_back(rv);

                // ? THE MAGIC: The new node becomes the ultimate DSU root for both components!
                dsu_parent[ru] = krt_nodes;
                dsu_parent[rv] = krt_nodes;
            }
        }

        // 3. Process the DFS for LCA and Subtree sizes
        // We loop through all nodes in case the graph is a disconnected forest
        for (int i = 1; i <= krt_nodes; i++)
        {
            if (find(i) == i)
            {
                // We found a supreme root! Run DFS downwards.
                dfs(i, 0);
            }
        }
    }

    // Returns the max edge weight on the path between U and V
    long long getMaxEdgeOnPath(int u, int v)
    {
        if (find(u) != find(v))
            return -1; // They are in disconnected components
        if (u == v)
            return 0; // Same node

        // Standard LCA Logic
        if (depth[u] < depth[v])
            swap(u, v);
        int diff = depth[u] - depth[v];

        for (int i = 0; i < LOG; i++)
        {
            if ((diff >> i) & 1)
                u = up[u][i];
        }

        if (u == v)
            return val[u]; // One is an ancestor of the other

        for (int i = LOG - 1; i >= 0; i--)
        {
            if (up[u][i] != up[v][i])
            {
                u = up[u][i];
                v = up[v][i];
            }
        }

        // The LCA node IS the edge that connected them!
        int lca = up[u][0];
        return val[lca];
    }

    // Returns how many original cities can be reached from U using edges <= max_w
    int getReachableLeaves(int u, long long max_w)
    {
        // Jump up the tree as long as the parent edge weight is valid
        for (int i = LOG - 1; i >= 0; i--)
        {
            int p = up[u][i];
            // Parent 0 means we hit the root boundaries.
            if (p != 0 && val[p] <= max_w)
            {
                u = p; // Safe to jump!
            }
        }
        // Once we reach the highest valid ancestor, its subtree contains ALL reachable cities!
        return leaves_count[u];
    }
};

void solve()
{
    int n, m, q;
    // Read N nodes, M edges, Q queries
    if (!(cin >> n >> m >> q))
        return;

    vector<Edge> edges(m);
    for (int i = 0; i < m; i++)
    {
        cin >> edges[i].u >> edges[i].v >> edges[i].w;
    }

    KRT krt(n, edges);

    while (q--)
    {
        int type;
        cin >> type;
        if (type == 1)
        {
            int u, v;
            cin >> u >> v;
            long long ans = krt.getMaxEdgeOnPath(u, v);
            if (ans == -1)
                cout << "Disconnected\n";
            else
                cout << ans << "\n";
        }
        else
        {
            int u;
            long long w;
            cin >> u >> w;
            cout << krt.getReachableLeaves(u, w) << "\n";
        }
    }
}

int main()
{
    ios_base::sync_with_stdio(false);
    cin.tie(NULL);
    // solve();
    return 0;
}
\end{lstlisting}
\subsubsection{Offline Dynamic connectivity}
\begin{lstlisting}[language=C++]
/**
 *  THE ULTIMATE CP TEMPLATE BLUEPRINT: DSU WITH ROLLBACKS & SEGMENT TREE OVER TIME
 *
 * 1.  WHAT PROBLEM DOES THIS SOLVE?
 * - Keywords: "Add and remove edges", "Dynamic Connectivity", "Offline queries", "Bipartite over time".
 * - Classic Scenarios: You have a graph and Q queries. Some queries ADD an edge, some REMOVE 
 *   an existing edge, and some ask for the number of connected components. Standard DSU 
 *   can only add edges, it CANNOT remove them!
 * - The Magic: "Time Travel using Segment Tree". Since we read all queries offline, we know the 
 *   exact "lifespan" of every edge (from time L when it was added, to time R when it was removed). 
 *   We treat a Segment Tree as a timeline. We insert each edge into the segment tree over the 
 *   range [L, R]. Then, we DFS through the segment tree. 
 *   Going DOWN: We add the edges using a modified DSU. 
 *   At the LEAVES: We are at a specific point in time! We answer the queries for that moment.
 *   Going UP: We UNDO (rollback) the edges we added. Because we strictly use Union-by-Size 
 *   (NO path compression), our DSU history is perfectly preserved and easily reversible!
 *
 * 2.  HOW TO USE IT (THE BLACK BOX)
 * - Initialization: Create the segment tree. 
 *       segmentTree st(N); 
 *
 * - Add Edges: For every edge, find its active time range [L, R) and add it.
 *       st.set(L, R, edge(u, v));
 *
 * - Execution: Call the main DFS function to traverse time and answer queries.
 *       st.get();
 *
 * - Complexity:
 *       Time: O(Q log Q log N)  Every edge is split into log Q segments, and each DSU operation takes log N.
 *       Space: O(Q log Q) to store the edges in the segment tree nodes.
 *
 * 3.  HOW TO ADAPT IT (THE DIALS & SWITCHES)
 * - The `N` vs `Q` Trap: In this current template, `segmentTree(int n)` uses the same `n` for 
 *   BOTH the number of nodes `d(n)` and the segment tree time span. Usually, Nodes (N) != Queries (Q). 
 *   You should change the constructor to: `segmentTree(int max_time, int max_nodes) : d(max_nodes)`.
 * - How to track L and R: Use a `map<pair<int,int>, int> active_edges`. When an edge is added at 
 *   time `t1`, store it. When it's removed at time `t2`, call `st.set(t1, t2, edge)`. Don't forget 
 *   to add the edges that are never removed until the end: `st.set(t1, Q, edge)`.
 * - Answering Queries: Go to the `if (rx - lx == 1)` block inside `get(...)`. At this exact line, 
 *   the DSU holds the exact state of the graph at time `lx`. You can print `d.get_forests()` 
 *   or check `d.same_set(u, v)` based on what query was asked at time `lx`.
 */

class DSU
{
private:
    vector<int> parent, size;
    stack<int> changes;
    int forests{};

    void link(int x, int y)
    {
        if (size[x] > size[y])
            swap(x, y);
        changes.push(x);
        parent[x] = y;
        size[y] += size[x];
    }

public:
    DSU(int n)
    {
        parent = vector<int>(n), size = vector<int>(n), changes = stack<int>();
        forests = n;
        for (int i = 0; i < n; i++)
            parent[i] = i, size[i] = 1;
    }

    int find_set(int x)
    {
        while (x != parent[x])
            x = parent[x];
        return x;
    }

    int union_sets(int x, int y)
    {
        x = find_set(x), y = find_set(y);
        if (x != y)
        {
            link(x, y);
            forests--;
        }
        return x != y;
    }

    void rollback()
    {
        if (changes.empty())
            return;
        int x = changes.top();
        changes.pop();
        size[parent[x]] -= size[x];
        parent[x] = x;
        forests++;
    }

    bool same_set(int x, int y)
    {
        x = find_set(x), y = find_set(y);
        return x == y;
    }

    int get_size(int x)
    {
        return size[find_set(x)];
    }

    int get_forests() const
    {
        return forests;
    }
};

struct edge
{
    ll from, to;

    edge()
    {
    }

    edge(ll from, ll to) : from(from), to(to)
    {
    }

    bool operator==(const edge &e)
    {
        return from == e.from && to == e.to;
    }
};

struct Node
{
    int val;

    Node()
    {
        val = 0;
    }

    Node(int val) : val(val)
    {
    }

    void change(int x)
    {
        val = x;
    }
};

struct segmentTree
{
    int sz;
    vector<Node> seg;
    vector<vector<edge>> edges;
    DSU d;

    segmentTree(int n) : d(n)
    {
        sz = 1;
        while (sz < n)
            sz *= 2;
        seg.assign(2 * sz, Node());
        edges.assign(2 * sz, {});
    }

    void set(int l, int r, int node, int lx, int rx, edge edge)
    {
        if (lx >= l && rx <= r)
        {
            edges[node].push_back(edge);
            return;
        }
        if (rx <= l || lx >= r)
            return;
        int mid = (lx + rx) / 2;
        set(l, r, 2 * node + 1, lx, mid, edge);
        set(l, r, 2 * node + 2, mid, rx, edge);
    }

    void set(int l, int r, edge edge)
    {
        set(l, r, 0, 0, sz, edge);
    }

    void get(int node, int lx, int rx)
    {
        int cnt{};
        for (auto &[x, y] : edges[node])
            cnt += d.union_sets(x, y);
        if (rx - lx == 1)
        {
            // answer queries
        }
        else
        {
            int mid = (lx + rx) / 2;
            get(2 * node + 1, lx, mid);
            get(2 * node + 2, mid, rx);
        }
        while (cnt--)
            d.rollback();
    }

    void get()
    {
        return get(0, 0, sz);
    }
};

\end{lstlisting}
\subsubsection{Partially persistent DSU}
\begin{lstlisting}[language=C++]
/**
 *  THE ULTIMATE CP TEMPLATE BLUEPRINT: PARTIALLY PERSISTENT DSU
 *
 * 1.  WHAT PROBLEM DOES THIS SOLVE?
 * - Keywords: "At what time did U and V connect?", "Were they connected after K edges?", "Historical connectivity".
 * - Classic Scenarios: You have a graph where edges are slowly being built over time.
 *   You get queries asking if two cities were connected at a specific year in the past.
 * - The Magic: Because we ONLY use Union-by-Size (and strictly ban path compression), the DSU
 *   naturally forms a tree with a maximum depth of O(log N). By simply recording the `time`
 *   an edge was formed on the child node pointing to the parent, we can "time travel".
 *   To find the root at time T, we just climb the parent pointers, but we STOP climbing if
 *   the edge above us was built AFTER time T!
 *
 * 2.  HOW TO USE IT (THE BLACK BOX)
 * - Initialization: Pass the number of nodes (1-indexed).
 *       PartiallyPersistentDSU pp_dsu(N);
 *
 * - Updates: Merge U and V at time T. (T must be strictly increasing across calls).
 *       pp_dsu.unite(U, V, T);
 *
 * - Queries:
 *       // Were U and V connected at time T?
 *       bool connected = pp_dsu.isConnected(U, V, T);
 *
 *       // What is the EXACT time U and V became connected?
 *       int time_joined = pp_dsu.timeConnected(U, V);
 *
 * - Complexity:
 *       Time: O(log N) per merge and per query.
 *       Space: O(N) memory.
 *
 * 3.  HOW TO ADAPT IT (THE DIALS & SWITCHES)
 * - 0-indexed nodes?
 *   The template naturally allocates `n + 1` space. It works flawlessly for both 0-indexed
 *   and 1-indexed inputs without changing any code!
 * - Getting the size of a component at time T?
 *   This requires a bit more work. You would need to make `sz` a vector of `vector<pair<int, int>>`
 *   to store the history of sizes: `sz[root].push_back({time, new_size})`, and then binary search it.
 *   If you just need standard connectivity, this template is perfectly lightweight.
 */

#include <bits/stdc++.h>
using namespace std;

struct PartiallyPersistentDSU
{
private:
    vector<int> parent;
    vector<int> sz;
    vector<int> time_joined; // The time this node became a child of its parent
    const int INF = 2e9 + 7;

public:
    PartiallyPersistentDSU(int n)
    {
        parent.resize(n + 1);
        sz.assign(n + 1, 1);
        time_joined.assign(n + 1, INF);
        iota(parent.begin(), parent.end(), 0);
    }

    // ? CRITICAL: Notice there is NO path compression `parent[i] = find(...)`!
    // Path compression would destroy the historical edges.
    int find(int i, int t)
    {
        // We stop climbing if we hit a root, or if the edge above us hasn't been built yet at time `t`
        if (i == parent[i] || time_joined[i] > t)
        {
            return i;
        }
        return find(parent[i], t);
    }

    // Merge u and v at time t. (Assumes t is strictly increasing as you feed edges)
    bool unite(int u, int v, int t)
    {
        // When merging, we only care about the present moment, so we pass INF as time
        int root_u = find(u, INF);
        int root_v = find(v, INF);

        if (root_u == root_v)
            return false;

        // Union by Size ensures the tree depth stays O(log N)
        if (sz[root_u] < sz[root_v])
        {
            swap(root_u, root_v);
        }

        parent[root_v] = root_u;
        sz[root_u] += sz[root_v];

        // ? THE MAGIC: We record EXACTLY when root_v surrendered to root_u
        time_joined[root_v] = t;

        return true;
    }

    // Check if u and v were in the same component at time t
    bool isConnected(int u, int v, int t)
    {
        return find(u, t) == find(v, t);
    }

    // Find the exact time u and v became connected. Returns INF if they never connect.
    int timeConnected(int u, int v)
    {
        if (!isConnected(u, v, INF))
            return INF; // Never connected

        // Binary search the timeline to find the exact moment of connection!
        int l = 0, r = INF - 1;
        int ans = INF;

        while (l <= r)
        {
            int mid = l + (r - l) / 2;
            if (isConnected(u, v, mid))
            {
                ans = mid;
                r = mid - 1; // Try to find an earlier time
            }
            else
            {
                l = mid + 1; // We need more time
            }
        }
        return ans;
    }
};

void solve()
{
    int n, m, q;
    // Example Input: N nodes, M edges added sequentially, Q queries
    if (!(cin >> n >> m >> q))
        return;

    PartiallyPersistentDSU pp_dsu(n);

    for (int i = 1; i <= m; i++)
    {
        int u, v;
        cin >> u >> v;
        // Assume edges are added at time `i`
        pp_dsu.unite(u, v, i);
    }

    while (q--)
    {
        int type, u, v;
        cin >> type >> u >> v;

        if (type == 1)
        {
            // Were u and v connected at time T?
            int t;
            cin >> t;
            cout << (pp_dsu.isConnected(u, v, t) ? "YES" : "NO") << "\n";
        }
        else
        {
            // At what time did u and v finally connect?
            int t = pp_dsu.timeConnected(u, v);
            if (t > m)
                cout << "Not connected\n";
            else
                cout << "Connected at edge " << t << "\n";
        }
    }
}

int main()
{
    ios_base::sync_with_stdio(false);
    cin.tie(NULL);
    // solve();
    return 0;
}
\end{lstlisting}
\subsubsection{Small to large}
\begin{lstlisting}[language=C++]
// Small to large merging

vector<int> colors;
GRAPH graph;

set<int> dfs(int u, int p)
{
    set<int> st{colors[u]};
    for (auto &ch : graph[u])
    {
        if (ch == p)
            continue;
        set<int> child_st = dfs(ch, u);
        if (int(child_st.size()) > int(st.size()))
            swap(st, child_st);
        st.merge(child_st);
    }
    return st;
}

\end{lstlisting}
\subsection{Fenwick Tree}
\subsubsection{2D bit}
\begin{lstlisting}[language=C++]
/**
 *  THE ULTIMATE CP TEMPLATE BLUEPRINT: 2D FENWICK TREE (POINT UPDATE & RANGE QUERY)
 *
 * 1.  WHAT PROBLEM DOES THIS SOLVE?
 * - Keywords: "2D Grid", "Update a single cell", "Sum of subgrid", "Rectangle Queries".
 * - Classic Scenarios: You have an N x M grid. You need to repeatedly change the value of
 *   a specific cell (X, Y). Then, you need to query the sum of all elements inside a
 *   subgrid bounded by top-left (X1, Y1) and bottom-right (X2, Y2).
 * - The Magic: Replaces a bulky, complex 2D Segment Tree. It computes 2D prefix sums dynamically
 *   in O(log N * log M) time using ridiculously short and fast bitwise loops.
 *
 * 2.  HOW TO USE IT (THE BLACK BOX)
 * - Initialization: Pass the dimensions (rows, cols) to the constructor.
 *   (Note: Fenwick Trees MUST be 1-indexed. The bounds are [1, N] and [1, M]).
 *       FenwickTree2D fenwick(N, M);
 *
 * - Point Updates (1-indexed (x, y)):
 *       // Add 'val' to the cell at row X, col Y
 *       fenwick.update(X, Y, val);
 *
 * - Range Queries (Inclusive 1-indexed [x1, x2] and [y1, y2]):
 *       // Get the sum of all cells in the rectangle
 *       long long total = fenwick.query(x1, y1, x2, y2);
 *
 * - Complexity:
 *       Time: O(log N * log M) per update and query.
 *       Space: O(N * M) memory.
 *
 * 3.  HOW TO ADAPT IT (THE DIALS & SWITCHES)
 * - "Point Set" instead of "Point Add"?
 *   Fenwick Trees can ONLY add values. If the problem asks you to REPLACE the value at (X, Y)
 *   with `V`, you must maintain a 2D array `grid` of the current values.
 *   Then, do:
 *       long long diff = V - grid[X][Y];
 *       fenwick.update(X, Y, diff);
 *       grid[X][Y] = V;
 *
 * - 0-Indexed Input?
 *   Simply add +1 to all coordinates before passing them into `update` and `query`.
 *   NEVER change the internal 1-based indexing of the BIT.
 */

#include <bits/stdc++.h>
using namespace std;

struct FenwickTree2D
{
private:
    int n, m;
    vector<vector<long long>> bit;

    // Get the sum of the rectangle from (1, 1) to (x, y)
    long long queryPrefix(int x, int y)
    {
        long long sum = 0;
        for (int i = x; i > 0; i -= i & -i)
        {
            for (int j = y; j > 0; j -= j & -j)
            {
                sum += bit[i][j];
            }
        }
        return sum;
    }

public:
    FenwickTree2D(int rows, int cols)
    {
        n = rows;
        m = cols;
        // 1-indexed, so allocate size + 1
        bit.assign(n + 1, vector<long long>(m + 1, 0));
    }

    // Add `val` to the cell at (x, y)
    void update(int x, int y, long long val)
    {
        for (int i = x; i <= n; i += i & -i)
        {
            for (int j = y; j <= m; j += j & -j)
            {
                bit[i][j] += val;
            }
        }
    }

    // Get the sum of the rectangle from (x1, y1) to (x2, y2)
    long long query(int x1, int y1, int x2, int y2)
    {
        // Standard 2D Inclusion-Exclusion Principle
        return queryPrefix(x2, y2) - queryPrefix(x1 - 1, y2) - queryPrefix(x2, y1 - 1) + queryPrefix(x1 - 1, y1 - 1);
    }
};

void solve()
{
    int n, m, q;
    // Example Input Format: N rows, M cols, Q queries
    if (!(cin >> n >> m >> q))
        return;

    FenwickTree2D fenwick(n, m);

    // Optional: Maintain original grid if "Point Set" operations are required
    // vector<vector<long long>> grid(n + 1, vector<long long>(m + 1, 0));

    while (q--)
    {
        int type;
        cin >> type;

        if (type == 1)
        { // Point Add
            int x, y;
            long long val;
            cin >> x >> y >> val;
            fenwick.update(x, y, val);
        }
        else if (type == 2)
        { // Rectangle Query
            int x1, y1, x2, y2;
            cin >> x1 >> y1 >> x2 >> y2;
            cout << fenwick.query(x1, y1, x2, y2) << "\n";
        }
    }
}

int main()
{
    // Fast I/O
    ios_base::sync_with_stdio(false);
    cin.tie(NULL);
    solve();
    return 0;
}
\end{lstlisting}
\subsubsection{2D bit rangeUpdate rangeQuery}
\begin{lstlisting}[language=C++]
/**
 *  THE ULTIMATE CP TEMPLATE BLUEPRINT: 2D FENWICK TREE (RANGE UPDATE & QUERY)
 *
 * 1.  WHAT PROBLEM DOES THIS SOLVE?
 * - Keywords: "2D Grid", "Add X to subgrid", "Sum of subgrid", "Rectangle Updates".
 * - Classic Scenarios: You are given an N x M grid (like a terrain map or matrix).
 *   You need to repeatedly add a value V to all cells in a rectangle from (X1, Y1) to (X2, Y2).
 *   Then, you need to query the sum of all cells in a rectangle from (X1, Y1) to (X2, Y2).
 * - The Magic: A 2D Segment tree with Lazy Propagation is notoriously difficult to code and uses
 *   gigantic amounts of memory. This 2D Fenwick Tree uses 4 parallel grids to mathematically
 *   track the inclusion-exclusion bounds of the updates. It uses a fraction of the memory
 *   and runs blazingly fast in O(log N * log M) time.
 *
 * 2.  HOW TO USE IT (THE BLACK BOX)
 * - Initialization: Pass the dimensions (rows, cols) to the constructor.
 *   (Note: Fenwick Trees MUST be 1-indexed. The bounds are [1, N] and [1, M]).
 *       FenwickTree2DRange fenwick(N, M);
 *
 * - Range Updates (Inclusive 1-indexed [x1, x2] and [y1, y2]):
 *       // Add 'val' to every cell in the rectangle bounded by top-left (x1, y1) and bottom-right (x2, y2)
 *       fenwick.update(x1, y1, x2, y2, val);
 *
 * - Range Queries (Inclusive 1-indexed [x1, x2] and [y1, y2]):
 *       // Get the sum of all cells in the rectangle
 *       long long total = fenwick.query(x1, y1, x2, y2);
 *
 * - Complexity:
 *       Time: O(log N * log M) per update and query.
 *       Space: O(N * M) memory.
 *
 * 3.  HOW TO ADAPT IT (THE DIALS & SWITCHES)
 * - 0-Indexed Input?
 *   If your problem gives coordinates as 0-indexed (e.g., (0, 0) to (N-1, M-1)), simply
 *   add +1 to all coordinates before passing them into `update` and `query`.
 *   NEVER change the internal 1-based indexing of the BIT, as `index & -index` fails on 0!
 * - Modulo Arithmetic?
 *   If the problem asks for the answer modulo 10^9+7, wrap every `+=` and `-=` inside the
 *   `add` and `queryPrefix` functions with modulo logic: `res = (res % MOD + MOD) % MOD;`
 */

#include <bits/stdc++.h>
using namespace std;

struct FenwickTree2DRange
{
private:
    int n, m;

    // We need 4 parallel 2D BITs to mathematically calculate the 2D Prefix Sums
    // D[i][j], D[i][j] * i, D[i][j] * j, D[i][j] * i * j
    vector<vector<long long>> bit1, bit2, bit3, bit4;

    // Add value to coordinate (x, y) across all 4 mathematical states
    void add(int x, int y, long long val)
    {
        for (int i = x; i <= n; i += i & -i)
        {
            for (int j = y; j <= m; j += j & -j)
            {
                bit1[i][j] += val;
                bit2[i][j] += val * x;
                bit3[i][j] += val * y;
                bit4[i][j] += val * x * y;
            }
        }
    }

    // Get the sum of the rectangle from (1, 1) to (x, y)
    long long queryPrefix(int x, int y)
    {
        long long res = 0;
        for (int i = x; i > 0; i -= i & -i)
        {
            for (int j = y; j > 0; j -= j & -j)
            {
                // The mathematical expansion of 2D Difference Array Prefix Sums:
                // Sum = D * (x+1)(y+1) - (D*x) * (y+1) - (D*y) * (x+1) + (D*x*y)
                res += bit1[i][j] * (x + 1) * (y + 1) - bit2[i][j] * (y + 1) - bit3[i][j] * (x + 1) + bit4[i][j];
            }
        }
        return res;
    }

public:
    FenwickTree2DRange(int rows, int cols)
    {
        n = rows;
        m = cols;
        // 1-indexed, so allocate size + 1
        bit1.assign(n + 1, vector<long long>(m + 1, 0));
        bit2.assign(n + 1, vector<long long>(m + 1, 0));
        bit3.assign(n + 1, vector<long long>(m + 1, 0));
        bit4.assign(n + 1, vector<long long>(m + 1, 0));
    }

    // Add `val` to the rectangle from (x1, y1) to (x2, y2)
    void update(int x1, int y1, int x2, int y2, long long val)
    {
        // 2D Difference Array Inclusion-Exclusion
        add(x1, y1, val);
        add(x2 + 1, y1, -val);
        add(x1, y2 + 1, -val);
        add(x2 + 1, y2 + 1, val);
    }

    // Get the sum of the rectangle from (x1, y1) to (x2, y2)
    long long query(int x1, int y1, int x2, int y2)
    {
        // 2D Prefix Sum Inclusion-Exclusion
        return queryPrefix(x2, y2) - queryPrefix(x1 - 1, y2) - queryPrefix(x2, y1 - 1) + queryPrefix(x1 - 1, y1 - 1);
    }
};

void solve()
{
    int n, m, q;
    // Example Input Format: N rows, M cols, Q queries
    if (!(cin >> n >> m >> q))
        return;

    FenwickTree2DRange fenwick(n, m);

    while (q--)
    {
        int type;
        cin >> type;

        if (type == 1)
        { // Range Add
            int x1, y1, x2, y2;
            long long val;
            cin >> x1 >> y1 >> x2 >> y2 >> val;
            fenwick.update(x1, y1, x2, y2, val);
        }
        else if (type == 2)
        { // Range Query
            int x1, y1, x2, y2;
            cin >> x1 >> y1 >> x2 >> y2;
            cout << fenwick.query(x1, y1, x2, y2) << "\n";
        }
    }
}

int main()
{
    // Crucial fast I/O setup for dense CP test cases
    ios_base::sync_with_stdio(false);
    cin.tie(NULL);
    solve();
    return 0;
}
\end{lstlisting}
\subsubsection{Bit rangeUpdate rangeQuery}
\begin{lstlisting}[language=C++]
/**
 *  THE ULTIMATE CP TEMPLATE BLUEPRINT: 1D FENWICK TREE (RANGE UPDATE & QUERY)
 *
 * 1.  WHAT PROBLEM DOES THIS SOLVE?
 * - Keywords: "Range addition", "Subarray sum", "Add X to range [L, R]".
 * - Classic Scenarios: You are given an array. You need to repeatedly add a value V
 *   to all elements from index L to R. Then, you need to query the sum of all elements
 *   from index L to R.
 * - The Magic: Normally, you need a Lazy Segment Tree for this. However, by using
 *   two parallel Fenwick Trees, we can track the difference array and its index-multiplier
 *   to mathematically simulate Lazy Propagation. It is ridiculously fast and takes zero
 *   extra memory.
 *
 * 2.  HOW TO USE IT (THE BLACK BOX)
 * - Initialization: Pass the size of the array to the constructor.
 *   (Note: Fenwick Trees MUST be 1-indexed. The bounds are [1, N]).
 *       FenwickTreeRange fenwick(N);
 *
 * - Range Updates (Inclusive 1-indexed [L, R]):
 *       // Add 'val' to every cell from index L to R
 *       fenwick.update(L, R, val);
 *
 * - Range Queries (Inclusive 1-indexed [L, R]):
 *       // Get the sum of all cells from index L to R
 *       long long total = fenwick.query(L, R);
 *
 * - Complexity:
 *       Time: O(log N) per update and query.
 *       Space: O(N) memory.
 *
 * 3.  HOW TO ADAPT IT (THE DIALS & SWITCHES)
 * - Initial Array Values?
 *   If the problem gives you an initial array `A`, you don't need a special build function.
 *   Simply loop through the array and call `fenwick.update(i, i, A[i]);` for each element.
 * - 0-Indexed Input?
 *   If the problem uses 0-based indices, add +1 to L and R before passing them to
 *   `update` or `query`. NEVER change the internal 1-based logic, as `idx & -idx` breaks on 0.
 * - Modulo Arithmetic?
 *   If the problem asks for the sum modulo 10^9+7, wrap the `+=` and `-=` assignments inside
 *   `add` and `queryPrefix` with modulo logic: `res = (res % MOD + MOD) % MOD;`.
 */

#include <bits/stdc++.h>
using namespace std;

struct FenwickTreeRange
{
private:
    int n;

    // bit1 tracks the standard difference array: D[i]
    // bit2 tracks the difference array multiplied by the index: D[i] * i
    vector<long long> bit1, bit2;

    // Point add to both parallel BITs
    void add(int idx, long long val)
    {
        // We use the actual `idx` where the value was inserted for the bit2 multiplier
        for (int i = idx; i <= n; i += i & -i)
        {
            bit1[i] += val;
            bit2[i] += val * idx;
        }
    }

    // Get the sum from index 1 to idx
    long long queryPrefix(int idx)
    {
        long long res = 0;
        for (int i = idx; i > 0; i -= i & -i)
        {
            // The magic algebraic expansion: Sum = (idx + 1) * D[i] - (D[i] * i)
            res += bit1[i] * (idx + 1) - bit2[i];
        }
        return res;
    }

public:
    FenwickTreeRange(int size)
    {
        n = size;
        bit1.assign(n + 1, 0);
        bit2.assign(n + 1, 0);
    }

    // Add `val` to the inclusive range [l, r]
    void update(int l, int r, long long val)
    {
        // Standard difference array inclusion-exclusion logic
        add(l, val);
        add(r + 1, -val);
    }

    // Get the sum of the inclusive range [l, r]
    long long query(int l, int r)
    {
        return queryPrefix(r) - queryPrefix(l - 1);
    }
};

void solve()
{
    int n, q;
    // Example Input Format: N elements, Q queries
    if (!(cin >> n >> q))
        return;

    FenwickTreeRange fenwick(n);

    while (q--)
    {
        int type;
        cin >> type;

        if (type == 1)
        { // Range Add
            int l, r;
            long long val;
            cin >> l >> r >> val;
            fenwick.update(l, r, val);
        }
        else if (type == 2)
        { // Range Query
            int l, r;
            cin >> l >> r;
            cout << fenwick.query(l, r) << "\n";
        }
    }
}

int main()
{
    // Fast I/O for competitive programming
    ios_base::sync_with_stdio(false);
    cin.tie(NULL);
    solve();
    return 0;
}
\end{lstlisting}
\subsubsection{Fenwick tree}
\begin{lstlisting}[language=C++]
template <typename T>
struct FenwickTree
{
    int n;
    vector<T> tree;

    // 1. Initialize an empty tree of size n (all zeros)
    FenwickTree(int n)
    {
        this->n = n;
        tree.assign(n + 1, 0);
    }

    // 2. Fast O(N) initialization from an existing array
    FenwickTree(const vector<T> &a)
    {
        n = a.size();
        tree.assign(n + 1, 0);

        // Step 1: Copy values in (shifting to 1-based indexing)
        for (int i = 0; i < n; i++)
        {
            tree[i + 1] = a[i];
        }

        // Step 2: Propagate the sums forward to the immediate parent
        for (int i = 1; i <= n; i++)
        {
            int parent = i + (i & -i);
            if (parent <= n)
            {
                tree[parent] += tree[i];
            }
        }
    }

    void add(int i, T delta)
    {
        for (; i <= n; i += (i & -i))
        {
            tree[i] += delta;
        }
    }

    T query(int i)
    {
        T sum = 0;
        for (; i > 0; i -= (i & -i))
        {
            sum += tree[i];
        }
        return sum;
    }

    T query(int l, int r)
    {
        return query(r) - query(l - 1);
    }
};

\end{lstlisting}
\subsubsection{Lower bound on bit}
\begin{lstlisting}[language=C++]
/**
 *  THE ULTIMATE CP TEMPLATE BLUEPRINT: LOWER BOUND ON BIT (BINARY LIFTING)
 *
 * 1.  WHAT PROBLEM DOES THIS SOLVE?
 * - Keywords: "Find the K-th smallest element", "First index with prefix sum >= X".
 * - Classic Scenarios: You have an array of frequencies (how many times a number appears).
 *   You need to dynamically add/remove numbers and frequently ask: "What is the K-th
 *   smallest number currently in my set?"
 * - The Magic: Normally, you'd use a C++ PBDS (Policy Based Data Structure) for this.
 *   If PBDS is banned, or you need raw execution speed, this algorithm does a binary
 *   search directly ON the Fenwick tree's internal structure. It takes exactly O(log N)
 *   steps instead of O(log^2 N).
 *
 * 2.  HOW TO USE IT (THE BLACK BOX)
 * - Initialization: Pass the size of the array to the constructor.
 *       FenwickTreeLB fenwick(N);
 *
 * - Standard Point Updates (1-indexed):
 *       // Add 'val' to index 'idx'
 *       fenwick.update(idx, val);
 *
 * - The K-th Element Query (Lower Bound):
 *       // Find the smallest index where the prefix sum is >= V
 *       int ans_idx = fenwick.lower_bound(V);
 *
 * - Complexity:
 *       Time: O(log N) per update and lower_bound query.
 *       Space: O(N) memory.
 *
 * 3.  HOW TO ADAPT IT (THE DIALS & SWITCHES)
 * -  CRITICAL RULE: MONOTONICITY
 *   This function ONLY works if the original array contains ONLY non-negative numbers!
 *   If you add negative numbers, the prefix sums go up and down randomly. Binary search
 *   (and therefore this binary lifting trick) will completely fail.
 * - Out of Bounds Check:
 *   If you search for a value `V` that is larger than the total sum of the entire array,
 *   the function will return `N + 1`. Always check `if (ans_idx > N)` if you aren't sure
 *   the sum exists!
 */

#include <bits/stdc++.h>
using namespace std;

struct FenwickTreeLB
{
private:
    int n;
    int LOGN;
    vector<long long> bit;

public:
    FenwickTreeLB(int size)
    {
        n = size;
        bit.assign(n + 1, 0);

        // Find the highest power of 2 that is <= n
        LOGN = 0;
        while ((1 << (LOGN + 1)) <= n)
        {
            LOGN++;
        }
    }

    // Add `val` to the cell at `idx`
    void update(int idx, long long val)
    {
        for (; idx <= n; idx += idx & -idx)
        {
            bit[idx] += val;
        }
    }

    // Get standard prefix sum up to `idx`
    long long query(int idx)
    {
        long long sum = 0;
        for (; idx > 0; idx -= idx & -idx)
        {
            sum += bit[idx];
        }
        return sum;
    }

    // ? THE MAGIC: Binary Lifting directly on the Fenwick Tree
    // Returns the 1-based index of the FIRST element whose prefix sum is >= v
    int lower_bound(long long v)
    {
        if (v <= 0)
            return 1; // Edge case: Target is already reached at the start

        int idx = 0;
        long long current_sum = 0;

        // Jump down from the highest power of 2 to 1
        for (int i = LOGN; i >= 0; i--)
        {
            int next_idx = idx + (1 << i);

            // If taking this jump keeps us inside the array bounds,
            // AND the sum AT this jump is still STRICTLY LESS than our target `v`
            if (next_idx <= n && current_sum + bit[next_idx] < v)
            {
                // We commit to the jump!
                idx = next_idx;
                current_sum += bit[next_idx];
            }
        }

        // `idx` is now the largest index where the sum is STRICTLY LESS than `v`.
        // Therefore, `idx + 1` is the first index where the sum is >= `v`!
        return idx + 1;
    }
};

void solve()
{
    int n, q;
    // Example Input Format: Array of size N, Q queries
    if (!(cin >> n >> q))
        return;

    FenwickTreeLB fenwick(n);

    while (q--)
    {
        int type;
        cin >> type;

        if (type == 1)
        { // Point Add
            int idx;
            long long val;
            cin >> idx >> val;
            fenwick.update(idx, val);
        }
        else if (type == 2)
        { // Lower Bound Query
            long long target_sum;
            cin >> target_sum;

            int ans = fenwick.lower_bound(target_sum);

            if (ans > n)
            {
                cout << "Target sum not reached!\n";
            }
            else
            {
                cout << "Index: " << ans << "\n";
            }
        }
    }
}

int main()
{
    // Fast I/O
    ios_base::sync_with_stdio(false);
    cin.tie(NULL);
    solve();
    return 0;
}
\end{lstlisting}
\subsection{HLD}
\subsubsection{HLD}
\begin{lstlisting}[language=C++]
/**
 *  THE ULTIMATE CP TEMPLATE BLUEPRINT: HEAVY-LIGHT DECOMPOSITION (HLD)
 *
 * 1.  WHAT PROBLEM DOES THIS SOLVE?
 * - Keywords: "Path queries on a tree", "Subtree queries with point updates", "O(log^2 N) tree path".
 * - Classic Scenarios: You have a tree with values on the nodes. You need to answer $Q$ queries 
 *   that either update a node's value OR ask for the maximum/sum on the simple path between $U$ and $V$. 
 *   Standard DFS/BFS takes $O(N)$ per query, which gives TLE. 
 * - The Magic: "Tree Flattening into 1D". HLD cuts the tree into a set of disjoint vertical paths 
 *   ("heavy chains"). It then maps these chains into a 1D array so we can build a Segment Tree over it. 
 *   The mathematical guarantee is beautiful: ANY path between two nodes in the tree will cross 
 *   at most $O(\log N)$ heavy chains. Thus, a path query becomes just $O(\log N)$ segment tree queries!
 *
 * 2.  HOW TO USE IT (THE BLACK BOX)
 * - Initialization: Pass the number of nodes, the adjacency list (0-indexed or 1-indexed safely), 
 *   and a 1-indexed vector of initial node values.
 *       // Note: Define GRAPH as vector<vector<int>> beforehand
 *       HLD hld(N, adj, initial_values);
 *
 * - Point Update: Change the value of node U to X.
 *       hld.point_update(U, X);
 *
 * - Path Query: Get the maximum value on the path from U to V.
 *       int max_val = hld.query_path(U, V);
 *
 * - Subtree Query: Get the maximum value in the entire subtree of U.
 *       int sub_max = hld.query_subtree(U);
 *
 * - Complexity:
 *       Time: Build $O(N)$, Updates $O(\log N)$, Subtree Queries $O(\log N)$, Path Queries $O(\log^2 N)$.
 *       Space: $O(N)$ for the arrays and Segment Tree.
 *
 * 3.  HOW TO ADAPT IT (THE DIALS & SWITCHES)
 * - Sum/Min instead of Max? 
 *   1. Change the identity value in the `Node` constructor (e.g., `-OO` to `0` for sum, or `OO` for min).
 *   2. Change `max(lf.val, rt.val)` in `merge()` to `+` or `min`.
 *   3. Inside `HLD::query_path`, change `int res{-OO};` and `res = max(...)` to match your operation.
 * - Edges instead of Nodes?
 *   If weights are on edges, push the edge weight down to the deeper node (the child). 
 *   Then, in `query_path`, when querying the final segment between the LCA and the other node, 
 *   you must EXCLUDE the LCA (because the LCA holds the weight of the edge above it, which is not on the path). 
 *   Change `seg.get(tin[v], tin[u] + 1)` to `seg.get(tin[v] + 1, tin[u] + 1)`.
 * - GRAPH Typedef: Ensure you have `typedef vector<vector<int>> GRAPH;` globally so the constructor compiles.
 */
struct Node
{
    int val;

    Node()
    {
        val = -OO;
    }

    Node(int val) : val(val)
    {
    }

    void change(int x)
    {
        val = x;
    }
};

struct segmentTree
{
    int sz;
    vector<Node> seg;

    segmentTree(int n)
    {
        sz = 1;
        while (sz < n)
            sz *= 2;
        seg.assign(2 * sz, Node());
    }

    void init(vector<int> &arr, int node, int lx, int rx)
    {
        if (rx - lx == 1)
        {
            if (lx < int(arr.size()))
                seg[node] = Node(arr[lx]);
            return;
        }
        int mid = (lx + rx) / 2;
        init(arr, 2 * node + 1, lx, mid);
        init(arr, 2 * node + 2, mid, rx);
        seg[node] = merge(seg[2 * node + 1], seg[2 * node + 2]);
    }

    void init(vector<int> &arr)
    {
        init(arr, 0, 0, sz);
    }

    Node merge(Node &lf, Node &rt)
    {
        Node ans = Node();
        ans.val = max(lf.val, rt.val);
        return ans;
    }

    void set(int idx, int val, int node, int lx, int rx)
    {
        if (rx - lx == 1)
        {
            seg[node].change(val);
            return;
        }
        int mid = (lx + rx) / 2;
        if (idx < mid)
            set(idx, val, 2 * node + 1, lx, mid);
        else
            set(idx, val, 2 * node + 2, mid, rx);
        seg[node] = merge(seg[2 * node + 1], seg[2 * node + 2]);
    }

    void set(int idx, int val)
    {
        set(idx, val, 0, 0, sz);
    }

    Node get(int l, int r, int node, int lx, int rx)
    {
        if (lx >= l && rx <= r)
            return seg[node];
        if (rx <= l || lx >= r)
            return Node();
        int mid = (lx + rx) / 2;
        Node lf = get(l, r, 2 * node + 1, lx, mid);
        Node ri = get(l, r, 2 * node + 2, mid, rx);
        return merge(lf, ri);
    }

    int get(int l, int r)
    {
        return get(l, r, 0, 0, sz).val;
    }
};

struct HLD
{
private:
    vector<int> sz, par, dep, tin, tout, head;
    vector<vector<int>> graph;
    vector<int> a;
    int timer{1};
    segmentTree seg;

public:
    HLD(int n, GRAPH graph, vector<int> &v) : graph(std::move(graph)), seg(n + 1)
    {
        sz = vector<int>(n + 1), par = vector<int>(n + 1), dep = vector<int>(n + 1), tin = vector<int>(n + 1), tout = vector<int>(n + 1), head = vector<int>(n + 1), a = vector<int>(n + 1);
        dfs1(1, 0);
        head[1] = 1;
        dfs2(1);
        for (int i = 1; i <= n; i++)
            a[tin[i]] = v[i];
        seg.init(a);
    }

    void dfs1(int u, int p)
    {
        dep[u] = dep[p] + 1;
        par[u] = p;
        sz[u] = 1;
        if (p != 0)
            graph[u].erase(find(graph[u].begin(), graph[u].end(), p));
        for (auto &x : graph[u])
        {
            dfs1(x, u);
            sz[u] += sz[x];
            if (sz[x] > sz[graph[u][0]])
                swap(x, graph[u][0]);
        }
    }

    void dfs2(int u)
    {
        tin[u] = timer++;
        for (auto &x : graph[u])
        {
            head[x] = x == graph[u][0] ? head[u] : x;
            dfs2(x);
        }
        tout[u] = timer;
    }

    int lca(int u, int v)
    {
        while (head[u] != head[v])
        {
            if (dep[head[u]] > dep[head[v]])
                u = par[head[u]];
            else
                v = par[head[v]];
        }
        if (dep[u] < dep[v])
            return u;
        return v;
    }

    void point_update(int u, int val)
    {
        seg.set(tin[u], val);
    }

    int query_subtree(int u)
    {
        return seg.get(tin[u], tout[u]);
    }

    // now querying for maximum
    int query_path(int u, int v)
    {
        int res{-OO}; // TO
        while (head[u] != head[v])
        {
            if (dep[head[u]] < dep[head[v]])
                swap(u, v);
            res = max(res, seg.get(tin[head[u]], tin[u] + 1)); // TO
            u = par[head[u]];
        }
        if (dep[u] < dep[v])
            swap(u, v);
        res = max(res, seg.get(tin[v], tin[u] + 1)); // TO
        return res;
    }
};

\end{lstlisting}
\subsection{Misc}
\subsubsection{CDQ}
\begin{lstlisting}[language=C++]
#include <bits/stdc++.h>
using namespace std;

/**
 *  THE ULTIMATE CP TEMPLATE BLUEPRINT: CDQ DIVIDE AND CONQUER
 *
 * 1.  WHAT PROBLEM DOES THIS SOLVE?
 * - Keywords: "3D Partial Order", "Dynamic 2D Range Queries", "Offline range counting".
 * - Classic Scenarios: You need to count how many points i have (x_i <= X, y_i <= Y, z_i <= Z).
 * - The Magic: CDQ D&C reduces dimensionality.
 *   - Dimension 1 (X): Handled by sorting the array.
 *   - Dimension 2 (Y): Handled by the Divide and Conquer recursion.
 *   - Dimension 3 (Z): Handled by a 1D Fenwick Tree (BIT).
 *
 * 2.  HOW TO USE IT (THE BLACK BOX)
 * - Define a struct Element {x, y, z, id}.
 * - Sort the elements initially by x.
 * - Call cdq(0, n-1) to process.
 *
 * - Complexity:
 *       Time: O(N log^2 N) where N is the number of elements.
 *       Space: O(N) memory.
 */

struct Element
{
    int x, y, z, id;
};

// 1D Fenwick Tree for the 3rd dimension
struct BIT
{
    int n;
    vector<int> tree;
    BIT(int n) : n(n), tree(n + 1, 0) {}

    void update(int i, int delta)
    {
        for (; i <= n; i += i & -i)
            tree[i] += delta;
    }

    int query(int i)
    {
        int sum = 0;
        for (; i > 0; i -= i & -i)
            sum += tree[i];
        return sum;
    }
};

int n;
vector<Element> arr;
vector<int> result;
BIT bit(1000005); // Adjust size to coordinate range

void cdq(int l, int r)
{
    if (l >= r)
        return;
    int mid = l + (r - l) / 2;

    cdq(l, mid);
    cdq(mid + 1, r);

    // Sort both halves by Y (Dimension 2)
    vector<Element> temp;
    int i = l, j = mid + 1;

    while (i <= mid && j <= r)
    {
        if (arr[i].y <= arr[j].y)
        {
            bit.update(arr[i].z, 1);
            temp.push_back(arr[i++]);
        }
        else
        {
            result[arr[j].id] += bit.query(arr[j].z);
            temp.push_back(arr[j++]);
        }
    }

    // Process remaining
    while (j <= r)
    {
        result[arr[j].id] += bit.query(arr[j].z);
        temp.push_back(arr[j++]);
    }

    // Clear BIT before moving on
    for (int k = l; k < i; k++)
        bit.update(arr[k].z, -1);

    while (i <= mid)
        temp.push_back(arr[i++]);
    for (int k = 0; k < temp.size(); k++)
        arr[l + k] = temp[k];
}

int main()
{
    ios_base::sync_with_stdio(false);
    cin.tie(NULL);

    // Example: Read N elements with (x, y, z) coordinates
    // Ensure all coordinates are mapped to 1-indexed values for BIT

    // cdq(0, n - 1);

    return 0;
}
\end{lstlisting}
\subsubsection{Leftist heap}
\begin{lstlisting}[language=C++]
#include <bits/stdc++.h>
using namespace std;

/**
 *  THE ULTIMATE CP TEMPLATE BLUEPRINT: LEFTIST HEAP (MELDABLE)
 *
 * 1.  WHAT PROBLEM DOES THIS SOLVE?
 * - Keywords: "Meldable Heaps", "Priority Queue on Trees", "Greedy Tree DP".
 * - Classic Scenarios: You need to maintain a set of values for each node in a tree
 *   and frequently merge them with children. std::priority_queue cannot merge.
 * - The Magic: It keeps the tree "left-heavy" to guarantee that the right spine is
 *   always short (O(log N)). Merging two heaps is just merging their right spines!
 *
 * 2.  HOW TO USE IT
 * - Initialization: Root = nullptr.
 * - Merge: root = merge(root1, root2);
 * - Complexity: O(log N) per merge, push, and pop.
 */

struct Node
{
    int val, npl;
    Node *left, *right;

    Node(int v) : val(v), npl(0), left(nullptr), right(nullptr) {}
};

int getNpl(Node *n)
{
    return n ? n->npl : -1;
}

// The core MELD operation
Node *merge(Node *a, Node *b)
{
    if (!a)
        return b;
    if (!b)
        return a;

    // To maintain a Min-Heap, ensure 'a' has the smaller root
    if (a->val > b->val)
        swap(a, b);

    // ! PERSISTENCE NOTE: If you need a fully persistent heap,
    // uncomment the line below to clone the node before modifying.
    // a = new Node(*a);

    a->right = merge(a->right, b);

    // Maintain Leftist Property: Left NPL >= Right NPL
    if (getNpl(a->left) < getNpl(a->right))
    {
        swap(a->left, a->right);
    }

    a->npl = getNpl(a->right) + 1;
    return a;
}

Node *push(Node *root, int val)
{
    return merge(root, new Node(val));
}

Node *pop(Node *root)
{
    Node *left = root->left;
    Node *right = root->right;
    delete root; // Clean up memory
    return merge(left, right);
}

int top(Node *root)
{
    return root->val;
}

int main()
{
    // Example: Standard Usage
    Node *root = nullptr;
    root = push(root, 10);
    root = push(root, 5);

    cout << "Top: " << top(root) << endl; // 5
    root = pop(root);
    cout << "Top after pop: " << top(root) << endl; // 10

    return 0;
}
\end{lstlisting}
\subsubsection{Monotonic queue}
\begin{lstlisting}[language=C++]
class Monotonic_queue
{
    deque<pair<int, int>> dq;
    int k;

public:
    Monotonic_queue(int k) : k(k) {}

    void push(int idx, int val)
    {
        if (!dq.empty() && dq.front().second <= idx - k)
            dq.pop_front();
        while (!dq.empty() && dq.back().second <= val)
            dq.pop_back();
        dq.push_back({idx, val});
    }

    int max() const
    {
        return dq.front().second;
    }
};
\end{lstlisting}
\subsubsection{Monotonic stack}
\begin{lstlisting}[language=C++]
class MonotonicStack
{
public:
    static vector<int> getNextGreater(const vector<int> &nums)
    {
        int n = nums.size();
        vector<int> nextGreater(n, n);
        stack<int> s;
        for (int i = 0; i < n; ++i)
        {
            while (!s.empty() && nums[i] > nums[s.top()])
            {
                nextGreater[s.top()] = i;
                s.pop();
            }
            s.push(i);
        }
        return nextGreater;
    }
};
\end{lstlisting}
\subsubsection{ODT}
\begin{lstlisting}[language=C++]
/**
 *  THE ULTIMATE CP TEMPLATE BLUEPRINT: OLD DRIVER TREE (ODT / CHTHOLLY TREE)
 *
 * 1.  WHAT PROBLEM DOES THIS SOLVE?
 * - Keywords: "Color blocks", "Range assign", "Randomized data", "Interval set".
 * - Classic Scenarios: You have an array that is mostly composed of long segments of identical
 *   values. You need to perform Range Assignments (A[i] = V for L to R) and other Range Queries.
 * - The Magic: An ODT stores the array as a set of intervals {l, r, v}.
 *   The "Assign" operation is the key: it erases all intervals in the set between [L, R]
 *   and inserts one single new interval. This keeps the number of intervals very small.
 *
 * 2.  HOW TO USE IT (THE BLACK BOX)
 * - Initialization: Insert the initial [1, N] interval into the `odt` set.
 * - Split: You MUST call `split(l)` and `split(r+1)` before ANY operation to ensure
 *   the boundaries are exact.
 * - Complexity:
 *       Time: Amortized O(N log log N) on random data. O(N^2) on worst-case non-random data.
 *       Space: O(N) memory.
 */

#include <bits/stdc++.h>
using namespace std;

struct Node
{
    int l, r;
    mutable long long v;

    Node(int l, int r = -1, long long v = 0) : l(l), r(r), v(v) {}

    // Set must be sorted by left endpoint
    bool operator<(const Node &o) const
    {
        return l < o.l;
    }
};

set<Node> odt;

// Splits the interval containing 'pos' into [l, pos-1] and [pos, r]
// Returns an iterator to the interval starting at 'pos'
auto split(int pos)
{
    auto it = odt.lower_bound({pos});
    if (it != odt.end() && it->l == pos)
        return it;

    --it;
    int l = it->l, r = it->r;
    long long v = it->v;

    odt.erase(it);
    odt.insert({l, pos - 1, v});
    return odt.insert({pos, r, v}).first;
}

//  HOW TO ADAPT IT (THE DIALS & SWITCHES)
// - If the problem asks for Range Sum: loop through the set from itl to itr and sum (r-l+1)*v.
// - If the problem asks for K-th smallest or Count: you must keep a secondary data
//   structure (like a Segment Tree) updated alongside the ODT.
// The core efficiency of ODT: replaces range [l, r] with one single interval
void assign(int l, int r, long long v)
{
    auto itr = split(r + 1), itl = split(l);
    odt.erase(itl, itr);
    odt.insert({l, r, v});
}

void solve()
{
    int n, m;
    cin >> n >> m;

    // Initially add [1, n] with value 0
    odt.insert({1, n, 0});
}

int main()
{
    ios_base::sync_with_stdio(false);
    cin.tie(NULL);
    solve();
    return 0;
}
\end{lstlisting}
\subsection{More Data structures}
\subsubsection{Wavelet tree}
\begin{lstlisting}[language=C++]
/**
 *  THE ULTIMATE CP TEMPLATE BLUEPRINT: WAVELET TREE
 *
 * 1.  WHAT PROBLEM DOES THIS SOLVE?
 * - Keywords: "K-th smallest in range", "Range frequency count", "Count elements < K".
 * - Classic Scenarios: You have an array and you need to answer queries like:
 *   "Find the k-th smallest element in subarray [L, R]" or "Count occurrences
 *   of values in [A, B] within index range [L, R]".
 * - The Magic: It is essentially a Merge Sort Tree, but instead of storing
 *   sorted vectors (which consume O(N log N) memory), it stores a bit-vector
 *   representation of the partitioning. This allows queries in O(log(MaxVal)) time.
 *
 * 2.  HOW TO USE IT (THE BLACK BOX)
 * - Initialization:
 *       // Pass the range of values [min_val, max_val] and the input vector
 *       WaveletTree wt(1, 1000000, array);
 *
 * - Queries:
 *       // How many numbers < K in subarray [L, R]
 *       int count = wt.count_less(L, R, K);
 *       // Find the k-th smallest number in subarray [L, R] (1-indexed k)
 *       int val = wt.kth_smallest(L, R, k);
 *
 * - Complexity:
 *       Time: O(N log(MaxVal)) to build. O(log(MaxVal)) per query.
 *       Space: O(N log(MaxVal)) memory.
 */

#include <bits/stdc++.h>
using namespace std;

struct WaveletTree
{
    int low, high;
    WaveletTree *left, *right;
    vector<int> pref; // Prefix sums of elements going to the left child

    WaveletTree(int l, int r, const vector<int> &a) : low(l), high(r), left(nullptr), right(nullptr)
    {
        if (low == high || a.empty())
            return;

        int mid = low + (high - low) / 2;

        // Define function to check if value goes to left child
        auto is_left = [&](int x)
        { return x <= mid; };

        pref.reserve(a.size() + 1);
        pref.push_back(0);

        vector<int> l_vec, r_vec;
        for (int x : a)
        {
            if (is_left(x))
            {
                l_vec.push_back(x);
                pref.push_back(pref.back() + 1);
            }
            else
            {
                r_vec.push_back(x);
                pref.push_back(pref.back());
            }
        }

        left = new WaveletTree(low, mid, l_vec);
        right = new WaveletTree(mid + 1, high, r_vec);
    }

    // Count numbers <= k in range [ql, qr] (0-indexed)
    int count_less_equal(int ql, int qr, int k)
    {
        if (ql > qr || low > k)
            return 0;
        if (high <= k)
            return qr - ql + 1;

        int l_count_before = pref[ql];
        int r_count_before = ql - l_count_before;
        int l_count_all = pref[qr + 1];
        int r_count_all = (qr + 1) - l_count_all;

        return left->count_less_equal(l_count_before, l_count_all - 1, k) +
               right->count_less_equal(r_count_before, r_count_all - 1, k);
    }

    // Find k-th smallest number in range [ql, qr]
    int kth_smallest(int ql, int qr, int k)
    {
        if (low == high)
            return low;

        int l_count_before = pref[ql];
        int r_count_before = ql - l_count_before;
        int l_count_all = pref[qr + 1];
        int r_count_all = (qr + 1) - l_count_all;

        int left_size = l_count_all - l_count_before;

        if (k <= left_size)
            return left->kth_smallest(l_count_before, l_count_all - 1, k);
        else
            return right->kth_smallest(r_count_before, r_count_all - 1, k - left_size);
    }
};
\end{lstlisting}
\subsection{SQRT decomposition MO}
\subsubsection{MO}
\begin{lstlisting}[language=C++]

// how many distinct values are there in a range [l,r]

int SQ = 400;
struct Query
{
    int l, r, idx;
    bool operator<(const Query &q) const
    {
        if (l / SQ != q.l / SQ)
        {
            return l / SQ < q.l / SQ;
        }
        return (l / SQ & 1) ? r < q.r : r > q.r;
    }
};

void solve()
{
    int n, q;
    cin >> n >> q;
    vector<int> a(n), frq(n + 10);
    map<int, int> mp;
    int id{};
    for (int i = 0; i < n; i++)
    {
        cin >> a[i];
        if (!mp.count(a[i]))
        {
            mp[a[i]] = id++;
        }
    }
    for (int i = 0; i < n; i++)
        a[i] = mp[a[i]];

    vector<Query> query(q);
    for (int i = 0; i < q; i++)
    {
        cin >> query[i].l >> query[i].r;
        query[i].l--, query[i].r--;
        query[i].idx = i;
    }
    sort(query.begin(), query.end());
    int ans = 0;
    auto add = [&](int i)
    {
        // if (frq[a[i]] == 0)
        // {
        //     ans++;
        // }
        // frq[a[i]]++;
    };
    auto remove = [&](int i)
    {
        // frq[a[i]]--;
        // if (frq[a[i]] == 0)
        // {
        //     ans--;
        // }
    };

    vector<int> res(q);
    int l = 0, r = -1;
    for (auto [ql, qr, idx] : query)
    {
        while (l > ql)
        {
            --l;
            add(l);
        }
        while (r < qr)
        {
            ++r;
            add(r);
        }
        while (l < ql)
        {
            remove(l);
            ++l;
        }
        while (r > qr)
        {
            remove(r);
            --r;
        }
        res[idx] = ans;
    }
    for (int i = 0; i < q; i++)
        cout << res[i] << "\n";
}
\end{lstlisting}
\subsubsection{MO on tree}
\begin{lstlisting}[language=C++]
/**
 *  THE ULTIMATE CP TEMPLATE BLUEPRINT: MO'S ALGORITHM ON TREES
 *
 * 1.  WHAT PROBLEM DOES THIS SOLVE?
 * - Keywords: "Path queries", "Distinct values on path", "Offline tree queries".
 * - Classic Scenarios: You are given a tree where each node has a color. You are asked Q queries:
 *   "How many distinct colors exist on the simple path between node U and node V?"
 * - The Magic: We flatten the tree using an Euler Tour. Every node appears twice in the array:
 *   once when we enter it (tin), and once when we leave it (tout).
 *   If we want the path between U and V (where tin[U] < tin[V]), we query the flat array:
 *   A) If U is an ancestor of V: Query range [tin[U], tin[V]].
 *   B) If U is NOT an ancestor: Query range [tout[U], tin[V]] PLUS the LCA(U, V).
 *   Because dead-end branches will appear twice in this range (entry and exit), we simply
 *   TOGGLE their state. 1st appearance = Add. 2nd appearance = Remove.
 *
 * 2.  HOW TO USE IT (THE BLACK BOX)
 * - Initialization: Pass your tree edges and node values, then run `dfs(1, 0)`.
 * - L and R Logic: We maintain an `in_path` boolean array. The `toggle()` function handles
 *   adding OR removing automatically based on the boolean!
 *
 * - Complexity:
 *       Time: O(N log N) to setup LCA + O((N + Q) * sqrt(2N)) for Mo's processing.
 *       Space: O(N log N) for LCA lifting table + O(N) arrays.
 *
 * 3.  HOW TO ADAPT IT (THE DIALS & SWITCHES)
 * - Changing the logic:
 *   Right now, the template counts "Distinct Colors" using a frequency array.
 *   Modify the `add()` and `remove()` functions inside `toggle()` to track whatever
 *   property your specific problem requires.
 */

#include <bits/stdc++.h>
using namespace std;

const int MAXN = 100005;
const int LOG = 20;

int n, q;
vector<int> adj[MAXN];
long long val[MAXN]; // The color or value of each node

// Euler Tour Variables
int timer_tour = 0;
int tin[MAXN], tout[MAXN];
int tour[2 * MAXN]; // The flattened 2N array

// LCA Variables
int up[MAXN][LOG];
int depth[MAXN];

void dfs(int u, int p)
{
    // 1. Record Entry Time
    tin[u] = ++timer_tour;
    tour[timer_tour] = u;

    // 2. Build Binary Lifting Table for LCA
    up[u][0] = p;
    for (int i = 1; i < LOG; i++)
    {
        up[u][i] = up[up[u][i - 1]][i - 1];
    }

    // 3. Traverse Children
    for (int v : adj[u])
    {
        if (v != p)
        {
            depth[v] = depth[u] + 1;
            dfs(v, u);
        }
    }

    // 4. Record Exit Time
    tout[u] = ++timer_tour;
    tour[timer_tour] = u;
}

int get_lca(int u, int v)
{
    if (depth[u] < depth[v])
        swap(u, v);

    int diff = depth[u] - depth[v];
    for (int i = 0; i < LOG; i++)
    {
        if ((diff >> i) & 1)
            u = up[u][i];
    }

    if (u == v)
        return u;

    for (int i = LOG - 1; i >= 0; i--)
    {
        if (up[u][i] != up[v][i])
        {
            u = up[u][i];
            v = up[v][i];
        }
    }
    return up[u][0];
}

int BLOCK_SIZE;

struct Query
{
    int l, r, id, lca;

    bool operator<(const Query &other) const
    {
        int block_own = l / BLOCK_SIZE;
        int block_other = other.l / BLOCK_SIZE;
        if (block_own != block_other)
        {
            return block_own < block_other;
        }
        // Parity optimization for slightly faster R pointer movement
        return (block_own & 1) ? r < other.r : r > other.r;
    }
};

bool in_path[MAXN];
long long freq[MAXN]; // Example: Frequency array for colors
long long distinct_colors = 0;

// ? THE MAGIC: Toggling a node.
// If it's entering our path, we add it. If we hit its exit time, it cancels out and we remove it!
void toggle(int node)
{
    long long c = val[node];

    if (in_path[node])
    {
        // It was in the path, so we are removing it
        freq[c]--;
        if (freq[c] == 0)
            distinct_colors--;
    }
    else
    {
        // It wasn't in the path, so we are adding it
        if (freq[c] == 0)
            distinct_colors++;
        freq[c]++;
    }

    in_path[node] ^= 1; // Flip the boolean state
}

void solve()
{
    if (!(cin >> n >> q))
        return;

    // Reset globals for multiple test cases (if needed)
    timer_tour = 0;
    distinct_colors = 0;
    memset(in_path, 0, sizeof(in_path));
    memset(freq, 0, sizeof(freq));

    // Coordinate compression for colors might be required here if values > 10^5!
    for (int i = 1; i <= n; i++)
        cin >> val[i];

    for (int i = 0; i < n - 1; i++)
    {
        int u, v;
        cin >> u >> v;
        adj[u].push_back(v);
        adj[v].push_back(u);
    }

    // 1. Flatten the tree!
    depth[1] = 0;
    dfs(1, 1);

    // Block size for an array of size 2N
    BLOCK_SIZE = max(1, (int)sqrt(2 * n));

    vector<Query> queries(q);
    for (int i = 0; i < q; i++)
    {
        int u, v;
        cin >> u >> v;

        // Ensure u is visited before v in the Euler Tour
        if (tin[u] > tin[v])
            swap(u, v);

        int lca = get_lca(u, v);

        if (lca == u)
        {
            // SCENARIO A: U is an ancestor of V.
            // The path is strictly the sequence of entries from U down to V.
            queries[i].l = tin[u];
            queries[i].r = tin[v];
            queries[i].lca = -1; // No extra LCA needed
        }
        else
        {
            // SCENARIO B: U and V are in different branches.
            // We take the path from the exit of U, up to the LCA, and down to the entry of V.
            queries[i].l = tout[u];
            queries[i].r = tin[v];
            queries[i].lca = lca; // We MUST manually include the LCA!
        }
        queries[i].id = i;
    }

    // 2. Sort Queries
    sort(queries.begin(), queries.end());

    vector<long long> ans(q);

    // 3. Run Mo's Pointers on the 2N Tour Array
    int L = 1, R = 0;

    for (const Query &qry : queries)
    {
        while (L > qry.l)
        {
            L--;
            toggle(tour[L]);
        }
        while (R < qry.r)
        {
            R++;
            toggle(tour[R]);
        }
        while (L < qry.l)
        {
            toggle(tour[L]);
            L++;
        }
        while (R > qry.r)
        {
            toggle(tour[R]);
            R--;
        }

        // If U and V were on different branches, the LCA itself was NOT included
        // in the [tout[U], tin[V]] range. We must temporarily add it!
        if (qry.lca != -1)
            toggle(qry.lca);

        // Record the answer
        ans[qry.id] = distinct_colors;

        // Revert the temporary LCA addition to keep the L/R pointers clean
        if (qry.lca != -1)
            toggle(qry.lca);
    }

    for (int i = 0; i < q; i++)
    {
        cout << ans[i] << "\n";
    }
}

int main()
{
    ios_base::sync_with_stdio(false);
    cin.tie(NULL);
    solve();
    return 0;
}
\end{lstlisting}
\subsubsection{Mo's with updates}
\begin{lstlisting}[language=C++]
/**
 *  THE ULTIMATE CP TEMPLATE BLUEPRINT: 3D MO'S ALGORITHM (WITH UPDATES)
 *
 * 1.  WHAT PROBLEM DOES THIS SOLVE?
 * - Keywords: "Point Updates", "Range Queries", "Offline", "Distinct elements with updates".
 * - Classic Scenarios: You are given an array of size N. You have two types of queries:
 *   Type 1: Query the number of distinct elements in range [L, R].
 *   Type 2: Change the element at index X to V.
 * - The Magic: A standard Segment Tree cannot count "Distinct Elements". Standard Mo's
 *   Algorithm cannot handle "Updates". 3D Mo combines them by adding a "Time" pointer (T).
 *   Queries are sorted into 3D blocks: [L_block, R_block, Time].
 *
 * 2.  HOW TO USE IT (THE BLACK BOX)
 * - Initialization:
 *   As you read the input, split it! If it's an update, push to the `updates` array.
 *   If it's a query, push to the `queries` array and record the CURRENT size of the
 *   `updates` array as its "Time" stamp!
 *
 * - Complexity:
 *       Time: O(N^{5/3}) optimally.
 *       Space: O(N + Q) memory.
 *
 * 3.  HOW TO ADAPT IT (THE DIALS & SWITCHES)
 * - The Block Size (B):
 *   Standard Mo uses B = N^{1/2}. 3D Mo MUST use B = N^{2/3}.
 *   If N = 100,000, N^{2/3} is roughly 2154. The template calculates this dynamically!
 * - Coordinate Compression:
 *   This template assumes array values are <= 10^6. If values are up to 10^9, you MUST
 *   coordinate compress the initial array AND all the update values before running Mo!
 */

#include <bits/stdc++.h>
using namespace std;

const int MAXN = 200005;
const int MAX_VAL = 1000005; // Maximum possible value in the array

int n, q;
int a[MAXN];         // The active array
int freq[MAX_VAL];   // Frequency array for distinct element counting
int current_ans = 0; // Tracks the current number of distinct elements
int BLOCK_SIZE;

struct Update
{
    int idx;
    int val; // When an update is applied, this holds the NEW value.
             // When reverted, it automatically swaps to hold the OLD value!
};

struct Query
{
    int l, r, t, id;
    int l_block, r_block;

    // Sort by L block, then R block, then Time.
    // Includes a 3D parity optimization for faster pointer movement!
    bool operator<(const Query &other) const
    {
        if (l_block != other.l_block)
            return l_block < other.l_block;

        if (r_block != other.r_block)
        {
            return (l_block & 1) ? r_block < other.r_block : r_block > other.r_block;
        }

        return (r_block & 1) ? t < other.t : t > other.t;
    }
};

vector<Update> updates;
vector<Query> queries;
vector<int> ans;

// Add an element at index 'idx' to our active window
inline void add(int idx)
{
    if (freq[a[idx]] == 0)
    {
        current_ans++;
    }
    freq[a[idx]]++;
}

// Remove an element at index 'idx' from our active window
inline void remove(int idx)
{
    freq[a[idx]]--;
    if (freq[a[idx]] == 0)
    {
        current_ans--;
    }
}

// ? THE MAGIC: The Swap Trick
// If we move forward in time, we swap the current array value with the update's value.
// Because we swap, the `updates[t_idx].val` now holds the OLD value!
// This means if we need to travel backward in time, we just call the exact same function!
inline void apply_time(int t_idx, int L, int R)
{
    int idx = updates[t_idx].idx;

    // If the index being updated is currently INSIDE our active [L, R] window,
    // we must remove its contribution using the old value.
    if (L <= idx && idx <= R)
    {
        remove(idx);
    }

    // Swap the array value and the update cache!
    swap(a[idx], updates[t_idx].val);

    // If the index is in our active window, add its contribution using the new value!
    if (L <= idx && idx <= R)
    {
        add(idx);
    }
}

void solve()
{
    if (!(cin >> n >> q))
        return;

    for (int i = 1; i <= n; i++)
    {
        cin >> a[i];
    }

    // 1-based indexing for updates is safer for the Time pointer
    updates.push_back({0, 0}); // Dummy update at T = 0

    // Parse queries and updates
    for (int i = 0; i < q; i++)
    {
        int type;
        cin >> type;
        if (type == 1)
        {
            // Type 1: Range Query (L, R)
            int l, r;
            cin >> l >> r;
            queries.push_back({l, r, (int)updates.size() - 1, (int)queries.size(), 0, 0});
        }
        else
        {
            // Type 2: Point Update (Index X becomes Value V)
            int idx, val;
            cin >> idx >> val;
            updates.push_back({idx, val});
        }
    }

    // ? THE MATH: Optimal block size for 3D Mo is N^(2/3)
    BLOCK_SIZE = max(1, (int)pow(n, 2.0 / 3.0));

    for (auto &qry : queries)
    {
        qry.l_block = qry.l / BLOCK_SIZE;
        qry.r_block = qry.r / BLOCK_SIZE;
    }

    sort(queries.begin(), queries.end());
    ans.resize(queries.size());

    int L = 1, R = 0, T = 0;

    for (const auto &qry : queries)
    {
        // Step 1: Expand the spatial window [L, R]
        while (L > qry.l)
        {
            L--;
            add(L);
        }
        while (R < qry.r)
        {
            R++;
            add(R);
        }

        // Step 2: Shrink the spatial window [L, R]
        while (L < qry.l)
        {
            remove(L);
            L++;
        }
        while (R > qry.r)
        {
            remove(R);
            R--;
        }

        // Step 3: Fast-forward or Rewind the Time window (T)
        // We do Time AFTER Space to ensure the active window bounds are correct
        // for the `apply_time` intersection check!
        while (T < qry.t)
        {
            T++;
            apply_time(T, L, R);
        }
        while (T > qry.t)
        {
            apply_time(T, L, R);
            T--;
        }

        // Record the answer for this specific snapshot in space and time
        ans[qry.id] = current_ans;
    }

    for (int i = 0; i < ans.size(); i++)
    {
        cout << ans[i] << "\n";
    }
}

int main()
{
    // Fast I/O is mandatory due to heavy pointer manipulation
    ios_base::sync_with_stdio(false);
    cin.tie(NULL);
    solve();
    return 0;
}
\end{lstlisting}
\subsubsection{Rollback mo}
\begin{lstlisting}[language=C++]
/**
 *  THE ULTIMATE CP TEMPLATE BLUEPRINT: ROLLBACK MO'S ALGORITHM
 *
 * 1.  WHAT PROBLEM DOES THIS SOLVE?
 * - Keywords: "Range queries", "Add is easy, Remove is impossible", "DSU on ranges".
 * - Classic Scenarios: You are given a graph with N edges. You receive queries asking:
 *   "If we only use edges from index L to R, how many connected components are there?"
 * - The Magic: Standard Mo's algorithm requires you to slide both the L and R pointers
 *   by adding AND removing elements. But you cannot "remove" an edge from a standard DSU!
 *   Rollback Mo sorts the queries so that the R pointer ONLY moves forward (add only).
 *   The L pointer moves backward temporarily, and we "undo" those L moves using a stack
 *   history (Rollback) before moving to the next query. NO `remove()` function needed!
 *
 * 2.  HOW TO USE IT (THE BLACK BOX)
 * - Define your state: You need a data structure that supports `add()`, `save_state()`,
 *   and `rollback(state)`. (I have included a Rollback DSU in this template).
 * - Read Queries: Push all queries into a vector and pass them to the Mo's runner.
 *
 * - Complexity:
 *       Time: O(Q * sqrt(N) * Cost(Add)). If DSU, it's O(Q * sqrt(N) * log N).
 *       Space: O(N + Q) memory.
 *
 * 3.  HOW TO ADAPT IT (THE DIALS & SWITCHES)
 * - The Block Size (B):
 *   The math for Mo's dictates the block size should be roughly N / sqrt(Q).
 *   If the problem has N = 10^5, setting `B = 350` or `B = 400` is the standard sweet spot.
 * - Changing the Data Structure:
 *   If the problem asks for "Max frequency in range", you don't need DSU. You just need
 *   an array tracking frequencies, and a variable `max_freq`. To rollback, you just
 *   remember the old `max_freq` and revert the array.
 */

#include <bits/stdc++.h>
using namespace std;

// 1. The Rollback Data Structure
struct RollbackDSU
{
    vector<int> parent, sz;
    int components;

    // We store the history of changes: {node_u, node_v, previous_components}
    struct Change
    {
        int u, v, comps;
    };
    stack<Change> history;

    RollbackDSU(int n)
    {
        parent.resize(n + 1);
        sz.assign(n + 1, 1);
        iota(parent.begin(), parent.end(), 0);
        components = n;
    }

    // ! CRITICAL: No Path Compression allowed! Only Union by Size.
    int find(int i)
    {
        while (i != parent[i])
            i = parent[i];
        return i;
    }

    void unite(int u, int v)
    {
        int root_u = find(u);
        int root_v = find(v);

        if (root_u == root_v)
        {
            // Even if they are already connected, we push a dummy change to keep
            // the stack height aligned with the number of `add()` calls.
            history.push({-1, -1, components});
            return;
        }

        if (sz[root_u] < sz[root_v])
            swap(root_u, root_v);

        // Record the exact state BEFORE we change it
        history.push({root_u, root_v, components});

        // Apply changes
        parent[root_v] = root_u;
        sz[root_u] += sz[root_v];
        components--;
    }

    // Returns the current "timestamp" or "height" of the history stack
    int save()
    {
        return history.size();
    }

    // Undos all operations until the stack shrinks back to `target_time`
    void rollback(int target_time)
    {
        while (history.size() > target_time)
        {
            auto [u, v, old_comps] = history.top();
            history.pop();

            if (u != -1)
            {
                sz[u] -= sz[v];
                parent[v] = v;
            }
            components = old_comps;
        }
    }

    void reset(int n)
    {
        while (!history.empty())
            history.pop();
        for (int i = 1; i <= n; i++)
        {
            parent[i] = i;
            sz[i] = 1;
        }
        components = n;
    }
};

// 2. Query Structure and Mo's Sorting logic
int BLOCK_SIZE;

struct Query
{
    int l, r, id;

    // Sort queries by their Left Block, then by their Right Endpoint
    bool operator<(const Query &other) const
    {
        int block_own = l / BLOCK_SIZE;
        int block_other = other.l / BLOCK_SIZE;
        if (block_own != block_other)
        {
            return block_own < block_other;
        }
        // If in the same block, sort R strictly ascending
        return r < other.r;
    }
};

struct Edge
{
    int u, v;
};

void solve()
{
    int n, m, q;
    // N vertices, M edges, Q queries
    if (!(cin >> n >> m >> q))
        return;

    // Standard block size for 10^5 is roughly 350
    BLOCK_SIZE = max(1, (int)sqrt(m));

    vector<Edge> edges(m + 1); // 1-indexed edges
    for (int i = 1; i <= m; i++)
    {
        cin >> edges[i].u >> edges[i].v;
    }

    vector<Query> queries(q);
    for (int i = 0; i < q; i++)
    {
        cin >> queries[i].l >> queries[i].r;
        queries[i].id = i;
    }

    // 1. Sort the queries
    sort(queries.begin(), queries.end());

    vector<int> ans(q);
    RollbackDSU dsu(n);

    int last_block = -1;
    int current_r = 0;

    for (const Query &query : queries)
    {
        int current_block = query.l / BLOCK_SIZE;

        // If we entered a brand new block, completely wipe the DSU clean!
        if (current_block != last_block)
        {
            dsu.reset(n);
            // Set our R pointer to the very edge of the current block
            current_r = min(m, (current_block + 1) * BLOCK_SIZE);
            last_block = current_block;
        }

        // SCENARIO 1: The query is incredibly tiny and fits entirely inside one block
        // We answer it naively in O(Sqrt(N)) and instantly roll it back.
        if (query.r < (current_block + 1) * BLOCK_SIZE)
        {
            int time_stamp = dsu.save();
            for (int i = query.l; i <= query.r; i++)
            {
                dsu.unite(edges[i].u, edges[i].v);
            }
            ans[query.id] = dsu.components;
            dsu.rollback(time_stamp);
            continue;
        }

        // SCENARIO 2: The query stretches outside the block.
        // Step A: Expand our R pointer forward until it hits the query's R limit.
        // Notice we NEVER rollback these operations while staying in the same block!
        while (current_r <= query.r)
        {
            dsu.unite(edges[current_r].u, edges[current_r].v);
            current_r++;
        }

        // Step B: We temporarily drag the L pointer backward from the block edge
        // to cover the left side of the query.
        int time_stamp = dsu.save();
        for (int i = query.l; i < (current_block + 1) * BLOCK_SIZE; i++)
        {
            dsu.unite(edges[i].u, edges[i].v);
        }

        // Step C: Record the answer!
        ans[query.id] = dsu.components;

        // Step D: Rollback the temporary L pointer movements.
        // The R pointer progress remains saved for the next queries!
        dsu.rollback(time_stamp);
    }

    for (int i = 0; i < q; i++)
    {
        cout << ans[i] << "\n";
    }
}

int main()
{
    // Fast I/O is mandatory for heavy Mo's algorithm problems
    ios_base::sync_with_stdio(false);
    cin.tie(NULL);
    solve();
    return 0;
}
\end{lstlisting}
\subsubsection{sqrt decomposition}
\begin{lstlisting}[language=C++]

struct SQRT
{
    vector<int> arr, blk;
    int SQ;
    SQRT(vector<int> &a, int sq = 450)
    {
        SQ = sq;
        arr = a;
        int n = arr.size();
        blk.resize(n / sq + 5);
        for (int i = 0; i < n; i++)
        {
            blk[i / sq] += arr[i];
        }
    }
    void update(int idx, int val)
    {
        int blkNum = idx / SQ;
        blk[blkNum] -= arr[idx];
        arr[idx] = val;
        blk[blkNum] += arr[idx];
    }

    int query(int r)
    {
        int res = 0;
        int blkNum = r / SQ;
        for (int i = 0; i < blkNum; i++)
        {
            res += blk[i];
        }
        for (int i = blkNum * SQ; i <= r; i++)
        {
            res += arr[i];
        }
        return res;
    }

    int query(int l, int r)
    {
        return query(r) - query(l - 1);
    }
    /*
    int query(int l,int r) {
        int res = oo;
        while (l<=r) {
            if (l % SQ == 0 && l+SQ <= r) {
                int blkNum =l/SQ;
             //   res = min(res, blk[blkNum]);
                l+=SQ;
            }
            else {
               // res = min(res,arr[l]);
                l++;
            }
        }
        return res;
    }
    */
};
\end{lstlisting}
\subsubsection{sqrt decomposition rangeUpdate rangeQuery}
\begin{lstlisting}[language=C++]
struct SqrtRangeUpdateRangeQuery
{
    int n, S;
    vector<long long> a, b, c; // b = lazy, c = sum of block

    SqrtRangeUpdateRangeQuery(int n, const vector<long long> &arr)
    {
        this->n = n;
        S = ceil(sqrt(n));
        a = arr;
        b.assign(S + 1, 0);
        c.assign(S + 1, 0);
        for (int i = 0; i < n; ++i)
        {
            c[i / S] += a[i];
        }
    }

    // Helper to get block actual sum safely
    long long get_block_sum(int k)
    {
        // Block sum = original sum + (lazy value * number of elements in block)
        long long elements_count = min(n - k * S, S);
        return c[k] + b[k] * elements_count;
    }

    // O(sqrt(N))
    void add(int l, int r, long long delta)
    {
        int c_l = l / S, c_r = r / S;

        if (c_l == c_r)
        {
            for (int i = l; i <= r; ++i)
            {
                a[i] += delta;
                c[c_l] += delta;
            }
        }
        else
        {
            for (int i = l, end = (c_l + 1) * S - 1; i <= end; ++i)
            {
                a[i] += delta;
                c[c_l] += delta;
            }
            for (int k = c_l + 1; k <= c_r - 1; ++k)
            {
                b[k] += delta; // Lazy update
            }
            for (int i = c_r * S; i <= r; ++i)
            {
                a[i] += delta;
                c[c_r] += delta;
            }
        }
    }

    // O(sqrt(N))
    long long query(int l, int r)
    {
        long long sum = 0;
        int c_l = l / S, c_r = r / S;

        if (c_l == c_r)
        {
            for (int i = l; i <= r; ++i)
                sum += a[i] + b[c_l];
        }
        else
        {
            for (int i = l, end = (c_l + 1) * S - 1; i <= end; ++i)
                sum += a[i] + b[c_l];

            for (int k = c_l + 1; k <= c_r - 1; ++k)
                sum += get_block_sum(k); // Full block sum

            for (int i = c_r * S; i <= r; ++i)
                sum += a[i] + b[c_r];
        }
        return sum;
    }
};
\end{lstlisting}
\subsection{Segment Tree}
\subsubsection{2D segment tree}
\begin{lstlisting}[language=C++]
/**
 *  THE ULTIMATE CP TEMPLATE BLUEPRINT: 2D SEGMENT TREE
 *
 * 1.  WHAT PROBLEM DOES THIS SOLVE?
 * - Keywords: "Subgrid queries", "Rectangle Sum/Max/Min", "Update a single cell in a grid".
 * - Classic Scenarios: You are given an N x M matrix. You need to repeatedly answer
 *   queries about a specific sub-rectangle (from top-left [R1, C1] to bottom-right [R2, C2]).
 *   Between queries, you can update the value of any specific cell [R, C].
 * - The Magic: A standard Segment Tree handles a 1D array. A 2D Segment Tree is literally
 *   a "Segment Tree of Segment Trees". The outer tree splits the rows, and every single
 *   node in the outer tree contains an inner tree that splits the columns for that row-range.
 *
 * 2.  HOW TO USE IT (THE BLACK BOX)
 * - Initialization: Pass the row and column sizes to the constructor.
 *   Then, (optionally) build it using an existing 0-indexed 2D vector.
 *       vector<vector<long long>> grid = {{1, 2}, {3, 4}};
 *       SegTree2D st(2, 2);
 *       st.build(grid);
 *
 * - Queries (All 0-indexed, Inclusive [r1, r2] and [c1, c2]):
 *       // Query the sum of the subgrid from top-left (0, 0) to bottom-right (1, 1)
 *       long long total = st.query(0, 0, 1, 1);
 *
 * - Updates (All 0-indexed):
 *       // Change the value at row R, column C to X
 *       st.update(R, C, X);
 *
 * - Complexity:
 *       Time: O(N * M) to build. O(log N * log M) per update and query.
 *       Space: O(N * M) memory. (Uses exactly 16 * N * M elements).
 *
 * 3.  HOW TO ADAPT IT (THE DIALS & SWITCHES)
 * - Changing from Range Sum to Range Min/Max?
 *   1. Change the `NEUTRAL` variable to `1e18` (for Min) or `-1e18` (for Max).
 *   2. Change the `combine()` function to return `min(a, b)` or `max(a, b)`.
 *
 * - Memory Limit Exceeded (MLE) Warning:
 *   If N = 1000 and M = 1000, this tree uses 16,000,000 long longs (~128 MB).
 *   This easily passes standard 256MB limits.
 *   BUT, if N, M > 2000, you will MLE. If that happens and the problem only asks for SUM,
 *   you must use a 2D Fenwick Tree (BIT) instead, which uses 1/16th of the memory!
 */

#include <bits/stdc++.h>
using namespace std;

struct SegTree2D
{
private:
    int n, m;
    vector<vector<long long>> seg;

    // ? Change these two dials to switch between Sum, Min, Max, GCD, etc.
    const long long NEUTRAL = 0;

    long long combine(long long a, long long b)
    {
        return a + b;
    }

    void build_y(int vx, int lx, int rx, int vy, int ly, int ry, const vector<vector<long long>> &a)
    {
        if (ly == ry)
        {
            // We are at a leaf in the column tree.
            if (lx == rx)
            {
                // We are ALSO at a leaf in the row tree. Base case!
                seg[vx][vy] = a[lx][ly];
            }
            else
            {
                // We are a combination of two row-segments at the SAME exact column.
                seg[vx][vy] = combine(seg[vx * 2][vy], seg[vx * 2 + 1][vy]);
            }
        }
        else
        {
            int my = ly + (ry - ly) / 2;
            build_y(vx, lx, rx, vy * 2, ly, my, a);
            build_y(vx, lx, rx, vy * 2 + 1, my + 1, ry, a);
            // Combine the left and right column-halves for the current row-segment
            seg[vx][vy] = combine(seg[vx][vy * 2], seg[vx][vy * 2 + 1]);
        }
    }

    void build_x(int vx, int lx, int rx, const vector<vector<long long>> &a)
    {
        if (lx != rx)
        {
            int mx = lx + (rx - lx) / 2;
            build_x(vx * 2, lx, mx, a);
            build_x(vx * 2 + 1, mx + 1, rx, a);
        }
        // Build the entire inner column-tree for the current row-segment
        build_y(vx, lx, rx, 1, 0, m - 1, a);
    }

    void update_y(int vx, int lx, int rx, int vy, int ly, int ry, int x, int y, long long val)
    {
        if (ly == ry)
        {
            if (lx == rx)
            {
                // Exact cell located!
                seg[vx][vy] = val;
            }
            else
            {
                // Update the combined row-segment value for this specific column
                seg[vx][vy] = combine(seg[vx * 2][vy], seg[vx * 2 + 1][vy]);
            }
        }
        else
        {
            int my = ly + (ry - ly) / 2;
            if (y <= my)
            {
                update_y(vx, lx, rx, vy * 2, ly, my, x, y, val);
            }
            else
            {
                update_y(vx, lx, rx, vy * 2 + 1, my + 1, ry, x, y, val);
            }
            seg[vx][vy] = combine(seg[vx][vy * 2], seg[vx][vy * 2 + 1]);
        }
    }

    void update_x(int vx, int lx, int rx, int x, int y, long long val)
    {
        if (lx != rx)
        {
            int mx = lx + (rx - lx) / 2;
            if (x <= mx)
            {
                update_x(vx * 2, lx, mx, x, y, val);
            }
            else
            {
                update_x(vx * 2 + 1, mx + 1, rx, x, y, val);
            }
        }
        // Force the inner Y-tree to update itself using the new X data
        update_y(vx, lx, rx, 1, 0, m - 1, x, y, val);
    }

    long long query_y(int vx, int vy, int ly, int ry, int qly, int qry)
    {
        if (qly > ry || qry < ly)
            return NEUTRAL;
        if (qly <= ly && ry <= qry)
            return seg[vx][vy];

        int my = ly + (ry - ly) / 2;
        return combine(
            query_y(vx, vy * 2, ly, my, qly, qry),
            query_y(vx, vy * 2 + 1, my + 1, ry, qly, qry));
    }

    long long query_x(int vx, int lx, int rx, int qlx, int qrx, int qly, int qry)
    {
        if (qlx > rx || qrx < lx)
            return NEUTRAL;

        // If the row segment is fully inside our query, hand it off to the column query!
        if (qlx <= lx && rx <= qrx)
        {
            return query_y(vx, 1, 0, m - 1, qly, qry);
        }

        int mx = lx + (rx - lx) / 2;
        return combine(
            query_x(vx * 2, lx, mx, qlx, qrx, qly, qry),
            query_x(vx * 2 + 1, mx + 1, rx, qlx, qrx, qly, qry));
    }

public:
    SegTree2D(int rows, int cols)
    {
        n = rows;
        m = cols;
        // Allocate 4N rows, each containing a 4M column segment tree
        seg.assign(4 * n, vector<long long>(4 * m, NEUTRAL));
    }

    void build(const vector<vector<long long>> &grid)
    {
        build_x(1, 0, n - 1, grid);
    }

    // Update cell at (row, col) to a new exact value
    void update(int row, int col, long long val)
    {
        update_x(1, 0, n - 1, row, col, val);
    }

    // Query inclusive rectangle from top-left (r1, c1) to bottom-right (r2, c2)
    long long query(int r1, int c1, int r2, int c2)
    {
        return query_x(1, 0, n - 1, r1, r2, c1, c2);
    }
};
\end{lstlisting}
\subsubsection{Dynamic segment tree}
\begin{lstlisting}[language=C++]
/**
 *  THE ULTIMATE CP TEMPLATE BLUEPRINT: DYNAMIC SEGMENT TREE
 *
 * 1.  WHAT PROBLEM DOES THIS SOLVE?
 * - Keywords: "Massive constraints", "10^9 or 10^18 indices", "Sparse updates", "Avoid Coordinate Compression".
 * - Classic Scenarios: You need a Segment Tree over a massive range of coordinates (e.g., from -10^9 to 10^9).
 *   A standard segment tree would need an array of size 4 * 10^9, which guarantees a Memory Limit Exceeded (MLE) crash.
 * - The Magic: Instead of creating all nodes at the start, we only create nodes *as we visit them*.
 *   Initially, only the Root exists. When an update goes down the tree, it allocates the exact $O(\log(\text{Range}))$
 *   children it needs to reach the bottom. Empty space simply doesn't exist in memory!
 *
 * 2.  HOW TO USE IT (THE BLACK BOX)
 * - Initialization: Pass the absolute minimum and maximum possible coordinates to the constructor.
 *       // Create a tree covering the range [-10^9, 10^9]
 *       DynamicSegTree st(-1e9, 1e9);
 *
 * - Updates (0-indexed point update):
 *       // Add X to the value at coordinate `idx`
 *       st.update(idx, X);
 *
 * - Queries (Inclusive [L, R]):
 *       // Query the sum of elements from coordinate L to R
 *       long long total = st.query(L, R);
 *
 * - Complexity:
 *       Time: $O(\log(\text{MAX\_VAL} - \text{MIN\_VAL}))$ per update and query.
 *       Space: $O(Q \log(\text{MAX\_VAL} - \text{MIN\_VAL}))$. (Creates ~30 nodes per update for 10^9 ranges).
 *
 * 3.  HOW TO ADAPT IT (THE DIALS & SWITCHES)
 * - Changing to Range Min/Max?
 *   1. Change `Node::sum` to `val` and initialize to `1e18` (Min) or `-1e18` (Max).
 *   2. Change the `NEUTRAL` variable to `1e18` or `-1e18`.
 *   3. In `query()`, return `min/max` instead of `+`.
 * - "My code crashes randomly during updates!":
 *   Never use references (`Node& curr = seg[node]`) before calling `.push_back()`. Vectors reallocate memory when they
 *   grow, which instantly invalidates all references to its elements, causing a terrifying segmentation fault.
 *   This template strictly uses indices (`seg[node]`) to completely avoid this trap.
 */

#include <bits/stdc++.h>
using namespace std;

struct DynamicSegTree
{
private:
    struct Node
    {
        long long sum;
        int lc; // Index of Left Child in the vector
        int rc; // Index of Right Child in the vector

        Node() : sum(0), lc(-1), rc(-1) {}
    };

    vector<Node> seg;
    long long MIN_VAL, MAX_VAL;
    const long long NEUTRAL = 0; // 0 for Sum. 1e18 for Min. -1e18 for Max.

    // ? Creates children on-demand if they don't exist yet
    void extend(int node)
    {
        if (seg[node].lc == -1)
        {
            seg[node].lc = seg.size();
            seg.emplace_back(); // Safely allocates a new blank node
        }
        if (seg[node].rc == -1)
        {
            seg[node].rc = seg.size();
            seg.emplace_back(); // Safely allocates a new blank node
        }
    }

    void update(int node, long long lx, long long rx, long long idx, long long val)
    {
        if (lx == rx)
        {
            seg[node].sum += val; // Point addition
            return;
        }

        // ! We are going deeper, so we MUST ensure the children exist
        extend(node);

        // ? This safe midpoint calculation gracefully handles negative ranges
        long long mid = lx + (rx - lx) / 2;

        if (idx <= mid)
        {
            update(seg[node].lc, lx, mid, idx, val);
        }
        else
        {
            update(seg[node].rc, mid + 1, rx, idx, val);
        }

        // Pull up the answer from the children
        seg[node].sum = seg[seg[node].lc].sum + seg[seg[node].rc].sum;
    }

    long long query(int node, long long lx, long long rx, long long ql, long long qr)
    {
        // ! If this branch doesn't exist, it means nothing was ever added here! Return Neutral.
        if (node == -1 || ql > rx || qr < lx)
        {
            return NEUTRAL;
        }

        // Total overlap
        if (ql <= lx && rx <= qr)
        {
            return seg[node].sum;
        }

        long long mid = lx + (rx - lx) / 2;

        // Partial overlap
        long long left_ans = query(seg[node].lc, lx, mid, ql, qr);
        long long right_ans = query(seg[node].rc, mid + 1, rx, ql, qr);

        return left_ans + right_ans;
    }

public:
    DynamicSegTree(long long min_v, long long max_v)
    {
        MIN_VAL = min_v;
        MAX_VAL = max_v;

        // Pre-allocate some memory to minimize slow vector reallocations
        seg.reserve(1000000);

        // Create the root node at index 0
        seg.emplace_back();
    }

    // Call this to add 'val' to coordinate 'idx'
    void update(long long idx, long long val)
    {
        update(0, MIN_VAL, MAX_VAL, idx, val);
    }

    // Call this to get the sum of the range [l, r]
    long long query(long long l, long long r)
    {
        return query(0, MIN_VAL, MAX_VAL, l, r);
    }
};
\end{lstlisting}
\subsubsection{Iterative segment tree}
\begin{lstlisting}[language=C++]
/**
 *  THE ULTIMATE CP TEMPLATE BLUEPRINT: ITERATIVE SEGMENT TREE
 *
 * 1.  WHAT PROBLEM DOES THIS SOLVE?
 * - Keywords: "Point Updates", "Range Queries", "Strict Time Limits", "No Lazy Propagation".
 * - Classic Scenarios: You need a Segment Tree, but the time limit is extremely tight (e.g., 1.0s for 5*10^5 queries),
 *   or the memory limit is small. Recursive segment trees have overhead from function calls and use 4*N memory.
 *   This iterative version uses EXACTLY 2*N memory and runs in a tight `while` loop, making it blindingly fast.
 * - The Magic: The leaves of the tree are stored in indices `N` to `2N-1`. The parents are stored in `1` to `N-1`.
 *   To find the parent of node `x`, you just do `x / 2` (or `x >> 1`).
 *   To find the sibling of node `x`, you just do `x ^ 1`.
 *
 * 2.  HOW TO USE IT (THE BLACK BOX)
 * - Initialization: Pass a 0-indexed vector to build the tree.
 *       vector<long long> a = {1, 5, 2, 8, 3};
 *       IterativeSegTree st(a);
 *
 * - Updates (0-indexed point SET):
 *       // Replace the value at index `idx` with `X`
 *       st.update(idx, X);
 *
 * - Queries (0-indexed, Inclusive [L, R]):
 *       // Query the combined answer from index L to R
 *       long long ans = st.query(L, R);
 *
 * - Complexity:
 *       Time: O(N) to build. O(log N) per update and query (with a tiny constant factor).
 *       Space: O(N) memory. Uses exactly 2*N elements.
 *
 * 3.  HOW TO ADAPT IT (THE DIALS & SWITCHES)
 * - Changing from Range Sum to Min/Max/GCD?
 *   1. Change the `NEUTRAL` variable to `1e18` (Min), `-1e18` (Max), or `0` (GCD).
 *   2. Update the `combine()` function to return `min(a, b)`, `max(a, b)`, or `gcd(a, b)`.
 * - Changing `update` from "Set" to "Add"?
 *   In `update()`, change `seg[idx += n] = val;` to `seg[idx += n] += val;`.
 */

#include <bits/stdc++.h>
using namespace std;

struct IterativeSegTree
{
private:
    int n;
    vector<long long> seg;

    // ? Dial 1: The Neutral Element (Identity)
    // 0 for Sum, 1e18 for Min, -1e18 for Max, 0 for GCD
    const long long NEUTRAL = 0;

    // ? Dial 2: The Combination Logic
    long long combine(long long a, long long b)
    {
        return a + b;
    }

public:
    IterativeSegTree(const vector<long long> &v)
    {
        n = v.size();
        // Allocate exactly 2 * N memory. The 0-th index is unused dummy space.
        seg.assign(2 * n, NEUTRAL);

        // 1. Insert the leaf nodes directly into the second half of the array
        for (int i = 0; i < n; i++)
        {
            seg[n + i] = v[i];
        }

        // 2. Build the parents from bottom to top
        for (int i = n - 1; i > 0; i--)
        {
            // A node `i` has children at `2*i` and `2*i + 1`
            seg[i] = combine(seg[i << 1], seg[i << 1 | 1]);
        }
    }

    // Replace the value at 0-indexed position `idx` with `val`
    void update(int idx, long long val)
    {
        // Jump straight to the leaf node
        idx += n;
        seg[idx] = val; // ! Change to += if problem requires Point Add instead of Point Set

        // Climb to the root, updating parents
        // `idx > 1` stops us before we update the unused 0-th index
        for (idx >>= 1; idx > 0; idx >>= 1)
        {
            // `idx << 1` is left child, `idx << 1 | 1` is right child
            seg[idx] = combine(seg[idx << 1], seg[idx << 1 | 1]);
        }
    }

    // Query the inclusive range [l, r]
    long long query(int l, int r)
    {
        long long res_left = NEUTRAL;
        long long res_right = NEUTRAL;

        // Jump straight to the leaves
        // Note: We use inclusive [l, r] so our loop condition is l <= r
        for (l += n, r += n; l <= r; l >>= 1, r >>= 1)
        {

            // ? If `l` is odd, it is a right child.
            // This means its parent covers data OUTSIDE our [l, r] range!
            // We must process this node independently, then move `l` to the right.
            if (l & 1)
            {
                res_left = combine(res_left, seg[l++]);
            }

            // ? If `r` is even, it is a left child.
            // This means its parent covers data OUTSIDE our [l, r] range!
            // We must process this node independently, then move `r` to the left.
            if (!(r & 1))
            {
                res_right = combine(seg[r--], res_right);
            }
        }

        // Combine the results from the left sweep and the right sweep.
        // Doing it this way preserves the strict Left-to-Right order for non-commutative operations!
        return combine(res_left, res_right);
    }
};

void solve()
{
    int n, q;
    if (!(cin >> n >> q))
        return;

    vector<long long> a(n);
    for (int i = 0; i < n; i++)
        cin >> a[i];

    IterativeSegTree st(a);

    while (q--)
    {
        int type;
        cin >> type;
        if (type == 1)
        {
            int idx;
            long long val;
            cin >> idx >> val;
            st.update(idx, val); // 0-indexed update
        }
        else
        {
            int l, r;
            cin >> l >> r;
            cout << st.query(l, r) << "\n"; // 0-indexed query
        }
    }
}

int main()
{
    // Standard fast I/O
    ios_base::sync_with_stdio(false);
    cin.tie(NULL);
    solve();
    return 0;
}
\end{lstlisting}
\subsubsection{Merge sort tree}
\begin{lstlisting}[language=C++]
struct Node
{
    vector<int> v;
};

struct segTree
{
    vector<Node> seg;
    int sz;

    segTree(int n)
    {
        sz = 1;
        while (sz < n)
            sz *= 2;
        seg.assign(2 * sz, Node());
    }

    Node merge(Node &l, Node &r)
    {
        Node res;
        res.v.reserve(l.v.size() + r.v.size());
        std::merge(all(l.v), all(r.v), back_inserter(res.v));
        return res;
    }

    void build(vector<int> &arr, int ni, int lx, int rx)
    {
        if (rx - lx == 1)
        {
            if (lx < (int)arr.size())
            {
                seg[ni].v = {arr[lx]};
            }
            return;
        }
        int mid = (lx + rx) / 2;
        build(arr, ni * 2 + 1, lx, mid);
        build(arr, ni * 2 + 2, mid, rx);
        seg[ni] = merge(seg[ni * 2 + 1], seg[ni * 2 + 2]);
    }

    void build(vector<int> &arr)
    {
        build(arr, 0, 0, sz);
    }

    // 1. Elements Strictly Less Than K (< K)
    int query_less(int l, int r, int k, int ni, int lx, int rx)
    {
        if (lx >= l && rx <= r)
        {
            return lower_bound(all(seg[ni].v), k) - seg[ni].v.begin();
        }
        if (lx >= r || rx <= l)
            return 0;
        int mid = (lx + rx) / 2;
        return query_less(l, r, k, ni * 2 + 1, lx, mid) + query_less(l, r, k, ni * 2 + 2, mid, rx);
    }

    // 2. Elements Less Than or Equal K (<= K)
    int query_less_equal(int l, int r, int k, int ni, int lx, int rx)
    {
        if (lx >= l && rx <= r)
        {
            return upper_bound(all(seg[ni].v), k) - seg[ni].v.begin();
        }
        if (lx >= r || rx <= l)
            return 0;
        int mid = (lx + rx) / 2;
        return query_less_equal(l, r, k, ni * 2 + 1, lx, mid) + query_less_equal(l, r, k, ni * 2 + 2, mid, rx);
    }

    // 3. Elements Strictly Greater Than K (> K)
    int query_greater(int l, int r, int k, int ni, int lx, int rx)
    {
        if (lx >= l && rx <= r)
        {
            return seg[ni].v.end() - upper_bound(all(seg[ni].v), k);
        }
        if (lx >= r || rx <= l)
            return 0;
        int mid = (lx + rx) / 2;
        return query_greater(l, r, k, ni * 2 + 1, lx, mid) + query_greater(l, r, k, ni * 2 + 2, mid, rx);
    }

    // 4. Elements Greater Than or Equal K (>= K)
    int query_greater_equal(int l, int r, int k, int ni, int lx, int rx)
    {
        if (lx >= l && rx <= r)
        {
            return seg[ni].v.end() - lower_bound(all(seg[ni].v), k);
        }
        if (lx >= r || rx <= l)
            return 0;
        int mid = (lx + rx) / 2;
        return query_greater_equal(l, r, k, ni * 2 + 1, lx, mid) + query_greater_equal(l, r, k, ni * 2 + 2, mid, rx);
    }
    // 5. Elements in Range [X, Y]
    int query_in_range(int l, int r, int x, int y)
    {
        // Total elements <= Y  MINUS  Total elements < X
        return get_less_equal(l, r, y) - get_less(l, r, x);
    }

    // Helpers
    int get_greater(int l, int r, int k) { return query_greater(l, r, k, 0, 0, sz); }
    int get_greater_equal(int l, int r, int k) { return query_greater_equal(l, r, k, 0, 0, sz); }
    int get_less(int l, int r, int k) { return query_less(l, r, k, 0, 0, sz); }
    int get_less_equal(int l, int r, int k)
    {
        return query_less_equal(l, r, k, 0, 0, sz);
    }
};
\end{lstlisting}
\subsubsection{PST}
\begin{lstlisting}[language=C++]
// Distinct elements in range
// kth element
// couting or frequency of numbers in range

struct PST
{
    struct Node
    {
        int l = 0;
        int r = 0;
        int val = 0;
    };

    vector<Node> tree;
    vector<int> roots;
    int n;
    int neutral;

    PST(int n, int neutral = 0) : n(n), neutral(neutral)
    {
        tree.reserve(1500000);
        tree.push_back({0, 0, neutral});
        roots.push_back(0);
    }

    int merge(int a, int b)
    {
        return a + b;
    }

    int build(const vector<int> &a, int lx, int rx)
    {
        int u = tree.size();
        tree.push_back({0, 0, neutral});

        if (rx - lx == 1)
        {
            if (lx < (int)a.size())
                tree[u].val = a[lx];
            return u;
        }

        int mid = lx + (rx - lx) / 2;
        tree[u].l = build(a, lx, mid);
        tree[u].r = build(a, mid, rx);
        tree[u].val = merge(tree[tree[u].l].val, tree[tree[u].r].val);
        return u;
    }

    void build_initial(const vector<int> &a)
    {
        roots.push_back(build(a, 0, n));
    }

    int update(int node, int idx, int val, int lx, int rx)
    {
        int u = tree.size();
        tree.push_back(tree[node]);

        if (rx - lx == 1)
        {
            tree[u].val += val;
            return u;
        }

        int mid = lx + (rx - lx) / 2;
        if (idx < mid)
        {
            int nxt = update(tree[node].l, idx, val, lx, mid);
            tree[u].l = nxt;
        }
        else
        {
            int nxt = update(tree[node].r, idx, val, mid, rx);
            tree[u].r = nxt;
        }

        tree[u].val = merge(tree[tree[u].l].val, tree[tree[u].r].val);
        return u;
    }

    int update(int version, int idx, int val)
    {
        int new_root = update(roots[version], idx, val, 0, n);
        roots.push_back(new_root);
        return (int)roots.size() - 1;
    }

    int query(int node, int l, int r, int lx, int rx)
    {
        if (node == 0)
            return neutral;
        if (lx >= r || rx <= l)
            return neutral;
        if (lx >= l && rx <= r)
            return tree[node].val;

        int mid = lx + (rx - lx) / 2;
        return merge(query(tree[node].l, l, r, lx, mid),
                     query(tree[node].r, l, r, mid, rx));
    }

    int query(int version, int l, int r)
    {
        return query(roots[version], l, r, 0, n);
    }

    int findKth(int l, int r, int k, int lx, int rx)
    {
        if (rx - lx == 1)
            return lx;
        int cnt = tree[tree[r].l].val - tree[tree[l].l].val;
        int mid = lx + (rx - lx) / 2;
        if (cnt >= k)
            return findKth(tree[l].l, tree[r].l, k, lx, mid);
        return findKth(tree[l].r, tree[r].r, k - cnt, mid, rx);
    }

    int findKth(int l, int r, int k)
    {
        return findKth(roots[l], roots[r], k, 0, n);
    }
};

\end{lstlisting}
\subsubsection{PST Lazy}
\begin{lstlisting}[language=C++]
/**
 *  THE ULTIMATE CP TEMPLATE BLUEPRINT: PERSISTENT LAZY SEGMENT TREE
 *
 * 1.  WHAT PROBLEM DOES THIS SOLVE?
 * - Keywords: "Time travel", "Historical versions", "Range updates on previous states".
 * - Classic Scenarios: You are given an array. You need to perform Range Additions or Range Sets.
 *   Each update creates a new "version" of the array. You must be able to query the Range Sum
 *   of ANY past version of the array efficiently.
 * - The Magic: Instead of modifying nodes in place, we create a clone of every node we touch.
 *   This means `update(v, L, R)` leaves version `v` completely untouched and returns the root ID
 *   of the brand new version. Because we only change O(log N) nodes per update, memory and time
 *   remain highly efficient!
 *
 * 2.  HOW TO USE IT (THE BLACK BOX)
 * - Initialization: Pass a 0-indexed vector<long long> to build Version 0.
 *       vector<long long> a = {1, 5, 2, 8, 3};
 *       PersistentLazySegTree pst(a);
 *
 * - Tracking Versions: Every update returns a NEW root index. You should store these!
 *       vector<int> versions;
 *       versions.push_back(pst.roots[0]); // Initial version
 *
 * - Updates (All 0-indexed, Inclusive [L, R]):
 *       // Add X to range [L, R] on a specific version, and save the new version ID
 *       int new_root = pst.updateAdd(versions.back(), L, R, X);
 *       versions.push_back(new_root);
 *
 * - Queries (All 0-indexed, Inclusive [L, R]):
 *       // Query the sum of range [L, R] exactly as it was in version 'V'
 *       long long sum = pst.querySum(versions[V], L, R);
 *
 * - Complexity:
 *       Time: O(N) to build. O(log N) per update and query.
 *       Space: O(N + Q log N) memory.
 *
 * 3.  HOW TO ADAPT IT (THE DIALS & SWITCHES)
 * - Memory Limit Exceeded (MLE)? Persistent Lazy creates ~4 new nodes per update.
 *   If N = 10^5 and Q = 10^5, you will generate roughly 10^5 + 400,000 nodes.
 *   The constructor already calls `seg.reserve(N * 40)` to prevent expensive vector reallocations.
 *   Increase this multiplier if Q is exceptionally large.
 * - Changing from Range Add to Range Set?
 *   Range Add is special because we can just accumulate `lazy` tags during the query without
 *   ever modifying the tree. If you want Range Set, you must pass down a `lazySet` flag.
 *   In `query`, if you hit a node with an active `lazySet`, simply return `lazySet * (r - l + 1)`
 *   immediately, ignoring its children.
 */

#include <bits/stdc++.h>
using namespace std;

struct PersistentLazySegTree
{
private:
    struct Node
    {
        int lc, rc;     // Indices of the left and right children
        long long sum;  // The range sum
        long long lazy; // The lazy tag (Range Add)

        Node() : lc(0), rc(0), sum(0), lazy(0) {}
        Node(int lc, int rc, long long sum, long long lazy)
            : lc(lc), rc(rc), sum(sum), lazy(lazy) {}
    };

    int n;
    vector<Node> seg;

    // ? Creates a brand new node that is an exact copy of the given node ID
    int clone(int node)
    {
        seg.push_back(seg[node]);
        return seg.size() - 1;
    }

    int build(int l, int r, const vector<long long> &a)
    {
        int curr = seg.size();
        seg.emplace_back(); // Allocate new node

        if (l == r)
        {
            if (l < (int)a.size())
            {
                seg[curr].sum = a[l];
            }
            return curr;
        }

        int mid = l + (r - l) / 2;
        int left_child = build(l, mid, a);
        int right_child = build(mid + 1, r, a);

        seg[curr].lc = left_child;
        seg[curr].rc = right_child;
        seg[curr].sum = seg[left_child].sum + seg[right_child].sum;

        return curr;
    }

    int updateAdd(int node, int l, int r, int ql, int qr, long long val)
    {
        // ? Core Persistence Logic: NEVER modify the original node. Always work on a clone!
        int curr = clone(node);

        // 1. Total Overlap
        if (ql <= l && r <= qr)
        {
            seg[curr].sum += val * (r - l + 1);
            seg[curr].lazy += val;
            return curr;
        }

        int mid = l + (r - l) / 2;

        // ? Lazy Propagation in Persistence:
        // Before we branch down, we must push the current node's lazy tag to its children.
        // Because pushing modifies the children, we MUST clone the children first!
        if (seg[curr].lazy != 0)
        {
            if (l != r)
            { // Don't push down if it's a leaf
                int left_child = clone(seg[curr].lc);
                seg[left_child].sum += seg[curr].lazy * (mid - l + 1);
                seg[left_child].lazy += seg[curr].lazy;
                seg[curr].lc = left_child;

                int right_child = clone(seg[curr].rc);
                seg[right_child].sum += seg[curr].lazy * (r - mid);
                seg[right_child].lazy += seg[curr].lazy;
                seg[curr].rc = right_child;
            }
            seg[curr].lazy = 0; // Cleared from the current node
        }

        // 2. Partial Overlap (Recursion)
        if (ql <= mid)
        {
            seg[curr].lc = updateAdd(seg[curr].lc, l, mid, ql, qr, val);
        }
        if (qr > mid)
        {
            seg[curr].rc = updateAdd(seg[curr].rc, mid + 1, r, ql, qr, val);
        }

        // 3. Update Current Node Sum
        seg[curr].sum = seg[seg[curr].lc].sum + seg[seg[curr].rc].sum;
        return curr;
    }

    // ? IMPORTANT: We CANNOT push down lazy tags during a query. Pushing clones nodes,
    // which would cause a read-only query to consume massive amounts of memory!
    // Instead, we just pass the accumulated lazy value down the recursion.
    long long querySum(int node, int l, int r, int ql, int qr, long long lazy_acc)
    {
        if (!node || ql > r || qr < l)
            return 0;

        // 1. Total Overlap
        if (ql <= l && r <= qr)
        {
            // Apply all the pending lazy additions from ancestors to this block
            return seg[node].sum + lazy_acc * (r - l + 1);
        }

        int mid = l + (r - l) / 2;

        // 2. Add the CURRENT node's lazy tag to the accumulator for its children
        long long current_lazy = lazy_acc + seg[node].lazy;

        // 3. Partial Overlap
        long long res = 0;
        if (ql <= mid)
        {
            res += querySum(seg[node].lc, l, mid, ql, qr, current_lazy);
        }
        if (qr > mid)
        {
            res += querySum(seg[node].rc, mid + 1, r, ql, qr, current_lazy);
        }

        return res;
    }

public:
    // Keeps track of the root node index for every version
    vector<int> roots;

    PersistentLazySegTree(const vector<long long> &a)
    {
        n = a.size();

        // Pre-allocate memory to prevent expensive vector resizing.
        // Formula: N + (Queries * 4 * log2(N)). 40 is a safe constant for 10^5.
        seg.reserve(n * 40);

        // Create a dummy node 0 (Null Node)
        seg.emplace_back();

        // Build Version 0
        roots.push_back(build(0, n - 1, a));
    }

    // Pass the root ID of the version you want to update, get the new root ID back
    int updateAdd(int root_version, int l, int r, long long val)
    {
        return updateAdd(root_version, 0, n - 1, l, r, val);
    }

    // Pass the root ID of the version you want to query
    long long querySum(int root_version, int l, int r)
    {
        return querySum(root_version, 0, n - 1, l, r, 0LL);
    }
};
\end{lstlisting}
\subsubsection{Segment tree}
\begin{lstlisting}[language=C++]
struct Node
{
    int val;

    Node()
    {
        val = 0;
    }

    Node(int val) : val(val)
    {
    }

    void change(int x)
    {
        val = x;
    }
};

struct segmentTree
{
    int sz;
    vector<Node> seg;

    segmentTree(int n)
    {
        sz = 1;
        while (sz < n)
            sz *= 2;
        seg.assign(2 * sz, Node());
    }

    void init(vector<int> &arr, int node, int lx, int rx)
    {
        if (rx - lx == 1)
        {
            if (lx < int(arr.size()))
                seg[node] = Node(arr[lx]);
            return;
        }
        int mid = (lx + rx) / 2;
        init(arr, 2 * node + 1, lx, mid);
        init(arr, 2 * node + 2, mid, rx);
        seg[node] = merge(seg[2 * node + 1], seg[2 * node + 2]);
    }

    void init(vector<int> &arr)
    {
        init(arr, 0, 0, sz);
    }

    Node merge(Node &lf, Node &rt)
    {
        Node ans = Node();
        ans.val = lf.val + rt.val;
        return ans;
    }

    void set(int idx, int val, int node, int lx, int rx)
    {
        if (rx - lx == 1)
        {
            seg[node].change(val);
            return;
        }
        int mid = (lx + rx) / 2;
        if (idx < mid)
            set(idx, val, 2 * node + 1, lx, mid);
        else
            set(idx, val, 2 * node + 2, mid, rx);
        seg[node] = merge(seg[2 * node + 1], seg[2 * node + 2]);
    }

    void set(int idx, int val)
    {
        set(idx, val, 0, 0, sz);
    }

    Node get(int l, int r, int node, int lx, int rx)
    {
        if (lx >= l && rx <= r)
            return seg[node];
        if (rx <= l || lx >= r)
            return Node();
        int mid = (lx + rx) / 2;
        Node lf = get(l, r, 2 * node + 1, lx, mid);
        Node ri = get(l, r, 2 * node + 2, mid, rx);
        return merge(lf, ri);
    }

    int get(int l, int r)
    {
        return get(l, r, 0, 0, sz).val;
    }
};

\end{lstlisting}
\subsubsection{Segment tree beats}
\begin{lstlisting}[language=C++]
/**
 *  THE ULTIMATE CP TEMPLATE BLUEPRINT: SEGMENT TREE BEATS
 *
 * 1.  WHAT PROBLEM DOES THIS SOLVE?
 * - Keywords: "Range chmin" (A[i] = min(A[i], X)), "Range chmax", "Range Sum with Min/Max updates".
 * - Classic Scenarios: Standard Segment Trees completely fail if you try to perform a range minimum
 *   update (e.g., "for all elements from L to R, replace the element with min(element, X)")
 *   while simultaneously asking for the Range Sum.
 * - The Magic: Segment Tree Beats tracks the *First Maximum* and *Second Maximum* in every node.
 *   If your update X is between the First and Second Max, it only affects the First Max!
 *   It uses this trick to update the sum mathematically in amortized O(log^2 N) time.
 *
 * 2.  HOW TO USE IT (THE BLACK BOX)
 * - Initialization: Pass a 0-indexed vector<long long> to the constructor.
 *       vector<long long> a = {1, 5, 2, 8, 3};
 *       SegmentTreeBeats seg(a);
 *
 * - Queries (All 0-indexed, Inclusive [L, R]):
 *       seg.updateMin(L, R, X) : Performs A[i] = min(A[i], X) for i in [L, R].
 *       seg.updateMax(L, R, X) : Performs A[i] = max(A[i], X) for i in [L, R].
 *       seg.updateAdd(L, R, X) : Performs A[i] = A[i] + X for i in [L, R].
 *       seg.updateSet(L, R, X) : Performs A[i] = X for i in [L, R].
 *       seg.querySum(L, R)     : Returns the sum of A[i] in [L, R].
 *       seg.queryMin(L, R)     : Returns the Minimum in [L, R].
 *       seg.queryMax(L, R)     : Returns the Maximum in [L, R].
 *
 * - Complexity:
 *       Time: O(N) to build. Amortized O(log^2 N) per query.
 *       Space: O(N) memory.
 *
 * 3.  HOW TO ADAPT IT (THE DIALS & SWITCHES)
 * - 1-based indexing? Instead of changing the recursive logic, simply add a dummy element `0`
 *   to the front of your vector before initializing: `v.insert(v.begin(), 0);`. Then query `[1, N]`.
 * - Numbers larger than 10^14? The `INF` value is set to 2 * 10^15. If your problem features numbers
 *   larger than that, change the `const long long INF` variable at the top of the struct.
 * - Need more speed? If the problem ONLY asks for `updateMin` and `querySum`, you can safely delete
 *   all variables related to `mn`, `secMn`, and `mnCnt` to cut memory and execution time in half.
 */

#include <bits/stdc++.h>
using namespace std;

// ! Assumes you are using long long to avoid any overflow bugs.
struct SegmentTreeBeats
{
private:
// ! FIX: Added strict parentheses to prevent operator precedence bugs in math
#define L (2 * node + 1)
#define R (2 * node + 2)
#define mid ((l + r) >> 1)

    // ! FIX: Safely typed INF to avoid macro expansion issues
    const long long INF = 2e15 + 5;

    struct Node
    {
        long long sum, mx, secMx, mxCnt, mn, mnCnt, secMn, lazyAdd, lazySet;

        Node(long long summ, long long mxx, long long secMxx, long long mxCntt,
             long long mnn, long long secMnn, long long mnCntt, long long add = 0, long long set = -1)
            : sum(summ), mx(mxx), secMx(secMxx), mxCnt(mxCntt),
              mn(mnn), secMn(secMnn), mnCnt(mnCntt),
              lazyAdd(add), lazySet(set)
        {
        }
    };

    int sz;
    vector<Node> seg;

    Node merge(Node &a, Node &b)
    {
        Node ret(0, -INF, -INF, 0, INF, INF, 0, 0, -1);
        ret.sum = a.sum + b.sum;

        // ? Merge Maximums
        ret.mx = max(a.mx, b.mx);
        ret.secMx = max(a.secMx, b.secMx);
        if (ret.mx == a.mx)
            ret.mxCnt += a.mxCnt;
        else
            ret.secMx = max(ret.secMx, a.mx);
        if (ret.mx == b.mx)
            ret.mxCnt += b.mxCnt;
        else
            ret.secMx = max(ret.secMx, b.mx);

        // ? Merge Minimums
        ret.mn = min(a.mn, b.mn);
        ret.secMn = min(a.secMn, b.secMn);
        if (ret.mn == a.mn)
            ret.mnCnt += a.mnCnt;
        else
            ret.secMn = min(ret.secMn, a.mn);
        if (ret.mn == b.mn)
            ret.mnCnt += b.mnCnt;
        else
            ret.secMn = min(ret.secMn, b.mn);

        return ret;
    }

    void pushSet(int l, int r, int node, long long val)
    {
        seg[node] = Node(val * (r - l + 1), val, -INF, r - l + 1, val, INF, r - l + 1, 0, val);
    }

    void pushMin(int l, int r, int node, long long mn)
    {
        if (seg[node].mn >= mn)
        {
            pushSet(l, r, node, mn);
            return;
        }
        if (seg[node].mx > mn)
        {
            if (seg[node].secMn == seg[node].mx)
                seg[node].secMn = mn;
            long long val = seg[node].mx - mn;
            seg[node].sum -= val * seg[node].mxCnt;
            seg[node].mx = mn;
        }
    }

    void pushMax(int l, int r, int node, long long mx)
    {
        if (seg[node].mx <= mx)
        {
            pushSet(l, r, node, mx);
            return;
        }
        if (seg[node].mn < mx)
        {
            if (seg[node].secMx == seg[node].mn)
                seg[node].secMx = mx;
            long long val = mx - seg[node].mn;
            seg[node].sum += val * seg[node].mnCnt;
            seg[node].mn = mx;
        }
    }

    void pushAdd(int l, int r, int node, long long val)
    {
        if (seg[node].mx == seg[node].mn)
        {
            pushSet(l, r, node, seg[node].mn + val);
            return;
        }
        seg[node].mx += val;
        if (seg[node].secMx != -INF)
            seg[node].secMx += val;
        seg[node].mn += val;
        if (seg[node].secMn != INF)
            seg[node].secMn += val;
        seg[node].sum += (r - l + 1) * val;
        seg[node].lazyAdd += val;
    }

    void propegate(int l, int r, int node)
    {
        if (l == r)
            return;
        if (seg[node].lazySet != -1)
        {
            pushSet(l, mid, L, seg[node].lazySet);
            pushSet(mid + 1, r, R, seg[node].lazySet);
            seg[node].lazySet = -1;
        }
        else
        {
            pushAdd(l, mid, L, seg[node].lazyAdd);
            pushAdd(mid + 1, r, R, seg[node].lazyAdd);
            seg[node].lazyAdd = 0;
            pushMin(l, mid, L, seg[node].mx);
            pushMin(mid + 1, r, R, seg[node].mx);
            pushMax(l, mid, L, seg[node].mn);
            pushMax(mid + 1, r, R, seg[node].mn);
        }
    }

    void build(int l, int r, int node, vector<long long> &v)
    {
        if (l == r)
        {
            if (l < (int)v.size())
            {
                seg[node] = Node(v[l], v[l], -INF, 1, v[l], INF, 1);
            }
            return;
        }
        build(l, mid, L, v);
        build(mid + 1, r, R, v);
        seg[node] = merge(seg[L], seg[R]);
    }

    void updateAdd(int l, int r, int node, int lx, int rx, long long val)
    {
        if (l > rx || r < lx)
            return;
        if (l >= lx && r <= rx)
        {
            pushAdd(l, r, node, val);
            return;
        }
        propegate(l, r, node);
        updateAdd(l, mid, L, lx, rx, val);
        updateAdd(mid + 1, r, R, lx, rx, val);
        seg[node] = merge(seg[L], seg[R]);
    }

    void updateSet(int l, int r, int node, int lx, int rx, long long val)
    {
        if (l > rx || r < lx)
            return;
        if (l >= lx && r <= rx)
        {
            pushSet(l, r, node, val);
            return;
        }
        propegate(l, r, node);
        updateSet(l, mid, L, lx, rx, val);
        updateSet(mid + 1, r, R, lx, rx, val);
        seg[node] = merge(seg[L], seg[R]);
    }

    void updateMin(int l, int r, int node, int lx, int rx, long long mn)
    {
        // ? Core BEATS logic: If the segment max is already <= mn, do nothing.
        if (r < lx || l > rx || seg[node].mx <= mn)
            return;
        // ? If the update affects the maximums, but NOT the second maximums, we can apply it safely!
        if (l >= lx && r <= rx && mn > seg[node].secMx)
        {
            pushMin(l, r, node, mn);
            return;
        }
        propegate(l, r, node);
        updateMin(l, mid, L, lx, rx, mn);
        updateMin(mid + 1, r, R, lx, rx, mn);
        seg[node] = merge(seg[L], seg[R]);
    }

    void updateMax(int l, int r, int node, int lx, int rx, long long mx)
    {
        if (l > rx || r < lx || mx <= seg[node].mn)
            return;
        if (l >= lx && r <= rx && mx < seg[node].secMn)
        {
            pushMax(l, r, node, mx);
            return;
        }
        propegate(l, r, node);
        updateMax(l, mid, L, lx, rx, mx);
        updateMax(mid + 1, r, R, lx, rx, mx);
        seg[node] = merge(seg[L], seg[R]);
    }

    long long querySum(int l, int r, int node, int lx, int rx)
    {
        if (l > rx || r < lx)
            return 0;
        if (l >= lx && r <= rx)
            return seg[node].sum;
        propegate(l, r, node);
        return querySum(l, mid, L, lx, rx) + querySum(mid + 1, r, R, lx, rx);
    }

    long long queryMin(int l, int r, int node, int lx, int rx)
    {
        if (l > rx || r < lx)
            return INF;
        if (l >= lx && r <= rx)
            return seg[node].mn;
        propegate(l, r, node);
        return min(queryMin(l, mid, L, lx, rx), queryMin(mid + 1, r, R, lx, rx));
    }

    long long queryMax(int l, int r, int node, int lx, int rx)
    {
        if (l > rx || r < lx)
            return -INF;
        if (l >= lx && r <= rx)
            return seg[node].mx;
        propegate(l, r, node);
        return max(queryMax(l, mid, L, lx, rx), queryMax(mid + 1, r, R, lx, rx));
    }

#undef L
#undef R
#undef mid

public:
    SegmentTreeBeats(vector<long long> &v)
    {
        int n = v.size() + 5;
        sz = 1;
        while (sz <= n)
            sz <<= 1;
        seg = vector<Node>(sz << 1, Node(0, -INF, -INF, 0, INF, INF, 0, 0, -1));
        build(0, sz - 1, 0, v);
    }

    void updateAdd(int l, int r, long long val) { updateAdd(0, sz - 1, 0, l, r, val); }
    void updateSet(int l, int r, long long val) { updateSet(0, sz - 1, 0, l, r, val); }
    void updateMin(int l, int r, long long mn) { updateMin(0, sz - 1, 0, l, r, mn); }
    void updateMax(int l, int r, long long mx) { updateMax(0, sz - 1, 0, l, r, mx); }
    long long queryMax(int l, int r) { return queryMax(0, sz - 1, 0, l, r); }
    long long queryMin(int l, int r) { return queryMin(0, sz - 1, 0, l, r); }
    long long querySum(int l, int r) { return querySum(0, sz - 1, 0, l, r); }
};
\end{lstlisting}
\subsubsection{Segment tree lazy}
\begin{lstlisting}[language=C++]
struct Node
{
    int val, lazy;
    bool is_lazy;

    Node()
    {
        val = 0, lazy = 0, is_lazy = false;
    }

    Node(int val) : val(val)
    {
        lazy = 0, is_lazy = false;
    }

    void change(int x, int lx, int rx)
    {
        val += x * (rx - lx);
        lazy += x;
        is_lazy = true;
    }

    void reset()
    {
        lazy = 0, is_lazy = false;
    }
};

struct segmentTree
{
    int sz;
    vector<Node> seg;

    segmentTree(int n)
    {
        sz = 1;
        while (sz < n)
            sz *= 2;
        seg.assign(2 * sz, Node());
    }

    void init(vector<int> &arr, int node, int lx, int rx)
    {
        if (rx - lx == 1)
        {
            if (lx < int(arr.size()))
                seg[node] = Node(arr[lx]);
            return;
        }
        int mid = (lx + rx) / 2;
        init(arr, 2 * node + 1, lx, mid);
        init(arr, 2 * node + 2, mid, rx);
        seg[node] = merge(seg[2 * node + 1], seg[2 * node + 2]);
    }

    void init(vector<int> &arr)
    {
        init(arr, 0, 0, sz);
    }

    Node merge(Node &lf, Node &rt)
    {
        Node ans = Node();
        ans.val = lf.val + rt.val;
        return ans;
    }

    void push(int node, int lx, int rx)
    {
        if (rx - lx == 1 || !seg[node].is_lazy)
            return;
        int mid = (lx + rx) / 2;
        seg[2 * node + 1].change(seg[node].lazy, lx, mid);
        seg[2 * node + 2].change(seg[node].lazy, mid, rx);
        seg[node].reset();
    }

    void update(int l, int r, int x, int node, int lx, int rx)
    {
        if (lx >= r || rx <= l)
            return;
        push(node, lx, rx);
        if (lx >= l && rx <= r)
        {
            seg[node].change(x, lx, rx);
            return;
        }
        int mid = (lx + rx) / 2;
        update(l, r, x, 2 * node + 1, lx, mid);
        update(l, r, x, 2 * node + 2, mid, rx);

        seg[node] = merge(seg[2 * node + 1], seg[2 * node + 2]);
    }

    void update(int l, int r, int x)
    {
        return update(l, r, x, 0, 0, sz);
    }

    Node get(int l, int r, int node, int lx, int rx)
    {
        if (rx <= l || lx >= r)
            return Node();
        push(node, lx, rx);
        if (lx >= l && rx <= r)
            return seg[node];
        int mid = (lx + rx) / 2;
        Node lf = get(l, r, 2 * node + 1, lx, mid);
        Node ri = get(l, r, 2 * node + 2, mid, rx);
        return merge(lf, ri);
    }

    int get(int l, int r)
    {
        return get(l, r, 0, 0, sz).val;
    }
};

\end{lstlisting}
\subsubsection{Segment tree merging}
\begin{lstlisting}[language=C++]
/**
 *  THE ULTIMATE CP TEMPLATE BLUEPRINT: SEGMENT TREE MERGING
 *
 * 1.  WHAT PROBLEM DOES THIS SOLVE?
 * - Keywords: "Subtree queries", "Colors in subtree", "Inversions in subtree", "Tree DP".
 * - Classic Scenarios: You are traversing a tree (or graph) bottom-up. Every node needs
 *   to track a frequency array or sums of its subtree.
 * - The Magic: Merging two Dynamic Segment Trees takes time proportional to their
 *   *intersecting* nodes. If you merge all children into their parents bottom-up,
 *   the total time and memory for the entire process is mathematically bounded to
 *   O(N log(Range)). It feels like magic because $10^5$ independent trees merge into
 *   one in a fraction of a second.
 *
 * 2.  HOW TO USE IT (THE BLACK BOX)
 * - Initialization:
 *       MergeableSegTree st(-1e9, 1e9);
 *       vector<int> roots(N + 1, -1); // Tracks the root ID for each node's personal tree
 *
 * - Updates:
 *       // Add 'val' at coordinate 'idx' for node U's personal tree
 *       roots[U] = st.update(roots[U], idx, val);
 *
 * - Merging:
 *       // Merge the tree of child V into the tree of parent U
 *       roots[U] = st.merge(roots[U], roots[V]);
 *
 * - Queries:
 *       // Query the sum in range [L, R] inside node U's tree
 *       long long ans = st.query(roots[U], L, R);
 *
 * 3.  HOW TO ADAPT IT (THE DIALS & SWITCHES)
 * -  DESTRUCTIVE MERGE WARNING:
 *   This template performs a "Destructive Merge". When you do `st.merge(roots[U], roots[V])`,
 *   the tree of V is absorbed into U. You CANNOT safely query `roots[V]` anymore!
 *   This is perfectly fine for 99% of Tree DP problems (you process bottom-up and don't look back).
 *   If a problem forces you to query V again later, you must change the `merge` function to
 *   create a clone instead of modifying `u` in place (Warning: clones cause massive memory usage).
 * - Changing to Min/Max:
 *   Change `seg[u].sum += seg[v].sum` to `= max(seg[u].max, seg[v].max)` in the base case of `merge()`.
 */

#include <bits/stdc++.h>
using namespace std;

struct MergeableSegTree
{
private:
    struct Node
    {
        long long sum = 0;
        int lc = -1; // Left child index
        int rc = -1; // Right child index
    };

    vector<Node> seg;
    long long MIN_VAL, MAX_VAL;

    int allocate()
    {
        seg.emplace_back();
        return seg.size() - 1;
    }

    int update(int node, long long lx, long long rx, long long idx, long long val)
    {
        // If the node doesn't exist yet, create it!
        if (node == -1)
            node = allocate();

        if (lx == rx)
        {
            seg[node].sum += val;
            return node;
        }

        long long mid = lx + (rx - lx) / 2;
        if (idx <= mid)
        {
            seg[node].lc = update(seg[node].lc, lx, mid, idx, val);
        }
        else
        {
            seg[node].rc = update(seg[node].rc, mid + 1, rx, idx, val);
        }

        // Pull up answers from children
        seg[node].sum = 0;
        if (seg[node].lc != -1)
            seg[node].sum += seg[seg[node].lc].sum;
        if (seg[node].rc != -1)
            seg[node].sum += seg[seg[node].rc].sum;

        return node;
    }

    // ? THE MAGIC: The Destructive Merge Function
    int merge(int u, int v, long long lx, long long rx)
    {
        // 1. If one of the branches is empty, just point to the other! (O(1) shortcut)
        if (u == -1)
            return v;
        if (v == -1)
            return u;

        // 2. Base case: Both trees have data at this leaf. Combine them!
        if (lx == rx)
        {
            seg[u].sum += seg[v].sum;
            return u;
        }

        // 3. Recursive case: Traverse overlapping paths and merge children
        long long mid = lx + (rx - lx) / 2;
        seg[u].lc = merge(seg[u].lc, seg[v].lc, lx, mid);
        seg[u].rc = merge(seg[u].rc, seg[v].rc, mid + 1, rx);

        // Pull up the new combined sum
        seg[u].sum = 0;
        if (seg[u].lc != -1)
            seg[u].sum += seg[seg[u].lc].sum;
        if (seg[u].rc != -1)
            seg[u].sum += seg[seg[u].rc].sum;

        return u; // Return the absorbed parent node
    }

    long long query(int node, long long lx, long long rx, long long ql, long long qr)
    {
        if (node == -1 || ql > rx || qr < lx)
            return 0;

        if (ql <= lx && rx <= qr)
            return seg[node].sum;

        long long mid = lx + (rx - lx) / 2;
        return query(seg[node].lc, lx, mid, ql, qr) +
               query(seg[node].rc, mid + 1, rx, ql, qr);
    }

public:
    MergeableSegTree(long long min_v, long long max_v)
    {
        MIN_VAL = min_v;
        MAX_VAL = max_v;
        // Pre-allocate to prevent slow resizing.
        // Formula: N * log(Range). Usually 40 * N is safe for 10^9 ranges.
        seg.reserve(4000000);
    }

    // Pass the root of the tree you want to update
    int update(int root, long long idx, long long val)
    {
        return update(root, MIN_VAL, MAX_VAL, idx, val);
    }

    // Pass the roots of the two trees you want to merge
    int merge(int root_u, int root_v)
    {
        return merge(root_u, root_v, MIN_VAL, MAX_VAL);
    }

    // Pass the root of the tree you want to query
    long long query(int root, long long l, long long r)
    {
        return query(root, MIN_VAL, MAX_VAL, l, r);
    }
};

void solve()
{
    int n;
    // ... Read tree structure ...
    // vector<vector<int>> adj(n + 1);

    // Create the global Mergeable Segment Tree
    MergeableSegTree st(1, 1e9);

    // Array to track the root of the segment tree belonging to each node
    vector<int> roots(n + 1, -1);

    // Example: Add value '1' to coordinate 'color[u]' for node U's tree
    // roots[u] = st.update(roots[u], color[u], 1);

    // Example Tree DP logic:
    // for (int v : adj[u]) {
    //     dfs(v, u);
    //     roots[u] = st.merge(roots[u], roots[v]);
    // }

    // Now roots[u] contains the fully merged frequency array of its entire subtree!
}

int main()
{
    ios_base::sync_with_stdio(false);
    cin.tie(NULL);
    // solve();
    return 0;
}
\end{lstlisting}
\subsection{Sparse Table}
\subsubsection{2D sparse table}
\begin{lstlisting}[language=C++]
/**
 *  THE ULTIMATE CP TEMPLATE BLUEPRINT: 2D SPARSE TABLE
 *
 * 1.  WHAT PROBLEM DOES THIS SOLVE?
 * - Keywords: "Static 2D Grid", "Maximum in subgrid", "Minimum in subgrid", "Huge number of queries".
 * - Classic Scenarios: You are given an N x M grid of numbers. The grid NEVER changes.
 *   You need to answer Q queries (up to 10^6) asking for the max/min in a rectangle bounded
 *   by top-left (r1, c1) and bottom-right (r2, c2).
 * - The Magic: A 2D Segment tree answers queries in O(log N * log M), which can Time Limit Exceed
 *   (TLE) if Q is massive. A 2D Sparse Table builds the answers for all powers of 2 in advance,
 *   allowing it to answer ANY rectangle query in absolute O(1) time by overlapping 4 smaller rectangles!
 *
 * 2.  HOW TO USE IT (THE BLACK BOX)
 * - Initialization: Pass a 0-indexed 2D vector to the constructor.
 *       vector<vector<int>> grid = {{1, 5}, {2, 8}};
 *       SparseTable2D st(grid);
 *
 * - Queries (All 0-indexed, Inclusive [r1, r2] and [c1, c2]):
 *       // Query the max in the subgrid from top-left (0, 0) to bottom-right (1, 1)
 *       int max_val = st.query(0, 0, 1, 1);
 *
 * - Complexity:
 *       Time: O(N * M * log N * log M) to build. Pure O(1) per query.
 *       Space: O(N * M * log N * log M) memory.
 *
 * 3.  HOW TO ADAPT IT (THE DIALS & SWITCHES)
 * - Changing from Max to Min or GCD?
 *   1. Change the `combine()` function to return `min(a, b)` or `std::gcd(a, b)`.
 * - Memory Limit Exceeded (MLE) Warning:
 *   If N = 1000 and M = 1000, log(N) = 10. The table requires 1000 * 1000 * 10 * 10 = 10^8 ints.
 *   This is ~400 MB of memory. This passes on most platforms (512MB limits), but if the limit
 *   is strictly 256MB, you will MLE. Always check the memory limit before choosing this!
 * - Idempotent Operations ONLY:
 *   Sparse tables ONLY work for operations where X + X = X (like Min, Max, GCD, Bitwise OR/AND).
 *   Do NOT use this for Sum queries! (Use a 2D Prefix Sum array for O(1) sums instead).
 */

#include <bits/stdc++.h>
using namespace std;

// struct SparseTable2D {
private:
int n, m;
int log_n, log_m;
// table[ir][ic][r][c]
// Represents the rectangle of size 2^ir x 2^ic starting at top-left (r, c)
vector<vector<vector<vector<int>>>> table;

// Precomputed logs for pure O(1) queries
vector<int> log_table;

// ? Dial 1: The Combination Logic (Change to min, gcd, |, &, etc.)
inline int combine(int a, int b)
{
    return max(a, b);
}

// public:
SparseTable2D(const vector<vector<int>> &grid)
{
    if (grid.empty() || grid[0].empty())
        return;
    n = grid.size();
    m = grid[0].size();

    // Precompute logarithms
    int max_dim = max(n, m);
    log_table.assign(max_dim + 1, 0);
    for (int i = 2; i <= max_dim; i++)
    {
        log_table[i] = log_table[i / 2] + 1;
    }

    log_n = log_table[n] + 1;
    log_m = log_table[m] + 1;

    // Allocate memory carefully: ir -> ic -> row -> col
    // This specific ordering prevents massive cache-miss penalties during the build phase
    table.assign(log_n, vector<vector<vector<int>>>(log_m,
                                                    vector<vector<int>>(n, vector<int>(m))));

    // 1. Base Case: 1x1 rectangles (2^0 x 2^0)
    for (int i = 0; i < n; i++)
    {
        for (int j = 0; j < m; j++)
        {
            table[0][0][i][j] = grid[i][j];
        }
    }

    // 2. Build purely horizontally (1D Sparse Table for each row)
    for (int ir = 0; ir < log_n; ir++)
    {
        for (int ic = 0; ic < log_m; ic++)
        {
            // Skip the base case we already filled
            if (ir == 0 && ic == 0)
                continue;

            for (int i = 0; i + (1 << ir) <= n; i++)
            {
                for (int j = 0; j + (1 << ic) <= m; j++)
                {

                    if (ir == 0)
                    {
                        // Merge two adjacent horizontal blocks
                        table[ir][ic][i][j] = combine(
                            table[ir][ic - 1][i][j],
                            table[ir][ic - 1][i][j + (1 << (ic - 1))]);
                    }
                    else
                    {
                        // Merge two adjacent vertical blocks
                        table[ir][ic][i][j] = combine(
                            table[ir - 1][ic][i][j],
                            table[ir - 1][ic][i + (1 << (ir - 1))][j]);
                    }
                }
            }
        }
    }
}

//    // Query inclusive rectangle from top-left (r1, c1) to bottom-right (r2, c2)
int query(int r1, int c1, int r2, int c2)
{
    // Find the largest power of 2 that fits inside the row and column lengths
    int k_r = log_table[r2 - r1 + 1];
    int k_c = log_table[c2 - c1 + 1];

    // ? THE MAGIC: We perfectly cover the requested rectangle using 4 smaller overlapping
    // rectangles of size (2^k_r) x (2^k_c).
    // Because the operation is Idempotent (max/min), overlapping data doesn't change the answer!

    int top_left = table[k_r][k_c][r1][c1];
    int top_right = table[k_r][k_c][r1][c2 - (1 << k_c) + 1];
    int bottom_left = table[k_r][k_c][r2 - (1 << k_r) + 1][c1];
    int bottom_right = table[k_r][k_c][r2 - (1 << k_r) + 1][c2 - (1 << k_c) + 1];

    return combine(combine(top_left, top_right), combine(bottom_left, bottom_right));
}
}
;

// void solve() {
int n, m, q;
// Example Input: N rows, M cols, Q queries
if (!(cin >> n >> m >> q))
    return;

vector<vector<int>> grid(n, vector<int>(m));
for (int i = 0; i < n; i++)
{
    for (int j = 0; j < m; j++)
    {
        cin >> grid[i][j];
    }
}

SparseTable2D st(grid);

while (q--)
{
    int r1, c1, r2, c2;
    // Assume 0-indexed coordinates are provided
    cin >> r1 >> c1 >> r2 >> c2;
    cout << st.query(r1, c1, r2, c2) << "\n";
}
}

int main()
{
    // Fast I/O is mandatory for heavy 2D queries
    ios_base::sync_with_stdio(false);
    cin.tie(NULL);
    solve();
    return 0;
}
\end{lstlisting}
\subsubsection{sparse table}
\begin{lstlisting}[language=C++]
struct SparseTable {
    int n, LG;
    vector<vector<int> > T;
    SparseTable(const vector<int> &arr) {
        n = arr.size();
        LG = __lg(n) + 2;
        T.assign(LG, vector<int>(n));
        for (int i = 0; i < n; i++) {
            T[0][i] = arr[i];
        }
        for (int pw = 1; (1 << pw) <= n; pw++) {
            for (int i = 0; i + (1 << pw) <= n; i++) {
                T[pw][i] = merge(T[pw - 1][i], T[pw - 1][i + (1 << (pw - 1))]);
            }
        }
    }

    int merge(int a, int b) {
        return min(a, b);
    }

    int query(int l, int r) {
        int len = r - l + 1;
        int pw = __lg(len);
        return merge(T[pw][l], T[pw][r - (1 << pw) + 1]);
    }

    int query1(int l, int r) { // o(log)
        int res = 1e9; // change
        int sz = r-l+1;
        for (int i = 0;l<=r;i++) {
            if (sz>>i&1) {
                res = merge(res,T[i][l]);
                l+=(1<<i);
            }
        }
        return res;
    }

};

\end{lstlisting}
\subsection{Treap}
\subsubsection{implicit treap}
\begin{lstlisting}[language=C++]
/**
 *  THE ULTIMATE CP TEMPLATE BLUEPRINT: IMPLICIT TREAP (CARTESIAN TREE)
 *
 * 1.  WHAT PROBLEM DOES THIS SOLVE?
 * - Keywords: "Dynamic Array", "Range Reverse", "Insert/Delete at index", "Cut and Paste".
 * - Classic Scenarios: You have an array and need to perform operations that a Segment
 *   Tree CANNOT do, such as:
 *      - Insert/Delete elements at specific indices (shifting the whole array).
 *      - Reversing a whole subsegment [L, R] dynamically.
 *      - Moving a block [L, R] to the front or end.
 * - The Magic: An Implicit Treap acts like a balanced BST, but instead of sorting by
 *   a "Key", it sorts by the "Implicit Index" (the size of the left subtree). This allows
 *   us to split and merge arbitrary subarrays in O(log N) time.
 *
 * 2.  HOW TO USE IT (THE BLACK BOX)
 * - Initialization: Root = nullptr. Use `insert` to populate.
 * - Range Operations:
 *       // Reverse range [L, R] (0-indexed)
 *       reverse_range(root, L, R);
 *       // Query sum of range [L, R]
 *       long long ans = query_sum(root, L, R);
 * - Complexity:
 *       Time: O(log N) per operation.
 *       Space: O(N) memory.
 *
 * 3.  HOW TO ADAPT IT (THE DIALS & SWITCHES)
 * - Need Range Add? Add a `lazy` field to the Node and push it in the `push()` function.
 * - Need Max/Min instead of Sum? Update `update()` and `get_sum()` to maintain
 *   `max`/`min` fields instead of `sum`.
 */

#include <bits/stdc++.h>
using namespace std;

struct Node
{
    int val, priority, size;
    long long sum;
    bool rev; // Flag for range reversal (Lazy Propagation)
    Node *left, *right;

    Node(int v) : val(v), priority(rand()), size(1), sum(v), rev(false), left(nullptr), right(nullptr) {}
};

int get_size(Node *t) { return t ? t->size : 0; }
long long get_sum(Node *t) { return t ? t->sum : 0; }

// Update subtree size and sum based on children
void update(Node *t)
{
    if (!t)
        return;
    t->size = 1 + get_size(t->left) + get_size(t->right);
    t->sum = t->val + get_sum(t->left) + get_sum(t->right);
}

// Push lazy propagation (reverse flag) down to children
void push(Node *t)
{
    if (t && t->rev)
    {
        t->rev = false;
        swap(t->left, t->right);
        if (t->left)
            t->left->rev ^= true;
        if (t->right)
            t->right->rev ^= true;
    }
}

// Split treap t into l (first k nodes) and r (remaining nodes)
void split(Node *t, int k, Node *&l, Node *&r)
{
    if (!t)
    {
        l = r = nullptr;
        return;
    }
    push(t);
    int left_size = get_size(t->left);
    if (left_size < k)
    {
        split(t->right, k - left_size - 1, t->right, r);
        l = t;
    }
    else
    {
        split(t->left, k, l, t->left);
        r = t;
    }
    update(t);
}

// Merge treap l and treap r
Node *merge(Node *l, Node *r)
{
    push(l);
    push(r);
    if (!l || !r)
        return l ? l : r;
    if (l->priority > r->priority)
    {
        l->right = merge(l->right, r);
        update(l);
        return l;
    }
    else
    {
        r->left = merge(l, r->left);
        update(r);
        return r;
    }
}

// Reverse subarray [l, r] (0-indexed)
void reverse_range(Node *&root, int l, int r)
{
    Node *t1, *t2, *t3;
    split(root, r + 1, t2, t3);
    split(t2, l, t1, t2);
    if (t2)
        t2->rev ^= true;
    root = merge(t1, merge(t2, t3));
}

// Query sum in subarray [l, r]
long long query_sum(Node *&root, int l, int r)
{
    Node *t1, *t2, *t3;
    split(root, r + 1, t2, t3);
    split(t2, l, t1, t2);
    long long res = get_sum(t2);
    root = merge(t1, merge(t2, t3));
    return res;
}

// Insert value at index pos
void insert(Node *&root, int pos, int val)
{
    Node *t1, *t2;
    split(root, pos, t1, t2);
    root = merge(merge(t1, new Node(val)), t2);
}

int main()
{
    srand(time(0));
    Node *root = nullptr;
    // Example: Insert values
    insert(root, 0, 10);
    insert(root, 1, 20);
    insert(root, 2, 30);

    // Example: Sum query [0, 1]
    // cout << query_sum(root, 0, 1) << endl;

    return 0;
}
\end{lstlisting}
\subsection{Trie}
\subsubsection{BinaryTrie Using Vector}
\begin{lstlisting}[language=C++]
class BinaryTrie {
    struct Node {
        int ch[2];
        int frq[2];
        Node() {
            ch[0] = ch[1] = -1;
            frq[0] = frq[1] = 0;
        }
    };
    vector<Node> tree;
public:
    BinaryTrie() {
        tree.emplace_back();
        insert(0);   
    }

    void insert(int n) {
        int cur = 0;
        for (int i = 30; i >= 0; i--) {
            int idx = (n >> i) & 1;

            if (tree[cur].ch[idx] == -1) {
                tree[cur].ch[idx] = tree.size();
                tree.emplace_back();
            }
 
            tree[cur].frq[idx]++;
            cur = tree[cur].ch[idx];
        }
    }
 
    void erase(int n) {
        int cur = 0;
        for (int i = 30; i >= 0; i--) {
            int idx = (n >> i) & 1;
            tree[cur].frq[idx]--;
            cur = tree[cur].ch[idx];
        }
    }

    bool search(int n) {
        int cur = 0;
        for (int i = 30; i >= 0; i--) {
            int idx = (n >> i) & 1;
            if (tree[cur].ch[idx] == -1 || tree[cur].frq[idx] == 0) {
                return false;
            }
            cur = tree[cur].ch[idx];
        }
        return true;
    }

    int maxXor(int x) {
        int cur = 0;
        int res = 0;
        for (int i = 30; i >= 0; i--) {
            int idx = (x >> i) & 1;
            int target = 1 - idx;  
 
            if (tree[cur].ch[target] != -1 && tree[cur].frq[target] > 0) {
                res |= (1 << i);
                cur = tree[cur].ch[target];
            } else {
                cur = tree[cur].ch[idx];
            }
        }
        return res;
    }
 
    void clear() {
        tree.clear();
        tree.emplace_back();
        insert(0);
    }
};



\end{lstlisting}
\subsubsection{Binary trie}
\begin{lstlisting}[language=C++]
struct BitTrie
{
private:
    BitTrie *child[2];
    int cnt{};

public:
    BitTrie()
    {
        child[0] = child[1] = nullptr;
    }

    ~BitTrie()
    {
        delete child[0];
        delete child[1];
    }

    void insert(int num)
    {
        BitTrie *curr = this;
        for (int i = 31; i >= 0; i--)
        {
            int bit = num >> i & 1;
            if (!curr->child[bit])
                curr->child[bit] = new BitTrie();
            curr = curr->child[bit];
            curr->cnt++;
        }
    }

    void erase(int num)
    {
        BitTrie *curr = this;
        for (int i = 31; i >= 0; i--)
        {
            int bit = num >> i & 1;
            curr = curr->child[bit];
            curr->cnt--;
        }
    }

    int getMaxXor(int num)
    {
        BitTrie *curr = this;
        int maxXor = 0;
        for (int i = 31; i >= 0; i--)
        {
            int need = 1 - (num >> i & 1);
            if (curr->child[need] && curr->child[need]->cnt > 0)
            {
                maxXor |= 1ll << i;
                curr = curr->child[need];
            }
            else
                curr = curr->child[1 - need];
        }
        return maxXor;
    }

    // count Xors with X less than K
    int countXorLessThan(int x, int k)
    {
        int ret{};
        BitTrie *curr = this;
        for (int i = 31; i >= 0; i--)
        {
            if (!curr)
                break;
            int bitX = x >> i & 1;
            int bitK = k >> i & 1;
            if (bitK == 0)
                curr = curr->child[bitX];
            else
            {
                if (curr->child[bitX])
                    ret += curr->child[bitX]->cnt;
                curr = curr->child[1 - bitX];
            }
        }
        return ret;
    }
};

// get the Max Xor of two numbers in the array
int getMaxXorOfTwoNumbersInArray(vector<int> &v)
{
    if (v.empty())
        return 0;
    int n = int(v.size());
    BitTrie trie;
    trie.insert(v[0]);
    int max{};
    for (int i = 1; i < n; i++)
    {
        max = std::max(max, trie.getMaxXor(v[i]));
        trie.insert(v[i]);
    }
    return max;
}

struct Query
{
    int x, a, id;

    bool operator<(const Query &other) const
    {
        return a < other.a;
    }
};

// max xor queries for a number with array with threshold A [The element i pick from array must be <= A]
void getMaxXorOfNumWithArrayLessThanThresholdQueries(vector<int> &v)
{
    int n = int(v.size());
    sort(v.begin(), v.end());
    int q;
    cin >> q;
    vector<Query> queries(q);
    for (int i = 0; i < q; i++)
    {
        cin >> queries[i].x >> queries[i].a;
        queries[i].id = i;
    }
    sort(queries.begin(), queries.end());
    BitTrie trie;
    int idx{};
    vector<int> ans(q);
    for (int i = 0; i < q; i++)
    {
        auto [x, a, id] = queries[i];
        while (idx < n && v[idx] <= a)
            trie.insert(v[idx++]);
        if (idx == 0)
            ans[id] = -1;
        else
            ans[id] = trie.getMaxXor(x);
    }
    for (int i = 0; i < q; i++)
        cout << ans[i] << "\n";
}

// Max xor subarray
int getMaxXorSubarray(vector<int> &a)
{
    int n = int(a.size());
    BitTrie trie;
    trie.insert(0);
    int pre{}, ret{};
    for (int i = 0; i < n; i++)
    {
        pre ^= a[i];
        ret = max(ret, trie.getMaxXor(pre));
        trie.insert(pre);
    }
    return ret;
}

int numberOfSubarraysWithXorLessThanK(vector<int> &a, int k)
{
    int n = int(a.size());
    BitTrie trie;
    trie.insert(0);
    int pre{}, ret{};
    for (int i = 0; i < n; i++)
    {
        pre ^= a[i];
        ret += trie.countXorLessThan(pre, k);
        trie.insert(pre);
    }
    return ret;
}

\end{lstlisting}
\subsubsection{Persistent Trie}
\begin{lstlisting}[language=C++]
/**
 *  THE ULTIMATE CP TEMPLATE BLUEPRINT: PERSISTENT BINARY TRIE (XOR MAXIMIZER)
 *
 * 1.  WHAT PROBLEM DOES THIS SOLVE?
 * - Keywords: "Maximum XOR in a subarray", "Range XOR", "Subsegment XOR with X".
 * - Classic Scenarios: You are given an array of N integers. You need to answer Q queries.
 *   Each query gives you a range [L, R] and a value X. You need to find an element Y
 *   in the subarray from L to R such that (Y ^ X) is maximized.
 * - The Magic: A standard Trie can find the max XOR in a whole array. By making it Persistent,
 *   we save a new "version" of the Trie after inserting every single element. To query the
 *   range [L, R], we take the Trie at version R and mathematically subtract the paths from
 *   version L-1. It does this in blindingly fast O(30) time!
 *
 * 2.  HOW TO USE IT (THE BLACK BOX)
 * - Initialization: Pass a 0-indexed vector to build the historical versions.
 *       vector<long long> a = {3, 10, 5, 25, 2};
 *       PersistentTrie ptrie(a);
 *
 * - Queries (All 0-indexed, Inclusive [L, R]):
 *       // Find the maximum XOR of 7 with any element in the subarray from index 1 to 3
 *       long long ans = ptrie.queryMaxXor(1, 3, 7);
 *
 * - Complexity:
 *       Time: O(N * BITS) to build. O(BITS) per query.
 *       Space: O(N * BITS) memory.
 *
 * 3.  HOW TO ADAPT IT (THE DIALS & SWITCHES)
 * - Dealing with huge numbers (up to 10^18)?
 *   Change the `BITS` constant from 30 to 60. (30 is enough for numbers up to 10^9).
 * - Minimum XOR instead of Maximum XOR?
 *   I have included a `queryMinXor` function right below the Max one! It simply tries to
 *   take the path of the SAME bit instead of the FLIPPED bit.
 * - Need to dynamically add elements to the back of the array?
 *   Just call `ptrie.insertBack(X);`. It will add the element and create a new version
 *   that you can instantly query.
 */

#include <bits/stdc++.h>
using namespace std;

struct PersistentTrie
{
private:
    // ? Dial 1: Set to 30 for <= 10^9. Set to 60 for <= 10^18.
    const int BITS = 30;

    struct Node
    {
        int child[2];
        int cnt; // How many elements exist in this subtree

        Node()
        {
            child[0] = child[1] = 0;
            cnt = 0;
        }
    };

    vector<Node> trie;

    // ? Creates a brand new node that is an exact copy of the given node ID
    int clone(int node)
    {
        trie.push_back(trie[node]);
        return trie.size() - 1;
    }

    int insert(int prev_root, long long val)
    {
        int curr = clone(prev_root);
        int new_root = curr;
        int prev = prev_root;

        trie[curr].cnt++; // Increment count for the root of this version

        for (int i = BITS; i >= 0; i--)
        {
            int b = (val >> i) & 1;

            // ? We are taking path 'b'. To safely modify it without destroying
            // the historical version, we MUST clone the 'b' child of the previous version!
            trie[curr].child[b] = clone(trie[prev].child[b]);

            // The OTHER child pointer (b ^ 1) was naturally copied over during the
            // clone() of the parent, so it still safely points to the historical branch!

            curr = trie[curr].child[b];
            prev = trie[prev].child[b];

            trie[curr].cnt++; // Increment the count of elements that pass through this path
        }

        return new_root;
    }

public:
    // roots[i] represents the Trie containing all elements from index 0 to i-1.
    // roots[0] is the empty Trie.
    vector<int> roots;

    PersistentTrie(const vector<long long> &a = {})
    {
        // Pre-allocate memory to prevent slow vector reallocations.
        // N elements * roughly (BITS + 1) nodes per element
        trie.reserve(a.size() * (BITS + 2) + 100);

        // Node 0 is the Dummy/Empty node
        trie.emplace_back();

        // Version 0 is the empty tree
        roots.push_back(0);

        for (long long x : a)
        {
            insertBack(x);
        }
    }

    // Add an element to the end of the array dynamically
    void insertBack(long long val)
    {
        int new_root = insert(roots.back(), val);
        roots.push_back(new_root);
    }

    // Find the MAXIMUM value of (val ^ Y) where Y is an element in a[l...r]
    long long queryMaxXor(int l, int r, long long val)
    {
        // Because roots[i] covers [0...i-1], the range [L, R] is extracted by
        // taking Version (R + 1) and subtracting Version L.
        int curr_R = roots[r + 1];
        int curr_L = roots[l];
        long long ans = 0;

        for (int i = BITS; i >= 0; i--)
        {
            int b = (val >> i) & 1;
            int flipped = b ^ 1; // To MAXIMIZE XOR, we want the opposite bit!

            // How many numbers in our specific range actually take the 'flipped' path?
            int count_flipped = trie[trie[curr_R].child[flipped]].cnt -
                                trie[trie[curr_L].child[flipped]].cnt;

            if (count_flipped > 0)
            {
                // The path exists! We take it.
                ans |= (1LL << i);
                curr_R = trie[curr_R].child[flipped];
                curr_L = trie[curr_L].child[flipped];
            }
            else
            {
                // The opposite bit doesn't exist in this range. We are forced to take the same bit.
                curr_R = trie[curr_R].child[b];
                curr_L = trie[curr_L].child[b];
            }
        }
        return ans;
    }

    // Find the MINIMUM value of (val ^ Y) where Y is an element in a[l...r]
    long long queryMinXor(int l, int r, long long val)
    {
        int curr_R = roots[r + 1];
        int curr_L = roots[l];
        long long ans = 0;

        for (int i = BITS; i >= 0; i--)
        {
            int b = (val >> i) & 1;

            // To MINIMIZE XOR, we want the EXACT SAME bit!
            int count_same = trie[trie[curr_R].child[b]].cnt -
                             trie[trie[curr_L].child[b]].cnt;

            if (count_same > 0)
            {
                // The path exists! We take it. (Ans bit remains 0 because b ^ b = 0)
                curr_R = trie[curr_R].child[b];
                curr_L = trie[curr_L].child[b];
            }
            else
            {
                // Forced to take the flipped bit, which increases our XOR penalty.
                int flipped = b ^ 1;
                ans |= (1LL << i);
                curr_R = trie[curr_R].child[flipped];
                curr_L = trie[curr_L].child[flipped];
            }
        }
        return ans;
    }
};

void solve()
{
    int n, q;
    if (!(cin >> n >> q))
        return;

    vector<long long> a(n);
    for (int i = 0; i < n; i++)
        cin >> a[i];

    PersistentTrie ptrie(a);

    while (q--)
    {
        int l, r;
        long long x;
        // 0-indexed L and R, and the XOR target X
        cin >> l >> r >> x;

        cout << ptrie.queryMaxXor(l, r, x) << "\n";
    }
}

int main()
{
    ios_base::sync_with_stdio(false);
    cin.tie(NULL);
    // solve();
    return 0;
}
\end{lstlisting}
\subsubsection{Trie (String)}
\begin{lstlisting}[language=C++]
class Trie {
    struct Node {
        map<string, int> ch;
        Node() {
        }
    };
    vector<Node>root;
public:
    Trie() {
        root.emplace_back();
    }
    void insert(vector<string> &words) {
        int cur = 0;
        for (auto s : words) {
            if (!root[cur].ch[s]) {
                root[cur].ch[s] = root.size();
                root.emplace_back();
            }
            cur = root[cur].ch[s];
        }
    }
};
\end{lstlisting}
\subsubsection{Trie Using Vector}
\begin{lstlisting}[language=C++]
class Trie {
    struct Node {
        int ch[26];
        int prefix, end;
        Node() {
            prefix = end = 0;
            memset(ch, -1, sizeof(ch));
        }
    };
    vector<Node> tree;
public:
    Trie() {
        tree.emplace_back();
    }

    void insert(string &s) {
        int cur = 0;  
        for (auto i : s) {
            int idx = i - 'a';
            
            if (tree[cur].ch[idx] == -1) {
                tree[cur].ch[idx] = tree.size();
                tree.emplace_back();
            }
            cur = tree[cur].ch[idx];
            tree[cur].prefix++;
        }
        tree[cur].end++;
    }

    bool search(string &s) {
        int cur = 0;
        for (auto i : s) {
            int idx = i - 'a';
            if (tree[cur].ch[idx] == -1) {
                return false;
            }
            cur = tree[cur].ch[idx];
        }
        return tree[cur].end > 0;
    }

    int countWords(string &s) {
        int cur = 0;
        for (auto i : s) {
            int idx = i - 'a';
            if (tree[cur].ch[idx] == -1) {
                return 0;
            }
            cur = tree[cur].ch[idx];
        }
        return tree[cur].end;
    }

    int countPrefix(string &s) {
        int cur = 0;
        for (auto i : s) {
            int idx = i - 'a';
            if (tree[cur].ch[idx] == -1) {
                return 0;
            }
            cur = tree[cur].ch[idx];
        }
        return tree[cur].prefix;
    }
    
    // DFS in Trie
    void dfs(int u) {
        for (int v = 0; v < 26; v++) {
            if (root[u].ch[v] != -1) {
                int x = root[u].ch[v];
                if (root[u].prefix != root[x].prefix) {
                    ans+=root[x].prefix;
                }
                dfs(x);
            }
        }
    }
 
    void clear() {
        tree.clear();
        tree.emplace_back();
    }
};


\end{lstlisting}
\subsubsection{Trie}
\begin{lstlisting}[language=C++]
vector<int> ans;

class Trie
{
private:
    Trie *child[26];
    int word_cnt{}, prefix_cnt{};

public:
    Trie()
    {
        memset(child, 0, sizeof child);
    }

    void insert(const string &str, int idx = 0)
    {
        prefix_cnt++;
        if (idx == int(str.size()))
        {
            word_cnt++;
            return;
        }
        int curr = str[idx] - 'a';
        if (!child[curr])
            child[curr] = new Trie();
        child[curr]->insert(str, idx + 1);
    }

    void erase(const string &str, int idx = 0)
    {
        if (idx == 0 && !wordSearch(str))
            return;
        prefix_cnt--;
        if (idx == int(str.size()))
        {
            word_cnt--;
            return;
        }
        int curr = str[idx] - 'a';
        child[curr]->erase(str, idx + 1);
    }

    void print(int idx = 0)
    {
        for (int i = 0; i < 26; i++)
        {
            if (!child[i] || !child[i]->prefix_cnt)
                continue;
            for (int k = 0; k < idx; k++)
                cout << " ";
            cout << char(i + 'a') << "\n";
            child[i]->print(idx + 1);
        }
    }

    // count number of chars needed to be typed by the user for each word [using auto-completion]
    // trie.insert("") at the beginning because the user needs to write the first letter
    void count(int cnt = 0)
    {
        if (word_cnt)
            ans.push_back(cnt);
        int num{};
        for (int i = 0; i < 26; i++)
        {
            if (child[i])
                num++;
        }
        for (int i = 0; i < 26; i++)
        {
            if (!child[i] || !child[i]->prefix_cnt)
                continue;
            if (num == 1 && word_cnt == 0)
                child[i]->count(cnt);
            else
                child[i]->count(cnt + 1);
        }
    }

    bool wordSearch(const string &str, int idx = 0)
    {
        if (idx == int(str.size()))
            return word_cnt;
        int curr = str[idx] - 'a';
        if (!child[curr] || child[curr]->prefix_cnt == 0)
            return false;
        return child[curr]->wordSearch(str, idx + 1);
    }

    bool prefixSearch(const string &str, int idx = 0)
    {
        if (idx == int(str.size()))
            return prefix_cnt;
        int curr = str[idx] - 'a';
        if (!child[curr] || child[curr]->prefix_cnt == 0)
            return false;
        return child[curr]->prefixSearch(str, idx + 1);
    }

    int countWordsEqualTo(const string &str, int idx = 0)
    {
        if (idx == int(str.size()))
            return word_cnt;
        int curr = str[idx] - 'a';
        if (!child[curr] || child[curr]->prefix_cnt == 0)
            return 0;
        return child[curr]->countWordsEqualTo(str, idx + 1);
    }

    int countWordsStartingWith(const string &str, int idx = 0)
    {
        if (idx == int(str.size()))
            return prefix_cnt;
        int curr = str[idx] - 'a';
        if (!child[curr] || child[curr]->prefix_cnt == 0)
            return 0;
        return child[curr]->countWordsStartingWith(str, idx + 1);
    }

    void printWordsStartingWith(const string &str, int idx = 0, string res = "")
    {
        if (idx < int(str.size()))
        {
            int curr = str[idx] - 'a';
            if (!child[curr] || child[curr]->prefix_cnt == 0)
                return;
            child[curr]->printWordsStartingWith(str, idx + 1, res + str[idx]);
        }
        else
        {
            if (word_cnt && idx != int(str.size()))
                // ans.push_back(res); //// uncomment this <-----
                for (int i = 0; i < 26; i++)
                {
                    if (!child[i] || child[i]->prefix_cnt == 0)
                        continue;
                    child[i]->printWordsStartingWith(str, idx + 1, res + char('a' + i));
                }
        }
    }

    // longest word with all prefixes problem
    bool checkCompleteness(const string &str, int idx = 0)
    {
        if (idx == int(str.size()))
            return word_cnt > 0;
        int curr = str[idx] - 'a';
        if (!child[curr] || child[curr]->word_cnt == 0)
            return false;
        return child[curr]->checkCompleteness(str, idx + 1);
    }

    string longestCompleteString(vector<string> &v)
    {
        int n = int(v.size());
        string longest{};
        for (int i = 0; i < n; i++)
        {
            if (checkCompleteness(v[i]))
            {
                if (v[i].length() > longest.length() || (v[i].length() == longest.length() && v[i] < longest))
                    longest = v[i];
            }
        }
        return longest.empty() ? "No Complete String" : longest;
    }

    int countNumberOfNodes()
    {
        int ret{};
        for (int i = 0; i < 26; i++)
        {
            if (child[i])
                ret += 1 + child[i]->countNumberOfNodes();
        }
        return ret;
    }
};

// count number of distinct substrings
int countNumberOfDistinctSubStrings(string &str)
{
    int n = int(str.size());
    Trie trie;
    for (int i = 0; i < n; i++)
        trie.insert(str.substr(i));
    return trie.countNumberOfNodes();
}

\end{lstlisting}
\section{Graphs}
\subsection{2-Sat}
\subsubsection{2sat}
\begin{lstlisting}[language=C++]
/**
 *  THE ULTIMATE CP TEMPLATE BLUEPRINT: 2-SAT SOLVER (VIA KOSARAJU'S ALGORITHM & SCC)
 *
 * 1.  WHAT PROBLEM DOES THIS SOLVE?
 * - Keywords: "2-SAT", "Boolean Satisfiability", "Either X or Y must happen", "Implication Graph".
 * - Classic Scenarios: You have N boolean variables (True/False) and M constraints. 
 *   Each constraint involves exactly two variables, usually in the form "At least one of X or Y must be true" 
 *   or "If X is chosen, Y cannot be chosen".
 * - The Magic: "The Implication Graph". We translate every logic clause (X OR Y) into two directed edges: 
 *   (!X -> Y) and (!Y -> X). We treat each variable as TWO nodes: one for True (2*i) and one for False (2*i+1). 
 *   Then we find the Strongly Connected Components (SCCs). 
 *   - If a variable X and its negation !X end up in the SAME component, a valid assignment is IMPOSSIBLE.
 *   - If possible, we greedily assign True to whichever version of the variable appears LATER in the 
 *     topological sort. Because Kosaraju's algorithm builds SCCs in reverse topological order, 
 *     `SCC_id[X] > SCC_id[!X]` mathematically guarantees a valid boolean assignment!
 *
 * 2.  HOW TO USE IT (THE BLACK BOX)
 * - Base State: Define `n` (variables) and `m` (clauses).
 * - Add Constraints: For every clause (X OR Y), call `add_clause(X, val_X, Y, val_Y)`. 
 *   (e.g., `add_clause(1, true, 2, false)` means "Variable 1 must be True OR Variable 2 must be False").
 * - Execution: 
 *       1. Call `SCC()` to build the components.
 *       2. Loop over 1 to N: check if `SCC_id[2*i] == SCC_id[2*i+1]`. If yes -> IMPOSSIBLE.
 *       3. Otherwise, `assignment[i] = (SCC_id[2*i] > SCC_id[2*i+1])`.
 *
 * - Complexity:
 *       Time: O(N + M)  Lightning fast linear time DFS traversals.
 *       Space: O(N + M) to store the implication graph and reverse graph.
 *
 * 3.  HOW TO ADAPT IT (THE DIALS & SWITCHES)
 * -  CRITICAL ARRAY SIZE WARNING: Because every variable splits into 2 nodes (positive and negative), 
 *   your graph MUST have space for `2 * N + 2` nodes! Make sure `cond_graph`, `graph`, and `vis` arrays 
 *   are initialized to `2 * n + 2`, NOT just `n + 1`.
 * - How to FORCE a variable to be True? 
 *   Add a clause stating (X OR X): `add_clause(x, true, x, true)`.
 * - How to handle XOR (Exactly one is True)?
 *   (X ^ Y) means they must be different. This is equivalent to TWO OR-clauses: 
 *   (X OR Y) AND (!X OR !Y). Call `add_clause` twice for these two conditions.
 * - How to handle XNOR (Both True or Both False)?
 *   (X == Y) means they must be the same. Equivalent to: (!X OR Y) AND (X OR !Y).
 */

const int N = 1e5 + 9;

int SCC_id[N];
GRAPH graph, graph_reverse, cond_graph;
int n, m;
vector<int> order, vis, vis_reverse, order_reverse;

void dfs(int node)
{
    vis[node] = true;
    for (auto &x : graph[node])
    {
        if (!vis[x])
            dfs(x);
    }
    order.emplace_back(node);
}

void dfs_reverse(int node)
{
    vis_reverse[node] = true;
    for (auto &x : graph_reverse[node])
    {
        if (!vis_reverse[x])
            dfs_reverse(x);
    }
    order_reverse.emplace_back(node);
}

void SCC()
{
    graph.assign(n + 1, {}), graph_reverse.assign(n + 1, {}), vis.assign(n + 1, {}), vis_reverse.assign(n + 1, {}), cond_graph.assign(n + 1, {});
    for (int i = 0; i < m; i++)
    {
        int x, y;
        cin >> x >> y;
        graph[x].emplace_back(y), graph_reverse[y].emplace_back(x);
    }
    order.clear();
    for (int i = 1; i <= n; i++)
    {
        if (!vis[i])
            dfs(i);
    }
    reverse(order.begin(), order.end());
    vector<vector<int>> scc;
    scc.push_back({});
    for (auto &x : order)
    {
        if (!vis_reverse[x])
        {
            order_reverse.clear();
            dfs_reverse(x);
            scc.emplace_back(order_reverse);
            for (auto &val : order_reverse)
                SCC_id[val] = int(scc.size()) - 1;
        }
    }
    for (int u = 1; u <= n; u++)
    {
        for (int v : graph[u])
        {
            if (SCC_id[u] != SCC_id[v])
                cond_graph[SCC_id[u]].emplace_back(SCC_id[v]);
        }
    }
}

// Add this helper to add edges easily
void add_clause(int i, bool val_i, int j, bool val_j)
{
    // i, j are variable indices (1 to n)
    // val_i, val_j are the boolean requirements (true for X, false for !X)

    // Logic: (X or Y) is equivalent to (!X -> Y) AND (!Y -> X)

    // Map to 2*i and 2*i+1
    int u = 2 * i + (val_i ? 0 : 1);
    int v = 2 * j + (val_j ? 0 : 1);

    // !u is u^1, !v is v^1
    graph[u ^ 1].emplace_back(v);
    graph[v ^ 1].emplace_back(u);

    // Keep reversed graph for Kosaraju
    graph_reverse[v].emplace_back(u ^ 1);
    graph_reverse[u].emplace_back(v ^ 1);
}

void solve()
{
    // 1. Read input
    cin >> n >> m;
    // Remember to adjust N to 2*N+1 if N was 1e5

    // 2. Clear graph from your SCC template
    // ... (Use your SCC setup code here)

    for (int k = 0; k < m; k++)
    {
        int i, j;
        bool val_i, val_j;
        // Example: Reading "1 True 2 False" means (X1 or !X2)
        cin >> i >> val_i >> j >> val_j;
        add_clause(i, val_i, j, val_j);
    }

    // 3. Run Kosaraju
    SCC();

    // 4. Check for Impossibility
    vector<int> assignment(n + 1);
    for (int i = 1; i <= n; i++)
    {
        if (SCC_id[2 * i] == SCC_id[2 * i + 1])
        {
            cout << "IMPOSSIBLE\n";
            return;
        }
        // 5. The Assignment Logic (The "Higher ID is True" strategy)
        assignment[i] = (SCC_id[2 * i] > SCC_id[2 * i + 1]);
    }

    // 6. Print result
    cout << "POSSIBLE\n";
    for (int i = 1; i <= n; i++)
    {
        cout << assignment[i] << " ";
    }
}

\end{lstlisting}
\subsection{DAG}
\subsubsection{Dag reachability(bitset trick)}
\begin{lstlisting}[language=C++]
#include <bits/stdc++.h>
using namespace std;

/**
 *  THE ULTIMATE CP TEMPLATE BLUEPRINT: BITSET REACHABILITY
 *
 * 1.  WHAT PROBLEM DOES THIS SOLVE?
 * - Keywords: "Transitive Closure", "DAG Reachability", "Bitset Optimization".
 * - Classic Scenarios: You have a Directed Acyclic Graph (DAG) and need to answer Q
 *   queries of the form "Can I reach node Y from node X?".
 * - The Magic: Since a DAG has a topological order, we process nodes in REVERSE
 *   topological order. For each node U, its reachability set is the Bitwise OR
 *   of the reachability sets of all its neighbors. Because bitsets are 64x faster than
 *   arrays, this brings an O(N^3) problem down to an O(N^3 / 64) complexity.
 *
 * 2.  HOW TO USE IT (THE BLACK BOX)
 * - Initialization:
 *       compute_reachability(N); // Run this once after building adjacency list
 *
 * - Query:
 *       if (reach[start].test(target)) { ... }
 *
 * - Complexity:
 *       Time: O(N * (N + M) / 64).
 *       Space: O(N^2 / 8) bytes. (For N=20,000, this is ~50MB).
 *
 * 3.  HOW TO ADAPT IT (THE DIALS & SWITCHES)
 * - Memory Limit Warning:
 *   The current `MAXN = 20005` uses ~50MB. If N > 20,000, you will exceed standard 256MB limits.
 *   If N > 20,000, you MUST use the "Block Reachability" approach (process bits in
 *   chunks of 1024 or 2048) rather than holding the entire N x N matrix in memory.
 */

// ? Dial 1: Adjust MAXN based on memory limits (N^2 / 8 bytes)
const int MAXN = 20005;

// reach[u][v] = 1 if there is a path from u to v
bitset<MAXN> reach[MAXN];
vector<int> adj[MAXN];
vector<int> topo_order;
int in_degree[MAXN];

// Standard Kahn's Algorithm for Topological Sort
void compute_topological_sort(int n)
{
    queue<int> q;
    for (int i = 1; i <= n; i++)
    {
        if (in_degree[i] == 0)
            q.push(i);
    }

    while (!q.empty())
    {
        int u = q.front();
        q.pop();
        topo_order.push_back(u);

        for (int v : adj[u])
        {
            in_degree[v]--;
            if (in_degree[v] == 0)
                q.push(v);
        }
    }
}

void compute_reachability(int n)
{
    compute_topological_sort(n);

    // ? THE MAGIC: Process in REVERSE topological order
    // This ensures that when we process node u, all its neighbors v
    // are already fully computed.
    for (int i = n - 1; i >= 0; i--)
    {
        int u = topo_order[i];

        reach[u].set(u); // A node can always reach itself

        for (int v : adj[u])
        {
            // Bitwise OR is the secret sauce:
            // This ORs all bits of reach[v] into reach[u] at once
            reach[u] |= reach[v];
        }
    }
}

int main()
{
    ios_base::sync_with_stdio(false);
    cin.tie(NULL);

    int n, m;
    cin >> n >> m;

    for (int i = 0; i < m; i++)
    {
        int u, v;
        cin >> u >> v;
        adj[u].push_back(v);
        in_degree[v]++;
    }

    compute_reachability(n);

    // Example Query: Can I reach node 'target' from node 'start'?
    int start, target;
    cin >> start >> target;

    if (reach[start].test(target))
    {
        cout << "Reachable!\n";
    }
    else
    {
        cout << "Not reachable.\n";
    }

    return 0;
}
\end{lstlisting}
\subsection{Eulerian path}
\subsubsection{Eulerian path}
\begin{lstlisting}[language=C++]
/**
 *  THE ULTIMATE CP TEMPLATE BLUEPRINT: EULERIAN PATH / CIRCUIT (HIERHOLZER'S ALGORITHM)
 *
 * 1.  WHAT PROBLEM DOES THIS SOLVE?
 * - Keywords: "Visit every edge exactly once", "Eulerian Tour", "Draw without lifting pen", "De Bruijn Sequence".
 * - Classic Scenarios: You are given a graph (directed or undirected) and you need to find a 
 *   path or cycle that traverses EVERY edge exactly once. Standard DFS/BFS won't work because 
 *   they visit *nodes* once, but here we care about *edges*.
 * - The Magic: "Hierholzer's Algorithm". We do a greedy DFS, destroying edges as we walk over them 
 *   using `.pop_back()` (to ensure we never process an edge twice, keeping time complexity strictly O(E)). 
 *   When we hit a "dead end" (a node with no outgoing edges left), we add it to our answer array. 
 *   Because of the mathematical properties of Eulerian graphs, this backtracking naturally constructs 
 *   the exact path in REVERSE!
 *
 * 2.  HOW TO USE IT (THE BLACK BOX)
 * - Initialization: Clear the `vis` array and the `ans` vector. 
 * - Build Graph: For every edge, assign it a unique `id` from 0 to M-1.
 *       graph[u].push_back({v, id});
 *       // If undirected, add the reverse edge with the SAME id!
 *       // graph[v].push_back({u, id}); 
 *
 * - Execution: Call `dfs(start_node)`. 
 *       dfs(start_node);
 *       reverse(ans.begin(), ans.end()); // ? CRITICAL: The path is generated backwards!
 *
 * - Complexity:
 *       Time: O(V + E)  Each edge is popped and visited exactly once.
 *       Space: O(V + E) for the adjacency list and recursion stack.
 *
 * 3.  HOW TO ADAPT IT (THE DIALS & SWITCHES)
 * -  CRITICAL PREREQUISITE (EXISTENCE CHECK): This DFS assumes the Eulerian path ALREADY EXISTS. 
 *   If you run it on a random graph, it will produce garbage. You MUST check the degrees first!
 *   - Undirected Circuit: All nodes must have an EVEN degree.
 *   - Undirected Path: Exactly 0 or 2 nodes must have an ODD degree. (Start DFS from an odd node!).
 *   - Directed Circuit: Every node's In-Degree == Out-Degree.
 *   - Directed Path: At most one node has (Out - In == 1) [Start Node], and one has (In - Out == 1) [End Node].
 * - Lexicographically Smallest Path?
 *   If the problem asks for the lexicographically smallest path, sort the adjacency list of every 
 *   node in DESCENDING order before starting the DFS. Since we use `.pop_back()`, the smallest 
 *   destinations will be popped first!
 * - Disconnected components: If the graph has isolated edges, a single DFS won't catch them. 
 *   Always check if `ans.size() == M + 1` at the end to ensure all edges were visited.
 */
const int N = 1e5, M = 1e5;

struct Edge
{
    int to, id;
};

vector<Edge> graph[N];
bool vis[M];
vector<int> ans;

void dfs(int u)
{
    while ((int)graph[u].size() > 0)
    {
        auto [v, id] = graph[u].back();
        graph[u].pop_back();
        if (!vis[id])
        {
            vis[id] = true;
            dfs(v);
        }
    }
    ans.emplace_back(u); // Eulerian path reversed
}

\end{lstlisting}
\subsection{Flows}
\subsubsection{Gomory-Hu Tree}
\begin{lstlisting}[language=C++]
#include <bits/stdc++.h>
using namespace std;

/**
 *  THE ULTIMATE CP TEMPLATE BLUEPRINT: GOMORY-HU TREE
 *
 * 1.  WHAT PROBLEM DOES THIS SOLVE?
 * - Keywords: "All-pairs Min-Cut", "Minimum cut tree", "Undirected Graph Cut".
 * - Classic Scenarios: You have an undirected graph and need to know the minimum
 *   cut between *every* possible pair of vertices.
 * - The Magic: Normally, finding the min-cut between all pairs requires O(N^2)
 *   Max Flow calls. The Gomory-Hu Tree (specifically Gusfield's algorithm implementation)
 *   builds a tree representing all min-cuts using only exactly N-1 Max Flow calls!
 *   The min-cut between any two nodes (u, v) is simply the minimum weight edge on
 *   the path between them in this resulting tree.
 *
 * 2.  HOW TO USE IT (THE BLACK BOX)
 * - Initialization:
 *       GomoryHu gh(N); // N is the number of nodes (0-indexed)
 * - Add Edge:
 *       gh.add_edge(u, v, cap); // Undirected edge
 * - Build:
 *       gh.build(); // Runs the N-1 Max Flow calls to construct the tree
 * - Query:
 *       long long ans = gh.min_cut(u, v);
 *
 * - Complexity:
 *       Time: O(N * MaxFlow(V, E)) to build. O(N) per query (or O(1) if you add LCA).
 *       Space: O(N + E)
 *
 * 3.  HOW TO ADAPT IT (THE DIALS & SWITCHES)
 * - Dinic Dependency: This template assumes you have your `Dinic` struct available
 *   in the same file or included. Ensure `Dinic` has `add_edge(u, v, cap)` and `max_flow(s, t)`.
 * - 1-based indexing: If your nodes are 1 to N, instantiate `GomoryHu gh(N + 1)` and
 *   just ignore index 0.
 */

struct GomoryHu
{
    int n;
    struct Edge
    {
        int u, v;
        long long cap;
    };
    vector<Edge> edges;

    // Parent array for the tree (Gomory-Hu Tree structure)
    vector<int> p;
    // Weight array (min-cut value between i and p[i])
    vector<long long> weight;
    // Depth array for simple O(N) queries without complex LCA
    vector<int> depth;

    GomoryHu(int n) : n(n), p(n), weight(n), depth(n, 0) {}

    void add_edge(int u, int v, long long cap)
    {
        edges.push_back({u, v, cap});
    }

    void build()
    {
        // 1. Initialize parents: everyone initially points to node 0
        fill(p.begin(), p.end(), 0);
        fill(weight.begin(), weight.end(), 0);

        // 2. Iterate to build the tree (Gusfield's Algorithm)
        for (int i = 1; i < n; i++)
        {
            // We rebuild the graph for Dinic each time.
            // (Assuming `Dinic` is your standard Dinic template)
            Dinic dinic(n);
            for (auto &e : edges)
            {
                dinic.add_edge(e.u, e.v, e.cap);
                dinic.add_edge(e.v, e.u, e.cap); // Gomory-Hu requires Undirected edges!
            }

            // Find min s-t cut between current node `i` and its parent `p[i]`
            long long flow = dinic.max_flow(i, p[i]);
            weight[i] = flow;

            // Find which nodes are on the source side of the cut
            // (reachable from s in the residual graph)
            vector<bool> reachable(n, false);
            dfs_residual(i, reachable, dinic);

            // Update parents for next iterations
            // If any node j > i points to p[i] but is actually on i's side of the cut,
            // we redirect its parent to i.
            for (int j = i + 1; j < n; j++)
            {
                if (reachable[j] && p[j] == p[i])
                {
                    p[j] = i;
                }
            }
        }

        // 3. Precalculate depths to allow simple climbing during queries
        for (int i = 1; i < n; i++)
        {
            depth[i] = depth[p[i]] + 1;
        }
    }

    void dfs_residual(int u, vector<bool> &reachable, Dinic &dinic)
    {
        reachable[u] = true;
        for (auto &e : dinic.adj[u])
        {
            // If capacity remains (cap - flow > 0), it is not saturated,
            // so we can walk across it in the residual graph.
            if (e.cap - e.flow > 0 && !reachable[e.to])
            {
                dfs_residual(e.to, reachable, dinic);
            }
        }
    }

    // Query min cut between any two nodes
    long long min_cut(int u, int v)
    {
        if (u == v)
            return 0; // Or INF depending on problem definition

        long long res = 2e18; // Use large INF

        // Walk up the tree from u and v simultaneously to find their LCA,
        // maintaining the minimum edge weight seen on the path.
        while (u != v)
        {
            if (depth[u] > depth[v])
            {
                res = min(res, weight[u]);
                u = p[u];
            }
            else
            {
                res = min(res, weight[v]);
                v = p[v];
            }
        }
        return res;
    }
};
\end{lstlisting}
\subsubsection{MCMF}
\begin{lstlisting}[language=C++]
#include <bits/stdc++.h>
using namespace std;

/**
 *  THE ULTIMATE CP TEMPLATE BLUEPRINT: MIN COST MAX FLOW (MCMF)
 *
 * 1.  WHAT PROBLEM DOES THIS SOLVE?
 * - Keywords: "Cheapest Flow", "Matching with Costs", "Assignment Problem".
 * - Classic Scenarios: You have workers and tasks with different costs. You need
 *   to assign every worker a task such that the total cost is minimized.
 * - The Magic: It works by finding the "shortest" path (where "shortest" is defined by
 *   edge cost) in the residual graph. It keeps augmenting flow along these shortest
 *   paths until no path from source to sink exists.
 *
 * 2.  HOW TO USE IT (THE BLACK BOX)
 * - Initialization:
 *       MCMF solver(N);
 * - Add Edge:
 *       solver.add_edge(u, v, capacity, cost);
 * - Query:
 *       pair<long long, long long> result = solver.min_cost_max_flow(source, sink);
 *       // result.first = max flow, result.second = min cost
 *
 * - Complexity:
 *       O(F * E * log V) using Dijkstra + Potentials.
 *       O(F * E * V) using SPFA (included below for simplicity with negative edges).
 */

const long long INF = 1e18;

struct Edge
{
    int to;
    long long cap, flow, cost;
    int rev; // Reverse edge index
};

struct MCMF
{
    int n;
    vector<vector<Edge>> adj;
    vector<long long> dist, parent_edge, parent_node;

    MCMF(int n) : n(n), adj(n), dist(n), parent_edge(n), parent_node(n) {}

    void add_edge(int from, int to, long long cap, long long cost)
    {
        adj[from].push_back({to, cap, 0, cost, (int)adj[to].size()});
        adj[to].push_back({from, 0, 0, -cost, (int)adj[from].size() - 1});
    }

    // SPFA is used here to handle negative costs.
    // If all costs are positive, Dijkstra with potentials is faster.
    bool spfa(int s, int t, long long &flow, long long &cost)
    {
        dist.assign(n, INF);
        parent_node.assign(n, -1);
        parent_edge.assign(n, -1);
        vector<bool> in_queue(n, false);
        queue<int> q;

        dist[s] = 0;
        q.push(s);
        in_queue[s] = true;

        while (!q.empty())
        {
            int u = q.front();
            q.pop();
            in_queue[u] = false;

            for (int i = 0; i < (int)adj[u].size(); i++)
            {
                Edge &e = adj[u][i];
                if (e.cap - e.flow > 0 && dist[e.to] > dist[u] + e.cost)
                {
                    dist[e.to] = dist[u] + e.cost;
                    parent_node[e.to] = u;
                    parent_edge[e.to] = i;
                    if (!in_queue[e.to])
                    {
                        q.push(e.to);
                        in_queue[e.to] = true;
                    }
                }
            }
        }

        if (dist[t] == INF)
            return false;

        // Push flow
        long long push = INF;
        int curr = t;
        while (curr != s)
        {
            int prev = parent_node[curr];
            int idx = parent_edge[curr];
            push = min(push, adj[prev][idx].cap - adj[prev][idx].flow);
            curr = prev;
        }

        flow += push;
        cost += push * dist[t];
        curr = t;
        while (curr != s)
        {
            int prev = parent_node[curr];
            int idx = parent_edge[curr];
            adj[prev][idx].flow += push;
            int rev_idx = adj[prev][idx].rev;
            adj[curr][rev_idx].flow -= push;
            curr = prev;
        }

        return true;
    }

    pair<long long, long long> min_cost_max_flow(int s, int t)
    {
        long long flow = 0, cost = 0;
        while (spfa(s, t, flow, cost))
            ;
        return {flow, cost};
    }
};
\end{lstlisting}
\subsubsection{Max flow}
\begin{lstlisting}[language=C++]
#include <bits/stdc++.h>
using namespace std;

/**
 *  THE ULTIMATE CP TEMPLATE BLUEPRINT: DINIC'S ALGORITHM
 *
 * 1.  WHAT PROBLEM DOES THIS SOLVE?
 * - Keywords: "Max Flow", "Min-Cut", "Bipartite Matching", "Maximum Closure".
 * - Classic Scenarios: You have a network of pipes with capacities and need to find
 *   the maximum amount of "stuff" (flow) that can move from a Source to a Sink.
 * - The Magic:
 *     a) BFS: Builds a Level Graph (only allows flow from distance d to d+1).
 *     b) DFS: Pushes flow along the Level Graph.
 *     c) ptr: The "pointer optimization" ensures we don't re-scan edges that are
 *        already saturated in the current DFS phase (the secret to O(V^2 E)).
 *
 * 2.  HOW TO USE IT (THE BLACK BOX)
 * - Initialization:
 *       Dinic flow(N);
 * - Add Edge:
 *       flow.add_edge(u, v, capacity);
 * - Query:
 *       long long max_f = flow.max_flow(source, sink);
 *
 * - Complexity:
 *       General: O(V^2 * E).
 *       Unit Networks (e.g., Bipartite Matching): O(E * sqrt(V)).
 */

struct Dinic
{
    struct Edge
    {
        int to;
        long long cap, flow;
        int rev; // Index of the reverse edge in adj[to]
    };

    int n;
    vector<vector<Edge>> adj;
    vector<int> level, ptr;

    Dinic(int n) : n(n), adj(n), level(n), ptr(n) {}

    void add_edge(int from, int to, long long cap)
    {
        // Forward edge: capacity 'cap', flow 0
        adj[from].push_back({to, cap, 0, (int)adj[to].size()});
        // Backward edge: capacity 0, flow 0
        adj[to].push_back({from, 0, 0, (int)adj[from].size() - 1});
    }

    // BFS to build the Level Graph
    bool bfs(int s, int t)
    {
        fill(level.begin(), level.end(), -1);
        level[s] = 0;
        queue<int> q;
        q.push(s);
        while (!q.empty())
        {
            int v = q.front();
            q.pop();
            for (auto &edge : adj[v])
            {
                if (edge.cap - edge.flow > 0 && level[edge.to] == -1)
                {
                    level[edge.to] = level[v] + 1;
                    q.push(edge.to);
                }
            }
        }
        return level[t] != -1;
    }

    // DFS to push flow through the Level Graph
    long long dfs(int v, int t, long long pushed)
    {
        if (pushed == 0)
            return 0;
        if (v == t)
            return pushed;

        // ? Dial 1: The Pointer Optimization
        // ptr[v] ensures we don't look at edges that we already know are dead ends
        for (int &cid = ptr[v]; cid < (int)adj[v].size(); cid++)
        {
            auto &edge = adj[v][cid];
            int tr = edge.to;
            if (level[v] + 1 != level[tr] || edge.cap - edge.flow == 0)
                continue;

            long long tr_pushed = dfs(tr, t, min(pushed, edge.cap - edge.flow));
            if (tr_pushed == 0)
                continue;

            edge.flow += tr_pushed;
            adj[tr][edge.rev].flow -= tr_pushed;
            return tr_pushed;
        }
        return 0;
    }

    long long max_flow(int s, int t)
    {
        long long flow = 0;
        while (bfs(s, t))
        {
            fill(ptr.begin(), ptr.end(), 0);
            while (long long pushed = dfs(s, t, 1e18))
            {
                flow += pushed;
            }
        }
        return flow;
    }
};
\end{lstlisting}
\subsection{Graph Traverse}
\subsubsection{Bipartite Graph}
\begin{lstlisting}[language=C++]
// Check Odd Cycle
// using in problems teams or groups or colors
vector<int>  teams;
Graph adj;
int n, m;
bool valid_team = true;
void bfs(int start) {
    queue<int> q;
    q.push(start);
    if (teams[start] == 0)
        teams[start] = 1;
    while (!q.empty())
    {
        int cur = q.front(); q.pop();
        for (auto ch : adj[cur])
        {
            if (!teams[ch])
            {
                q.push(ch);
                teams[ch] = 3 - teams[cur];
            }
            else if (teams[ch] == teams[cur])
            {
                valid_team = false; 
                return;
            }
        }
    }

}
\end{lstlisting}
\subsubsection{Euler tour}
\begin{lstlisting}[language=C++]
const int N = 1e5 + 5;
GRAPH graph(N);
int tin[N], tout[N], value[N], flat[N];
int timer{1};

void dfs(int node, int p)
{
    flat[timer] = value[node];
    tin[node] = timer++;
    for (auto &x : graph[node])
    {
        if (x != p)
            dfs(x, node);
    }
    tout[node] = timer;
    // [tin, tout)
}

\end{lstlisting}
\subsubsection{Top sort}
\begin{lstlisting}[language=C++]
vector<int> res;

void topo_DFS(int node)
{
    vis[node] = true;
    for (auto &neighbor : graph[node])
    {
        if (!vis[neighbor])
            topo_DFS(neighbor);
    }
    res.push_back(node);
}

// Then reverse

vector<int> topSort()
{
    int n = int(graph.size());
    vector<int> inDegree(n);
    for (int i = 0; i < n; i++)
    {
        for (auto &neighbor : graph[i])
            inDegree[neighbor]++;
    }
    queue<int> ready;
    vector<int> ordering;
    for (int i = 0; i < n; i++)
    {
        if (!inDegree[i])
            ready.push(i);
    }
    while (!ready.empty())
    {
        int node = ready.front();
        ready.pop();
        ordering.push_back(node);
        for (auto &neighbor : graph[node])
        {
            if (!--inDegree[neighbor])
                ready.push(neighbor);
        }
    }
    if ((int)ordering.size() != int(graph.size()))
        ordering.clear();
    return ordering;
}

\end{lstlisting}
\subsection{Graph connectivity}
\subsubsection{Articulations}
\begin{lstlisting}[language=C++]
/**
 *  THE ULTIMATE CP TEMPLATE BLUEPRINT: BRIDGES & ARTICULATION POINTS (CUT VERTICES & EDGES)
 *
 * 1.  WHAT PROBLEM DOES THIS SOLVE?
 * - Keywords: "Critical connections", "Single point of failure", "Vulnerable roads/servers", "Network disconnection".
 * - Classic Scenarios: You have a network of cities and bidirectional roads. You need to find which 
 *   specific road (Bridge) or which specific city (Articulation Point), if destroyed, would split the 
 *   network into two or more disconnected components. Naively removing each edge/node and running DFS 
 *   takes O(E * (V + E)), which yields TLE for large graphs.
 * - The Magic: "Discovery Time (disc) and Lowest Reachable Time (low)". As we do a DFS, we timestamp 
 *   when we visit each node (`disc`). We also track the "oldest" ancestor we can reach via a back-edge (`low`).
 *   - BRIDGE: If a child `v` has `low[v] > disc[u]`, it means the subtree at `v` has NO back-edges reaching 
 *     above `u` or even to `u`. The edge `u-v` is its ONLY lifeline!
 *   - ARTICULATION POINT (AP): If `low[v] >= disc[u]`, the subtree at `v` can at best reach `u`, but nothing 
 *     above it. Destroying `u` traps `v`'s subtree!
 *
 * 2.  HOW TO USE IT (THE BLACK BOX)
 * - Initialization: Clear the `adj`, `bridges`, `articulation_points`, `disc`, and `low` arrays. 
 *   Set `timer = 0`.
 * - Build Graph: Add undirected edges to the `adj` list.
 * - Execution: Call `dfs(i, -1)` for all unvisited nodes (to handle disconnected graphs).
 *       for(int i = 1; i <= n; i++) {
 *           if(!disc[i]) dfs(i, -1);
 *       }
 *
 * - Complexity:
 *       Time: O(V + E)  Just a single DFS traversal!
 *       Space: O(V + E) for adjacency list, recursion stack, and tracking arrays.
 *
 * 3.  HOW TO ADAPT IT (THE DIALS & SWITCHES)
 * -  MULTIPLE EDGES (MULTIGRAPH) TRAP: If there are multiple direct edges between `u` and `v`, 
 *   the condition `if(v == p)` is dangerous because it blocks ALL edges back to the parent, even valid 
 *   parallel edges. FIX: Pass the `edge_index` instead of the parent `node` to the DFS, and skip only 
 *   the specific edge that brought you to `u`.
 * - Why a `set` for Articulation Points? A node might be an AP for multiple different branches of its 
 *   subtree. Pushing to a vector would cause duplicates. `std::set` handles this automatically, but if 
 *   you need more speed, use a boolean array `is_AP[MAXN]` instead.
 * - Bridge Tree (2-Edge-Connected Components): After finding all bridges, if you remove them from the graph, 
 *   the remaining connected components can be "shrunk" into single super-nodes. The resulting graph 
 *   will always be a perfect Tree!
 */
#include <bits/stdc++.h>
using namespace std;

const int MAXN = 1e5 + 5;
vector<int> adj[MAXN];
int disc[MAXN], low[MAXN], timer;
vector<pair<int, int>> bridges;
set<int> articulation_points;

void dfs(int u, int p = -1)
{
    disc[u] = low[u] = ++timer;
    int children = 0; // Needed for the root case

    for (int v : adj[u])
    {
        if (v == p)
            continue; // Don't go back to parent

        if (disc[v])
        {
            // Back-edge found!
            low[u] = min(low[u], disc[v]);
        }
        else
        {
            // Tree-edge
            dfs(v, u);
            low[u] = min(low[u], low[v]);

            // BRIDGE CONDITION
            if (low[v] > disc[u])
                bridges.push_back({u, v});

            // ARTICULATION POINT CONDITION
            if (low[v] >= disc[u] && p != -1)
                articulation_points.insert(u);

            children++;
        }
    }

    // ROOT SPECIAL CASE: If root has > 1 child, it's an AP
    if (p == -1 && children > 1)
        articulation_points.insert(u);
}

\end{lstlisting}
\subsubsection{SCC}
\begin{lstlisting}[language=C++]
const int N = 1e5 + 9;

int SCC_id[N];
GRAPH graph, graph_reverse, cond_graph;
int n, m;
vector<int> order, vis, vis_reverse, order_reverse;

void dfs(int node)
{
    vis[node] = true;
    for (auto &x : graph[node])
    {
        if (!vis[x])
            dfs(x);
    }
    order.emplace_back(node);
}

void dfs_reverse(int node)
{
    vis_reverse[node] = true;
    for (auto &x : graph_reverse[node])
    {
        if (!vis_reverse[x])
            dfs_reverse(x);
    }
    order_reverse.emplace_back(node);
}

void SCC()
{
    graph.assign(n + 1, {}), graph_reverse.assign(n + 1, {}), vis.assign(n + 1, {}), vis_reverse.assign(n + 1, {}),
        cond_graph.assign(n + 1, {});
    for (int i = 0; i < m; i++)
    {
        int x, y;
        cin >> x >> y;
        graph[x].emplace_back(y), graph_reverse[y].emplace_back(x);
    }
    order.clear();
    for (int i = 1; i <= n; i++)
    {
        if (!vis[i])
            dfs(i);
    }
    reverse(order.begin(), order.end());
    vector<vector<int>> scc;
    scc.push_back({});
    for (auto &x : order)
    {
        if (!vis_reverse[x])
        {
            order_reverse.clear();
            dfs_reverse(x);
            scc.emplace_back(order_reverse);
            for (auto &val : order_reverse)
                SCC_id[val] = int(scc.size()) - 1;
        }
    }
    for (int u = 1; u <= n; u++)
    {
        for (int v : graph[u])
        {
            if (SCC_id[u] != SCC_id[v])
                cond_graph[SCC_id[u]].emplace_back(SCC_id[v]);
        }
    }
}

\end{lstlisting}
\subsection{Graph counting}
\subsubsection{BEST}
\begin{lstlisting}[language=C++]
#include <bits/stdc++.h>
using namespace std;

/**
 *  THE ULTIMATE CP TEMPLATE BLUEPRINT: BEST THEOREM
 *
 * 1.  WHAT PROBLEM DOES THIS SOLVE?
 * - Keywords: "Count Eulerian Circuits", "Directed Graph Cycles".
 * - Classic Scenarios: You have a directed graph where every node's in-degree equals
 *   its out-degree, and you need to know how many distinct Eulerian circuits it has.
 * - The Magic: The number of Eulerian circuits = t_w(G) * prod((deg(v) - 1)!)
 *   where t_w(G) is the number of arborescences rooted at any vertex w.
 */

const int MOD = 1e9 + 7;

// Gaussian Elimination for Determinant (MOD)
long long power(long long base, long long exp)
{
    long long res = 1;
    base %= MOD;
    while (exp > 0)
    {
        if (exp % 2 == 1)
            res = (res * base) % MOD;
        base = (base * base) % MOD;
        exp /= 2;
    }
    return res;
}

long long modInverse(long long n) { return power(n, MOD - 2); }

long long determinant(vector<vector<long long>> &mat, int n)
{
    long long det = 1;
    for (int i = 0; i < n; i++)
    {
        int pivot = i;
        for (int j = i + 1; j < n; j++)
            if (abs(mat[j][i]) > abs(mat[pivot][i]))
                pivot = j;
        swap(mat[i], mat[pivot]);
        if (i != pivot)
            det = (MOD - det) % MOD;
        if (mat[i][i] == 0)
            return 0;
        det = (det * mat[i][i]) % MOD;
        long long inv = modInverse(mat[i][i]);
        for (int j = i + 1; j < n; j++)
        {
            long long factor = (mat[j][i] * inv) % MOD;
            for (int k = i + 1; k < n; k++)
                mat[j][k] = (mat[j][k] - factor * mat[i][k] % MOD + MOD) % MOD;
        }
    }
    return det;
}

long long best_theorem(int n, vector<vector<int>> &adj)
{
    vector<int> out_degree(n, 0);
    vector<vector<long long>> L(n, vector<long long>(n, 0));

    // 1. Build Directed Laplacian
    for (int u = 0; u < n; ++u)
    {
        for (int v : adj[u])
        {
            out_degree[u]++;
            L[u][v]--;
        }
    }
    for (int i = 0; i < n; ++i)
        L[i][i] += out_degree[i];

    // 2. Compute arborescences (t_w)
    // Delete row 0 and col 0 (root at 0)
    vector<vector<long long>> minor(n - 1, vector<long long>(n - 1));
    for (int i = 1; i < n; i++)
        for (int j = 1; j < n; j++)
            minor[i - 1][j - 1] = (L[i][j] % MOD + MOD) % MOD;

    long long t_w = determinant(minor, n - 1);

    // 3. Multiply by prod((deg(v) - 1)!)
    long long fact_prod = 1;
    vector<long long> fact(n + 1, 1);
    for (int i = 2; i <= n; i++)
        fact[i] = (fact[i - 1] * i) % MOD;

    for (int i = 0; i < n; i++)
    {
        if (out_degree[i] > 0)
            fact_prod = (fact_prod * fact[out_degree[i] - 1]) % MOD;
    }

    return (t_w * fact_prod) % MOD;
}

int main()
{
    ios_base::sync_with_stdio(false);
    cin.tie(NULL);
    // int n; cin >> n; ... adj list setup ...
    return 0;
}
\end{lstlisting}
\subsubsection{Prufer code}
\begin{lstlisting}[language=C++]
#include <bits/stdc++.h>
using namespace std;

/**
 *  THE ULTIMATE CP TEMPLATE BLUEPRINT: PRUFER CODE
 *
 * 1.  WHAT PROBLEM DOES THIS SOLVE?
 * - Keywords: "Labeled Trees", "Cayley's Formula", "Tree Counting".
 * - Classic Scenarios: Bijectively mapping a labeled tree to a sequence of length N-2.
 * - The Magic: It allows you to count trees with specific constraints (like degrees)
 *   by counting sequences instead, which is mathematically trivial.
 */

struct Prufer
{
    // Converts a labeled tree (adj list) to a Prufer sequence of length N-2
    // Complexity: O(N log N)
    static vector<int> tree_to_prufer(int n, vector<vector<int>> &adj)
    {
        if (n == 2)
            return {}; // A 2-node tree has no Prufer sequence

        vector<int> degree(n + 1);
        priority_queue<int, vector<int>, greater<int>> leaves;

        for (int i = 1; i <= n; i++)
        {
            degree[i] = adj[i].size();
            if (degree[i] == 1)
                leaves.push(i);
        }

        vector<int> prufer;
        vector<bool> removed(n + 1, false);

        for (int i = 0; i < n - 2; i++)
        {
            int leaf = leaves.top();
            leaves.pop();
            removed[leaf] = true;

            // Find the neighbor of this leaf
            int neighbor = -1;
            for (int v : adj[leaf])
            {
                if (!removed[v])
                {
                    neighbor = v;
                    break;
                }
            }

            prufer.push_back(neighbor);
            degree[neighbor]--;
            if (degree[neighbor] == 1)
                leaves.push(neighbor);
        }

        return prufer;
    }

    // Converts a Prufer sequence of length N-2 back to an edge list of the tree
    // Complexity: O(N log N)
    static vector<pair<int, int>> prufer_to_tree(int n, const vector<int> &prufer)
    {
        vector<int> degree(n + 1, 1);
        for (int x : prufer)
            degree[x]++;

        priority_queue<int, vector<int>, greater<int>> leaves;
        for (int i = 1; i <= n; i++)
        {
            if (degree[i] == 1)
                leaves.push(i);
        }

        vector<pair<int, int>> edges;
        for (int x : prufer)
        {
            int leaf = leaves.top();
            leaves.pop();
            edges.push_back({leaf, x});

            degree[x]--;
            if (degree[x] == 1)
                leaves.push(x);
        }

        // Connect the last two remaining nodes
        int u = leaves.top();
        leaves.pop();
        int v = leaves.top();
        leaves.pop();
        edges.push_back({u, v});

        return edges;
    }
};

void solve()
{
    int n;
    cin >> n;
    // Example: Tree to Prufer
    // vector<vector<int>> adj(n+1);
    // ... fill adj ...
    // vector<int> p = Prufer::tree_to_prufer(n, adj);
}

int main()
{
    ios_base::sync_with_stdio(false);
    cin.tie(NULL);
    return 0;
}
\end{lstlisting}
\subsection{LCA}
\subsubsection{LCA(Hassan)}
\begin{lstlisting}[language=C++]
struct LCA {
    vector<vector<pair<int,int>> > adj;
    vector<vector<int>> up,cost[2];
    vector<int> depth;
    int LOG = 30;
    int ign = 1e9;
    void build(int n, vector<vector<pair<int,int>> > &G, int root = 1) {
        adj = G;
        up = vector<vector<int> >(LOG + 5, vector<int>(n + 5));
        cost[0] = vector<vector<int> >(LOG + 5, vector<int>(n + 5));
        cost[1] = vector<vector<int> >(LOG + 5, vector<int>(n + 5));
        depth = vector<int>(n + 5);
        dfs(root, -1);
    }
    void dfs(int u, int p) {
        for (auto [v,w]: adj[u]) {
            if (v == p) continue;
            depth[v] = depth[u] + 1;

            up[0][v] = u;
            cost[0][0][v] = w;
            cost[1][0][v] = w;
            for (int i = 1; i < LOG; i++) {
                up[i][v] = up[i - 1][up[i - 1][v]];
                cost[0][i][v] = min(cost[0][i - 1][v], cost[0][i - 1][up[i - 1][v]]);
                cost[1][i][v]  = max(cost[1][i - 1][v], cost[1][i - 1][up[i - 1][v]]);

            }
            dfs(v, u);
        }
    }

    int kth_ancestor(int u, int k) {
        for (int i = 0; i < LOG; i++) {
            if (k >> i & 1) {
                u = up[i][u];
            }
        }
        return u;
    }

    int getLca(int u, int v) {
        if (depth[u] < depth[v]) swap(u, v);
        int k = depth[u] - depth[v];
        u = kth_ancestor(u, k);
        if (u == v) return u;

        for (int i = LOG - 1; i >= 0; --i) {
            if (up[i][u] != up[i][v]) {
                u = up[i][u];
                v = up[i][v];
            }
        }
        return up[0][v];
    }
    pair<int,int> query(int u,int v) {
        int mn = ign,mx = -ign;
        int lca = getLca(u, v);
        int k = depth[u] - depth[lca];
        for (int i=0;i<LOG;i++) {
            if (k>>i&1) {
                mn = min(mn,cost[0][i][u]);
                mx = max(mx,cost[1][i][u]);
                u = up[i][u];
            }
        }
        k = depth[v] - depth[lca];
        for (int i=0;i<LOG;i++) {
            if (k>>i&1) {
                mn = min(mn,cost[0][i][v]);
                mx = max(mx,cost[1][i][v]);
                v = up[i][v];
            }
        }
        return {mn,mx};
    }


};

\end{lstlisting}
\subsubsection{LCA}
\begin{lstlisting}[language=C++]
struct LCA
{
    int n, LG, SKIP;
    vector<vector<int>> T, Cost;
    vector<int> par, dep;

    LCA(int nodes_count, const GRAPH &graph, int root = 1, int skip_val = OO)
    {
        n = nodes_count;
        LG = __lg(n) + 2;
        SKIP = skip_val;
        T.assign(LG, vector<int>(n + 1, -1));
        Cost.assign(LG, vector<int>(n + 1, SKIP));
        par.assign(n + 1, -1);
        dep.assign(n + 1, 0);

        Cost[0][root] = SKIP;
        dfs(root, -1, 0, graph);

        for (int i = 1; i <= n; i++)
            T[0][i] = par[i];

        for (int pw = 1; (1 << pw) <= n; pw++)
        {
            for (int i = 1; i <= n; i++)
            {
                if (T[pw - 1][i] != -1)
                {
                    T[pw][i] = T[pw - 1][T[pw - 1][i]];
                    Cost[pw][i] = merge(Cost[pw - 1][i], Cost[pw - 1][T[pw - 1][i]]);
                }
            }
        }
    }

    int merge(int a, int b)
    {
        return min(a, b);
    }

    void dfs(int node, int p, int d, const GRAPH &graph)
    {
        dep[node] = d;
        par[node] = p;
        for (auto &[i, w] : graph[node])
        {
            if (i != p)
            {
                Cost[0][i] = w;
                dfs(i, node, d + 1, graph);
            }
        }
    }

    int get_kth(int x, int k)
    {
        while (k)
        {
            if (x == -1)
                break;
            int pw = __builtin_ctz(k);
            x = T[pw][x];
            k &= (k - 1);
        }
        return x;
    }

    int get_lca(int u, int v)
    {
        if (dep[u] < dep[v])
            swap(u, v);
        u = get_kth(u, dep[u] - dep[v]);
        if (u == v)
            return u;
        for (int j = LG - 1; j >= 0; j--)
        {
            if (T[j][u] != T[j][v])
            {
                u = T[j][u];
                v = T[j][v];
            }
        }
        return par[u];
    }

    int get_cost(int u, int v)
    {
        if (dep[u] < dep[v])
            swap(u, v);
        int k = dep[u] - dep[v], x = u;
        int cost = SKIP;
        while (k)
        {
            if (x == -1)
                break;
            int pw = __builtin_ctz(k);
            cost = merge(cost, Cost[pw][x]);
            x = T[pw][x];
            k &= (k - 1);
        }
        return cost;
    }

    int query_path(int u, int v)
    {
        int lca = get_lca(u, v);
        return merge(get_cost(u, lca), get_cost(v, lca));
    }

    int get_dist(int u, int v)
    {
        int lca = get_lca(u, v);
        return dep[u] + dep[v] - 2 * dep[lca];
    }
};

\end{lstlisting}
\subsubsection{Lca o1}
\begin{lstlisting}[language=C++]
  vector<int> dep(n), in(n), out(n), seq, pa(n, -1);

  auto dfs = [&](auto &&self, int u) -> void {
	in[u] = seq.size();
	seq.push_back(u);
	for (auto v : G[u])
	  if (v != pa[u]) {
		pa[v] = u, dep[v] = dep[u] + 1;
		self(self, v);
	  }
	out[u] = seq.size();
  };
  seq.reserve(n);
  dfs(dfs, 0);

  const int M = __lg(n);
  vector dp(M + 1, vector<int>(n));
  auto cmp = [&](int a, int b) { return dep[a] < dep[b] ? a : b; };

  dp[0] = seq;
  for (int i = 0; i < M; i++)
	for (int j = 0; j + (2 << i) <= n; j++)
	  dp[i + 1][j] = cmp(dp[i][j], dp[i][j + (1 << i)]);

  auto lca = [&](int a, int b) {
	if (a == b)return a;
	a = in[a] + 1, b = in[b] + 1;
	if (a > b)swap(a, b);
	int h = __lg(b - a);
	return pa[cmp(dp[h][a], dp[h][b - (1 << h)])];
  };

\end{lstlisting}
\subsection{MST}
\subsubsection{Directed MST}
\begin{lstlisting}[language=C++]
#include <bits/stdc++.h>
using namespace std;

/**
 *  THE ULTIMATE CP TEMPLATE BLUEPRINT: DIRECTED MST (CHU-LIU/EDMONDS)
 *
 * 1.  WHAT PROBLEM DOES THIS SOLVE?
 * - Keywords: "Minimum Spanning Arborescence", "Directed MST".
 * - Classic Scenarios: Finding the cheapest way to connect all nodes in a DIRECTED graph
 *   rooted at a specific node, such that every node (except the root) has exactly 1 incoming edge.
 * - The Magic: It greedily picks the minimum incoming edge for every node. If this creates
 *   cycles, it "contracts" those cycles into a single super-node, updates the weights of edges
 *   entering the cycle, and recurses.
 *
 * 2.  HOW TO USE IT
 * - edge: {u, v, cost}
 * - roots: The node that acts as the root of the arborescence.
 * - Complexity: O(V * E)
 */

struct Edge
{
    int u, v;
    long long cost;
};

long long edmonds(int n, int root, vector<Edge> &edges)
{
    long long total_cost = 0;
    vector<int> min_edge(n + 1);
    vector<int> parent(n + 1);
    vector<int> id(n + 1);
    vector<int> vis(n + 1);

    while (true)
    {
        fill(min_edge.begin(), min_edge.end(), 2e18); // Use large INF
        for (const auto &e : edges)
        {
            if (e.u != e.v && e.cost < min_edge[e.v])
            {
                min_edge[e.v] = e.cost;
                parent[e.v] = e.u;
            }
        }

        for (int i = 1; i <= n; i++)
        {
            if (i != root && min_edge[i] == 2e18)
                return -1; // Impossible
        }

        int cnt_nodes = 0;
        fill(id.begin(), id.end(), -1);
        fill(vis.begin(), vis.end(), -1);
        min_edge[root] = 0;

        for (int i = 1; i <= n; i++)
        {
            total_cost += min_edge[i];
            int v = i;
            while (vis[v] != i && id[v] == -1 && v != root)
            {
                vis[v] = i;
                v = parent[v];
            }
            if (v != root && id[v] == -1)
            {
                for (int u = parent[v]; u != v; u = parent[u])
                    id[u] = cnt_nodes;
                id[v] = cnt_nodes++;
            }
        }

        if (cnt_nodes == 0)
            break;
        for (int i = 1; i <= n; i++)
        {
            if (id[i] == -1)
                id[i] = cnt_nodes++;
        }

        for (auto &e : edges)
        {
            long long v_cost = min_edge[e.v];
            e.u = id[e.u];
            e.v = id[e.v];
            if (e.u != e.v)
                e.cost -= v_cost;
        }
        n = cnt_nodes;
        root = id[root];
    }
    return total_cost;
}
\end{lstlisting}
\subsubsection{Kirchoff}
\begin{lstlisting}[language=C++]
#include <bits/stdc++.h>
using namespace std;

/**
 *  THE ULTIMATE CP TEMPLATE BLUEPRINT: KIRCHHOFF'S THEOREM
 *
 * 1.  WHAT PROBLEM DOES THIS SOLVE?
 * - Keywords: "Count spanning trees", "Matrix Tree Theorem".
 * - Classic Scenarios: You have a graph and need to know exactly how many different
 *   spanning trees can be formed.
 * - The Magic: It converts a graph problem into a linear algebra problem.
 *   It constructs the Laplacian Matrix and computes a determinant using
 *   Gaussian Elimination.
 *
 * 2.  HOW TO USE IT
 * - Build the Laplacian matrix (L = Degree - Adjacency).
 * - Remove the last row and last column.
 * - Compute the determinant modulo MOD.
 * - Complexity: O(N^3) (due to Gaussian Elimination).
 */

const int MOD = 1e9 + 7;

long long power(long long base, long long exp)
{
    long long res = 1;
    base %= MOD;
    while (exp > 0)
    {
        if (exp % 2 == 1)
            res = (res * base) % MOD;
        base = (base * base) % MOD;
        exp /= 2;
    }
    return res;
}

long long modInverse(long long n)
{
    return power(n, MOD - 2);
}

//  Gaussian Elimination to find Determinant
long long determinant(vector<vector<long long>> &mat, int n)
{
    long long det = 1;
    for (int i = 0; i < n; i++)
    {
        int pivot = i;
        for (int j = i + 1; j < n; j++)
        {
            if (abs(mat[j][i]) > abs(mat[pivot][i]))
                pivot = j;
        }
        swap(mat[i], mat[pivot]);
        if (i != pivot)
            det = (MOD - det) % MOD;
        if (mat[i][i] == 0)
            return 0;

        det = (det * mat[i][i]) % MOD;
        long long inv = modInverse(mat[i][i]);
        for (int j = i + 1; j < n; j++)
        {
            long long factor = (mat[j][i] * inv) % MOD;
            for (int k = i + 1; k < n; k++)
            {
                mat[j][k] = (mat[j][k] - factor * mat[i][k] % MOD + MOD) % MOD;
            }
        }
    }
    return det;
}

void solve()
{
    int n, m;
    cin >> n >> m;

    // Laplacian Matrix
    vector<vector<long long>> L(n, vector<long long>(n, 0));

    for (int i = 0; i < m; i++)
    {
        int u, v;
        cin >> u >> v;
        --u;
        --v; // 0-indexed

        L[u][u]++;
        L[v][v]++;
        L[u][v]--;
        L[v][u]--;
    }

    // Convert to positive modulo
    for (int i = 0; i < n; i++)
        for (int j = 0; j < n; j++)
            L[i][j] = (L[i][j] % MOD + MOD) % MOD;

    // Remove last row and column to get the Minor
    vector<vector<long long>> minor(n - 1, vector<long long>(n - 1));
    for (int i = 0; i < n - 1; i++)
    {
        for (int j = 0; j < n - 1; j++)
        {
            minor[i][j] = L[i][j];
        }
    }

    cout << determinant(minor, n - 1) << "\n";
}

int main()
{
    ios_base::sync_with_stdio(false);
    cin.tie(NULL);
    solve();
    return 0;
}
\end{lstlisting}
\subsubsection{Kruskal}
\begin{lstlisting}[language=C++]
class DSU
{
private:
    vector<int> parent, rank;
    int forests{};

    void link(int x, int y)
    {
        if (rank[x] > rank[y])
            swap(x, y);
        parent[x] = y;
        if (rank[x] == rank[y])
            rank[y]++;
    }

public:
    DSU(int n)
    {
        parent = vector<int>(n), rank = vector<int>(n);
        forests = n;
        for (int i = 0; i < n; i++)
            parent[i] = i, rank[i] = 1;
    }

    int find_set(int x)
    {
        if (x == parent[x])
            return x;
        return parent[x] = find_set(parent[x]);
    }

    int union_sets(int x, int y)
    {
        x = find_set(x), y = find_set(y);
        if (x != y)
        {
            link(x, y);
            forests--;
        }
        return x != y;
    }
};

struct edge
{
    int from, to, w;

    edge(int from, int to, int w) : from(from), to(to), w(w)
    {
    }

    bool operator<(const edge &e) const
    {
        return w < e.w;
    }
};

int Kruskal(vector<edge> &edgeList, int n)
{
    DSU uf(n + 1);
    vector<edge> edges;
    int mnCost{};
    sort(edgeList.begin(), edgeList.end());
    for (auto &[from, to, cost] : edgeList)
    {
        if (uf.union_sets(from, to))
        {
            edges.emplace_back(from, to, cost);
            mnCost += cost;
        }
    }
    return mnCost;
}

\end{lstlisting}
\subsubsection{Prim}
\begin{lstlisting}[language=C++]
int n;
GRAPH graph;

void prim(int node)
{
    int total{};
    priority_queue<pair<int, int>, vector<pair<int, int>>, greater<pair<int, int>>> pq;
    pq.emplace(0, node);
    vector<bool> visited(n + 1, 0);
    vector<int> dist(n + 1, OO);
    dist[node] = 0;
    while (!pq.empty())
    {
        auto [w, node] = pq.top();
        pq.pop();
        if (visited[node])
            continue;
        visited[node] = true;
        total += w;
        for (auto &[to, c] : graph[node])
        {
            if (!visited[to] && dist[to] > c)
            {
                dist[to] = c;
                pq.push({dist[to], to});
            }
        }
    }
}

\end{lstlisting}
\subsection{Matching}
\subsubsection{Bipartite matching}
\begin{lstlisting}[language=C++]
#include <bits/stdc++.h>
using namespace std;

/**
 *  THE ULTIMATE CP TEMPLATE BLUEPRINT: HOPCROFT-KARP ALGORITHM
 *
 * 1.  WHAT PROBLEM DOES THIS SOLVE?
 * - Keywords: "Maximum Bipartite Matching", "Maximum Independent Set", "Minimum Vertex Cover".
 * - Classic Scenarios: You have two distinct groups (e.g., Workers and Jobs, or Boys and Girls).
 *   Each worker can only do specific jobs. You want to pair up as many workers with jobs as possible
 *   such that no worker gets two jobs, and no job gets two workers.
 * - The Magic: It finds the maximum matching in a bipartite graph. It works by using BFS to find
 *   multiple shortest augmenting paths simultaneously, and then uses DFS to push flow along them.
 *
 * 2.  HOW TO USE IT (THE BLACK BOX)
 * - Initialization: Pass the size of the Left partition (N) and Right partition (M).
 *   (1-indexed nodes!).
 *       HopcroftKarp hk(N, M);
 *
 * - Add Edge:
 *       hk.add_edge(u, v); // Worker `u` (1 to N) can do Job `v` (1 to M).
 *
 * - Query:
 *       int max_pairs = hk.max_matching();
 *
 * - Check Pairs:
 *       // After running max_matching(), pair_u[u] gives the job assigned to worker u.
 *       // If pair_u[u] == 0, worker u is unmatched.
 *
 * - Complexity:
 *       Time: O(E * sqrt(V)). Blazing fast for up to 10^5 nodes and edges.
 *       Space: O(V + E) memory.
 *
 * 3.  HOW TO ADAPT IT (THE DIALS & SWITCHES)
 * - Knig's Theorem (Min Vertex Cover):
 *   To find the minimum number of nodes needed to "cover" all edges:
 *   1. Run Max Matching. Min Vertex Cover = Max Matching.
 *   2. To find the ACTUAL nodes: Run a DFS from all UNMATCHED nodes in the Left partition
 *      using alternating paths. The Vertex Cover = (Unvisited nodes in Left) + (Visited nodes in Right).
 *
 * - Maximum Independent Set:
 *   The maximum number of nodes where no two nodes share an edge.
 *   Formula: Max Independent Set = (Total Nodes) - (Min Vertex Cover).
 */

struct HopcroftKarp
{
    int n; // Size of the Left partition (U)
    int m; // Size of the Right partition (V)

    vector<vector<int>> adj;

    // pair_u[i] = the node in V matched with node i in U
    // pair_v[j] = the node in U matched with node j in V
    vector<int> pair_u, pair_v;

    // distance array used by BFS to find shortest augmenting paths
    vector<int> dist;

    const int INF = 1e9;

    HopcroftKarp(int n, int m) : n(n), m(m)
    {
        adj.resize(n + 1);
        pair_u.assign(n + 1, 0);
        pair_v.assign(m + 1, 0);
        dist.assign(n + 1, 0);
    }

    void add_edge(int u, int v)
    {
        adj[u].push_back(v); // u is in Left (1...N), v is in Right (1...M)
    }

    // BFS builds the level graph, finding the shortest paths to unmatched nodes in V
    bool bfs()
    {
        queue<int> q;

        for (int u = 1; u <= n; u++)
        {
            // If u is unmatched, it's a starting point for an augmenting path
            if (pair_u[u] == 0)
            {
                dist[u] = 0;
                q.push(u);
            }
            else
            {
                dist[u] = INF; // Otherwise, mark as unvisited
            }
        }

        // Node 0 acts as a dummy node for all unmatched nodes in V
        dist[0] = INF;

        while (!q.empty())
        {
            int u = q.front();
            q.pop();

            // If this node has a distance less than the dummy node, we can explore it
            if (dist[u] < dist[0])
            {
                for (int v : adj[u])
                {
                    // If the node matched to v is unvisited, we can walk there
                    if (dist[pair_v[v]] == INF)
                    {
                        dist[pair_v[v]] = dist[u] + 1;
                        q.push(pair_v[v]);
                    }
                }
            }
        }

        // If dist[0] is no longer INF, we found at least one augmenting path!
        return dist[0] != INF;
    }

    // DFS greedily pushes "flow" along the paths discovered by BFS
    bool dfs(int u)
    {
        if (u != 0)
        {
            for (int v : adj[u])
            {
                // Ensure we only walk along the shortest path DAG built by BFS
                if (dist[pair_v[v]] == dist[u] + 1)
                {
                    if (dfs(pair_v[v]))
                    {
                        // Augment the path! Flip the matching status.
                        pair_v[v] = u;
                        pair_u[u] = v;
                        return true;
                    }
                }
            }
            // If no augmenting path could be found from u, mark it as dead end
            dist[u] = INF;
            return false;
        }
        // Base case: Reached the dummy node 0 (an unmatched node in V)
        return true;
    }

    // Repeatedly runs BFS and DFS to build the maximum matching
    int max_matching()
    {
        int matching = 0;

        // While there is an augmenting path...
        while (bfs())
        {
            // Try to augment from every unmatched node in U
            for (int u = 1; u <= n; u++)
            {
                if (pair_u[u] == 0 && dfs(u))
                {
                    matching++;
                }
            }
        }
        return matching;
    }
};

void solve()
{
    int n, m, edges;
    // Example: N workers, M jobs, and E compatible worker-job pairs
    if (!(cin >> n >> m >> edges))
        return;

    HopcroftKarp hk(n, m);

    for (int i = 0; i < edges; i++)
    {
        int u, v;
        cin >> u >> v;
        hk.add_edge(u, v);
    }

    cout << "Maximum Pairs: " << hk.max_matching() << "\n";

    // Print the actual assignments
    for (int u = 1; u <= n; u++)
    {
        if (hk.pair_u[u] != 0)
        {
            cout << "Worker " << u << " is assigned Job " << hk.pair_u[u] << "\n";
        }
    }
}

int main()
{
    ios_base::sync_with_stdio(false);
    cin.tie(NULL);
    // solve();
    return 0;
}
\end{lstlisting}
\subsubsection{Hungarian algorithm}
\begin{lstlisting}[language=C++]
#include <bits/stdc++.h>
using namespace std;

/**
 *  THE ULTIMATE CP TEMPLATE BLUEPRINT: HUNGARIAN ALGORITHM
 *
 * 1.  WHAT PROBLEM DOES THIS SOLVE?
 * - Keywords: "Optimal Assignment", "Bipartite Weighted Matching", "Min-Cost Perfect Matching".
 * - Classic Scenarios: You have N workers and M jobs. Every worker-job pair has a specific
 *   cost or time to complete. You want to assign exactly one job to each worker such that
 *   the total cost is minimized.
 * - The Magic: While MCMF works, Hungarian runs in a strict, highly-optimized O(N * M^2) time,
 *   making it the undisputed king for dense assignment matrices.
 *
 * 2.  HOW TO USE IT (THE BLACK BOX)
 * - Initialization:
 *       // Create a 1-indexed cost matrix of size (N+1) x (M+1)
 *       // Important: N must be <= M! (Rows <= Cols)
 *       Hungarian solver(n, m, cost_matrix);
 *
 * - Query:
 *       long long min_cost = solver.solve();
 *
 * - Extract Assignments:
 *       // solver.ans[i] will contain the job (1 to M) assigned to worker i.
 *
 * 3.  HOW TO ADAPT IT (THE DIALS & SWITCHES)
 * - Maximum Profit Assignment?
 *   If you want to MAXIMIZE the score, simply negate all values in your matrix (cost = -cost)
 *   before passing it in. Then negate the final answer returned by `solve()`.
 * - Forbidden Assignments?
 *   If worker `i` cannot do job `j`, set `cost[i][j] = 1e15` (Infinity).
 */

struct Hungarian
{
    int n, m;
    vector<vector<long long>> a;
    vector<long long> u, v, minv;
    vector<int> p, way, ans;
    const long long INF = 1e15;

    Hungarian(int n, int m, const vector<vector<long long>> &cost_matrix)
    {
        this->n = n;
        this->m = m;
        // 1-based indexing logic
        a = cost_matrix;
        u.assign(n + 1, 0);
        v.assign(m + 1, 0);
        p.assign(m + 1, 0);
        way.assign(m + 1, 0);
        ans.assign(n + 1, 0);
    }

    /* STREAMING_CHUNK: Core O(N * M^2) solver using augmenting paths */
    long long solve()
    {
        // Iterate over every worker (row) to find them a job
        for (int i = 1; i <= n; ++i)
        {
            p[0] = i;
            int j0 = 0; // Dummy column 0

            minv.assign(m + 1, INF);
            vector<bool> used(m + 1, false);

            do
            {
                used[j0] = true;
                int i0 = p[j0], j1 = 0;
                long long delta = INF;

                // Update minv array and find the next column to explore
                for (int j = 1; j <= m; ++j)
                {
                    if (!used[j])
                    {
                        long long cur = a[i0][j] - u[i0] - v[j];
                        if (cur < minv[j])
                        {
                            minv[j] = cur;
                            way[j] = j0;
                        }
                        if (minv[j] < delta)
                        {
                            delta = minv[j];
                            j1 = j;
                        }
                    }
                }

                // Update potentials (dual variables)
                for (int j = 0; j <= m; ++j)
                {
                    if (used[j])
                    {
                        u[p[j]] += delta;
                        v[j] -= delta;
                    }
                    else
                    {
                        minv[j] -= delta;
                    }
                }
                j0 = j1;
            } while (p[j0] != 0); // Stop when we reach an unmatched column

            do
            {
                int j1 = way[j0];
                p[j0] = p[j1];
                j0 = j1;
            } while (j0 != 0);
        }

        // The optimal job assigned to worker `i`
        for (int j = 1; j <= m; ++j)
        {
            if (p[j] != 0)
            {
                ans[p[j]] = j;
            }
        }

        // The final minimum cost is stored negated in the dummy column potential v[0]
        return -v[0];
    }
};

void example_usage()
{
    int n, m;
    // Example: N workers, M jobs (ensure N <= M)
    if (!(cin >> n >> m))
        return;

    // Build the 1-indexed cost matrix
    vector<vector<long long>> cost(n + 1, vector<long long>(m + 1, 0));
    for (int i = 1; i <= n; i++)
    {
        for (int j = 1; j <= m; j++)
        {
            cin >> cost[i][j];
        }
    }

    Hungarian solver(n, m, cost);
    cout << "Minimum Total Cost: " << solver.solve() << "\n";

    for (int i = 1; i <= n; i++)
    {
        cout << "Worker " << i << " assigned to Job " << solver.ans[i] << "\n";
    }
}
\end{lstlisting}
\subsubsection{Stable Marriage Problem}
\begin{lstlisting}[language=C++]
#include <bits/stdc++.h>
using namespace std;

/**
 *  THE ULTIMATE CP TEMPLATE BLUEPRINT: STABLE MARRIAGE (GALE-SHAPLEY)
 *
 * 1.  WHAT PROBLEM DOES THIS SOLVE?
 * - Keywords: "Stable Matching", "Preference Lists", "Gale-Shapley".
 * - Classic Scenarios: Matching N men to N women, or N students to N colleges,
 *   based on ranked preference lists. A matching is "stable" if there is no
 *   pair (Man A, Woman B) who both prefer each other over their current partners.
 * - The Magic: The Gale-Shapley algorithm guarantees a stable matching in O(N^2) time.
 *
 * 2.  HOW TO USE IT (THE BLACK BOX)
 * - Initialization:
 *       StableMarriage sm(N);
 *       sm.set_preferences(men_prefs, women_prefs);
 * - Query:
 *       sm.solve();
 * - Extract Assignments:
 *       // match_proposer[i] contains the receiver assigned to proposer i.
 *       // match_receiver[j] contains the proposer assigned to receiver j.
 *
 * 3.  HOW TO ADAPT IT (THE DIALS & SWITCHES)
 * - The Proposer Advantage:
 *   Gale-Shapley strongly favors the group doing the proposing!
 *   If you want the optimal outcome for Group X, pass Group X as the "Proposers"
 *   (the first argument in set_preferences).
 * - Unequal sizes: If Proposers > Receivers, add "dummy" receivers that rank
 *   everyone equally at the bottom.
 */

struct StableMarriage
{
    int n;

    // pref_proposer[i][k] = the k-th preferred receiver by proposer i
    vector<vector<int>> pref_proposer;

    // To make rejection O(1), we invert the receiver's preferences.
    // rank_receiver[j][i] = the rank position (1 to N) of proposer i in receiver j's list.
    // Lower rank number means highly preferred!
    vector<vector<int>> rank_receiver;

    // The current matches. 0 means unassigned.
    vector<int> match_proposer;
    vector<int> match_receiver;

    StableMarriage(int n) : n(n)
    {
        pref_proposer.assign(n + 1, vector<int>(n + 1, 0));
        rank_receiver.assign(n + 1, vector<int>(n + 1, 0));
        match_proposer.assign(n + 1, 0);
        match_receiver.assign(n + 1, 0);
    }

    // Both 2D vectors should be 1-indexed: [1..n][1..n]
    void set_preferences(const vector<vector<int>> &proposers, const vector<vector<int>> &receivers)
    {
        pref_proposer = proposers;

        // Invert the receiver's preference list for O(1) comparison later
        for (int j = 1; j <= n; ++j)
        {
            for (int rank = 1; rank <= n; ++rank)
            {
                int proposer_id = receivers[j][rank];
                rank_receiver[j][proposer_id] = rank;
            }
        }
    }

    void solve()
    {
        queue<int> free_proposers;
        // next_proposal[i] tracks the index of the next receiver that proposer i will ask
        vector<int> next_proposal(n + 1, 1);

        for (int i = 1; i <= n; ++i)
        {
            free_proposers.push(i);
        }

        // Loop until every proposer is matched
        while (!free_proposers.empty())
        {
            int p = free_proposers.front();

            // The receiver this proposer is going to ask next
            int r = pref_proposer[p][next_proposal[p]];
            next_proposal[p]++; // Increment pointer so he doesn't ask her again if rejected

            if (match_receiver[r] == 0)
            {
                // Scenario 1: The receiver is entirely free. She happily accepts!
                match_receiver[r] = p;
                match_proposer[p] = r;
                free_proposers.pop(); // The proposer is no longer free
            }
            else
            {
                // Scenario 2: The receiver is already engaged to someone else (current_p).
                int current_p = match_receiver[r];

                // Does she prefer the new proposer (p) over her current partner (current_p)?
                if (rank_receiver[r][p] < rank_receiver[r][current_p])
                {
                    // She dumps her current partner!
                    match_proposer[current_p] = 0;
                    free_proposers.push(current_p); // The dumped partner is now free and goes back in line

                    // She accepts the new proposer
                    match_receiver[r] = p;
                    match_proposer[p] = r;
                    free_proposers.pop();
                }
                // If she prefers her current partner, she rejects 'p', and 'p' stays in the queue.
            }
        }
    }
};

void solve_example()
{
    int n;
    // Read N pairs
    if (!(cin >> n))
        return;

    // 1-indexed preference lists
    vector<vector<int>> men_prefs(n + 1, vector<int>(n + 1));
    vector<vector<int>> women_prefs(n + 1, vector<int>(n + 1));

    // Read Men's preferences (Proposers)
    for (int i = 1; i <= n; ++i)
    {
        for (int j = 1; j <= n; ++j)
        {
            cin >> men_prefs[i][j];
        }
    }

    // Read Women's preferences (Receivers)
    for (int i = 1; i <= n; ++i)
    {
        for (int j = 1; j <= n; ++j)
        {
            cin >> women_prefs[i][j];
        }
    }

    StableMarriage sm(n);
    sm.set_preferences(men_prefs, women_prefs);
    sm.solve();

    cout << "Stable Matches (Man -> Woman):\n";
    for (int i = 1; i <= n; ++i)
    {
        cout << "Man " << i << " matches with Woman " << sm.match_proposer[i] << "\n";
    }
}

int main()
{
    ios_base::sync_with_stdio(false);
    cin.tie(NULL);
    // solve_example();
    return 0;
}
\end{lstlisting}
\subsection{Misc}
\subsubsection{Check Cycle }
\begin{lstlisting}[language=C++]
// Check Cycle in Undirected
Graph graph ;
vector<bool>  vis;
bool isCycle = 0;
void dfs(int u,int prv )
{
    vis[u] = 1;
    for (auto& n : graph[u])
    {
        if (n == prv) continue;
        if (!vis[n])
            dfs(n, u);
        else
            isCycle = 1;
    }
}

// Check Cycle in directed
Graph graph ;
vector<int>  vis;
bool isCycle = 0;
void dfs(int u )
{
    vis[u] = 2;  // mark node that exsit in recurion stack
    for (auto& n : graph[u])
    {
        if (!vis[n])  // if node not vis & not exist inf recursion stack
            dfs(n);
        else if (vis[n] == 2) // node is visited and exsit in recursion stack
            isCycle = 1;
    }

    vis[u] = 1;  // mark node that visited and not exsit in recursion stack

}
\end{lstlisting}
\subsubsection{Degree Sequences}
\begin{lstlisting}[language=C++]
#include <bits/stdc++.h>
using namespace std;

/**
 *  THE ULTIMATE CP TEMPLATE BLUEPRINT: DEGREE SEQUENCES
 *
 * 1.  WHAT PROBLEM DOES THIS SOLVE?
 * - Keywords: "Given an array of degrees", "Can we form a simple graph", "Construct a graph".
 * - Classic Scenarios: You are given N integers representing the number of friends each person has.
 *   You need to verify if this network is mathematically possible, and if so, pair them up.
 *
 * 2.  HOW TO USE IT (THE BLACK BOX)
 * - Validation Only (Fast):
 *       bool possible = DegreeSequence::erdos_gallai(degrees);
 *       // Complexity: O(N log N). Use this when N = 10^5 and you only need YES/NO.
 *
 * - Construction (Slower, but gives edges):
 *       vector<pair<int, int>> edges = DegreeSequence::havel_hakimi(degrees);
 *       // Complexity: O(N^2 log N). Use this when you must PRINT the graph.
 *       // Returns an empty vector if construction is impossible.
 *
 * 3.  HOW TO ADAPT IT (THE DIALS & SWITCHES)
 * - 0-based vs 1-based indexing for Construction:
 *   The `havel_hakimi` function returns edges using 1-based indexing by default.
 *   If you need 0-based indexing, change `nodes.push_back({d[i], i + 1});` to `i`.
 */

struct DegreeSequence
{

    static bool erdos_gallai(vector<long long> d)
    {
        int n = d.size();
        long long sum = 0;

        // A valid graph MUST have an even sum of degrees (Handshaking Lemma)
        for (long long x : d)
            sum += x;
        if (sum % 2 != 0)
            return false;

        // Sort degrees in descending order: d_1 >= d_2 >= ... >= d_n
        sort(d.begin(), d.end(), greater<long long>());

        // Build 1-indexed prefix sums for O(1) range queries
        vector<long long> pref(n + 1, 0);
        for (int i = 0; i < n; i++)
        {
            pref[i + 1] = pref[i] + d[i];
        }

        // Evaluate the Erds-Gallai inequality for every k from 1 to n:
        // Sum(i=1 to k) d_i <= k(k-1) + Sum(i=k+1 to n) min(d_i, k)
        for (int k = 1; k <= n; k++)
        {
            long long left_sum = pref[k];

            // Binary search to find the split point in the suffix where d_i drops below k.
            // This avoids an O(N) loop and keeps the overall complexity at O(N log N).
            int p = max(k, (int)(lower_bound(d.begin(), d.end(), k, greater<long long>()) - d.begin()));

            long long right_sum = 1LL * k * (k - 1);

            // Elements from k to p-1 are >= k, so min(d_i, k) is always exactly k
            right_sum += 1LL * k * (p - k);

            // Elements from p to n-1 are < k, so min(d_i, k) is just d_i
            right_sum += (pref[n] - pref[p]);

            if (left_sum > right_sum)
            {
                return false; // Mathematical contradiction found!
            }
        }

        return true;
    }

    static vector<pair<int, int>> havel_hakimi(const vector<int> &d)
    {
        int n = d.size();
        long long sum = 0;
        for (int x : d)
            sum += x;
        if (sum % 2 != 0)
            return {}; // Fails Handshaking Lemma

        // Store {degree, original_node_id}
        vector<pair<int, int>> nodes;
        for (int i = 0; i < n; i++)
        {
            if (d[i] > 0)
            {
                nodes.push_back({d[i], i + 1}); // 1-based indexing for nodes
            }
        }

        vector<pair<int, int>> edges;

        while (!nodes.empty())
        {
            // Sort descending by degree
            sort(nodes.begin(), nodes.end(), greater<pair<int, int>>());

            // The node with the highest degree demands to be connected
            int current_degree = nodes[0].first;
            int u = nodes[0].second;

            // Remove it from the pool of available nodes
            nodes.erase(nodes.begin());

            // If it demands more connections than there are nodes left, it's impossible
            if (current_degree > nodes.size())
                return {};

            // Connect 'u' greedily to the next 'current_degree' highest nodes
            for (int i = 0; i < current_degree; i++)
            {
                edges.push_back({u, nodes[i].second});
                nodes[i].first--; // They used up one of their connections

                // If a degree drops below 0, it means we forced a connection that broke the graph
                if (nodes[i].first < 0)
                    return {};
            }

            // Purge fully satisfied nodes (degree == 0) to speed up sorting
            vector<pair<int, int>> remaining_nodes;
            for (auto &node : nodes)
            {
                if (node.first > 0)
                {
                    remaining_nodes.push_back(node);
                }
            }
            nodes = remaining_nodes;
        }

        return edges;
    }
};

void solve()
{
    int n;
    // Example: Read N, then N degrees
    if (!(cin >> n))
        return;

    vector<int> d_int(n);
    vector<long long> d_ll(n);
    for (int i = 0; i < n; i++)
    {
        cin >> d_int[i];
        d_ll[i] = d_int[i];
    }

    // 1. O(N log N) Check (Best for large N)
    if (DegreeSequence::erdos_gallai(d_ll))
    {
        cout << "YES\n";

        // 2. O(N^2 log N) Construction (Only run if you actually need the edges!)
        // In CF/ICPC, if N > 10^4, they usually don't ask you to print the edges.
        vector<pair<int, int>> edges = DegreeSequence::havel_hakimi(d_int);
        for (auto &edge : edges)
        {
            cout << edge.first << " " << edge.second << "\n";
        }
    }
    else
    {
        cout << "NO\n";
    }
}

int main()
{
    ios_base::sync_with_stdio(false);
    cin.tie(NULL);
    // solve();
    return 0;
}
\end{lstlisting}
\subsubsection{Tree hashing}
\begin{lstlisting}[language=C++]
#include <bits/stdc++.h>
using namespace std;

/**
 *  THE ULTIMATE CP TEMPLATE BLUEPRINT: TREE ISOMORPHISM (HASHING)
 *
 * Solves: Checks if two unrooted trees are structurally identical.
 * Time Complexity: O(N)
 * Space Complexity: O(N)
 */

struct TreeIsomorphism
{
    // Generate a random seed once per program execution to prevent anti-hash tests
    static uint64_t get_seed()
    {
        mt19937_64 rng(chrono::steady_clock::now().time_since_epoch().count());
        return rng();
    }

    // A highly robust non-linear permutation function (SplitMix64)
    static uint64_t shift_hash(uint64_t x)
    {
        static const uint64_t SEED = get_seed();
        x += SEED;
        x ^= (x >> 33);
        x *= 0xff51afd7ed558ccdULL;
        x ^= (x >> 33);
        x *= 0xc4ceb9fe1a85ec53ULL;
        x ^= (x >> 33);
        return x;
    }

    static vector<int> get_centers(int n, const vector<vector<int>> &adj)
    {
        vector<int> degree(n + 1);
        queue<int> q;

        for (int i = 1; i <= n; i++)
        {
            degree[i] = adj[i].size();
            // Leaves have degree 1 (or 0 if n=1)
            if (degree[i] <= 1)
            {
                q.push(i);
            }
        }

        int remaining_nodes = n;

        // Peel away the leaves layer by layer like an onion
        while (remaining_nodes > 2)
        {
            int layer_size = q.size();
            remaining_nodes -= layer_size;

            for (int i = 0; i < layer_size; i++)
            {
                int u = q.front();
                q.pop();

                for (int v : adj[u])
                {
                    if (--degree[v] == 1)
                    {
                        q.push(v);
                    }
                }
            }
        }

        // The remaining 1 or 2 nodes are the absolute centers
        vector<int> centers;
        while (!q.empty())
        {
            centers.push_back(q.front());
            q.pop();
        }
        return centers;
    }

    static uint64_t compute_hash(int u, int p, const vector<vector<int>> &adj)
    {
        uint64_t current_hash = 1; // Base value for a node

        for (int v : adj[u])
        {
            if (v != p)
            {
                // To combine children, we hash the child's subtree, run it through
                // our non-linear shift function, and sum them up.
                // Using SUM guarantees that the order of children does not matter!
                current_hash += shift_hash(compute_hash(v, u, adj));
            }
        }
        return current_hash;
    }

    static bool are_isomorphic(int n, const vector<vector<int>> &adj1, const vector<vector<int>> &adj2)
    {
        if (n == 1)
            return true; // Trivial case

        // 1. Find centers of both trees
        vector<int> centers1 = get_centers(n, adj1);
        vector<int> centers2 = get_centers(n, adj2);

        // If the number of centers doesn't match, they physically cannot be isomorphic
        if (centers1.size() != centers2.size())
        {
            return false;
        }

        // 2. Hash Tree 1 from its first center
        uint64_t hash1 = compute_hash(centers1[0], 0, adj1);

        // 3. Check Tree 2 against Tree 1's hash
        for (int c2 : centers2)
        {
            uint64_t hash2 = compute_hash(c2, 0, adj2);
            if (hash1 == hash2)
            {
                return true; // Match found!
            }
        }

        return false;
    }
};

void solve()
{
    int n;
    // Example: Read N, then N-1 edges for Tree 1, then N-1 edges for Tree 2
    if (!(cin >> n))
        return;

    vector<vector<int>> tree1(n + 1);
    for (int i = 0; i < n - 1; i++)
    {
        int u, v;
        cin >> u >> v;
        tree1[u].push_back(v);
        tree1[v].push_back(u);
    }

    vector<vector<int>> tree2(n + 1);
    for (int i = 0; i < n - 1; i++)
    {
        int u, v;
        cin >> u >> v;
        tree2[u].push_back(v);
        tree2[v].push_back(u);
    }

    if (TreeIsomorphism::are_isomorphic(n, tree1, tree2))
    {
        cout << "Isomorphic (They are the exact same tree)\n";
    }
    else
    {
        cout << "Not Isomorphic\n";
    }
}

int main()
{
    ios_base::sync_with_stdio(false);
    cin.tie(NULL);
    // solve();
    return 0;
}
\end{lstlisting}
\subsection{Shortest paths}
\subsubsection{BFS01}
\begin{lstlisting}[language=C++]
#include <bits/stdc++.h>
using namespace std;

void zero_one_bfs(int start, int n, vector<vector<pair<int, int>>> &adj, vector<int> &dist)
{
    dist.assign(n, 1e9);
    deque<int> dq;
    dist[start] = 0;
    dq.push_front(start);

    while (!dq.empty())
    {
        int u = dq.front();
        dq.pop_front();

        for (auto &edge : adj[u])
        {
            int v = edge.first;
            int w = edge.second;
            if (dist[v] > dist[u] + w)
            {
                dist[v] = dist[u] + w;
                if (w == 0)
                    dq.push_front(v);
                else
                    dq.push_back(v);
            }
        }
    }
}
\end{lstlisting}
\subsubsection{Bellman ford}
\begin{lstlisting}[language=C++]
#include <bits/stdc++.h>
using namespace std;

const long long INF = 1e18;

struct Edge
{
    int u, v, w;
};

bool bellman_ford(int n, int start, vector<Edge> &edges, vector<long long> &dist)
{
    dist.assign(n, INF);
    dist[start] = 0;

    for (int i = 0; i < n - 1; ++i)
    {
        for (auto &e : edges)
        {
            if (dist[e.u] != INF && dist[e.v] > dist[e.u] + e.w)
                dist[e.v] = dist[e.u] + e.w;
        }
    }

    for (auto &e : edges)
    {
        if (dist[e.u] != INF && dist[e.v] > dist[e.u] + e.w)
            return true;
    }
    return false;
}
\end{lstlisting}
\subsubsection{Dijkstra}
\begin{lstlisting}[language=C++]
ll n;
GRAPH graph;

vector<ll> dijkstra(ll x)
{
    vector<ll> dist(n + 1, OO);
    dist[x] = 0;
    priority_queue<pair<ll, ll>, vector<pair<ll, ll>>, greater<pair<ll, ll>>> pq;
    pq.push({0, x});
    while (!pq.empty())
    {
        auto [l, node] = pq.top();
        pq.pop();
        if (l > dist[node])
            continue;
        for (auto &[to, w] : graph[node])
        {
            if (dist[to] > dist[node] + w)
            {
                dist[to] = dist[node] + w;
                pq.push({dist[to], to});
            }
        }
    }
    return dist;
}


// Path
ll n;
GRAPH graph;
vector<int> parent;
deque<int> path;

vector<ll> dijkstra(ll x)
{
    parent.assign(n + 1, -1);
    vector<ll> dist(n + 1, OO);
    dist[x] = 0;
    priority_queue<pair<ll, ll>, vector<pair<ll, ll>>, greater<pair<ll, ll>>> pq;
    pq.push({0, x});
    while (!pq.empty())
    {
        auto [l, node] = pq.top();
        pq.pop();
        if (l > dist[node])
            continue;
        for (auto &[to, w] : graph[node])
        {
            if (dist[to] > dist[node] + w)
            {
                dist[to] = dist[node] + w;
                parent[to] = node;
                pq.push({dist[to], to});
            }
        }
    }
    return dist;
}

void build(int node)
{
    path.push_front(node);
    if (parent[node] == -1)
        return;
    build(parent[node]);
}

// Destination
ll n;
GRAPH graph;

vector<ll> dijkstra(ll x, ll d)
{
    vector<ll> dist(n + 1, OO);
    dist[x] = 0;
    priority_queue<pair<ll, ll>, vector<pair<ll, ll>>, greater<pair<ll, ll>>> pq;
    pq.push({0, x});
    while (!pq.empty())
    {
        auto [l, node] = pq.top();
        pq.pop();
        if (node == d)
            break;
        if (l > dist[node])
            continue;
        for (auto &[to, w] : graph[node])
        {
            if (dist[to] > dist[node] + w)
            {
                dist[to] = dist[node] + w;
                pq.push({dist[to], to});
            }
        }
    }
    return dist;
}


// Dijkstra indexed PQ
#include <bits/stdc++.h>

#include <ext/pb_ds/priority_queue.hpp>
using namespace std;

using namespace __gnu_pbds;
typedef long long ll;
const ll OO = 1e18;
typedef vector<vector<pair<ll, ll>>> GRAPH;
typedef __gnu_pbds::priority_queue<pair<ll, ll>, greater<pair<ll, ll>>, pairing_heap_tag> indexedPQ;

ll n;
GRAPH graph;

vector<ll> dijkstra(ll x)
{
    vector<ll> dist(n + 1, OO);
    indexedPQ pq;
    vector<indexedPQ::point_iterator> key(n + 1, nullptr);
    dist[x] = 0;
    key[x] = pq.push({0, x});
    while (!pq.empty())
    {
        auto [l, node] = pq.top();
        pq.pop();

        for (auto &[to, w] : graph[node])
        {
            if (dist[to] > dist[node] + w)
            {
                dist[to] = dist[node] + w;
                if (key[to] == nullptr)
                    key[to] = pq.push({dist[to], to});
                else
                    pq.modify(key[to], {dist[to], to});
            }
        }
    }
    return dist;
}

// Dijkstra remove
ll n;
GRAPH graph;

vector<ll> dijkstra(ll x)
{
    vector<ll> dist(n + 1, OO);
    dist[x] = 0;
    set<pair<ll, ll>> s;
    s.insert({0, x});
    while (!s.empty())
    {
        auto [l, node] = *s.begin();
        s.erase(s.begin());
        if (l > dist[node])
            continue;
        for (auto &[to, w] : graph[node])
        {
            if (dist[to] > dist[node] + w)
            {
                dist[to] = dist[node] + w;
                s.insert({dist[to], to});
            }
        }
    }
    return dist;
}

\end{lstlisting}
\subsubsection{Floyd}
\begin{lstlisting}[language=C++]
int n;
GRAPH graph;

vector<vector<int>> floyd()
{
    vector<vector<int>> dist = graph;
    for (int k = 0; k < n; k++)
    {
        for (int i = 0; i < n; i++)
        {
            for (int j = 0; j < n; j++)
            {
                if (dist[i][k] < OO && dist[k][j] < OO)
                    dist[i][j] = min(dist[i][j], dist[i][k] + dist[k][j]);
            }
        }
    }
    return dist;
}

// real
int n;
GRAPH graph;

vector<vector<double>> floyd()
{
    vector<vector<double>> dist = graph;
    for (int k = 0; k < n; k++)
    {
        for (int i = 0; i < n; i++)
        {
            for (int j = 0; j < n; j++)
            {
                if (dist[i][k] < OO && dist[k][j] < OO)
                {
                    if (dist[i][k] + dist[k][j] < dist[i][j] - DBL_EPSILON)
                        dist[i][j] = min(dist[i][j], dist[i][k] + dist[k][j]);
                }
            }
        }
    }
    return dist;
}

// path
int n;
GRAPH graph;
vector<int> path;
vector<vector<int>> phase;

vector<vector<int>> floyd()
{
    phase.assign(n + 1, vector<int>(n + 1, -1));
    vector<vector<int>> dist = graph;
    for (int k = 0; k < n; k++)
    {
        for (int i = 0; i < n; i++)
        {
            for (int j = 0; j < n; j++)
            {
                if (dist[i][k] < OO && dist[k][j] < OO)
                {
                    if (dist[i][j] > dist[i][k] + dist[k][j])
                    {
                        dist[i][j] = min(dist[i][j], dist[i][k] + dist[k][j]);
                        phase[i][j] = k;
                    }
                }
            }
        }
    }
    return dist;
}

// push the start node manually
void build(int i, int j)
{
    if (phase[i][j] == -1)
    {
        path.push_back(j);
        return;
    }
    int k = phase[i][j];
    build(i, k);
    build(k, j);
}

\end{lstlisting}
\subsubsection{Floyd Warshall(Hassan)}
\begin{lstlisting}[language=C++]
void hassan() {
    int n,m;cin >>n>>m;
    vector<vector<int>> dist(n+1, vector<int>(n+1, oo)); // oo >= n * maxCost

    for (int i=0;i<m;i++) {
        int u,v,w;cin>>u>>v>>w;
        dist[u][v] = min(dist[u][v], w);
    }
    for (int i=1;i<=n;i++) {
        dist[i][i] = 0;
    }

    for (int k = 1; k <= n; k++) {
        for (int from = 1; from <= n; from++) {
            for (int to = 1; to <= n; to++) {
                if (dist[from][k] < oo && dist[k][to] < oo) {
                    dist[from][to] = min(dist[from][to], dist[from][k] + dist[k][to]);
                }
            }
        }

    for (int i=1;i<=n;i++) {
        for (int j=1;j<=n;j++) {
            cout << "from " << i << " to " << j << " Dist " <<dist[i][j] << "\n";
        }
    }
}

void reachability(int n, int m) {
    vector<vector<bool>> canReach(n+1, vector<bool>(n+1, false));
    // ... input u, v ... canReach[u][v] = true;
    for (int i=1; i<=n; i++) canReach[i][i] = true;

    for (int k=1; k<=n; k++) {
        for (int i=1; i<=n; i++) {
            for (int j=1; j<=n; j++) {     
                canReach[i][j] = canReach[i][j] || (canReach[i][k] && canReach[k][j]);
            }
        }
    }
}

///  Find Minimax Or Maxmini Path ///////
void hassan() {
    int n,m;cin >>n>>m;
    vector<vector<int>> dist(n+1, vector<int>(n+1, oo)); // oo >= n * maxCost
    for (int i=1;i<=n;i++) {
        dist[i][i] = 0;
     }
    for (int i=0;i<m;i++) {
        int u,v,w;cin>>u>>v>>w;
        dist[u][v] = min(dist[u][v], w);
    }
  
    for (int k = 1; k <= n; k++) {
        for (int from = 1; from <= n; from++) {
            for (int to = 1; to <= n; to++) {
                if (dist[from][k] < oo && dist[k][to] < oo) {
                    dist[from][to] = min(dist[from][to], max(dist[from][k],dist[k][to]));
                    //dist[from][to] = max(dist[from][to], min(dist[from][k],dist[k][to]));
                }
            }
        }
        for (int i=1;i<=n;i++) {
            for (int j=1;j<=n;j++) {
                cout << "from " << i << " to " << j << " Dist " <<dist[i][j] << "\n";
            }
        }
    }
}

///// Counting Shortest Paths ///////
void hassan_counting() {
    int n, m; cin >> n >> m;
    long long oo = 1e15;

    vector<vector<long long>> dist(n + 1, vector<long long>(n + 1, oo));
    vector<vector<long long>> pathsCount(n + 1, vector<long long>(n + 1, 0));

    for (int i = 1; i <= n; i++) {
        dist[i][i] = 0;
        pathsCount[i][i] = 1;   
    }

    for (int i = 0; i < m; i++) {
        int u, v, w; cin >> u >> v >> w;
        if (w < dist[u][v]) {
            dist[u][v] = w;
            pathsCount[u][v] = 1;
        } else if (w == dist[u][v]) {
            pathsCount[u][v]++; 
        }
    }

    for (int k = 1; k <= n; k++) {
        for (int i = 1; i <= n; i++) {
            for (int j = 1; j <= n; j++) {
                if (dist[i][k] + dist[k][j] < dist[i][j]) {
                    dist[i][j] = dist[i][k] + dist[k][j];
                    pathsCount[i][j] = pathsCount[i][k] * pathsCount[k][j];
                } 
                else if (dist[i][k] + dist[k][j] == dist[i][j] && i != k && k != j) {
                    pathsCount[i][j] += (pathsCount[i][k] * pathsCount[k][j]);
                }
            }
        }
    }

    for (int i = 1; i <= n; i++) {
        for (int j = 1; j <= n; j++) {
            if (dist[i][j] == oo) continue;
            cout << "From " << i << " to " << j << ": Min Dist = " << dist[i][j] 
                 << " | Ways = " << pathsCount[i][j] << "\n";
        }
    }
}


/////// Diameter and Radius  /////////
int diameter = 0;
int radius = oo;

for (int i=1; i<=n; i++) {
    int max_dist_for_i = 0;
    for (int j=1; j<=n; j++) {
        if (dist[i][j] != oo) 
            max_dist_for_i = max(max_dist_for_i, dist[i][j]);
    }
    // Diameter:  " "   
    diameter = max(diameter, max_dist_for_i);
    // Radius:  " "     
    radius = min(radius, max_dist_for_i);
}

\end{lstlisting}
\subsubsection{MaxMin Cost Using Dijkstra}
\begin{lstlisting}[language=C++]
/*
* Problem Description:
* The goal is to find a path from the source node (Node 1) to all other nodes in the graph
* such that the minimum edge weight along this path is maximized (Maximize the minimum weight).
*/
struct node {
    int u,w;

    bool operator<(const node &b) const {
        return w > b.w;
    }
};

void hassan() {
    int n, m;
    cin >> n >> m;
    vector<vector<pair<int, int> > > adj(n + 1);
    for (int i = 1; i <= m; i++) {
        int u, v, w;
        cin >> u >> v >> w;
        adj[u].push_back({v, w});
        adj[v].push_back({u, w});
    }
    vector<int> maxMin(n + 1, -1);
    priority_queue<node> pq;
    pq.push({1, oo});
    maxMin[1] = oo;
    while (!pq.empty()) {
        auto [u,w] = pq.top();
        pq.pop();

        if (w < maxMin[u]) continue;

        for (auto [v,c]: adj[u]) {
            if (min(maxMin[u], c) > maxMin[v]) {
                maxMin[v] = min(maxMin[u], c);
                pq.push({v, maxMin[v]});
            }
        }
    }
    for (int i = 1; i <= n; i++) {
        cout << maxMin[i] << " ";
    }
}
\end{lstlisting}
\subsubsection{MinMax Cost Using Dijkstra}
\begin{lstlisting}[language=C++]
/*
 * Problem Description:
 * The goal is to find a path from the source node (Node 1) to all other nodes in the graph
 * such that the Maximum edge weight along this path is Minimized (Minimize the maximum weight).
*/
struct node {
    int u, w;

    bool operator<(const node &b) const {
        return w > b.w;
    }
};

void hassan() {
    int n, m;
    cin >> n >> m;
    vector<vector<pair<int, int> > > adj(n + 1);
    for (int i = 1; i <= m; i++) {
        int u, v, w;
        cin >> u >> v >> w;
        adj[u].push_back({v, w});
        adj[v].push_back({u, w});
    }
    priority_queue<node> pq;
    pq.push({1, 0});
    vector<int> minMax(n + 1, oo);
    minMax[1] = 0;
    while (!pq.empty()) {
        auto [u,w] = pq.top();
        pq.pop();

        if (w > minMax[u]) continue;

        for (auto [v,c]: adj[u]) {
            if (max(minMax[u], c) < minMax[v]) {
                minMax[v] = max(minMax[u], c);
                pq.push({v, minMax[v]});
            }
        }
    }
    for (int i = 1; i <= n; i++) {
        cout << minMax[i] << " ";
    }
}
\end{lstlisting}
\subsection{Trees}
\subsubsection{Rerooting}
\begin{lstlisting}[language=C++]
/**
 *  THE ULTIMATE CP TEMPLATE BLUEPRINT: TREE REROOTING DP (IN-OUT DP)
 *
 * 1.  WHAT PROBLEM DOES THIS SOLVE?
 * - Keywords: "Sum of distances to all nodes", "Tree Rerooting", "Tree DP", "Calculate for EVERY node".
 * - Classic Scenarios: You are given a tree and asked to find the sum of distances from node U to 
 *   ALL other nodes, but you need to output this for EVERY node U in the tree. A naive BFS/DFS 
 *   from every single node takes O(N^2), which gives TLE.
 * - The Magic: "Rerooting Technique". We solve it in two linear passes:
 *   1. Bottom-Up (dfs): Root the tree arbitrarily at node 1. Calculate the subtree size (`sz`) 
 *      and the sum of distances from node 1 to all nodes in its subtree (`dist`).
 *   2. Top-Down (dfs2): Shift the root from parent `u` to child `v` in O(1) time! 
 *      When we move the root from `u` to `v`, all nodes in `v`'s subtree become 1 step CLOSER 
 *      (so we subtract `sz[v]`), and all other nodes in the tree become 1 step FURTHER 
 *      (so we add `n - sz[v]`).
 *
 * 2.  HOW TO USE IT (THE BLACK BOX)
 * - Initialization: 
 *       Ensure `GRAPH` is defined (e.g., `typedef vector<vector<int>> GRAPH;`).
 *       The `solve()` function automatically initializes all vectors to size `N + 1`.
 *
 * - Execution: 
 *       1. Build the undirected graph.
 *       2. Call `dfs(1, 0)` to precalculate subtree sizes and the answer for node 1.
 *       3. Set the base case: `ans[1] = dist[1]`.
 *       4. Call `dfs2(1, 0)` to magically propagate the answer to all other nodes.
 *
 * - Complexity:
 *       Time: O(N)  Just two simple DFS traversals.
 *       Space: O(N) for the tree, sizes, and distances.
 *
 * 3.  HOW TO ADAPT IT (THE DIALS & SWITCHES)
 * - Weighted Edges?
 *   If edges have weights `W`, change `graph` to store pairs `{neighbor, weight}`.
 *   In `dfs1`: `dist[u] += dist[v] + sz[v] * W`.
 *   In `dfs2`: `ans[v] = ans[u] - sz[v] * W + (n - sz[v]) * W`.
 * - Max Distance instead of Sum?
 *   Rerooting can also find the farthest node from every node! However, the transition is harder: 
 *   you must keep track of the MAXIMUM and SECOND MAXIMUM depths in `dfs1` so that when you reroot 
 *   into the branch that contained the maximum, you can fall back to the second maximum.
 */
int n;
vector<int> sz, dist, ans;
GRAPH graph;

void dfs(int u, int p)
{
    sz[u] = 1;
    for (auto &v : graph[u])
    {
        if (v == p)
            continue;
        dfs(v, u);
        sz[u] += sz[v];
        dist[u] += dist[v] + sz[v];
    }
}

void dfs2(int u, int p)
{
    for (auto &v : graph[u])
    {
        if (v == p)
            continue;
        ans[v] = ans[u] - sz[v] + (n - sz[v]);
        dfs2(v, u);
    }
}

void solve()
{
    cin >> n;
    sz.assign(n + 1, {}), dist.assign(n + 1, {}), graph.assign(n + 1, {}), ans.assign(n + 1, {});
    for (int i = 1; i < n; i++)
    {
        int a, b;
        cin >> a >> b;
        graph[a].emplace_back(b), graph[b].emplace_back(a);
    }
    dfs(1, 0);
    ans[1] = dist[1];
    dfs2(1, 0);
    for (int i = 1; i <= n; i++)
        cout << ans[i] << " ";
}

\end{lstlisting}
\subsubsection{Virtual tree}
\begin{lstlisting}[language=C++]
#include <bits/stdc++.h>
using namespace std;

/**
 *  THE ULTIMATE CP TEMPLATE BLUEPRINT: VIRTUAL TREE (AUXILIARY TREE)
 *
 * 1.  WHAT PROBLEM DOES THIS SOLVE?
 * - Keywords: "Paths between subset nodes", "DP on subset of tree", "Sparse subset queries".
 * - Classic Scenarios: A tree has N=10^5 nodes. You have Q queries. Each query gives you a subset
 *   of K nodes. You need to perform a Tree DP or path calculation on these K nodes, but running
 *   an O(N) DFS per query is impossible.
 * - The Magic: The Virtual Tree builds a smaller tree containing only the K subset nodes
 *   AND their necessary LCAs to maintain tree structure. This new tree has O(K) nodes.
 *
 * 2.  HOW TO USE IT (THE BLACK BOX)
 * - Initialization:
 *       1. Build the full tree using standard DFS.
 *       2. For each query, collect the subset of K nodes.
 *       3. Call `build_virtual_tree(subset)`.
 *       4. Perform your DP on `virtual_adj` (the edges of the virtual tree).
 * - Complexity:
 *       Time: O(K log K) per query to build.
 *       Space: O(N) to store the full tree + O(K) for the virtual tree.
 */

const int MAXN = 2e5 + 5;
const int LOG = 20;

// adj: Original tree, virtual_adj: The sparse tree we build for the subset
vector<int> adj[MAXN], virtual_adj[MAXN];
int tin[MAXN], tout[MAXN], depth[MAXN];
int up[MAXN][LOG];
int timer;

// Standard DFS to precalculate depths, entry times, and binary lifting table
void dfs(int u, int p)
{
    tin[u] = ++timer;
    depth[u] = depth[p] + 1;
    up[u][0] = p;
    for (int i = 1; i < LOG; i++)
        up[u][i] = up[up[u][i - 1]][i - 1];

    for (int v : adj[u])
    {
        if (v != p)
            dfs(v, u);
    }
    tout[u] = timer;
}

// Find LCA using binary lifting
int get_lca(int u, int v)
{
    if (depth[u] < depth[v])
        swap(u, v);
    for (int i = LOG - 1; i >= 0; i--)
    {
        if (depth[u] - (1 << i) >= depth[v])
            u = up[u][i];
    }
    if (u == v)
        return u;
    for (int i = LOG - 1; i >= 0; i--)
    {
        if (up[u][i] != up[v][i])
        {
            u = up[u][i], v = up[v][i];
        }
    }
    return up[u][0];
}

// Ancestor check using Entry/Exit times
bool is_ancestor(int u, int v)
{
    return tin[u] <= tin[v] && tout[u] >= tout[v];
}

// Main function to build the virtual tree from a subset of nodes
void build_virtual_tree(vector<int> &nodes)
{
    // 1. Sort by entry time so we can process the tree in DFS order
    sort(nodes.begin(), nodes.end(), [](int a, int b)
         { return tin[a] < tin[b]; });

    // 2. Inject LCAs: To keep the tree connected, we must add the LCA of every
    // adjacent pair in our sorted list.
    int sz = nodes.size();
    for (int i = 0; i < sz - 1; i++)
    {
        nodes.push_back(get_lca(nodes[i], nodes[i + 1]));
    }

    // 3. Sort again by entry time and remove duplicates (LCAs might have been added multiple times)
    sort(nodes.begin(), nodes.end(), [](int a, int b)
         { return tin[a] < tin[b]; });
    nodes.erase(unique(nodes.begin(), nodes.end()), nodes.end());

    // 4. Connect the nodes using a Stack
    // We maintain a stack representing the current path from the root.
    vector<int> st;
    st.push_back(nodes[0]);

    for (size_t i = 1; i < nodes.size(); i++)
    {
        // Pop nodes that are not ancestors of the current node
        while (st.size() > 1 && !is_ancestor(st.back(), nodes[i]))
        {
            virtual_adj[st[st.size() - 2]].push_back(st.back());
            st.pop_back();
        }
        st.push_back(nodes[i]);
    }
    // Connect remaining nodes in the stack
    while (st.size() > 1)
    {
        virtual_adj[st[st.size() - 2]].push_back(st.back());
        st.pop_back();
    }
}

void solve()
{
    int n;
    cin >> n;

    // Clear graph before building for multiple test cases
    for (int i = 0; i <= n; ++i)
    {
        adj[i].clear();
        virtual_adj[i].clear();
    }

    // Example: Reading N-1 edges
    for (int i = 0; i < n - 1; ++i)
    {
        int u, v;
        cin >> u >> v;
        adj[u].push_back(v);
        adj[v].push_back(u);
    }

    // Precompute full tree properties
    timer = 0;
    dfs(1, 1);

    int k;
    cin >> k;
    vector<int> subset(k);
    for (int &x : subset)
        cin >> x;

    // Build the virtual tree for the K nodes
    build_virtual_tree(subset);

    // Now perform DP on `virtual_adj`!
    // subset now contains all nodes in the virtual tree.
}

int main()
{
    ios_base::sync_with_stdio(false);
    cin.tie(NULL);
    solve();
    return 0;
}
\end{lstlisting}
\section{Math}
\subsection{Calculus}
\subsubsection{Differentiation integration}
\begin{lstlisting}[language=C++]
/**
 *  THE ULTIMATE CP TEMPLATE BLUEPRINT: ADAPTIVE SIMPSON's RULE
 *
 * 1.  WHAT PROBLEM DOES THIS SOLVE?
 * - Keywords: "Area under curve", "Evaluate integral", "Continuous probability".
 * - Classic Scenarios: You are given a mathematical function f(x) and need the
 *   integral from A to B with a precision of 10^-6.
 * - The Magic: Instead of chopping the interval into 10^6 fixed pieces (which can
 *   TLE or lose precision), Adaptive Simpson's evaluates the error dynamically.
 *   It recursively splits the interval only in areas where the curve is highly
 *   irregular, guaranteeing extreme precision in logarithmic time.
 *
 * 2.  HOW TO USE IT (THE BLACK BOX)
 * - Query:
 *       double area = integrate(A, B, 1e-7, my_function);
 *
 * - Complexity:
 *       Time: O(Depends on curve complexity, practically instant).
 *       Space: O(Recursion depth).
 */

double simpson(double l, double r, double f_l, double f_r, double f_mid)
{
    return (f_l + 4 * f_mid + f_r) * (r - l) / 6.0;
}

double adaptive_simpson(double l, double r, double f_l, double f_r, double f_mid, double S, double eps, const function<double(double)> &f)
{
    double mid = (l + r) / 2.0;
    double l_mid = (l + mid) / 2.0;
    double r_mid = (mid + r) / 2.0;

    double f_l_mid = f(l_mid);
    double f_r_mid = f(r_mid);

    double left_S = simpson(l, mid, f_l, f_mid, f_l_mid);
    double right_S = simpson(mid, r, f_mid, f_r, f_r_mid);

    if (abs(left_S + right_S - S) <= 15.0 * eps)
    {
        return left_S + right_S + (left_S + right_S - S) / 15.0;
    }

    return adaptive_simpson(l, mid, f_l, f_mid, f_l_mid, left_S, eps / 2.0, f) +
           adaptive_simpson(mid, r, f_mid, f_r, f_r_mid, right_S, eps / 2.0, f);
}

double integrate(double a, double b, double eps, const function<double(double)> &f)
{
    double f_a = f(a);
    double f_b = f(b);
    double f_mid = f((a + b) / 2.0);
    double S = simpson(a, b, f_a, f_b, f_mid);
    return adaptive_simpson(a, b, f_a, f_b, f_mid, S, eps, f);
}
\end{lstlisting}
\subsection{Convolution}
\subsubsection{FWHT}
\begin{lstlisting}[language=C++]
/**
 *  THE ULTIMATE CP TEMPLATE BLUEPRINT: FAST WALSH-HADAMARD TRANSFORM
 *
 * 1.  WHAT PROBLEM DOES THIS SOLVE?
 * - Keywords: "Bitwise convolution", "XOR sum", "AND sum", "OR sum".
 * - Classic Scenarios: Find array C where C[k] = sum(A[i] * B[j]) for all
 *   pairs (i, j) such that (i XOR j) == k.
 * - The Magic: It alters the NTT butterfly structure to map binary bitwise
 *   operations into addition/subtraction, resolving the convolution in O(N log N).
 *
 * 2.  HOW TO USE IT (THE BLACK BOX)
 * - Query:
 *       vector<long long> C = fwht_xor(A, B);
 *       // Swap 'fwht_xor' for 'fwht_and' or 'fwht_or' as needed.
 *
 * - Complexity:
 *       Time: O(N log N). Note that N MUST be a power of 2.
 *       Space: O(N).
 *
 * 3.  HOW TO ADAPT IT (THE DIALS & SWITCHES)
 * - Padding: The input arrays MUST be padded with zeros until their size
 *   is exactly a power of 2 before passing them into these functions.
 */

using ll = long long;
const ll MOD = 998244353;

ll binpow_fwht(ll a, ll b)
{
    ll res = 1;
    a %= MOD;
    while (b > 0)
    {
        if (b & 1)
            res = res * a % MOD;
        a = a * a % MOD;
        b >>= 1;
    }
    return res;
}

ll modInv(ll n) { return binpow_fwht(n, MOD - 2); }

void fwht_xor_base(vector<ll> &a, bool inv)
{
    int n = a.size();
    for (int len = 1; 2 * len <= n; len <<= 1)
    {
        for (int i = 0; i < n; i += 2 * len)
        {
            for (int j = 0; j < len; j++)
            {
                ll u = a[i + j];
                ll v = a[i + len + j];
                a[i + j] = (u + v) % MOD;
                a[i + len + j] = (u - v + MOD) % MOD;
            }
        }
    }
    if (inv)
    {
        ll invN = modInv(n);
        for (ll &x : a)
            x = x * invN % MOD;
    }
}

void fwht_and_base(vector<ll> &a, bool inv)
{
    int n = a.size();
    for (int len = 1; 2 * len <= n; len <<= 1)
    {
        for (int i = 0; i < n; i += 2 * len)
        {
            for (int j = 0; j < len; j++)
            {
                ll u = a[i + j];
                ll v = a[i + len + j];
                if (!inv)
                {
                    a[i + j] = (u + v) % MOD;
                }
                else
                {
                    a[i + j] = (u - v + MOD) % MOD;
                }
            }
        }
    }
}

void fwht_or_base(vector<ll> &a, bool inv)
{
    int n = a.size();
    for (int len = 1; 2 * len <= n; len <<= 1)
    {
        for (int i = 0; i < n; i += 2 * len)
        {
            for (int j = 0; j < len; j++)
            {
                ll u = a[i + j];
                ll v = a[i + len + j];
                if (!inv)
                {
                    a[i + len + j] = (u + v) % MOD;
                }
                else
                {
                    a[i + len + j] = (v - u + MOD) % MOD;
                }
            }
        }
    }
}

vector<ll> fwht_xor(vector<ll> a, vector<ll> b)
{
    fwht_xor_base(a, false);
    fwht_xor_base(b, false);
    for (size_t i = 0; i < a.size(); i++)
        a[i] = a[i] * b[i] % MOD;
    fwht_xor_base(a, true);
    return a;
}

vector<ll> fwht_and(vector<ll> a, vector<ll> b)
{
    fwht_and_base(a, false);
    fwht_and_base(b, false);
    for (size_t i = 0; i < a.size(); i++)
        a[i] = a[i] * b[i] % MOD;
    fwht_and_base(a, true);
    return a;
}

vector<ll> fwht_or(vector<ll> a, vector<ll> b)
{
    fwht_or_base(a, false);
    fwht_or_base(b, false);
    for (size_t i = 0; i < a.size(); i++)
        a[i] = a[i] * b[i] % MOD;
    fwht_or_base(a, true);
    return a;
}
\end{lstlisting}
\subsubsection{GDC LCM Convolution}
\begin{lstlisting}[language=C++]
/**
 *  THE ULTIMATE CP TEMPLATE BLUEPRINT: GCD & LCM CONVOLUTION
 *
 * 1.  WHAT PROBLEM DOES THIS SOLVE?
 * - Keywords: "gcd(i, j) = k", "lcm(i, j) = k", "Divisor transformation".
 * - Classic Scenarios: Find array C where C[k] = sum(A[i] * B[j]) over all pairs
 *   (i, j) such that gcd(i, j) == k (or lcm(i, j) == k).
 * - The Magic: Bypasses polynomial rings entirely. It transforms the arrays into
 *   prefix sums over their multiples (for GCD) or divisors (for LCM), multiplies
 *   the point values, and reverses the transformation using Mobius-style subtraction.
 *
 * 2.  HOW TO USE IT (THE BLACK BOX)
 * - Query:
 *       vector<long long> C_gcd = gcd_convolution(A, B);
 *       vector<long long> C_lcm = lcm_convolution(A, B);
 *
 * - Complexity:
 *       Time: O(N log log N) utilizing a prime sieve optimization.
 *       Space: O(N).
 *
 * 3.  HOW TO ADAPT IT (THE DIALS & SWITCHES)
 * - 1-Based Indexing: These arrays MUST be 1-indexed (i.e., A[1] is the first
 *   element, A[0] is ignored). Ensure A and B are the same size before passing.
 */

using ll = long long;
const ll MOD = 998244353;

vector<int> get_primes(int n)
{
    vector<bool> is_prime(n + 1, true);
    vector<int> primes;
    for (int p = 2; p <= n; p++)
    {
        if (is_prime[p])
        {
            primes.push_back(p);
            for (int i = p * 2; i <= n; i += p)
                is_prime[i] = false;
        }
    }
    return primes;
}

vector<ll> gcd_convolution(vector<ll> a, vector<ll> b)
{
    int n = a.size() - 1;
    vector<int> primes = get_primes(n);

    for (int p : primes)
    {
        for (int i = n / p; i > 0; i--)
        {
            a[i] = (a[i] + a[i * p]) % MOD;
            b[i] = (b[i] + b[i * p]) % MOD;
        }
    }

    vector<ll> c(n + 1);
    for (int i = 1; i <= n; i++)
    {
        c[i] = a[i] * b[i] % MOD;
    }

    for (int p : primes)
    {
        for (int i = 1; i * p <= n; i++)
        {
            c[i] = (c[i] - c[i * p] + MOD) % MOD;
        }
    }
    return c;
}

vector<ll> lcm_convolution(vector<ll> a, vector<ll> b)
{
    int n = a.size() - 1;
    vector<int> primes = get_primes(n);

    for (int p : primes)
    {
        for (int i = 1; i * p <= n; i++)
        {
            a[i * p] = (a[i * p] + a[i]) % MOD;
            b[i * p] = (b[i * p] + b[i]) % MOD;
        }
    }

    vector<ll> c(n + 1);
    for (int i = 1; i <= n; i++)
    {
        c[i] = a[i] * b[i] % MOD;
    }

    for (int p : primes)
    {
        for (int i = n / p; i > 0; i--)
        {
            c[i * p] = (c[i * p] - c[i] + MOD) % MOD;
        }
    }
    return c;
}
\end{lstlisting}
\subsubsection{NTT}
\begin{lstlisting}[language=C++]
/**
 *  THE ULTIMATE CP TEMPLATE BLUEPRINT: NUMBER THEORETIC TRANSFORM (NTT)
 *
 * 1.  WHAT PROBLEM DOES THIS SOLVE?
 * - Keywords: "Polynomial multiplication", "Convolution", "Sum of pairs".
 * - Classic Scenarios: You are given two arrays A and B and need to find C
 *   where C[k] = sum(A[i] * B[k-i]) modulo 998244353.
 * - The Magic: A naive double loop takes O(N^2) and TLEs for N = 10^5.
 *   NTT evaluates the polynomials at primitive roots of unity, multiplies
 *   the point values in O(N), and interpolates them back in O(N log N).
 *
 * 2.  HOW TO USE IT (THE BLACK BOX)
 * - Query:
 *       vector<long long> C = multiply(A, B);
 *
 * - Complexity:
 *       Time: O(N log N)
 *       Space: O(N)
 *
 * 3.  HOW TO ADAPT IT (THE DIALS & SWITCHES)
 * - Modulo: This template is strictly hardcoded for MOD = 998244353 and
 *   its primitive root 3. If the problem has a different modulo, you MUST
 *   use an Arbitrary Modulo NTT (MTT) or find the primitive root of the new modulo.
 */

using ll = long long;
const ll MOD = 998244353;
const ll ROOT = 3;

ll binpow(ll a, ll b)
{
    ll res = 1;
    a %= MOD;
    while (b > 0)
    {
        if (b & 1)
            res = res * a % MOD;
        a = a * a % MOD;
        b >>= 1;
    }
    return res;
}

ll modInverse(ll n)
{
    return binpow(n, MOD - 2);
}

void ntt(vector<ll> &a, bool invert)
{
    int n = a.size();
    for (int i = 1, j = 0; i < n; i++)
    {
        int bit = n >> 1;
        for (; j & bit; bit >>= 1)
            j ^= bit;
        j ^= bit;
        if (i < j)
            swap(a[i], a[j]);
    }
    for (int len = 2; len <= n; len <<= 1)
    {
        ll wlen = binpow(ROOT, (MOD - 1) / len);
        if (invert)
            wlen = modInverse(wlen);
        for (int i = 0; i < n; i += len)
        {
            ll w = 1;
            for (int j = 0; j < len / 2; j++)
            {
                ll u = a[i + j];
                ll v = a[i + j + len / 2] * w % MOD;
                a[i + j] = (u + v < MOD ? u + v : u + v - MOD);
                a[i + j + len / 2] = (u - v >= 0 ? u - v : u - v + MOD);
                w = w * wlen % MOD;
            }
        }
    }
    if (invert)
    {
        ll n_inv = modInverse(n);
        for (ll &x : a)
            x = x * n_inv % MOD;
    }
}

vector<ll> multiply(vector<ll> const &a, vector<ll> const &b)
{
    vector<ll> fa(a.begin(), a.end()), fb(b.begin(), b.end());
    int n = 1;
    while (n < a.size() + b.size())
        n <<= 1;
    fa.resize(n);
    fb.resize(n);
    ntt(fa, false);
    ntt(fb, false);
    for (int i = 0; i < n; i++)
        fa[i] = fa[i] * fb[i] % MOD;
    ntt(fa, true);
    while (fa.size() > 1 && fa.back() == 0)
        fa.pop_back();
    return fa;
}
\end{lstlisting}
\subsection{FFT}
\subsubsection{bitIntMultiply}
\begin{lstlisting}[language=C++]
string mul_two_big_int(const string &s1, const string &s2) {
    int n = s1.size(), m = s2.size();

    vector<int> poly1(n), poly2(m);
    for (int i = 0; i < n; ++i) {
        poly1[n-i-1] = s1[i] - '0';
    }

    for (int i = 0; i < m; ++i) {
        poly2[m-i-1] = s2[i] - '0';
    }

    vector<int> ans = multiply(poly1, poly2);
    int k = ans.size();

    for (int i = 0; i < k - 1; ++i) {
        ans[i + 1] += ans[i] / 10;
        ans[i] = ans[i] % 10;
    }

    string final = to_string(ans[k - 1]);
    for (int i = k - 2; i >= 0; --i) {
        final += (char)(ans[i] + '0');
    }

    for (int i = 0; i < k; ++i) {
        if(final[i] != '0') return final.substr(i);
    }
    return "0";
}

string power_of_big_int(string s, int p) {
    string ans = "1";
    while (p) {
        if(p&1) ans = mul_two_big_int(ans, s);
        s = mul_two_big_int(s, s);
        p >>= 1;
    }
    return ans;
}
\end{lstlisting}
\subsubsection{fft}
\begin{lstlisting}[language=C++]
// Multiplies two polynomials (or large numbers) in O(N log N) time
// Uses the Fast Fourier Transform (FFT) to convert coefficients to point-values and back

#include <bits/stdc++.h>
using namespace std;
#define int long long

using cd = complex<double>;
const double PI = acos(-1);

// Fast Fourier Transform: 
// invert = false converts coefficients to point-values (Forward FFT)
// invert = true converts point-values back to coefficients (Inverse FFT)
void fft(vector<cd> &a, bool invert) {
    int n = a.size();

    // Bit-reversal permutation: Rearranges elements so we can build the FFT bottom-up iteratively
    for (int i = 1, j = 0; i < n; i++) {
        int bit = n >> 1;
        for (; j & bit; bit >>= 1)
            j ^= bit;
        j ^= bit;
        if (i < j)
            swap(a[i], a[j]);
    }

    // Iterative FFT: Merge smaller subproblems into larger ones
    for (int len = 2; len <= n; len <<= 1) {
        double ang = 2 * PI / len * (invert ? -1 : 1);
        cd wlen(cos(ang), sin(ang));
        for (int i = 0; i < n; i += len) {
            cd w(1.0, 0.0);
            for (int j = 0; j < len / 2; j++) {
                // Butterfly operation
                cd u = a[i + j], v = a[i + j + len / 2] * w;
                a[i + j] = u + v;
                a[i + j + len / 2] = u - v;
                w *= wlen;
            }
        }
    }

    // If doing Inverse FFT, divide by N to complete the interpolation
    if (invert) {
        for (cd &x: a)
            x /= n;
    }
}

// Multiply two polynomials using FFT
vector<int> multiply(const vector<int> &a, const vector<int> &b) {
    vector<cd> fa(a.begin(), a.end()), fb(b.begin(), b.end());
    
    // Pad the polynomial sizes to the nearest power of 2 (required for Radix-2 FFT)
    int n = 1;
    while (n < (int) (a.size() + b.size()))
        n <<= 1;
    fa.resize(n);
    fb.resize(n);

    // 1. Evaluate polynomials at complex roots of unity
    fft(fa, false);
    fft(fb, false);
    
    // 2. Point-wise multiplication in O(N) time
    for (int i = 0; i < n; i++)
        fa[i] *= fb[i];
        
    // 3. Interpolate the results back into polynomial coefficients
    fft(fa, true);

    vector<int> result(n);
    for (int i = 0; i < n; i++)
        // Round to nearest integer to fix minor floating-point inaccuracies
        result[i] = round(fa[i].real());

    // Optional: trim trailing zeros
    while (!result.empty() && result.back() == 0)
        result.pop_back();

    return result;
}

// Optional: test usage
void solve() {
    int n;
    cin >> n;

    // Read coefficients from degree 0 up to degree n
    vector<int> a(n + 1), b(n + 1);
    for (int i = 0; i <= n; i++) cin >> a[i];
    for (int i = 0; i <= n; i++) cin >> b[i];

    // Get the multiplied polynomial
    vector<int> c = multiply(a, b);
    
    // The resulting polynomial has degree 2N, so we print 2N + 1 coefficients
    for (int i = 0; i <= 2*n; i++)
        // Handle cases where trailing zeros were trimmed
        cout << (i < c.size() ? c[i] : 0) << " ";
    cout << endl;
}

/*
If
a.size() = A (i.e. degree = A - 1)
b.size() = B (i.e. degree = B - 1)

Then:
 result.size() = A + B - 1 (before trimming)
 */

signed main() {
    ios::sync_with_stdio(false);
    cin.tie(nullptr);
    int t = 1;
    cin >> t;
    while (t--) solve();
    return 0;
}
\end{lstlisting}
\subsubsection{fftMod}
\begin{lstlisting}[language=C++]
// Arbitrary Modulo Polynomial Convolution (MTT) in O(N log N) time
// Multiplies two polynomials modulo ANY number (does not require NTT primes)

#include <bits/stdc++.h>
using namespace std;

#define rep(aa, bb, cc) for(int aa = bb; aa < cc;aa++)
#define sz(a) (int)a.size()
typedef long long ll;
typedef vector<int> vi;
typedef complex<double> C;
typedef vector<double> vd;

// Highly optimized Fast Fourier Transform
void fft(vector<C>& a) {
    int n = sz(a), L = 31 - __builtin_clz(n);
    static vector<complex<long double>> R(2, 1);
    static vector<C> rt(2, 1);  // Caches roots of unity for speed
    
    // Precompute roots of unity
    for (static int k = 2; k < n; k *= 2) {
        R.resize(n); rt.resize(n);
        auto x = polar(1.0L, acos(-1.0L) / k);
        rep(i,k,2*k) rt[i] = R[i] = i&1 ? R[i/2] * x : R[i/2];
    }
    
    // Bit-reversal permutation
    vi rev(n);
    rep(i,0,n) rev[i] = (rev[i / 2] | (i & 1) << L) / 2;
    rep(i,0,n) if (i < rev[i]) swap(a[i], a[rev[i]]);
    
    // Iterative FFT
    for (int k = 1; k < n; k *= 2)
        for (int i = 0; i < n; i += 2 * k) rep(j,0,k) {
                // Hand-rolled complex multiplication for a ~25% speed boost
                auto x = (double *)&rt[j+k], y = (double *)&a[i+j+k];        
                C z(x[0]*y[0] - x[1]*y[1], x[0]*y[1] + x[1]*y[0]);           
                a[i + j + k] = a[i + j] - z;
                a[i + j] += z;
            }
}
 
// Arbitrary Modulo Convolution
template<int M> vi convMod(const vi &a, const vi &b) {
    if (a.empty() || b.empty()) return {};
    vi res(sz(a) + sz(b) - 1);
    
    // Pad to nearest power of 2
    int B=32-__builtin_clz(sz(res)), n=1<<B;
    
    // Split values into chunks of roughly sqrt(M) to prevent floating-point overflow
    int cut=int(sqrt(M)); 
    vector<C> L(n), R(n), outs(n), outl(n);
    
    // Encode polynomial A as (A_div + i * A_mod) and B as (B_div + i * B_mod)
    rep(i,0,sz(a)) L[i] = C((int)a[i] / cut, (int)a[i] % cut);
    rep(i,0,sz(b)) R[i] = C((int)b[i] / cut, (int)b[i] % cut);
    fft(L), fft(R);
    
    // Clever math trick to extract the transforms of the real/imaginary parts using conjugates
    rep(i,0,n) {
        int j = -i & (n - 1); // Equivalent to (n - i) % n
        outl[j] = (L[i] + conj(L[j])) * R[i] / (2.0 * n);
        outs[j] = (L[i] - conj(L[j])) * R[i] / (2.0 * n) / 1i;
    }
    
    // Inverse FFT to get the cross-multiplications
    fft(outl), fft(outs);
    
    // Recombine the pieces modulo M
    rep(i,0,sz(res)) {
        ll av = ll(real(outl[i])+.5), cv = ll(imag(outs[i])+.5);
        ll bv = ll(imag(outl[i])+.5) + ll(real(outs[i])+.5);
        res[i] = ((av % M * cut + bv) % M * cut + cv) % M;
    }
    return res;
}

// 
const int MOD = 1e9 + 7;
vector<int> a = {1, 2, 3};
vector<int> b = {4, 5, 6};

vector<int> result = convMod<MOD>(a, b);
\end{lstlisting}
\subsubsection{ntt}
\begin{lstlisting}[language=C++]



// This template implements the Number Theoretic Transform (NTT),
//  which is used to multiply two polynomials in $O(N \log N)$ time
//  modulo a specific prime. It also includes the building blocks for
//  Arbitrary Modulo
//  Convolution (MTT) using the Chinese Remainder Theorem (CRT).

#include <bits/stdc++.h>
using namespace std;
typedef long long ll;

// Standard NTT-friendly prime (998244353) and its primitive root
const ll mod = (119 << 23) + 1, root = 62; 
// For p < 2^30 there is also e.g. 5 << 25, 7 << 26, 479 << 21
// and 483 << 21 (same root). The last two are > 10^9.

// Computes (b^e) % mod in O(log e) time
ll modpow(ll b, ll e) {
    ll ans = 1;
    for (; e; b = b * b % mod, e /= 2)
        if (e & 1) ans = ans * b % mod;
    return ans;
}

// Finds a primitive root for the given modulo dynamically.
// This is useful if you change the 'mod' variable and don't know its root.
int generator () {
    vector<int> fact;
    int phi = mod-1,  n = phi;
    
    // Prime factorize phi
    for (int i=2; i*i<=n; ++i)
        if (n % i == 0) {
            fact.push_back (i);
            while (n % i == 0) n /= i;
        }
    if (n > 1) fact.push_back (n);

    // Test candidates to find the primitive root
    for (int res=2; res<=mod; ++res) {
        bool ok = true;
        for (size_t i=0; i<fact.size() && ok; ++i)
            ok &= modpow (res, phi / fact[i]) != 1;
        if (ok)  return res;
    }
    return -1;
}

// Computes (b^e) % m with a custom modulo 'm'
int modpow(int b, int e, int m) {
    int ans = 1;
    for (; e; b = (ll)b * b % m, e /= 2)
        if (e & 1) ans = (ll)ans * b % m;
    return ans;
}

// The core Number Theoretic Transform (NTT) algorithm
// Converts polynomial coefficients to point-values and vice versa
void ntt(vector<int> &a) {
    int n = (int)a.size(), L = 31 - __builtin_clz(n);
    vector<int> rt(2, 1); 
    
    // Precompute roots of unity
    for (int k = 2, s = 2; k < n; k *= 2, s++) { 
        rt.resize(n);
        int z[] = {1, modpow(root, mod >> s, mod)};
        for (int i = k; i < 2*k; ++i) rt[i] = (ll)rt[i / 2] * z[i & 1] % mod;
    }
    
    // Bit-reversal permutation
    vector<int> rev(n);
    for (int i = 0; i < n; ++i) rev[i] = (rev[i / 2] | (i & 1) << L) / 2;
    for (int i = 0; i < n; ++i) if (i < rev[i]) swap(a[i], a[rev[i]]);
    
    // Iterative butterfly operations
    for (int k = 1; k < n; k *= 2) {
        for (int i = 0; i < n; i += 2 * k) {
            for (int j = 0; j < k; ++j) {
                int z = (ll)rt[j + k] * a[i + j + k] % mod, &ai = a[i + j];
                a[i + j + k] = ai - z + (z > ai ? mod : 0);
                ai += (ai + z >= mod ? z - mod : z);
            }
        }
    }
}

// Multiplies two polynomials (a and b) modulo 'mod'
vector<int> conv(const vector<int> &a, const vector<int> &b) {
    if (a.empty() || b.empty()) return {};
    int s = (int)a.size() + (int)b.size() - 1, B = 32 - __builtin_clz(s), n = 1 << B;
    int inv = modpow(n, mod - 2, mod);
    vector<int> L(a), R(b), out(n);
    L.resize(n), R.resize(n);
    
    ntt(L), ntt(R); // Convert to point-values
    for (int i = 0; i < n; ++i) out[-i & (n - 1)] = (ll)L[i] * R[i] % mod * inv % mod; // Multiply
    ntt(out); // Convert back to coefficients
    
    return {out.begin(), out.begin() + s};
}

// Chinese Remainder Theorem for 2 variables
// Solves: x = a (mod m1) and x = b (mod m2)
ll CRT(ll a, ll m1, ll b, ll m2) {
    __int128 m = m1*m2;
    ll ans = a*m2%m*modpow(m2, m1-2, m1)%m + m1*b%m*modpow(m1, m2-2, m2)%m;
    return ans % m;
}

/*
// --- ARBITRARY MODULO MULTIPLICATION SETUP ---
// This commented section uses 3 specific NTT-friendly primes.
// By computing the polynomial multiplication separately for mod1, mod2, and mod3,
// you can use the 3-variable CRT function below to find the exact answer modulo ANYTHING (like 10^9+7).

int mod, root, desired_mod = 1000000007;
const int mod1 = 167772161;
const int mod2 = 469762049;
const int mod3 = 754974721;
const int root1 = 3;
const int root2 = 3;
const int root3 = 11;

// Chinese Remainder Theorem for 3 variables
int CRT(int a, int b, int c, int m1, int m2, int m3) {
    __int128 M = (__int128)m1*m2*m3;
    ll M1 = (ll)m2*m3;
    ll M2 = (ll)m1*m3;
    ll M3 = (ll)m2*m1;

    int M_1 = modpow(M1%m1, m1 - 2, m1);
    int M_2 = modpow(M2%m2, m2 - 2, m2);
    int M_3 = modpow(M3%m3, m3 - 2, m3);

    __int128 ans = (__int128)a*M1*M_1;
    ans += (__int128)b*M2*M_2;
    ans += (__int128)c*M3*M_3;

    return (ans % M) % desired_mod;
}
*/
\end{lstlisting}
\subsubsection{pow}
\begin{lstlisting}[language=C++]
vector<int> poly_pow(vector<int> poly, int p) {
    vector<int> ans{1};
    while (p) {
        if(p&1) ans = multiply(ans, poly);
        poly = multiply(poly, poly);
        p >>= 1;
    }
    return ans;
}
\end{lstlisting}
\subsection{Linear algebra}
\subsubsection{Determinant(composite mod)}
\begin{lstlisting}[language=C++]
/**
 *  DETERMINANT WITH COMPOSITE MODULO
 *
 * WHAT PROBLEM DOES THIS SOLVE?
 * Standard matrix determinants require division (modular inverse). If the modulo
 * is composite, the modular inverse might not exist. This algorithm bypasses
 * division entirely by using Euclidean-like steps (similar to GCD) to eliminate
 * row elements.
 *
 * Complexity: O(N^3 log M)
 */
long long det_composite(vector<vector<long long>> a, long long mod)
{
    int n = a.size();
    long long det = 1;

    for (int i = 0; i < n; i++)
    {
        for (int j = i + 1; j < n; j++)
        {
            // Use Euclidean-like steps to eliminate a[j][i] without using fractions
            while (a[j][i] != 0)
            {
                long long t = a[i][i] / a[j][i];

                for (int k = i; k < n; k++)
                {
                    // Subtract a multiple of row j from row i
                    a[i][k] = (a[i][k] - t * a[j][k]) % mod;

                    // Swap row i and row j to keep the algorithm progressing
                    swap(a[i][k], a[j][k]);
                }
                // Every time we swap two rows, the sign of the determinant flips
                det = -det;
            }
        }

        // If a diagonal element becomes 0, the matrix is singular (determinant is 0)
        if (a[i][i] == 0)
            return 0;

        // Multiply the diagonal elements to build the final determinant
        det = det * a[i][i] % mod;
    }

    // Ensure the final determinant is positive before returning
    return (det % mod + mod) % mod;
}
\end{lstlisting}
\subsubsection{Determinant(prime mod)}
\begin{lstlisting}[language=C++]
/**
 *  THE ULTIMATE CP TEMPLATE BLUEPRINT: DETERMINANT (PRIME MODULO)
 *
 * 1.  WHAT PROBLEM DOES THIS SOLVE?
 * - Keywords: "Matrix Tree Theorem", "Number of Spanning Trees".
 * - Classic Scenarios: You have an N x N matrix and need its determinant modulo
 *   a prime number like 1e9+7 or 998244353.
 * - The Magic: Standard Gaussian elimination. It tracks row swaps because every
 *   swap multiplies the final determinant by -1.
 *
 * 2.  HOW TO USE IT (THE BLACK BOX)
 * - Query:
 *       long long ans = det_prime(matrix, MOD);
 *
 * - Complexity:
 *       Time: O(N^3)
 *       Space: O(N^2)
 */

using ll = long long;

ll binpow_det(ll a, ll b, ll mod)
{
    ll res = 1;
    a %= mod;
    while (b > 0)
    {
        if (b & 1)
            res = res * a % mod;
        a = a * a % mod;
        b >>= 1;
    }
    return res;
}

ll modInverse_det(ll n, ll mod)
{
    return binpow_det(n, mod - 2, mod);
}

ll det_prime(vector<vector<ll>> a, ll mod)
{
    int n = a.size();
    ll ans = 1;
    for (int i = 0; i < n; ++i)
    {
        int pivot = i;
        for (int j = i + 1; j < n; ++j)
        {
            if (a[j][i] > a[pivot][i])
                pivot = j;
        }
        if (a[pivot][i] == 0)
            return 0;
        if (i != pivot)
        {
            swap(a[i], a[pivot]);
            ans = (mod - ans) % mod;
        }
        ans = ans * a[i][i] % mod;
        ll inv = modInverse_det(a[i][i], mod);
        for (int j = i + 1; j < n; ++j)
        {
            ll factor = a[j][i] * inv % mod;
            for (int k = i; k < n; ++k)
            {
                a[j][k] = (a[j][k] - factor * a[i][k]) % mod;
                if (a[j][k] < 0)
                    a[j][k] += mod;
            }
        }
    }
    return ans;
}
\end{lstlisting}
\subsubsection{Gaussian elemination}
\begin{lstlisting}[language=C++]
/**
 *  THE ULTIMATE CP TEMPLATE BLUEPRINT: GAUSSIAN ELIMINATION
 *
 * 1.  WHAT PROBLEM DOES THIS SOLVE?
 * - Keywords: "System of equations", "Probabilities", "Markov Chains".
 * - Classic Scenarios: Solving Ax = B for floating-point values.
 * - The Magic: It reduces the matrix to row-echelon form to find the exact
 *   values of the variables.
 *
 * 2.  HOW TO USE IT (THE BLACK BOX)
 * - Query:
 *       vector<double> ans;
 *       int solutions = gauss(matrix, ans);
 * - Result Breakdown:
 *       Returns 1 if there is exactly one solution (stored in 'ans').
 *       Returns 0 if there are no solutions.
 *       Returns 2 if there are infinitely many solutions.
 *
 * - Complexity:
 *       Time: O(N^3)
 *       Space: O(N^2)
 */

const double EPS = 1e-9;

int gauss(vector<vector<double>> a, vector<double> &ans)
{
    int n = a.size();
    if (n == 0)
        return 0;
    int m = a[0].size() - 1;
    vector<int> where(m, -1);

    for (int col = 0, row = 0; col < m && row < n; ++col)
    {
        int sel = row;
        for (int i = row; i < n; ++i)
        {
            if (abs(a[i][col]) > abs(a[sel][col]))
                sel = i;
        }
        if (abs(a[sel][col]) < EPS)
            continue;

        for (int i = col; i <= m; ++i)
            swap(a[sel][i], a[row][i]);
        where[col] = row;

        for (int i = 0; i < n; ++i)
        {
            if (i != row)
            {
                double c = a[i][col] / a[row][col];
                for (int j = col; j <= m; ++j)
                {
                    a[i][j] -= a[row][j] * c;
                }
            }
        }
        ++row;
    }

    ans.assign(m, 0);
    for (int i = 0; i < m; ++i)
    {
        if (where[i] != -1)
            ans[i] = a[where[i]][m] / a[where[i]][i];
    }

    for (int i = 0; i < n; ++i)
    {
        double sum = 0;
        for (int j = 0; j < m; ++j)
        {
            sum += ans[j] * a[i][j];
        }
        if (abs(sum - a[i][m]) > EPS)
            return 0;
    }

    for (int i = 0; i < m; ++i)
    {
        if (where[i] == -1)
            return 2;
    }
    return 1;
}
\end{lstlisting}
\subsubsection{Gaussian elemination mod2}
\begin{lstlisting}[language=C++]
/**
 *  THE ULTIMATE CP TEMPLATE BLUEPRINT: GAUSSIAN ELIMINATION MOD 2
 *
 * 1.  WHAT PROBLEM DOES THIS SOLVE?
 * - Keywords: "XOR subset", "Boolean equations", "Lights out game".
 * - Classic Scenarios: Finding if a specific XOR sum can be formed from an array,
 *   or solving Ax = B where all values and operations are modulo 2 (XOR).
 * - The Magic: Replaces standard arrays with `std::bitset`, shrinking the time
 *   constant by a factor of 64. Instead of division and subtraction, it simply
 *   uses XOR.
 *
 * 2.  HOW TO USE IT (THE BLACK BOX)
 * - Query:
 *       bitset<2005> ans;
 *       int solutions = gauss_mod2(matrix, N, M, ans);
 *
 * - Complexity:
 *       Time: O(N^3 / 64)
 *       Space: O(N^2 / 64)
 */

const int MAX_V = 2005;

int gauss_mod2(vector<bitset<MAX_V>> a, int n, int m, bitset<MAX_V> &ans)
{
    vector<int> where(m, -1);
    for (int col = 0, row = 0; col < m && row < n; ++col)
    {
        for (int i = row; i < n; ++i)
        {
            if (a[i][col])
            {
                swap(a[i], a[row]);
                break;
            }
        }
        if (!a[row][col])
            continue;
        where[col] = row;
        for (int i = 0; i < n; ++i)
        {
            if (i != row && a[i][col])
            {
                a[i] ^= a[row];
            }
        }
        ++row;
    }

    ans.reset();
    for (int i = 0; i < m; ++i)
    {
        if (where[i] != -1)
            ans[i] = a[where[i]][m];
    }

    for (int i = 0; i < n; ++i)
    {
        int sum = 0;
        for (int j = 0; j < m; ++j)
        {
            if (ans[j] && a[i][j])
                sum ^= 1;
        }
        if (sum != a[i][m])
            return 0;
    }

    for (int i = 0; i < m; ++i)
    {
        if (where[i] == -1)
            return 2;
    }
    return 1;
}
\end{lstlisting}
\subsubsection{XOR basis}
\begin{lstlisting}[language=C++]
/**
 *  THE ULTIMATE CP TEMPLATE BLUEPRINT: XOR BASIS (LINEAR BASIS)
 *
 * 1.  WHAT PROBLEM DOES THIS SOLVE?
 * - Keywords: "Maximum XOR subset", "K-th XOR sum", "Can we form X".
 * - Classic Scenarios: You are given an array of size 10^5 and need to find the
 *   maximum possible XOR sum of any subset.
 * - The Magic: It reduces the entire array into a set of at most B "basis" vectors
 *   (where B is the number of bits, usually 60 for long long). These basis vectors
 *   span the exact same XOR space as the original array.
 *
 * 2.  HOW TO USE IT (THE BLACK BOX)
 * - Initialization:
 *       XORBasis basis;
 *       for (ll x : arr) basis.insert(x);
 * - Queries:
 *       long long max_val = basis.max_xor();
 *       bool exists = basis.contains(X);
 *
 * 3.  HOW TO ADAPT IT (THE DIALS & SWITCHES)
 * - Reduced Row Echelon Form (RREF) / K-th Smallest: To answer "K-th smallest XOR sum"
 *   (Topic 383/384), you MUST call `basis.build_rref()` first. Then use `basis.kth(K)`.
 *   (Note: K is 1-indexed. If K=1, it returns the smallest possible XOR sum > 0).
 */

using ll = long long;

struct XORBasis
{
    static const int BITS = 60; // Max bits for long long
    vector<ll> basis;
    vector<ll> rref; // For K-th queries
    int sz;
    bool is_rref_built;

    XORBasis()
    {
        basis.assign(BITS, 0);
        sz = 0;
        is_rref_built = false;
    }

    bool insert(ll x)
    {
        for (int i = BITS - 1; i >= 0; i--)
        {
            if ((x >> i) & 1)
            {
                if (!basis[i])
                {
                    basis[i] = x;
                    sz++;
                    is_rref_built = false;
                    return true;
                }
                x ^= basis[i];
            }
        }
        return false;
    }

    bool contains(ll x)
    {
        for (int i = BITS - 1; i >= 0; i--)
        {
            if ((x >> i) & 1)
            {
                if (!basis[i])
                    return false;
                x ^= basis[i];
            }
        }
        return true;
    }

    ll max_xor(ll current_xor = 0)
    {
        ll res = current_xor;
        for (int i = BITS - 1; i >= 0; i--)
        {
            res = max(res, res ^ basis[i]);
        }
        return res;
    }

    // Call this strictly BEFORE querying kth()
    void build_rref()
    {
        if (is_rref_built)
            return;
        rref.clear();
        vector<ll> temp = basis;
        for (int i = BITS - 1; i >= 0; i--)
        {
            if (!temp[i])
                continue;
            for (int j = i + 1; j < BITS; j++)
            {
                if ((temp[j] >> i) & 1)
                    temp[j] ^= temp[i];
            }
        }
        for (int i = 0; i < BITS; i++)
        {
            if (temp[i])
                rref.push_back(temp[i]);
        }
        is_rref_built = true;
    }

    // Returns K-th smallest XOR sum (1-indexed, excluding 0)
    // Returns -1 if K is out of bounds
    ll kth(ll k)
    {
        build_rref();
        if (k > (1LL << sz) - 1)
            return -1;
        ll res = 0;
        for (int i = 0; i < rref.size(); i++)
        {
            if ((k >> i) & 1)
                res ^= rref[i];
        }
        return res;
    }
};
\end{lstlisting}
\subsection{Linear recurrence}
\subsubsection{Berlekamp}
\begin{lstlisting}[language=C++]
/**
 *  THE ULTIMATE CP TEMPLATE BLUEPRINT: BERLEKAMP-MASSEY
 *
 * 1.  WHAT PROBLEM DOES THIS SOLVE?
 * - Keywords: "Find the recurrence", "Sequence guesser".
 * - Classic Scenarios: You brute-forced the first 50 terms of a DP or counting
 *   problem, but you cannot figure out the actual mathematical recurrence relation.
 * - The Magic: It iteratively builds a polynomial that perfectly models the
 *   input sequence. If the true recurrence has length N, giving BM the first
 *   2N terms guarantees it will perfectly reconstruct the formula.
 *
 * 2.  HOW TO USE IT (THE BLACK BOX)
 * - Query:
 *       vector<ll> C = berlekamp_massey(S, MOD);
 *       // S is the sequence you brute-forced.
 *       // C contains the coefficients: S[i] = C[0]*S[i-1] + C[1]*S[i-2] + ...
 *
 * - Complexity:
 *       Time: O(|S|^2)
 *       Space: O(|S|)
 */

using ll = long long;

ll binpow_bm(ll a, ll b, ll mod)
{
    ll res = 1;
    a %= mod;
    while (b > 0)
    {
        if (b & 1)
            res = res * a % mod;
        a = a * a % mod;
        b >>= 1;
    }
    return res;
}

ll modInverse_bm(ll n, ll mod)
{
    return binpow_bm(n, mod - 2, mod);
}

vector<ll> berlekamp_massey(const vector<ll> &s, ll mod)
{
    vector<ll> c = {1}, b = {1};
    ll L = 0, m = 1, b_val = 1;
    for (int i = 0; i < s.size(); i++)
    {
        ll d = 0;
        for (int j = 0; j <= L; j++)
        {
            d = (d + c[j] * s[i - j]) % mod;
        }
        if (d == 0)
        {
            m++;
        }
        else
        {
            vector<ll> temp = c;
            ll c_val = (d * modInverse_bm(b_val, mod)) % mod;
            while (c.size() <= b.size() + m)
                c.push_back(0);
            for (int j = 0; j < b.size(); j++)
            {
                c[j + m] = (c[j + m] - c_val * b[j] % mod + mod) % mod;
            }
            if (2 * L <= i)
            {
                L = i + 1 - L;
                b = temp;
                b_val = d;
                m = 1;
            }
            else
            {
                m++;
            }
        }
    }
    c.erase(c.begin());
    for (ll &x : c)
        x = (mod - x) % mod;
    return c;
}
\end{lstlisting}
\subsection{Matrix}
\subsubsection{Frievald's algorithm}
\begin{lstlisting}[language=C++]
/**
 *  THE ULTIMATE CP TEMPLATE BLUEPRINT: FREIVALD'S ALGORITHM
 *
 * 1.  WHAT PROBLEM DOES THIS SOLVE?
 * - Keywords: "Verify Matrix Multiplication", "A * B = C".
 * - Classic Scenarios: You are given three massive N x N matrices (A, B, and C).
 *   You need to determine if A * B equals C without TLEing.
 * - The Magic: A full multiplication takes O(N^3). Freivald's generates a random
 *   1D vector (r) of 0s and 1s, and checks if A * (B * r) == C * r.
 *   Since matrix-vector multiplication is O(N^2), this is drastically faster.
 *   Repeating it K times drops the probability of a false positive to (1/2)^K.
 *
 * 2.  HOW TO USE IT (THE BLACK BOX)
 * - Query:
 *       bool is_correct = freivalds(A, B, C, 15); // 15 iterations is perfectly safe
 *
 * - Complexity:
 *       Time: O(K * N^2)
 *       Space: O(N)
 */

#include <random>
#include <chrono>

using ll = long long;
const ll MOD = 1e9 + 7;

vector<ll> multiply_mat_vec(const vector<vector<ll>> &mat, const vector<ll> &vec)
{
    int n = mat.size();
    vector<ll> res(n, 0);
    for (int i = 0; i < n; i++)
    {
        for (int j = 0; j < n; j++)
        {
            res[i] = (res[i] + mat[i][j] * vec[j]) % MOD;
        }
    }
    return res;
}

bool freivalds(const vector<vector<ll>> &a, const vector<vector<ll>> &b, const vector<vector<ll>> &c, int k_iterations = 15)
{
    int n = a.size();
    mt19937 rng(chrono::steady_clock::now().time_since_epoch().count());
    uniform_int_distribution<int> dist(0, 1);

    for (int k = 0; k < k_iterations; k++)
    {
        vector<ll> r(n);
        for (int i = 0; i < n; i++)
        {
            r[i] = dist(rng);
        }

        vector<ll> br = multiply_mat_vec(b, r);
        vector<ll> abr = multiply_mat_vec(a, br);
        vector<ll> cr = multiply_mat_vec(c, r);

        for (int i = 0; i < n; i++)
        {
            if (abr[i] != cr[i])
            {
                return false;
            }
        }
    }
    return true;
}
\end{lstlisting}
\subsubsection{Matrix expo}
\begin{lstlisting}[language=C++]
typedef vector<vector<ll>> Matrix;

Matrix matrix_mul(Matrix &A, Matrix &B)
{
    Matrix C;
    ll n = ll(A.size());
    ll m = ll(A[0].size());
    ll l = ll(B[0].size());
    C.resize(n, vector<ll>(l));
    for (ll i = 0; i < n; i++)
    {
        for (ll k = 0; k < l; k++)
        {
            for (ll j = 0; j < m; j++)
            {
                C[i][k] += A[i][j] % MOD * B[j][k] % MOD;
                C[i][k] %= MOD;
            }
        }
    }
    return C;
}

Matrix matrix_expo(Matrix &matrix, ll pow)
{
    ll n = ll(matrix.size());
    Matrix res;
    res.resize(n, vector<ll>(n));
    for (int i = 0; i < n; i++)
        res[i][i] = 1;
    while (pow)
    {
        if (pow & 1)
            res = matrix_mul(matrix, res);
        matrix = matrix_mul(matrix, matrix);
        pow >>= 1;
    }
    return res;
}

Matrix I(int n)
{
    Matrix T(n, vector<int>(n));
    for (int i = 0; i < n; i++)
        T[i][i] = 1;
    return T;
}

// 2d Grid represent as graph
// count paths
short dx[8] = {1,1,-1,-1,2,2,-2,-2};
short dy[8] = {2,-2,2,-2,1,-1,1,-1};
void solve() {
    int k;cin >>k;
    Matrix T(65,Row(65));
    for (int i = 0;i<8;i++) {
        for (int j = 0;j<8;j++) {
            int cur = i*8 + j;
            for (int k = 0;k<8;k++) {
                short ni = i+dx[k];
                short nj = j+dy[k];
                if (ni < 0 || nj < 0 || ni > 7 || nj>7) continue;
                int nCur = ni*8+nj;
                T[nCur][cur] = 1;
            }
        }
    }
    for (int i = 0;i<65;i++) {
        T[i][64] = 1;
    }
    T = fastPower(T,k+1);
    int ans = T[0][64];
    cout << ans;
}
 

\end{lstlisting}
\subsubsection{Matrix expo Queries}
\begin{lstlisting}[language=C++]

 *  THE ULTIMATE CP TEMPLATE BLUEPRINT: FAST PATH COUNTING (PRECOMPUTED MATRIX EXPONENTIATION)
 * 1.  WHAT PROBLEM DOES THIS SOLVE?
 * - Keywords: "Number of paths of length K", "Multiple Queries (Q)", "State Vector multiplication".
 * - Classic Scenarios: You need to find the number of valid paths of exactly length K between 
 *   node U and node V in a directed graph, modulo some number. BUT, you have Q queries! 
 *   Standard Matrix Exponentiation per query takes O(Q * N^3 * log K), which will Time Limit Exceed (TLE).
 * - The Magic: "The Vector-Matrix Trick". We precompute the binary powers of the adjacency 
 *   matrix A^(2^0), A^(2^1), ..., A^(2^29) ONCE in O(N^3 * log K). 
 *   Then, for each query, we don't multiply N x N matrices together. Instead, we start with 
 *   a 1 x N state vector (where only the starting node U is 1) and multiply this VECTOR 
 *   by the precomputed N x N matrices. This slashes the query time down to just O(N^2 * log K)!
 *
 * 2.  HOW TO USE IT (THE BLACK BOX)
 * - Base State: Build your adjacency matrix `t`. If there is a directed edge u -> v, `t[u][v] = 1`.
 * - Precomputation: 
 *       pre[0] = t;
 *       pre[i] = mul(pre[i-1], pre[i-1]); // Computes A^(2^i)
 * - Queries: 
 *       Create a 1 x N matrix `s` (state vector). Set `s[0][u] = 
 *       For each set bit i in K, update `s = mul(s, pre[i])`.
 *       The answer is simply `s[0][v]`.
 *
 * - Complexity:
 *       Time: Precomputation O(N^3 * log K) + Queries O(Q * N^2 * log K).
 *       Space: O(log K * N^2) to store the 30 precomputed matrices.

const int mod = 1e9 + 7;
#define Row vector<int>
#define mat vector<Row>
mat mul(mat &a,mat &b) {
    int n = a.size(), m = b[0].size(),k = a[0].size();
    mat res(n,Row(m));
    for (int i = 0; i < n; i++) {
        for (int j = 0; j < m; j++) {
            for (int o = 0; o < k; o++) {
                res[i][j] += a[i][o]%mod * b[o][j]%mod;
                res[i][j]%=mod;
            }
        }
    }
    return res;
}

mat power(mat a,int b) {
    int n = a.size();
    mat res(n,Row(n));
    for (int i = 0; i < n; i++) res[i][i] = 1;
    while (b) {
        if (b&1) {
            res = mul(res,a);
        }
        b>>=1;
        a = mul(a,a);
    }
    return res;
}
void solve() {
    int n,m,q;cin >>n>>m>>q;
    mat t(n,Row(n));
    for (int i = 0; i < m; i++) {
        int u,v;cin >>u>>v;
        u--; v--;
        t[u][v] = 1;
    }
    vector<mat> pre(30,mat(n,Row(n)));
    pre[0] = t;
    for (int i = 1; i < 30; i++) {
        pre[i] = mul(pre[i-1],pre[i-1]);
    }
    while (q--) {
        int u,v,k;cin>>u>>v >>k;--u,--v;
        mat s(1,Row(n));
        s[0][u] = 1;
        for (int i = 0; i < 30; i++) {
            if (1 << i & k) {
                s = mul(s,pre[i]);
            }
        }
        cout << s[0][v] <<'\n';
    }
}

\end{lstlisting}
\subsubsection{Min Path UsingMatrix expo}
\begin{lstlisting}[language=C++]
*  THE ULTIMATE CP TEMPLATE BLUEPRINT: MIN-PLUS MATRIX EXPONENTIATION 
 * 1.  WHAT PROBLEM DOES THIS SOLVE?
 * - Keywords: "Shortest path of exactly K edges", "Paths of length K", "K is up to 10^9 or 10^18", "Min-Plus Algebra".
 * - Classic Scenarios: You are given a graph and need to find the minimum cost to travel between 
 *   nodes using EXACTLY (or at most) K edges. Since K is huge, standard algorithms like Dijkstra 
 *   or Bellman-Ford will Time Limit Exceed (TLE).
 * - The Magic: We treat the adjacency matrix as paths of length 1. If we multiply the matrix by itself, 
 *   we get paths of length 2. But instead of standard multiplication (+ and *), we use the 
 *   Min-Plus Semiring (min instead of +, and + instead of *). 
 *   Using binary exponentiation, we can compute A^K to find paths of length K in just O(N^3 log K)!
 *
 * 2.  HOW TO USE IT (THE BLACK BOX)
 * - Initialization: Create an N x N matrix filled with infinity (`oo`).
 *       Matrix t(n, Row(n, oo));
 * - Build Base State: Set the weight of directed edges. 
 *       t[u][v] = w; // Cost of 1 edge from u to v
 * - Execution: Call `fastPower` to raise the matrix to the power of K.
 *       t = fastPower(t, k);


#define Row vector<int>
#define Matrix vector<Row>
const int mod = 1e9 + 7;
const int oo = 2e18+5;

Matrix mul(Matrix &a,Matrix &b) {
    int n = a.size();
    int m = b[0].size();
    int sz = a[0].size();
    Matrix res(n,Row(m,oo));
    for (int i = 0; i < n; i++) {
        for (int j = 0; j < m; j++) {
            for (int o = 0; o < sz; o++) {
                res[i][j] = min(res[i][j],a[i][o] + b[o][j]);
            }
        }
    }
    return res;
}

Matrix fastPower(Matrix &a,int b) {
    int n = a.size();
 
    Matrix res(n, Row(n, oo));
    for (int i = 0; i < n; i++) {
        res[i][i] = 0; 
    }
    
    while (b) {
        if (b & 1)
            res = mul(res, a);
        a = mul(a, a);
        b /= 2;
    }
    return res;
}

void solve() {
    int n, m, k;
    cin >> n >> m >> k;
    Matrix t(n,Row(n,oo));
    for (int i = 0; i < m; i++) {
        int u, v, w;
        cin >> u >> v >> w;
        --u;
        --v;
        t[u][v] = w;
    }
    t = fastPower(t,k);

    int ans = oo;
    for (int i = 0; i < n; i++) {
        for (int j=0;j<n;j++) {
            ans = min(ans, t[i][j]);
        }
    }
    if (ans > 1e18) {
        cout << "IMPOSSIBLE\n";
    } else {
        cout << ans << "\n";
    }
}

\end{lstlisting}
\subsubsection{Permanent matrix}
\begin{lstlisting}[language=C++]
/**
 *  THE ULTIMATE CP TEMPLATE BLUEPRINT: PERMANENT OF A MATRIX (RYSER'S FORMULA)
 *
 * 1.  WHAT PROBLEM DOES THIS SOLVE?
 * - Keywords: "Number of perfect matchings", "Bipartite graph matchings".
 * - Classic Scenarios: Given an N x N adjacency matrix of a bipartite graph,
 *   find exactly how many distinct perfect matchings exist.
 * - The Magic: Uses Inclusion-Exclusion (Ryser's Formula) to bypass the
 *   factorial complexity of permutations, solving it in exactly O(N * 2^N).
 *
 * 2.  HOW TO USE IT (THE BLACK BOX)
 * - Query:
 *       long long ways = permanent(adj_matrix);
 *
 * - Complexity:
 *       Time: O(N * 2^N)  Generally safe for N <= 20.
 *       Space: O(N)
 */

using ll = long long;
const ll MOD = 1e9 + 7;

ll permanent(const vector<vector<ll>> &a)
{
    int n = a.size();
    if (n == 0)
        return 0;

    ll total_perm = 0;
    int max_mask = 1 << n;

    for (int mask = 1; mask < max_mask; mask++)
    {
        ll current_prod = 1;
        int set_bits = 0;

        for (int i = 0; i < n; i++)
        {
            ll row_sum = 0;
            int bit_count = 0;
            for (int j = 0; j < n; j++)
            {
                if ((mask >> j) & 1)
                {
                    row_sum = (row_sum + a[i][j]) % MOD;
                    if (i == 0)
                        bit_count++;
                }
            }
            if (i == 0)
                set_bits = bit_count;
            current_prod = (current_prod * row_sum) % MOD;
        }

        if ((n - set_bits) % 2 == 1)
        {
            total_perm = (total_perm - current_prod + MOD) % MOD;
        }
        else
        {
            total_perm = (total_perm + current_prod) % MOD;
        }
    }
    return total_perm;
}
\end{lstlisting}
\subsection{Optimization}
\subsubsection{Simplex}
\begin{lstlisting}[language=C++]
/**
 *  THE ULTIMATE CP TEMPLATE BLUEPRINT: SIMPLEX ALGORITHM
 *
 * 1.  WHAT PROBLEM DOES THIS SOLVE?
 * - Keywords: "Maximize profit", "Minimize cost", "Subject to constraints".
 * - Classic Scenarios: A generic resource allocation problem where you need to
 *   maximize an objective function C^T X subject to a system of inequalities AX <= B.
 * - The Magic: Bypasses the need for complex min-cost max-flow graph modeling
 *   or impossible greedy logic. You literally just pass the coefficients of your
 *   equations into the matrix.
 *
 * 2.  HOW TO USE IT (THE BLACK BOX)
 * - Query:
 *       vector<double> ans_vars;
 *       double max_profit = simplex(A, b, c, ans_vars);
 * - Result Breakdown:
 *       Returns numeric_limits<double>::infinity() if the problem is unbounded.
 *       Returns -numeric_limits<double>::infinity() if there is no feasible solution.
 *       Otherwise, returns the maximum possible value of the objective function,
 *       and populates `ans_vars` with the optimal values for variables X_1 to X_n.
 *
 * 3.  HOW TO ADAPT IT (THE DIALS & SWITCHES)
 * - Handling >= constraints: If an equation is 5x + 3y >= 10, multiply the entire
 *   row by -1 to make it -5x - 3y <= -10 before pushing it to matrix A and vector b.
 * - Handling == constraints: Split it into two inequalities:
 *   (5x + 3y <= 10) AND (-5x - 3y <= -10).
 */

const double EPS = 1e-9;
const double INF = numeric_limits<double>::infinity();

double simplex(vector<vector<double>> A, vector<double> b, vector<double> c, vector<double> &x)
{
    int m = b.size(), n = c.size();
    vector<int> N(n + 1), B(m);
    vector<vector<double>> D(m + 2, vector<double>(n + 2));

    for (int i = 0; i < m; i++)
    {
        for (int j = 0; j < n; j++)
            D[i][j] = A[i][j];
        B[i] = n + i;
        D[i][n] = -1;
        D[i][n + 1] = b[i];
    }
    for (int j = 0; j < n; j++)
    {
        N[j] = j;
        D[m][j] = -c[j];
    }
    N[n] = -1;
    D[m + 1][n] = 1;

    auto pivot = [&](int r, int s)
    {
        double inv = 1.0 / D[r][s];
        for (int i = 0; i < m + 2; i++)
        {
            if (i != r && abs(D[i][s]) > EPS)
            {
                double inv2 = D[i][s] * inv;
                for (int j = 0; j < n + 2; j++)
                {
                    D[i][j] -= D[r][j] * inv2;
                }
                D[i][s] = D[r][s] * inv2;
            }
        }
        for (int j = 0; j < n + 2; j++)
        {
            if (j != s)
                D[r][j] *= inv;
        }
        D[r][s] = inv;
        swap(B[r], N[s]);
    };

    auto phase = [&](int row)
    {
        while (true)
        {
            int s = -1;
            for (int j = 0; j <= n; j++)
            {
                if (N[j] != -1 && (s == -1 || D[row][j] < D[row][s] - EPS))
                {
                    s = j;
                }
            }
            if (D[row][s] > -EPS)
                return true;

            int r = -1;
            for (int i = 0; i < m; i++)
            {
                if (D[i][s] > EPS && (r == -1 || D[i][n + 1] / D[i][s] < D[r][n + 1] / D[r][s] - EPS))
                {
                    r = i;
                }
            }
            if (r == -1)
                return false;
            pivot(r, s);
        }
    };

    int r = 0;
    for (int i = 1; i < m; i++)
    {
        if (D[i][n + 1] < D[r][n + 1])
            r = i;
    }

    if (D[r][n + 1] < -EPS)
    {
        pivot(r, n);
        if (!phase(m + 1) || D[m + 1][n + 1] < -EPS)
            return -INF;
        for (int i = 0; i < m; i++)
        {
            if (B[i] == -1)
            {
                int s = -1;
                for (int j = 0; j <= n; j++)
                {
                    if (N[j] != -1 && (s == -1 || D[i][j] < D[i][s] - EPS))
                        s = j;
                }
                pivot(i, s);
            }
        }
    }

    if (!phase(m))
        return INF;

    x.assign(n, 0);
    for (int i = 0; i < m; i++)
    {
        if (B[i] < n)
            x[B[i]] = D[i][n + 1];
    }

    return D[m][n + 1];
}
\end{lstlisting}
\subsection{Polynomial}
\subsubsection{EXP POW SQRT}
\begin{lstlisting}[language=C++]
/**
 *  THE ULTIMATE CP TEMPLATE BLUEPRINT: POLYNOMIAL EXP, POW, & SQRT
 *
 * 1.  WHAT PROBLEM DOES THIS SOLVE?
 * - Keywords: "Exponentiate Generating Function", "e^{A(x)}", "A(x)^K", "sqrt(A(x))".
 * - Classic Scenarios: Solving combinatorial problems involving unlabelled sets
 *   or Euler transforms (Topic 335, 337, 338).
 * - The Magic: Polynomial Exp uses Newton's method on logarithms:
 *   F_{k+1} = F_k * (1 - ln(F_k) + A). Polynomial Pow uses the identity
 *   A^k = exp(k * ln(A)).
 *
 * 2.  HOW TO USE IT (THE BLACK BOX)
 * - Query:
 *       vector<ll> expA = poly_exp(A, N);
 *       vector<ll> powA = poly_pow(A, K, N);
 *       vector<ll> sqrtA = poly_sqrt(A, N);
 *
 * - Complexity:
 *       Time: O(N log N)
 *       Space: O(N)
 *
 * 3.  HOW TO ADAPT IT (THE DIALS & SWITCHES)
 * - Dependencies: Requires all functions from FPS Toolkit Part 1.
 * - Constant Term Strictness:
 *     - poly_exp: A[0] MUST be 0.
 *     - poly_pow: A[0] MUST be 1 (if not, factor out the lowest degree term first).
 *     - poly_sqrt: A[0] MUST be 1.
 */

vector<ll> poly_exp(vector<ll> a, int deg)
{
    if (deg == 1)
        return {1};
    vector<ll> b = poly_exp(a, (deg + 1) / 2);
    b.resize(deg);
    vector<ll> lnb = poly_log(b, deg);
    vector<ll> diff(deg);
    for (int i = 0; i < deg; i++)
    {
        diff[i] = (a.size() > i ? a[i] : 0) - lnb[i];
        if (diff[i] < 0)
            diff[i] += MOD;
    }
    diff[0] = (diff[0] + 1) % MOD;
    vector<ll> res = multiply(b, diff);
    res.resize(deg);
    return res;
}

vector<ll> poly_pow(vector<ll> a, ll k, int deg)
{
    vector<ll> lna = poly_log(a, deg);
    for (int i = 0; i < lna.size(); i++)
    {
        lna[i] = lna[i] * (k % MOD) % MOD;
    }
    return poly_exp(lna, deg);
}

vector<ll> poly_sqrt(vector<ll> a, int deg)
{
    if (deg == 1)
        return {1};
    vector<ll> b = poly_sqrt(a, (deg + 1) / 2);
    b.resize(deg);
    vector<ll> inv_b = poly_inv(b, deg);
    vector<ll> res = multiply(a, inv_b);
    ll inv2 = modInverse(2);
    res.resize(deg);
    for (int i = 0; i < deg; i++)
    {
        res[i] = (res[i] + b[i]) % MOD * inv2 % MOD;
    }
    return res;
}
\end{lstlisting}
\subsubsection{Lagrange}
\begin{lstlisting}[language=C++]
/**
 *  THE ULTIMATE CP TEMPLATE BLUEPRINT: LAGRANGE INTERPOLATION O(K)
 *
 * 1.  WHAT PROBLEM DOES THIS SOLVE?
 * - Keywords: "N-th term of a sequence", "Sum of i^p", "Polynomial fitting".
 * - Classic Scenarios: You know a sequence is generated by a polynomial of degree K.
 *   You need the exact value at a massive index X (e.g., X = 10^18), but you can
 *   only easily compute the first K+1 values (at x = 0, 1, 2, ..., K).
 * - The Magic: Standard Lagrange takes O(K^2). Because our evaluation points are
 *   strictly contiguous integers (0 to K), the products in the numerator and
 *   denominator become factorials and prefix/suffix products. This collapses the
 *   complexity to pure O(K), allowing you to evaluate degree 10^6 polynomials in
 *   less than a second.
 *
 * 2.  HOW TO USE IT (THE BLACK BOX)
 * - Query:
 *       long long ans = lagrange_eval(Y, X);
 *       // Y is a vector of size K+1 containing P(0), P(1), ..., P(K).
 *       // X is the massive point you want to evaluate at.
 *
 * - Complexity:
 *       Time: O(K) where K is the degree of the polynomial (size of Y - 1).
 *       Space: O(K) for prefix/suffix and factorial arrays.
 *
 * 3.  HOW TO ADAPT IT (THE DIALS & SWITCHES)
 * - Modulo: Adjust the global MOD variable if the problem uses 1e9+7 instead
 *   of 998244353.
 */

using ll = long long;
const ll MOD = 998244353;

ll binpow(ll base, ll exp)
{
    ll res = 1;
    base %= MOD;
    while (exp > 0)
    {
        if (exp % 2 == 1)
            res = (res * base) % MOD;
        base = (base * base) % MOD;
        exp /= 2;
    }
    return res;
}

ll modInverse(ll n)
{
    return binpow(n, MOD - 2);
}

ll lagrange_eval(const vector<ll> &y, ll x)
{
    int k = y.size() - 1;
    if (x <= k)
        return y[x];

    vector<ll> pref(k + 1, 1), suff(k + 1, 1);
    vector<ll> fact(k + 1, 1), invFact(k + 1, 1);

    for (int i = 1; i <= k; i++)
        fact[i] = fact[i - 1] * i % MOD;
    invFact[k] = modInverse(fact[k]);
    for (int i = k - 1; i >= 0; i--)
        invFact[i] = invFact[i + 1] * (i + 1) % MOD;

    pref[0] = x % MOD;
    for (int i = 1; i <= k; i++)
        pref[i] = pref[i - 1] * (x - i) % MOD;

    suff[k] = (x - k) % MOD;
    for (int i = k - 1; i >= 0; i--)
        suff[i] = suff[i + 1] * (x - i) % MOD;

    ll ans = 0;
    for (int i = 0; i <= k; i++)
    {
        ll num = (i == 0 ? 1 : pref[i - 1]) * (i == k ? 1 : suff[i + 1]) % MOD;
        ll den = invFact[i] * invFact[k - i] % MOD;
        ll term = y[i] * num % MOD * den % MOD;
        if ((k - i) % 2 == 1)
        {
            ans = (ans - term + MOD) % MOD;
        }
        else
        {
            ans = (ans + term) % MOD;
        }
    }
    return ans;
}
\end{lstlisting}
\subsubsection{Multipoint evaluation}
\begin{lstlisting}[language=C++]
/**
 *  THE ULTIMATE CP TEMPLATE BLUEPRINT: MULTIPOINT EVALUATION
 *
 * 1.  WHAT PROBLEM DOES THIS SOLVE?
 * - Keywords: "Evaluate polynomial at multiple points", "P(x_i)".
 * - Classic Scenarios: You have a massive polynomial P(x) of degree N, and you need
 *   to find its exact value at M different points (Topic 340).
 * - The Magic: Horner's method takes O(N * M) which TLEs. This uses a Subproduct Tree
 *   (a segment tree of polynomials) and polynomial division with remainder to process
 *   all queries simultaneously in O(N log^2 N).
 *
 * 2.  HOW TO USE IT (THE BLACK BOX)
 * - Query:
 *       vector<ll> results = multipoint_eval(P, X);
 *       // results[i] contains P(X[i]) mod 998244353.
 *
 * - Complexity:
 *       Time: O(N log^2 N)
 *       Space: O(N log N)
 *
 * 3.  HOW TO ADAPT IT (THE DIALS & SWITCHES)
 * - Polynomial Modulo: This struct includes `poly_mod(A, B)` which finds the remainder
 *   of A(x) / B(x). This is an incredibly powerful tool on its own.
 */

vector<ll> poly_rev(vector<ll> a)
{
    reverse(a.begin(), a.end());
    return a;
}

vector<ll> poly_mod(vector<ll> a, vector<ll> b)
{
    while (a.size() > 1 && a.back() == 0)
        a.pop_back();
    while (b.size() > 1 && b.back() == 0)
        b.pop_back();
    if (a.size() < b.size())
        return a;

    int n = a.size(), m = b.size();
    vector<ll> rev_a = poly_rev(a);
    vector<ll> rev_b = poly_rev(b);

    rev_a.resize(n - m + 1);
    vector<ll> inv_rev_b = poly_inv(rev_b, n - m + 1);
    vector<ll> q = multiply(rev_a, inv_rev_b);
    q.resize(n - m + 1);
    q = poly_rev(q);

    vector<ll> qb = multiply(q, b);
    vector<ll> r(m - 1, 0);
    for (int i = 0; i < m - 1; i++)
    {
        r[i] = (a[i] - (i < qb.size() ? qb[i] : 0)) % MOD;
        if (r[i] < 0)
            r[i] += MOD;
    }
    while (r.size() > 1 && r.back() == 0)
        r.pop_back();
    return r;
}

vector<vector<ll>> tree;

void build_subproduct_tree(int node, int l, int r, const vector<ll> &x)
{
    if (l == r)
    {
        tree[node] = {(MOD - x[l]) % MOD, 1};
        return;
    }
    int mid = l + (r - l) / 2;
    build_subproduct_tree(2 * node, l, mid, x);
    build_subproduct_tree(2 * node + 1, mid + 1, r, x);
    tree[node] = multiply(tree[2 * node], tree[2 * node + 1]);
}

void eval_subproduct_tree(int node, int l, int r, vector<ll> p, vector<ll> &ans)
{
    p = poly_mod(p, tree[node]);
    if (l == r)
    {
        ans[l] = p.empty() ? 0 : p[0];
        return;
    }
    int mid = l + (r - l) / 2;
    eval_subproduct_tree(2 * node, l, mid, p, ans);
    eval_subproduct_tree(2 * node + 1, mid + 1, r, p, ans);
}

vector<ll> multipoint_eval(const vector<ll> &p, const vector<ll> &x)
{
    if (x.empty())
        return {};
    int m = x.size();
    tree.resize(4 * m);
    build_subproduct_tree(1, 0, m - 1, x);
    vector<ll> ans(m);
    eval_subproduct_tree(1, 0, m - 1, p, ans);
    return ans;
}
\end{lstlisting}
\subsubsection{Poly inv log}
\begin{lstlisting}[language=C++]
/**
 *  THE ULTIMATE CP TEMPLATE BLUEPRINT: POLYNOMIAL INV & LOG
 *
 * 1.  WHAT PROBLEM DOES THIS SOLVE?
 * - Keywords: "Generating Functions", "Formal Power Series (FPS)", "Divide Polynomials".
 * - Classic Scenarios: You have a generating function expressed as a fraction A(x)/B(x)
 *   and need the first N coefficients.
 * - The Magic: Polynomial Inverse (Topic 334) uses Newton's Method to find B^{-1}(x)
 *   such that B(x) * B^{-1}(x) = 1 mod x^n in O(N log N). Polynomial Log (Topic 336)
 *   uses calculus: ln(A(x)) = integral( A'(x) * A^{-1}(x) ).
 *
 * 2.  HOW TO USE IT (THE BLACK BOX)
 * - Query:
 *       vector<ll> invA = poly_inv(A, N);
 *       vector<ll> logA = poly_log(A, N);
 *
 * - Complexity:
 *       Time: O(N log N)
 *       Space: O(N)
 *
 * 3.  HOW TO ADAPT IT (THE DIALS & SWITCHES)
 * - Dependencies: Requires `ntt`, `multiply`, and `modInverse` from your FFT/NTT snippet.
 * - Constant Term: For poly_inv, A[0] MUST NOT be 0. For poly_log, A[0] MUST be 1.
 */

using ll = long long;
const ll MOD = 998244353;

vector<ll> poly_inv(vector<ll> a, int deg)
{
    if (deg == 1)
        return {modInverse(a[0])};
    vector<ll> b = poly_inv(a, (deg + 1) / 2);
    int n = 1;
    while (n < deg * 2)
        n <<= 1;
    vector<ll> fa(a.begin(), a.begin() + min((int)a.size(), deg));
    fa.resize(n);
    vector<ll> fb = b;
    fb.resize(n);
    ntt(fa, false);
    ntt(fb, false);
    for (int i = 0; i < n; i++)
    {
        fb[i] = fb[i] * (2 - fa[i] * fb[i] % MOD + MOD) % MOD;
    }
    ntt(fb, true);
    fb.resize(deg);
    return fb;
}

vector<ll> poly_deriv(const vector<ll> &a)
{
    if (a.empty())
        return {};
    vector<ll> res(max(1, (int)a.size() - 1));
    for (int i = 1; i < a.size(); i++)
    {
        res[i - 1] = a[i] * i % MOD;
    }
    return res;
}

vector<ll> poly_integr(const vector<ll> &a)
{
    if (a.empty())
        return {0};
    vector<ll> res(a.size() + 1, 0);
    for (int i = 0; i < a.size(); i++)
    {
        res[i + 1] = a[i] * modInverse(i + 1) % MOD;
    }
    return res;
}

vector<ll> poly_log(vector<ll> a, int deg)
{
    vector<ll> deriv = poly_deriv(a);
    vector<ll> inv = poly_inv(a, deg);
    vector<ll> res = poly_integr(multiply(deriv, inv));
    res.resize(deg);
    return res;
}
\end{lstlisting}
\subsection{Probability}
\subsubsection{Expected occurence time of substring}
\begin{lstlisting}[language=C++]
/**
 *  THE ULTIMATE CP TEMPLATE BLUEPRINT: EXPECTED OCCURRENCE TIME OF SUBSTRING
 *
 * 1.  WHAT PROBLEM DOES THIS SOLVE?
 * - Keywords: "Expected number of coin flips", "Typing monkeys", "Expected time to see string".
 * - Classic Scenarios: A random generator picks characters from an alphabet of size A.
 *   What is the expected number of generated characters until string S appears?
 * - The Magic: Bypasses O(N^3) Gaussian elimination. By setting up a fair casino
 *   Martingale where gamblers bet on consecutive characters, the expected time simplifies
 *   to a summation over the borders (prefix == suffix) of the string.
 *
 * 2.  HOW TO USE IT (THE BLACK BOX)
 * - Query:
 *       long long expected_time = expected_string_occurrence(S, ALPHABET_SIZE, MOD);
 *
 * - Complexity:
 *       Time: O(|S|)
 *       Space: O(|S|)
 *
 * 3.  HOW TO ADAPT IT (THE DIALS & SWITCHES)
 * - Alphabet Size: If flipping a coin, alphabet_size = 2. If rolling a die, 6.
 * - Array Form: This strictly uses the KMP prefix function. The condition `j == pi[j-1]`
 *   efficiently finds all prefix-suffix borders of the entire string.
 */

using ll = long long;

vector<int> compute_pi(const string &s)
{
    int n = s.length();
    vector<int> pi(n);
    for (int i = 1; i < n; i++)
    {
        int j = pi[i - 1];
        while (j > 0 && s[i] != s[j])
            j = pi[j - 1];
        if (s[i] == s[j])
            j++;
        pi[i] = j;
    }
    return pi;
}

ll expected_string_occurrence(const string &s, ll alphabet_size, ll mod)
{
    vector<int> pi = compute_pi(s);
    int n = s.length();

    auto binpow = [&](ll base, ll exp)
    {
        ll res = 1;
        base %= mod;
        while (exp > 0)
        {
            if (exp % 2 == 1)
                res = (res * base) % mod;
            base = (base * base) % mod;
            exp /= 2;
        }
        return res;
    };

    ll expected_time = 0;
    int current_border = n;

    while (current_border > 0)
    {
        expected_time = (expected_time + binpow(alphabet_size, current_border)) % mod;
        current_border = pi[current_border - 1];
    }

    return expected_time;
}
\end{lstlisting}
\section{Number Theory}
\subsection{Fibonacci}
\subsubsection{Fast doubling fibonacci}
\begin{lstlisting}[language=C++]
/**
 *  THE ULTIMATE CP TEMPLATE BLUEPRINT: FAST DOUBLING FIBONACCI
 *
 * 1.  WHAT PROBLEM DOES THIS SOLVE?
 * - Keywords: "N-th Fibonacci", "Massive N".
 * - Classic Scenarios: You need F_N modulo M where N = 10^18.
 * - The Magic: Derives F_{2k} and F_{2k+1} directly from F_k and F_{k+1}.
 *   It completely sidesteps matrix overhead, making it drastically faster
 *   and much shorter to type than a robust matrix exponentiation struct.
 *
 * 2.  HOW TO USE IT (THE BLACK BOX)
 * - Query:
 *       pair<ll, ll> ans = fast_doubling(N, M);
 * - Result Breakdown:
 *       ans.first = F_N
 *       ans.second = F_{N+1}
 *
 * - Complexity:
 *       Time: Strictly O(log N).
 *       Space: O(log N) for the recursion stack.
 */

using ll = long long;

pair<ll, ll> fast_doubling(ll n, ll mod)
{
    if (n == 0)
        return {0, 1};
    auto p = fast_doubling(n >> 1, mod);
    ll c = p.first * (2 * p.second - p.first + mod) % mod;
    ll d = (p.first * p.first + p.second * p.second) % mod;
    if (n & 1)
        return {d, (c + d) % mod};
    return {c, d};
}
\end{lstlisting}
\subsubsection{PHI field}
\begin{lstlisting}[language=C++]
/**
 *  THE ULTIMATE CP TEMPLATE BLUEPRINT: PHI FIELD (Z[sqrt(5)])
 *
 * 1.  WHAT PROBLEM DOES THIS SOLVE?
 * - Keywords: "Fibonacci sum of powers", "Binet's formula modulo M".
 * - Classic Scenarios: You need to evaluate algebraic expressions involving
 *   Fibonacci numbers, such as sum(F_i^K) mod M.
 * - The Magic: It operates exactly like complex numbers, but instead of i^2 = -1,
 *   it maps a + b*sqrt(5) under modulo M. This allows you to directly use Binet's
 *   formula natively even if 5 is a quadratic non-residue modulo M!
 *
 * 2.  HOW TO USE IT (THE BLACK BOX)
 * - Query:
 *       Phi res = power_phi({A, B}, exp, MOD);
 *
 * - Complexity:
 *       Time: O(log(exp)) per multiplication.
 *       Space: O(1)
 */

using ll = long long;

struct Phi
{
    ll a, b;
};

Phi mult_phi(Phi x, Phi y, ll mod)
{
    return {
        (x.a * y.a + 5 * x.b % mod * y.b) % mod,
        (x.a * y.b + x.b * y.a) % mod};
}

Phi power_phi(Phi base, ll exp, ll mod)
{
    Phi res = {1, 0};
    base.a %= mod;
    base.b %= mod;
    while (exp > 0)
    {
        if (exp % 2 == 1)
            res = mult_phi(res, base, mod);
        base = mult_phi(base, base, mod);
        exp /= 2;
    }
    return res;
}
\end{lstlisting}
\subsection{Linear diophantine equation}
\subsubsection{Floor sum}
\begin{lstlisting}[language=C++]
/**
 *  THE ULTIMATE CP TEMPLATE BLUEPRINT: FLOOR SUM
 *
 * 1.  WHAT PROBLEM DOES THIS SOLVE?
 * - Keywords: "ax + by <= c", "Lattice points under line", "Sum of arithmetic progression mod M".
 * - Classic Scenarios: Compute sum_{i=0}^{n-1} floor((a * i + b) / m).
 * - The Magic: It operates exactly like the Euclidean algorithm, trading dimensions
 *   and folding the geometry to evaluate the entire sum logarithmically.
 *
 * 2.  HOW TO USE IT (THE BLACK BOX)
 * - Query:
 *       long long ans = floor_sum(n, m, a, b);
 *
 * - Complexity:
 *       Time: O(log(min(m, a)))
 *       Space: O(log(min(m, a))) due to recursion.
 *
 * 3.  HOW TO ADAPT IT (THE DIALS & SWITCHES)
 * - Equivalent to ax + by <= c: For a line equation in the first quadrant, rearrange
 *   it to y <= (c - ax) / b, which directly translates to summing floor((c - a*x) / b)
 *   over all valid x.
 */

using ll = long long;

ll floor_sum(ll n, ll m, ll a, ll b)
{
    ll ans = 0;
    if (a >= m)
    {
        ans += (n - 1) * n / 2 * (a / m);
        a %= m;
    }
    if (b >= m)
    {
        ans += n * (b / m);
        b %= m;
    }
    ll y_max = (a * n + b) / m;
    ll x_max = y_max * m - b;
    if (y_max == 0)
        return ans;
    ans += (n - (x_max + a - 1) / a) * y_max;
    ans += floor_sum(y_max, a, m, (a - x_max % a) % a);
    return ans;
}
\end{lstlisting}
\subsubsection{Generalized floor sum}
\begin{lstlisting}[language=C++]
/**
 *  THE ULTIMATE CP TEMPLATE BLUEPRINT: GENERALIZED FLOOR SUM
 *
 * 1.  WHAT PROBLEM DOES THIS SOLVE?
 * - Keywords: "Sum of i * floor", "Sum of squares of floors".
 * - Classic Scenarios: A step above standard floor sum, used when evaluating
 *   complex grid summations or polynomial constraints under a modulo line.
 * - The Magic: Through a complex series of algebraic expansions and Euclidean
 *   domain folding, it simultaneously resolves three interconnected sums in
 *   pure logarithmic time.
 *
 * 2.  HOW TO USE IT (THE BLACK BOX)
 * - Query:
 *       FSR ans = gen_floor_sum(n, a, b, c);
 * - Result Breakdown:
 *       ans.f = sum_{i=0}^n floor((a*i+b)/c)
 *       ans.g = sum_{i=0}^n i * floor((a*i+b)/c)
 *       ans.h = sum_{i=0}^n floor((a*i+b)/c)^2
 *
 * - Complexity:
 *       Time: O(log(min(a, c)))
 *       Space: O(log(min(a, c)))
 *
 * 3.  HOW TO ADAPT IT (THE DIALS & SWITCHES)
 * - Inclusive Bounds: This specifically calculates for i = 0 to N (inclusive).
 *   If your problem is 1-based, evaluate for N and subtract the i=0 term, or
 *   use N-1 if the upper bound is exclusive.
 */

using ll = long long;

struct FSR
{
    ll f, g, h;
};

FSR gen_floor_sum(ll n, ll a, ll b, ll c)
{
    ll ac = a / c, bc = b / c;
    if (a == 0)
    {
        ll f = (n + 1) * bc;
        ll g = bc * n * (n + 1) / 2;
        ll h = (n + 1) * bc * bc;
        return {f, g, h};
    }
    if (a >= c || b >= c)
    {
        FSR res = gen_floor_sum(n, a % c, b % c, c);
        ll f = res.f + n * (n + 1) / 2 * ac + (n + 1) * bc;
        ll g = res.g + ac * n * (n + 1) * (2 * n + 1) / 6 + bc * n * (n + 1) / 2;
        ll h = res.h + ac * ac * n * (n + 1) * (2 * n + 1) / 6 + bc * bc * (n + 1) + 2 * ac * bc * n * (n + 1) / 2 + 2 * ac * res.g + 2 * bc * res.f;
        return {f, g, h};
    }
    ll m = (a * n + b) / c;
    FSR res = gen_floor_sum(m - 1, c, c - b - 1, a);
    ll f = n * m - res.f;
    ll g = (m * n * (n + 1) - res.f - res.h) / 2;
    ll h = n * m * m - 2 * res.g - res.f;
    return {f, g, h};
}
\end{lstlisting}
\subsubsection{LDE}
\begin{lstlisting}[language=C++]
/**
 *  THE ULTIMATE CP TEMPLATE BLUEPRINT: BOUNDED LDE2
 *
 * 1.  WHAT PROBLEM DOES THIS SOLVE?
 * - Keywords: "Solutions in range", "ax + by = c", "Bounded variables".
 * - Classic Scenarios: You need the exact count of integer pairs (x, y) satisfying
 *   ax + by = c such that min_x <= x <= max_x and min_y <= y <= max_y.
 * - The Magic: It finds a base solution using Extended Euclidean Algorithm, then
 *   mathematically shifts the solution using (b/g) and (a/g) to find the absolute
 *   minimum and maximum valid boundaries, intersecting them safely.
 *
 * 2.  HOW TO USE IT (THE BLACK BOX)
 * - Query:
 *       long long total = count_solutions(a, b, c, min_x, max_x, min_y, max_y);
 *
 * - Complexity:
 *       Time: O(log(min(a, b)))
 *       Space: O(1)
 *
 * 3.  HOW TO ADAPT IT (THE DIALS & SWITCHES)
 * - Zero Handling: Safely handles edge cases where a = 0 or b = 0, which often
 *   crash naive implementations.
 */

using ll = long long;

ll extgcd(ll a, ll b, ll &x, ll &y)
{
    if (b == 0)
    {
        x = 1;
        y = 0;
        return a;
    }
    ll x1, y1;
    ll d = extgcd(b, a % b, x1, y1);
    x = y1;
    y = x1 - y1 * (a / b);
    return d;
}

void shift_solution(ll &x, ll &y, ll a, ll b, ll cnt)
{
    x += cnt * b;
    y -= cnt * a;
}

ll count_solutions(ll a, ll b, ll c, ll min_x, ll max_x, ll min_y, ll max_y)
{
    if (a == 0 && b == 0)
    {
        if (c == 0)
            return (max_x - min_x + 1) * (max_y - min_y + 1);
        return 0;
    }
    if (a == 0)
    {
        if (c % b == 0 && c / b >= min_y && c / b <= max_y)
            return max_x - min_x + 1;
        return 0;
    }
    if (b == 0)
    {
        if (c % a == 0 && c / a >= min_x && c / a <= max_x)
            return max_y - min_y + 1;
        return 0;
    }

    ll x, y;
    ll g = extgcd(abs(a), abs(b), x, y);
    if (c % g != 0)
        return 0;

    x *= c / g;
    y *= c / g;
    if (a < 0)
        x = -x;
    if (b < 0)
        y = -y;

    ll a_g = a / g, b_g = b / g;
    ll sign_a = a > 0 ? +1 : -1;
    ll sign_b = b > 0 ? +1 : -1;

    shift_solution(x, y, a_g, b_g, (min_x - x) / b_g);
    if (x < min_x)
        shift_solution(x, y, a_g, b_g, sign_b);
    if (x > max_x)
        return 0;
    ll lx1 = x;

    shift_solution(x, y, a_g, b_g, (max_x - x) / b_g);
    if (x > max_x)
        shift_solution(x, y, a_g, b_g, -sign_b);
    ll rx1 = x;

    shift_solution(x, y, a_g, b_g, -(min_y - y) / a_g);
    if (y < min_y)
        shift_solution(x, y, a_g, b_g, -sign_a);
    if (y > max_y)
        return 0;
    ll lx2 = x;

    shift_solution(x, y, a_g, b_g, -(max_y - y) / a_g);
    if (y > max_y)
        shift_solution(x, y, a_g, b_g, sign_a);
    ll rx2 = x;

    if (lx2 > rx2)
        swap(lx2, rx2);
    ll lx = max(lx1, lx2);
    ll rx = min(rx1, rx2);

    if (lx > rx)
        return 0;
    return (rx - lx) / abs(b_g) + 1;
}
\end{lstlisting}
\subsubsection{Pell equation}
\begin{lstlisting}[language=C++]
/**
 *  THE ULTIMATE CP TEMPLATE BLUEPRINT: PELL'S EQUATION
 *
 * 1.  WHAT PROBLEM DOES THIS SOLVE?
 * - Keywords: "x^2 - d*y^2 = 1", "Fundamental Solution".
 * - Classic Scenarios: Find the smallest positive integer solution (x, y)
 *   satisfying Pell's equation where d is a non-square positive integer.
 * - The Magic: It generates the continued fraction expansion of sqrt(d).
 *   The convergents (h_i / k_i) of this fraction are guaranteed to eventually
 *   satisfy the equation.
 *
 * 2.  HOW TO USE IT (THE BLACK BOX)
 * - Query:
 *       pair<__int128, __int128> ans = solve_pell(d);
 *       // Returns {x, y}. If d is a perfect square, returns {-1, -1}.
 *
 * - Complexity:
 *       Time: O(sqrt(d) * log(d))
 *       Space: O(1)
 *
 * 3.  HOW TO ADAPT IT (THE DIALS & SWITCHES)
 * - 128-bit Return Type: Fundamental solutions to Pell's equation grow
 *   exponentially fast. Even for small d (e.g., d = 61), x and y can easily
 *   overflow 64-bit integers. __int128 is highly recommended here.
 * - Generating More Solutions: If (x1, y1) is the fundamental solution,
 *   the next solution (x_k, y_k) is (x1 + y1*sqrt(d))^k.
 */

using u128 = __int128;
using ll = long long;

pair<u128, u128> solve_pell(ll d)
{
    ll sq = sqrt(d);
    if (sq * sq == d)
        return {-1, -1};

    ll m = 0, d_den = 1, a = sq;
    u128 num1 = 1, num0 = a;
    u128 den1 = 0, den0 = 1;

    while (num0 * num0 - (u128)d * den0 * den0 != 1)
    {
        m = d_den * a - m;
        d_den = (d - m * m) / d_den;
        a = (sq + m) / d_den;

        u128 num2 = num1;
        num1 = num0;
        num0 = (u128)a * num1 + num2;

        u128 den2 = den1;
        den1 = den0;
        den0 = (u128)a * den1 + den2;
    }
    return {num0, den0};
}
\end{lstlisting}
\subsection{Misc}
\subsubsection{Min X in Mod range}
\begin{lstlisting}[language=C++]
/**
 *  THE ULTIMATE CP TEMPLATE BLUEPRINT: MINIMUM X IN MODULO RANGE
 *
 * 1.  WHAT PROBLEM DOES THIS SOLVE?
 * - Keywords: "Smallest x", "L <= ax mod m <= R", "Min of Mod of AP".
 * - Classic Scenarios: Finding the first index where an arithmetic progression
 *   wraps around a modulo into a specific target window.
 * - The Magic: It uses a recursive Euclidean-style domain folding. Instead
 *   of simulating the progression, it inverts the bounds and steps to find
 *   the exact minimum x in O(log M) time.
 *
 * 2.  HOW TO USE IT (THE BLACK BOX)
 * - Query:
 *       long long x = min_x_in_mod_range(A, M, L, R);
 *       // Returns -1 if no such x exists.
 *
 * - Complexity:
 *       Time: O(log M)
 *       Space: O(log M) for recursion.
 *
 * 3.  HOW TO ADAPT IT (THE DIALS & SWITCHES)
 * - Min of Mod of AP (Topic 295): If you just need the absolute minimum of
 *   (A*x) mod M for x in [0, N], you can binary search the answer V in [0, M-1]
 *   and use this function to check if a valid x <= N exists for L=0, R=V.
 */

using ll = long long;

ll min_x_in_mod_range(ll a, ll m, ll l, ll r)
{
    if (l == 0)
        return 0;
    if (a == 0)
        return -1;
    ll step = (l + a - 1) / a;
    if (step * a <= r)
        return step;
    ll b = m % a;
    ll res = min_x_in_mod_range(b, a, a - r % a, a - l % a);
    if (res == -1)
        return -1;
    return (l + res * b + a - 1) / a;
}
\end{lstlisting}
\subsubsection{SIMPLEST FRACTION}
\begin{lstlisting}[language=C++]
/**
 *  THE ULTIMATE CP TEMPLATE BLUEPRINT: STERN-BROCOT TREE (SIMPLEST FRACTION)
 *
 * 1.  WHAT PROBLEM DOES THIS SOLVE?
 * - Keywords: "Stern-Brocot", "Simplest Fraction", "Rational Approximation".
 * - Classic Scenarios: Find the fraction P/Q with the absolute smallest
 *   denominator Q that strictly lies between two other fractions A/B and C/D.
 * - The Magic: Instead of simulating the binary search descent down the
 *   Stern-Brocot tree (which can TLE in O(max(P, Q))), it uses a continued
 *   fraction / Euclidean acceleration to jump down the branches.
 *
 * 2.  HOW TO USE IT (THE BLACK BOX)
 * - Query:
 *       pair<long long, long long> ans = simplest_fraction(A, B, C, D);
 *       // Returns {P, Q} such that A/B < P/Q < C/D.
 *
 * - Complexity:
 *       Time: O(log(max(A, B, C, D)))
 *       Space: O(log(max)) for recursion.
 *
 * 3.  HOW TO ADAPT IT (THE DIALS & SWITCHES)
 * - Fraction format: Ensure your input fractions A/B and C/D are simplified
 *   and that A/B is strictly less than C/D before passing them in.
 */

using ll = long long;

pair<ll, ll> simplest_fraction(ll a, ll b, ll c, ll d)
{
    if (a / b < c / d)
        return {a / b + 1, 1};
    if (a == 0)
        return {1, d / c + 1};
    auto res = simplest_fraction(d, c % d, b, a % b);
    return {res.second + (a / b) * res.first, res.first};
}
\end{lstlisting}
\subsection{Modular arithmetic}
\subsubsection{Baby step giant step}
\begin{lstlisting}[language=C++]
/**
 *  THE ULTIMATE CP TEMPLATE BLUEPRINT: BABY-STEP GIANT-STEP (DISCRETE LOG)
 *
 * 1.  WHAT PROBLEM DOES THIS SOLVE?
 * - Keywords: "Discrete Logarithm", "Find the exponent", "A^x = B mod M".
 * - Classic Scenarios: You are given A, B, and M, and need to find the minimum
 *   non-negative integer x such that A^x  B (mod M).
 * - The Magic: It trades space for time. By rewriting x = n*p - q (where n = sqrt(M)),
 *   the equation becomes A^(n*p)  B * A^q (mod M). We precompute the right side
 *   (Baby Steps) and store them in a hash map, then calculate the left side
 *   (Giant Steps) and check for collisions.
 *
 * 2.  HOW TO USE IT (THE BLACK BOX)
 * - Query:
 *       long long ans = solve_bsgs(A, B, M);
 *       // Returns -1 if no solution exists.
 *
 * - Complexity:
 *       Time: O(sqrt(M) * log(sqrt(M))) due to the map. O(sqrt(M)) if using unordered_map.
 *       Space: O(sqrt(M)) to store the baby steps.
 *
 * 3.  HOW TO ADAPT IT (THE DIALS & SWITCHES)
 * - Coprime Requirement: This standard version assumes gcd(A, M) == 1. If A and M
 *   are not coprime, you need the Extended BSGS algorithm (which involves dividing
 *   out the GCDs until they are coprime).
 */

using ll = long long;

ll solve_bsgs(ll a, ll b, ll m)
{
    a %= m, b %= m;
    if (b == 1 || m == 1)
        return 0;

    ll n = sqrt(m) + 1;
    map<ll, ll> vals;
    ll cur = b;
    for (ll q = 0; q <= n; ++q)
    {
        vals[cur] = q;
        cur = (cur * a) % m;
    }

    ll a_n = 1;
    for (ll i = 0; i < n; ++i)
    {
        a_n = (a_n * a) % m;
    }

    cur = 1;
    for (ll p = 1; p <= n; ++p)
    {
        cur = (cur * a_n) % m;
        if (vals.count(cur))
        {
            ll ans = n * p - vals[cur];
            return ans;
        }
    }
    return -1;
}
\end{lstlisting}
\subsubsection{Big Mod Power}
\begin{lstlisting}[language=C++]
/// if we need to calculate (n % m) if  EX: n = 10^10000 and m = 10^9
ll Big_Mod_Power(string n, ll m) 
{
    ll x = 0;
    for(ll i = 0; i < n.size(); i++)
    {
        x = ((x * 10)) + (n[i] - '0');
        x %= m;   
    }
    return x;  /// return n % m
}

\end{lstlisting}
\subsubsection{CRT}
\begin{lstlisting}[language=C++]
/**
 *  THE ULTIMATE CP TEMPLATE BLUEPRINT: CHINESE REMAINDER THEOREM (GENERAL)
 *
 * 1.  WHAT PROBLEM DOES THIS SOLVE?
 * - Keywords: "System of congruences", "x = a1 mod m1, x = a2 mod m2".
 * - The Magic: Unlike basic CRT which fails if the moduli are not coprime,
 *   this robust version iteratively combines the equations using the Extended
 *   Euclidean algorithm. It safely catches impossible systems.
 *
 * 2.  HOW TO USE IT (THE BLACK BOX)
 * - Query:
 *       pair<ll, ll> ans = crt(A, M);
 * - Result: Returns {remainder, modulo}. If no solution exists, returns {-1, -1}.
 */

using ll = long long;

ll extgcd_crt(ll a, ll b, ll &x, ll &y)
{
    if (b == 0)
    {
        x = 1;
        y = 0;
        return a;
    }
    ll x1, y1;
    ll d = extgcd_crt(b, a % b, x1, y1);
    x = y1;
    y = x1 - y1 * (a / b);
    return d;
}

pair<ll, ll> crt(const vector<ll> &a, const vector<ll> &m)
{
    if (a.empty())
        return {0, 1};
    ll rem = a[0], mod = m[0];
    for (size_t i = 1; i < a.size(); i++)
    {
        ll x, y;
        ll g = extgcd_crt(mod, m[i], x, y);
        if ((a[i] - rem) % g != 0)
            return {-1, -1};
        ll target = (a[i] - rem) / g;
        ll m_i_g = m[i] / g;
        x = (x % m_i_g + m_i_g) % m_i_g;
        x = (x * (target % m_i_g + m_i_g)) % m_i_g;
        rem += mod * x;
        mod *= m_i_g;
        rem = (rem % mod + mod) % mod;
    }
    return {rem, mod};
}
\end{lstlisting}
\subsubsection{Discrete root}
\begin{lstlisting}[language=C++]
/**
 *  THE ULTIMATE CP TEMPLATE BLUEPRINT: DISCRETE ROOT (K-TH ROOT)
 *
 * 1.  WHAT PROBLEM DOES THIS SOLVE?
 * - Keywords: "K-th Root Modulo P", "x^K = A mod P".
 * - Classic Scenarios: You need to reverse an exponentiation where the modulo is prime.
 * - The Magic: It finds a primitive root g of P. Since x can be represented as g^y,
 *   the equation becomes (g^y)^K  A (mod P) => (g^K)^y  A (mod P). This perfectly
 *   reduces the problem into a Discrete Logarithm (BSGS) problem!
 *
 * 2.  HOW TO USE IT (THE BLACK BOX)
 * - Query:
 *       long long root = discrete_root(K, A, P);
 *       // Returns -1 if no root exists.
 *
 * - Complexity:
 *       Time: O(sqrt(P) * log(sqrt(P))).
 *       Space: O(sqrt(P)).
 *
 * 3.  HOW TO ADAPT IT (THE DIALS & SWITCHES)
 * - Dependencies: This snippet requires both `custom_pow` and `solve_bsgs` to be
 *   available in your code (from the snippets above).
 */

using ll = long long;

ll primitive_root(ll p)
{
    vector<ll> fact;
    ll phi = p - 1, n = phi;
    for (ll i = 2; i * i <= n; ++i)
    {
        if (n % i == 0)
        {
            fact.push_back(i);
            while (n % i == 0)
                n /= i;
        }
    }
    if (n > 1)
        fact.push_back(n);

    for (ll res = 2; res <= p; ++res)
    {
        bool ok = true;
        for (size_t i = 0; i < fact.size() && ok; ++i)
        {
            ok &= custom_pow(res, phi / fact[i], p) != 1;
        }
        if (ok)
            return res;
    }
    return -1;
}

ll discrete_root(ll k, ll a, ll p)
{
    if (a == 0)
        return 0;
    if (k == 0)
        return (a == 1) ? 1 : -1;

    ll g = primitive_root(p);
    if (g == -1)
        return -1;

    ll base = custom_pow(g, k, p);
    ll y = solve_bsgs(base, a, p);
    if (y == -1)
        return -1;

    return custom_pow(g, y, p);
}
\end{lstlisting}
\subsubsection{Distinct kth power}
\begin{lstlisting}[language=C++]
/**
 *  THE ULTIMATE CP TEMPLATE BLUEPRINT: DISTINCT K-TH POWERS & X^2 = 1 MOD M
 *
 * 1.  WHAT PROBLEM DOES THIS SOLVE?
 * - Keywords: "How many perfect K-th powers modulo P", "Square roots of 1 modulo M".
 * - The Magic: These are closed-form formula generators. Calculating the number
 *   of valid values for x^2 = 1 mod M depends heavily on the prime factorization
 *   of M (specifically treating the powers of 2 differently than odd primes).
 *
 * 2.  HOW TO USE IT (THE BLACK BOX)
 * - Query:
 *       ll ans1 = distinct_kth_powers_prime_mod(K, P); // Requires P to be a prime
 *       ll ans2 = solutions_x2_1_mod_m(M);
 */

using ll = long long;

ll distinct_kth_powers_prime_mod(ll k, ll p)
{
    return (p - 1) / __gcd(k, p - 1) + 1;
}

ll solutions_x2_1_mod_m(ll m)
{
    if (m <= 1)
        return 1;
    ll count = 1;
    ll temp = m;
    int e2 = 0;
    while (temp % 2 == 0)
    {
        e2++;
        temp /= 2;
    }
    if (e2 == 2)
        count *= 2;
    else if (e2 >= 3)
        count *= 4;

    for (ll i = 3; i * i <= temp; i += 2)
    {
        if (temp % i == 0)
        {
            count *= 2;
            while (temp % i == 0)
                temp /= i;
        }
    }
    if (temp > 1)
        count *= 2;
    return count;
}
\end{lstlisting}
\subsubsection{Linear congurence}
\begin{lstlisting}[language=C++]
/**
 *  THE ULTIMATE CP TEMPLATE BLUEPRINT: LINEAR CONGRUENCE EQUATION
 *
 * 1.  WHAT PROBLEM DOES THIS SOLVE?
 * - Keywords: "Solve ax = b mod m".
 * - Classic Scenarios: You are given a, b, and m, and need to find x.
 * - The Magic: It uses the Extended Euclidean Algorithm to solve ax + my = b.
 *   It handles cases where gcd(a, m) != 1. If a solution exists, it returns
 *   the base solution x0 and the number of distinct solutions modulo m.
 *
 * 2.  HOW TO USE IT (THE BLACK BOX)
 * - Query:
 *       ll x, num_solutions;
 *       bool possible = linear_congruence(a, b, m, x, num_solutions);
 */

using ll = long long;

ll extgcd(ll a, ll b, ll &x, ll &y)
{
    if (b == 0)
    {
        x = 1;
        y = 0;
        return a;
    }
    ll x1, y1;
    ll d = extgcd(b, a % b, x1, y1);
    x = y1;
    y = x1 - y1 * (a / b);
    return d;
}

bool linear_congruence(ll a, ll b, ll m, ll &x, ll &num_solutions)
{
    ll x0, y0;
    ll g = extgcd(abs(a), m, x0, y0);
    if (b % g != 0)
    {
        x = 0;
        num_solutions = 0;
        return false;
    }
    x0 = (x0 % (m / g) + (m / g)) % (m / g);
    x = (x0 * (b / g)) % (m / g);
    num_solutions = g;
    return true;
}
\end{lstlisting}
\subsubsection{MULTIPLICATIVE ORDER CARMICHAEL'S LAMBDA}
\begin{lstlisting}[language=C++]
/**
 *  MULTIPLICATIVE ORDER & CARMICHAEL'S LAMBDA
 *
 * WHAT DOES THIS DO?
 * Finds the smallest integer k such that (a^k) % n == 1.
 * Using Euler's Totient phi(n) is too slow to find the exact order because
 * the order divides Carmichael's Lambda function, which divides phi(n).
 *
 * MATH BACKGROUND:
 * lambda(p^k) = p^(k-1) * (p - 1)  [except for p=2, k>=3 where it's divided by 2]
 * lambda(a * b) = lcm(lambda(a), lambda(b))
 * The multiplicative order of 'a' modulo 'n' strictly divides lambda(n).
 *
 * Complexity: O(sqrt(N)) to factorize N.
 */

using ll = long long;

ll gcd(ll a, ll b)
{
    while (b)
    {
        a %= b;
        swap(a, b);
    }
    return a;
}

ll lcm(ll a, ll b)
{
    return (a / gcd(a, b)) * b;
}

// Modular exponentiation using __int128 to prevent overflow when multiplying two long longs
ll power(ll base, ll exp, ll mod)
{
    ll res = 1;
    base %= mod;
    while (exp > 0)
    {
        if (exp % 2 == 1)
            res = (ll)((unsigned __int128)res * base % mod);
        base = (ll)((unsigned __int128)base * base % mod);
        exp /= 2;
    }
    return res;
}

ll carmichael_lambda(ll n)
{
    ll lambda = 1;
    ll temp = n;

    // Step 1: Prime factorization of n
    for (ll i = 2; i * i <= temp; i++)
    {
        if (temp % i == 0)
        {
            ll p_pow = 1;
            while (temp % i == 0)
            {
                p_pow *= i;
                temp /= i;
            }

            // term = phi(p^k) = p^k - p^(k-1)
            ll term = p_pow - p_pow / i;

            // Special case for Carmichael: if p = 2 and k >= 3, lambda is phi/2
            if (i == 2 && p_pow >= 8)
            {
                term /= 2;
            }
            lambda = lcm(lambda, term);
        }
    }

    // If there is a prime factor greater than sqrt(n)
    if (temp > 1)
    {
        lambda = lcm(lambda, temp - 1);
    }

    return lambda;
}

ll multiplicative_order(ll a, ll n)
{
    // If a and n are not coprime, order doesn't exist
    if (gcd(a, n) != 1)
        return -1;

    ll lambda = carmichael_lambda(n);
    ll order = lambda;
    ll temp_lambda = lambda;

    // Step 2: Find the exact order by removing prime factors from lambda
    // and checking if it still satisfies a^order % n == 1
    for (ll i = 2; i * i <= temp_lambda; i++)
    {
        if (temp_lambda % i == 0)
        {
            while (temp_lambda % i == 0)
            {
                temp_lambda /= i;
            }
            // Try dividing the order by the prime factor 'i'
            while (order % i == 0 && power(a, order / i, n) == 1)
            {
                order /= i;
            }
        }
    }

    // Check the remaining prime factor
    if (temp_lambda > 1)
    {
        while (order % temp_lambda == 0 && power(a, order / temp_lambda, n) == 1)
        {
            order /= temp_lambda;
        }
    }

    return order;
}
\end{lstlisting}
\subsubsection{discrete sqrt}
\begin{lstlisting}[language=C++]
/**
 *  THE ULTIMATE CP TEMPLATE BLUEPRINT: TONELLI-SHANKS ALGORITHM
 *
 * 1.  WHAT PROBLEM DOES THIS SOLVE?
 * - Keywords: "Quadratic Residue", "Modular Square Root", "x^2 = N mod P".
 * - Classic Scenarios: You are given an integer N and an odd prime P, and you
 *   need to find x such that x^2  N (mod P).
 * - The Magic: It factors P-1 as Q * 2^S. By finding a quadratic non-residue Z,
 *   it iteratively manipulates the exponents to isolate the square root. It is
 *   a beautiful but mathematically dense algorithm that you do not want to debug live.
 *
 * 2.  HOW TO USE IT (THE BLACK BOX)
 * - Query:
 *       long long root = tonelli_shanks(N, P);
 *       // Returns -1 if no square root exists.
 *       // If 'root' is a solution, 'P - root' is the other valid solution.
 *
 * - Complexity:
 *       Time: O(log^2 P) worst case, O(log P) on average. Extremely fast.
 *       Space: O(1)
 *
 * 3.  HOW TO ADAPT IT (THE DIALS & SWITCHES)
 * - Prime Modulo Only: This specifically requires P to be an odd prime. If you
 *   need to find roots modulo a composite number, you must use Tonelli-Shanks
 *   on its prime factors and combine the results using the Chinese Remainder Theorem (CRT).
 */

using ll = long long;

ll custom_pow(ll base, ll exp, ll mod)
{
    ll res = 1;
    base %= mod;
    while (exp > 0)
    {
        if (exp % 2 == 1)
            res = (res * base) % mod;
        base = (base * base) % mod;
        exp /= 2;
    }
    return res;
}

ll tonelli_shanks(ll n, ll p)
{
    if (n == 0)
        return 0;
    if (custom_pow(n, (p - 1) / 2, p) != 1)
        return -1;

    ll q = p - 1, s = 0;
    while (q % 2 == 0)
    {
        s++;
        q /= 2;
    }

    if (s == 1)
        return custom_pow(n, (p + 1) / 4, p);

    ll z = 2;
    while (custom_pow(z, (p - 1) / 2, p) != p - 1)
    {
        z++;
    }

    ll c = custom_pow(z, q, p);
    ll r = custom_pow(n, (q + 1) / 2, p);
    ll t = custom_pow(n, q, p);
    ll m = s;

    while (t != 1)
    {
        ll t2 = (t * t) % p;
        ll i = 1;
        for (; i < m; i++)
        {
            if (t2 == 1)
                break;
            t2 = (t2 * t2) % p;
        }

        ll b = c;
        for (int j = 0; j < m - i - 1; j++)
        {
            b = (b * b) % p;
        }

        r = (r * b) % p;
        c = (b * b) % p;
        t = (t * c) % p;
        m = i;
    }
    return r;
}
\end{lstlisting}
\subsection{Multiplicative functions}
\subsubsection{Min 25 sieve}
\begin{lstlisting}[language=C++]
/**
 *  MIN-25 SIEVE
 *
 * WHAT DOES THIS DO?
 * Computes the prefix sums of a multiplicative function f(x) in sub-linear time.
 * Specifically, it calculates sum(f(i)) for 1 <= i <= N in O(N^(3/4) / log N) or O(N^(1 - eps)).
 *
 * CONDITIONS FOR f(x):
 * 1. f(p) must be a polynomial (e.g., f(p) = 1, f(p) = p, f(p) = p^2).
 * 2. f(p^c) can be calculated quickly in O(1).
 *
 * This template computes the sum of primes and count of primes as the base layer,
 * which is the starting point for 90% of Min-25 problems.
 * Modify S(x, y) to match the f(p^c) of your specific problem.
 */

using ll = long long;

struct Min25
{
    ll n, sq, m;
    vector<int> primes;
    vector<bool> is_prime;

    // Arrays for the DP state
    vector<ll> w;         // Stores the distinct values of floor(n / i)
    vector<ll> g0;        // DP table for f(p) = 1 (Count of primes)
    vector<ll> g1;        // DP table for f(p) = p (Sum of primes)
    vector<int> id1, id2; // id1 for x <= sq, id2 for x > sq

    // Constructor initializes the structures and runs the prime sieve up to sqrt(N)
    Min25(ll _n) : n(_n)
    {
        sq = sqrt(n);
        is_prime.assign(sq + 1, true);
        is_prime[0] = is_prime[1] = false;

        for (int p = 2; p <= sq; p++)
        {
            if (is_prime[p])
            {
                primes.push_back(p);
                for (int i = p * p; i <= sq; i += p)
                    is_prime[i] = false;
            }
        }

        id1.assign(sq + 1, 0);
        id2.assign(sq + 1, 0);
        m = 0;

        // Step 1: Initialize all distinct values of floor(n / i)
        for (ll i = 1, j; i <= n; i = n / j + 1)
        {
            j = n / i;
            w.push_back(j);

            // Assign DP base cases: treating all numbers as if they were prime
            // g0 = j - 1 (count of numbers from 2 to j)
            // g1 = (j * (j + 1) / 2) - 1 (sum of numbers from 2 to j)
            g0.push_back(j - 1);
            g1.push_back(j * (j + 1) / 2 - 1);

            if (j <= sq)
                id1[j] = m;
            else
                id2[n / j] = m;
            m++;
        }

        // Step 2: The DP over primes (Filtering out composites)
        for (int p : primes)
        {
            ll p2 = 1LL * p * p;
            for (int i = 0; i < m && w[i] >= p2; i++)
            {
                ll next_val = w[i] / p;
                int id = (next_val <= sq) ? id1[next_val] : id2[n / next_val];

                // Subtract the contribution of composite numbers that have 'p' as their smallest prime factor
                // g(w[i]) = g(w[i]) - f(p) * (g(w[i] / p) - sum(f(p_j)) for j < current_prime)

                // For g0: f(p) = 1
                g0[i] -= 1LL * (g0[id] - g0[id1[p - 1]]);

                // For g1: f(p) = p
                g1[i] -= 1LL * p * (g1[id] - g1[id1[p - 1]]);
            }
        }
    }

    /**
     * Step 3: Compute the actual sum of f(x)
     * @param x: Current value (n / i)
     * @param y: Index of the minimum prime factor we are considering
     */
    ll S(ll x, int y)
    {
        if (x <= 1 || primes[y] > x)
            return 0;

        int id = (x <= sq) ? id1[x] : id2[n / x];

        // Initial value: sum of f(p) for all primes in the range (primes[y], x]
        // Example for f(p) = p: ans = g1[id] - g1[id1[primes[y]]]
        // MODIFICATION: Change this formula based on your f(p)
        ll ans = (g1[id] - g1[id1[primes[y]]]);

        // Add the contribution of composite numbers
        for (int i = y; i < primes.size() && 1LL * primes[i] * primes[i] <= x; i++)
        {
            ll p = primes[i];
            ll p_pow = p;
            ll p_next = p * p;

            for (int e = 1; p_next <= x; e++, p_pow = p_next, p_next *= p)
            {
                // ans += f(p^e) * S(x / p^e, i + 1) + f(p^(e+1))
                // MODIFICATION: Change the f(p^e) multipliers here
                ans += p_pow * S(x / p_pow, i + 1) + p_next;
            }
        }
        return ans;
    }

    ll solve()
    {
        // Adding f(1) at the end, which is usually 1
        return S(n, 0) + 1;
    }
};
\end{lstlisting}
\subsubsection{Power tower}
\begin{lstlisting}[language=C++]
/**
 *  THE ULTIMATE CP TEMPLATE BLUEPRINT: POWER TOWER (GENERALIZED EULER)
 *
 * 1.  WHAT PROBLEM DOES THIS SOLVE?
 * - Keywords: "Power Tower", "A^B^C^D mod M".
 * - Classic Scenarios: You are given an array of numbers and asked to compute their
 *   power tower (a[0] ^ (a[1] ^ (a[2] ^ ... ))) modulo M.
 * - The Magic: It applies the Generalized Euler Theorem to safely reduce exponents
 *   using the Phi function, effectively capping the recursion depth at O(log M).
 *
 * 2.  HOW TO USE IT (THE BLACK BOX)
 * - Query:
 *       long long ans = power_tower(arr, 0, M);
 *
 * - Complexity:
 *       Time: O(log M * sqrt(M))
 *       Space: O(log M) for the recursion stack.
 */

using ll = long long;

ll phi(ll n)
{
    ll res = n;
    for (ll i = 2; i * i <= n; i++)
    {
        if (n % i == 0)
        {
            while (n % i == 0)
                n /= i;
            res -= res / i;
        }
    }
    if (n > 1)
        res -= res / n;
    return res;
}

ll custom_pow(ll base, ll exp, ll mod)
{
    ll res = 1;
    base = (base >= mod) ? (base % mod + mod) : base;
    while (exp > 0)
    {
        if (exp % 2 == 1)
        {
            res *= base;
            res = (res >= mod) ? (res % mod + mod) : res;
        }
        base *= base;
        base = (base >= mod) ? (base % mod + mod) : base;
        exp /= 2;
    }
    return res;
}

ll power_tower(const vector<ll> &arr, int idx, ll mod)
{
    if (idx == arr.size() || mod == 1)
        return 1;
    if (arr[idx] == 1)
        return 1;
    ll next_mod = phi(mod);
    ll exponent = power_tower(arr, idx + 1, next_mod);
    return custom_pow(arr[idx], exponent, mod);
}
\end{lstlisting}
\subsubsection{phi}
\begin{lstlisting}[language=C++]
// computes the number of coprimes number with n up to n

const int N = 5000006;
ll phi[N];

void sieve()
{
    for (int i = 0; i < N; i++)
    {
        phi[i] = i;
    }
    for (int i = 2; i < N; i++)
    {
        if (phi[i] == i)
        {
            for (int j = i; j < N; j += i)
            {
                phi[j] -= phi[j] / i;
            }
        }
    }
}

ll phi(ll n)
{
    ll ret = n;
    for (ll i = 2; i * i <= n; i++)
    {
        if (n % i == 0)
        {
            while (n % i == 0)
                n /= i;
            ret -= ret / i;
        }
    }
    if (n > 1)
        ret -= ret / n;
    return ret;
}

\end{lstlisting}
\subsection{Primes And Divisors}
\subsubsection{Basic}
\begin{lstlisting}[language=C++]
struct NumberTheory {
    int n;
    vector<int> p, spf, phi_arr, fac, inv, finv;
    vector<bool> isPrime;
    vector<vector<int> > divisor;

    NumberTheory(int max_n = 1000000, bool pre_divs = false) {
        n = max_n;
        isPrime.assign(n + 1, true);
        spf.assign(n + 1, 0);
        phi_arr.assign(n + 1, 0);
        isPrime[0] = isPrime[1] = false;

        for (int i = 0; i <= n; i++) phi_arr[i] = i;

        for (int i = 2; i <= n; i++) {
            if (isPrime[i]) {
                p.push_back(i);
                spf[i] = i;
                phi_arr[i] -= phi_arr[i] / i;
            }
            for (int j = 0; j < (int) p.size() && i * p[j] <= n; j++) {
                isPrime[i * p[j]] = false;
                spf[i * p[j]] = p[j];
                if (i % p[j] == 0) {
                    phi_arr[i * p[j]] = phi_arr[i] * p[j];
                    break;
                } else {
                    phi_arr[i * p[j]] = phi_arr[i] * phi_arr[p[j]];
                }
            }
        }

        fac.assign(n + 1, 1);
        inv.assign(n + 1, 1);
        finv.assign(n + 1, 1);
        for (int i = 1; i <= n; i++) fac[i] = mul(fac[i - 1], i);
        for (int i = 2; i <= n; i++) inv[i] = (MOD - (MOD / i) * inv[MOD % i] % MOD) % MOD;
        for (int i = 2; i <= n; i++) finv[i] = mul(finv[i - 1], inv[i]);

        if (pre_divs) {
            divisor.resize(n + 1);
            for (int i = 1; i <= n; i++) {
                for (int j = i; j <= n; j += i) {
                    divisor[j].push_back(i);
                }
            }
        }
    }

    int add(int a, int b) { return ((a % MOD) + (b % MOD)) % MOD; }
    int sub(int a, int b) { return (((a - b) % MOD) + MOD) % MOD; }
    int mul(int a, int b) { return ((a % MOD) * (b % MOD)) % MOD; }

    int power(int b, int pow) {
        b %= MOD;
        int res = 1;
        while (pow) {
            if (pow & 1) res = mul(res, b);
            b = mul(b, b);
            pow >>= 1;
        }
        return res;
    }

    int modInv(int x) { return power(x, MOD - 2); }
    int divide(int a, int b) { return mul(a, modInv(b)); }

    int nCr(int x, int y) {
        if (x < 0 || y > x || y < 0) return 0;
        return mul(fac[x], mul(finv[y], finv[x - y]));
    }

    int nPr(int x, int y) {
        if (x < 0 || y > x || y < 0) return 0;
        return mul(fac[x], finv[x - y]);
    }

    int catalan(int x) { return mul(nCr(2 * x, x), modInv(x + 1)); }
    int StarsAndPars(int x, int k) { return nCr(x + k - 1, k - 1); }

    int nCr_no_mod(int x, int r) {
        if (r < 0 || x < 0 || r > x) return 0;
        int ret = 1;
        for (int i = 1; i <= min(r, x - r); i++) {
            ret *= x - i + 1;
            ret /= i;
        }
        return ret;
    }

    int nPr_no_mod(int x, int r) {
        if (r < 0 || x < 0 || r > x) return 0;
        int ret = 1;
        for (int i = 0; i < r; i++) ret *= (x - i);
        return ret;
    }

    int numberOfDivisors1ToN(int limit) {
        int sum = 0;
        for (int i = 1; i <= limit; i++) sum += limit / i;
        return sum;
    }

    bool is_prime(int x) {
        if (x <= 1) return false;
        if (x % 2 == 0) return x == 2;
        for (int i = 3; i * i <= x; i += 2) {
            if (x % i == 0) return false;
        }
        return true;
    }

    vector<int> factorization(int x) {
        vector<int> primes_factors;
        for (int i = 2; i * i <= x; i++) {
            while (x % i == 0) {
                primes_factors.push_back(i);
                x /= i;
            }
        }
        if (x > 1) primes_factors.push_back(x);
        return primes_factors;
    }

    vector<int> spf_factorization(int x) {
        vector<int> primes_factors;
        while (x != 1) {
            primes_factors.push_back(spf[x]);
            x /= spf[x];
        }
        return primes_factors;
    }

    map<int, int> spf_factorization_freq(int x) {
        map<int, int> frq;
        while (x != 1) {
            frq[spf[x]]++;
            x /= spf[x];
        }
        return frq;
    }

    int phi(int x) {
        int ret = x;
        for (int i = 2; i * i <= x; i++) {
            if (x % i == 0) {
                while (x % i == 0) x /= i;
                ret -= ret / i;
            }
        }
        if (x > 1) ret -= ret / x;
        return ret;
    }
};
    
\end{lstlisting}
\subsubsection{FastSumitionDivisor 1e16}
\begin{lstlisting}[language=C++]
*  THE ULTIMATE CP TEMPLATE BLUEPRINT: FAST SUM OF PROPER DIVISORS
 *
 * 1.  WHAT PROBLEM DOES THIS SOLVE?
 * - Keywords: "Sum of divisors", "Proper divisors", "Aliquot sum", "N up to 10^16".
 * - Classic Scenarios: You are given a massive number N (up to 10^16) and you need to find 
 *   the sum of all its divisors EXCEPT the number itself.

vector<int> Primes;
	
	void Sieve(int num)
	{
	    vector<bool> is_prime(num + 1, true);
	    is_prime[0], is_prime[1] = false;
	    for (int i = 2; i * i <= num; i++)
	    {
	        if (is_prime[i])
	        {
	            for (int j = i * i; j <= num; j += i)
	            is_prime[j] = false;
	        }
	    }
	    Primes.push_back(2);
	    for (int i = 3; i <= num; i += 2)
	    {
	        if (is_prime[i])
	        Primes.push_back(i);
	    }
	
	}
	
	void Hassan ()
	{
	//one integer between 1 and 1e16 inclusive
	    ll n; cin >> n;
	    ll ans = 1;
	    ll copy = n;
	    for (auto p : Primes)
	    {
	        if (1LL * p * p > n) 
	        break;
	        if (n % p == 0)
	        {
	            ll sum = 1, pow = 1;
	            while (n % p == 0)
	            {
	                pow *= p;
	                sum += pow;
	                n /= p;
	            }
	            ans *= sum;
	        }
	    }
	    if (n > 1) ans *= (n + 1);
	    cout << ans - copy << '\n'; /// note 
	}
	

\end{lstlisting}
\subsubsection{Linear sieve}
\begin{lstlisting}[language=C++]
const int N = 1e5 + 5;
vector<bool> isPrime(N, true);
vector<int> p;

void linear_sieve()
{
    isPrime[0] = isPrime[1] = 0;
    for (ll i = 2; i < N; i++)
    {
        if (isPrime[i] == 1)
            p.push_back(i);
        for (ll o = 0; o < p.size(); o++)
        {
            if (i * p[o] >= N)
                break;
            isPrime[i * p[o]] = 0;
            if (i % p[o] == 0)
                break;
        }
    }
}

\end{lstlisting}
\subsubsection{Miller pollard}
\begin{lstlisting}[language=C++]
/**
 *  THE ULTIMATE CP TEMPLATE BLUEPRINT: MILLER-RABIN & POLLARD'S RHO
 *
 * 1.  WHAT PROBLEM DOES THIS SOLVE?
 * - Keywords: "Primality Test", "Prime Factorization", "Count Divisors", "Massive N".
 * - Classic Scenarios: You need to check if N is prime, find all prime factors of N,
 *   or count the total number of divisors where N is up to 10^18. Standard O(sqrt(N))
 *   trial division will result in a Time Limit Exceeded (TLE) verdict.
 * - The Magic: Miller-Rabin deterministically checks primality for 64-bit integers
 *   in O(log N) time using 7 specific mathematical bases. If it's composite, Pollard's Rho
 *   steps in to find a prime factor in expected O(N^(1/4)) time. Together, they instantly
 *   break down 10^18 numbers.
 *
 * 2.  HOW TO USE IT (THE BLACK BOX)
 * - Primality Check:
 *       bool is_prime = MillerRabin(n);
 * - Prime Factorization:
 *       map<uint64_t, int> prime_factors;
 *       factorize(n, prime_factors); // Populates the map with {prime, frequency}
 * - Count Total Divisors:
 *       uint64_t total = count_divisors(n);
 *
 * - Complexity:
 *       Time: O(log N) for Miller-Rabin. Expected O(N^(1/4)) for Pollard's Rho.
 *       Space: O(log N) for recursion depth and the frequency map.
 *
 * 3.  HOW TO ADAPT IT (THE DIALS & SWITCHES)
 * - Type Aliases: This template strictly uses `u64` (uint64_t) and `u128` (__uint128_t)
 *   to prevent silent overflows during modular multiplication of massive numbers.
 *   Do not change these to standard `long long`.
 * - Randomness: `rand()` is used for the Rho polynomial `x^2 + c`. This is usually
 *   sufficient, but if you hit a malicious anti-hash test case, you can swap it
 *   for a faster, custom `mt19937_64` implementation.
 */

using u64 = uint64_t;
using u128 = __uint128_t;

u64 binpower(u64 base, u64 e, u64 mod)
{
    u64 result = 1;
    base %= mod;
    while (e)
    {
        if (e & 1)
            result = (u128)result * base % mod;
        base = (u128)base * base % mod;
        e >>= 1;
    }
    return result;
}

bool check_composite(u64 n, u64 a, u64 d, int s)
{
    u64 x = binpower(a, d, n);
    if (x == 1 || x == n - 1)
        return false;
    for (int r = 1; r < s; r++)
    {
        x = (u128)x * x % n;
        if (x == n - 1)
            return false;
    }
    return true;
}

bool MillerRabin(u64 n)
{
    if (n < 2)
        return false;
    int r = 0;
    u64 d = n - 1;
    while ((d & 1) == 0)
    {
        d >>= 1;
        r++;
    }
    // These 7 bases are deterministically sufficient for all n < 2^64
    for (u64 a : {2, 325, 9375, 28178, 450775, 9780504, 1795265022})
    {
        if (n == a)
            return true;
        if (check_composite(n, a, d, r))
            return false;
    }
    return true;
}

u64 pollard_rho(u64 n)
{
    if (n % 2 == 0)
        return 2;
    if (MillerRabin(n))
        return n;

    u64 x = 2, y = 2, d = 1, c = 1;
    auto f = [&](u64 x, u64 n, u64 c)
    {
        return (u64)(((u128)x * x + c) % n);
    };

    while (d == 1)
    {
        x = f(x, n, c);
        y = f(f(y, n, c), n, c);
        d = __gcd(x > y ? x - y : y - x, n);
        if (d == n)
        {
            x = rand() % (n - 2) + 2;
            y = x;
            c = rand() % (n - 1) + 1;
            d = 1;
        }
    }
    return d;
}

void factorize(u64 n, map<u64, int> &prime_factors)
{
    if (n == 1)
        return;
    if (MillerRabin(n))
    {
        prime_factors[n]++;
        return;
    }
    u64 divisor = pollard_rho(n);
    factorize(divisor, prime_factors);
    factorize(n / divisor, prime_factors);
}

u64 count_divisors(u64 n)
{
    map<u64, int> prime_factors;
    factorize(n, prime_factors);
    u64 total_divisors = 1;
    for (auto &pf : prime_factors)
    {
        total_divisors *= (pf.second + 1);
    }
    return total_divisors;
}
\end{lstlisting}
\subsubsection{Number of divisors 1e18}
\begin{lstlisting}[language=C++]
/**
 *  THE ULTIMATE CP TEMPLATE BLUEPRINT: FAST DIVISOR COUNTING (WHEEL FACTORIZATION)
 *
 * 1.  WHAT PROBLEM DOES THIS SOLVE?
 * - Keywords: "Number of divisors of N", "N up to 10^18", "Count divisors without full factorization".
 * - Classic Scenarios: You are given a massive number N (up to 10^14 or 10^18) and you need to
 *   find exactly how many divisors it has, but O(sqrt(N)) is too slow, and you can't build a Sieve
 *   because N is too large.
 * - The Magic: "Wheel Factorization". Instead of checking every number up to sqrt(N), we manually
 *   factor out the first 8 primes (2, 3, 5, 7, 11, 13, 17, 19). The product of these primes is ~9.6 million.
 *   For the remaining part of N, we ONLY check divisors that are coprime to these 8 primes.
 *   This eliminates about ~83% of the useless checks in the O(sqrt(N)) loop!
 *
 * 2.  HOW TO USE IT (THE BLACK BOX)
 * - Initialization: Create an instance globally or in `main`. It will precompute the "wheel spokes" ONCE.
 *       FastDivisorCounter fdc;
 *
 * - Queries: Pass your number N (up to 10^18) to get the number of divisors.
 *       ll total_divisors = fdc.count(N);
 *
 * - Complexity:
 *       Time (Initialization): Precomputing the wheel takes O(P) where P is ~9.6 * 10^6 (Done once in ~15ms).
 *       Time (Per Query): roughly O(sqrt(N) / 9.6*10^6). Lightning fast for huge queries!
 *       Space: ~6 MB for the precomputed wheel.
 *
 * 3.  HOW TO ADAPT IT (THE DIALS & SWITCHES)
 * - If Memory Limit is EXTREMELY tight (<10 MB), you can remove the `spokes` vector and calculate the
 *   coprime numbers inside the `count` function loop like your original code, but only if Q=1.
 *   For Q > 1, the precomputed `spokes` vector is vastly superior.
 */

#include <bits/stdc++.h>
using namespace std;
#define ll long long
struct FastDivisorCounter
{
private:
    vector<int> primes;
    ll wheel_size;
    vector<ll> spokes; // Stores numbers from 1 to wheel_size that are coprime to the base primes
public:
    FastDivisorCounter()
    {
        primes = {2, 3, 5, 7, 11, 13, 17, 19};
        wheel_size = 1;

        for (int p : primes)
        {
            wheel_size *= p;
        }
        // ? PRECOMPUTATION MAGIC: We find all valid "spokes" ONCE.
        // There are exactly 1,658,880 such numbers. By saving them, we skip millions
        // of useless iterations in future queries.
        for (ll i = 1; i < wheel_size; i++)
        {
            bool is_coprime = true;
            for (int p : primes)
            {
                if (i % p == 0)
                {
                    is_coprime = false;
                    break;
                }
            }
            if (is_coprime)
            {
                spokes.push_back(i);
            }
        }
    }
    ll count(ll n)
    {
        ll ans = 1;
        // Step 1: Remove all factors of the base primes and count their contribution
        for (int p : primes)
        {
            int c = 0;
            while (n % p == 0)
            {
                n /= p;
                c++;
            }
            // Standard divisor formula: multiply by (power + 1)
            ans *= (c + 1);
        }
        // Step 2: The remaining part of 'n' is now completely coprime to our base primes.
        // We only check potential divisors using our precomputed wheel spokes.
        ll all = 0;
        for (ll spoke : spokes)
        {
            ll o = spoke;
            // ? THE SPEED JUMP: We jump by `wheel_size` (9,699,690) instead of +1 or +2!
            for (; o * o < n; o += wheel_size)
            {
                if (n % o == 0)
                {
                    all += 2; // Found pair of divisors: `o` and `n/o`
                }
            }

            // Check perfect square
            if (o * o == n)
            {
                all++;
            }
        }
        // The total divisors of the original N is the product of the divisors
        // of the base-prime part and the divisors of the remaining coprime part.
        return ans * all;
    }
};
void solve()
{
    int t;
    if (!(cin >> t))
        return;
    // Initialize OUTSIDE the query loop so the wheel is built only once!
    FastDivisorCounter fdc;
    while (t--)
    {
        ll n;
        cin >> n;
        cout << fdc.count(n) << "\n";
    }
}

\end{lstlisting}
\subsubsection{Prime coutning}
\begin{lstlisting}[language=C++]
/**
 *  THE ULTIMATE CP TEMPLATE BLUEPRINT: PRIME COUNTING FUNCTION (LUCY-HEDGEHOG)
 *
 * 1.  WHAT PROBLEM DOES THIS SOLVE?
 * - Keywords: "Count primes up to N", "Pi(x)", "Lucy-Hedgehog", "Min-25 Sieve".
 * - Classic Scenarios: You are given N up to 10^11 and asked exactly how many
 *   prime numbers are less than or equal to N.
 * - The Magic: A standard sieve up to 10^11 will instantly crash with Memory Limit
 *   Exceeded and Time Limit Exceeded. The Lucy-Hedgehog dynamic programming approach
 *   computes this mathematically without iterating over or storing the actual primes!
 *
 * 2.  HOW TO USE IT (THE BLACK BOX)
 * - Query:
 *       uint64_t total_primes = prime_count(N);
 *
 * - Complexity:
 *       Time: O(N^(3/4))  Runs almost instantly for N = 10^10, and easily
 *       passes 1-second limits for N = 10^11.
 *       Space: O(sqrt(N))  Extremely memory-efficient.
 *
 * 3.  HOW TO ADAPT IT (THE DIALS & SWITCHES)
 * - Array Sizes: It only creates vectors of size sqrt(N). For N = 10^11, this is
 *   just an array of size 3.16 * 10^5, taking mere megabytes of RAM.
 */

using u64 = uint64_t;

u64 prime_count(u64 N)
{
    if (N <= 1)
        return 0;
    u64 v = sqrt(N);
    vector<u64> S0(v + 1), S1(v + 1);
    for (u64 i = 1; i <= v; i++)
    {
        S0[i] = i - 1;
        S1[i] = N / i - 1;
    }
    for (u64 p = 2; p <= v; p++)
    {
        if (S0[p] == S0[p - 1])
            continue;
        u64 p_cnt = S0[p - 1];
        u64 q = p * p;
        u64 end = min(v, N / q);
        for (u64 i = 1; i <= end; i++)
        {
            u64 d = i * p;
            if (d <= v)
                S1[i] -= S1[d] - p_cnt;
            else
                S1[i] -= S0[N / d] - p_cnt;
        }
        for (u64 i = v; i >= q; i--)
        {
            S0[i] -= S0[i / p] - p_cnt;
        }
    }
    return S1[1];
}
\end{lstlisting}
\subsubsection{Sieve 1e9}
\begin{lstlisting}[language=C++]
#include <bits/stdc++.h>
using namespace std;

// credit: min_25
// takes 0.5s for n = 1e9
vector<int> sieve(const int N, const int Q = 17, const int L = 1 << 15)
{
    static const int rs[] = {1, 7, 11, 13, 17, 19, 23, 29};
    struct P
    {
        P(int p) : p(p)
        {
        }

        int p;
        int pos[8];
    };
    auto approx_prime_count = [](const int N) -> int
    {
        return N > 60184
                   ? N / (log(N) - 1.1)
                   : max(1., N / (log(N) - 1.11)) + 1;
    };

    const int v = sqrt(N), vv = sqrt(v);
    vector<bool> isp(v + 1, true);
    for (int i = 2; i <= vv; ++i)
        if (isp[i])
        {
            for (int j = i * i; j <= v; j += i)
                isp[j] = false;
        }

    const int rsize = approx_prime_count(N + 30);
    vector<int> primes = {2, 3, 5};
    int psize = 3;
    primes.resize(rsize);

    vector<P> sprimes;
    size_t pbeg = 0;
    int prod = 1;
    for (int p = 7; p <= v; ++p)
    {
        if (!isp[p])
            continue;
        if (p <= Q)
            prod *= p, ++pbeg, primes[psize++] = p;
        auto pp = P(p);
        for (int t = 0; t < 8; ++t)
        {
            int j = (p <= Q) ? p : p * p;
            while (j % 30 != rs[t])
                j += p << 1;
            pp.pos[t] = j / 30;
        }
        sprimes.push_back(pp);
    }

    vector<unsigned char> pre(prod, 0xFF);
    for (size_t pi = 0; pi < pbeg; ++pi)
    {
        auto pp = sprimes[pi];
        const int p = pp.p;
        for (int t = 0; t < 8; ++t)
        {
            const unsigned char m = ~(1 << t);
            for (int i = pp.pos[t]; i < prod; i += p)
                pre[i] &= m;
        }
    }

    const int block_size = (L + prod - 1) / prod * prod;
    vector<unsigned char> block(block_size);
    unsigned char *pblock = block.data();
    const int M = (N + 29) / 30;

    for (int beg = 0; beg < M; beg += block_size, pblock -= block_size)
    {
        int end = min(M, beg + block_size);
        for (int i = beg; i < end; i += prod)
        {
            copy(pre.begin(), pre.end(), pblock + i);
        }
        if (beg == 0)
            pblock[0] &= 0xFE;
        for (size_t pi = pbeg; pi < sprimes.size(); ++pi)
        {
            auto &pp = sprimes[pi];
            const int p = pp.p;
            for (int t = 0; t < 8; ++t)
            {
                int i = pp.pos[t];
                const unsigned char m = ~(1 << t);
                for (; i < end; i += p)
                    pblock[i] &= m;
                pp.pos[t] = i;
            }
        }
        for (int i = beg; i < end; ++i)
        {
            for (int m = pblock[i]; m > 0; m &= m - 1)
            {
                primes[psize++] = i * 30 + rs[__builtin_ctz(m)];
            }
        }
    }
    assert(psize <= rsize);
    while (psize > 0 && primes[psize - 1] > N)
        --psize;
    primes.resize(psize);
    return primes;
}

int32_t main()
{
    ios_base::sync_with_stdio(0);
    cin.tie(0);
    int n, a, b;
    cin >> n >> a >> b;
    auto primes = sieve(n);
    vector<int> ans;
    for (int i = b; i < primes.size() && primes[i] <= n; i += a)
        ans.push_back(primes[i]);
    cout << primes.size() << ' ' << ans.size() << '\n';
    for (auto x : ans)
        cout << x << ' ';
    cout << '\n';
    return 0;
}

\end{lstlisting}
\subsubsection{segmented sieve}
\begin{lstlisting}[language=C++]
/**
 *  THE ULTIMATE CP TEMPLATE BLUEPRINT: SEGMENTED SIEVE
 *
 * 1.  WHAT PROBLEM DOES THIS SOLVE?
 * - Keywords: "Primes in Range [L, R]", "Massive Bounds".
 * - Classic Scenarios: Find all primes or the number of primes between L and R.
 *   L and R can be as massive as 10^12, but the difference (R - L) is small (<= 10^6).
 * - The Magic: A regular sieve cannot allocate an array up to 10^12. The segmented
 *   sieve only sieves up to sqrt(R), caches those base primes, and then "shifts"
 *   the sieving logic directly to the target window [L, R], ignoring the void before L.
 *
 * 2.  HOW TO USE IT (THE BLACK BOX)
 * - Query:
 *       vector<uint64_t> primes = segmented_sieve(L, R);
 *
 * - Complexity:
 *       Time: O((R - L + sqrt(R)) * log(log(R)))
 *       Space: O(sqrt(R) + (R - L))
 *
 * 3.  HOW TO ADAPT IT (THE DIALS & SWITCHES)
 * - Edge Cases Guarded: Safely handles constraints where L = 0 or L = 1
 *   (which are not prime but often trip up basic segmented implementations).
 */

using u64 = uint64_t;

vector<u64> segmented_sieve(u64 L, u64 R)
{
    if (L > R)
        return {};
    u64 lim = sqrt(R);
    vector<bool> mark(lim + 1, false);
    vector<u64> primes;
    for (u64 i = 2; i <= lim; ++i)
    {
        if (!mark[i])
        {
            primes.push_back(i);
            for (u64 j = i * i; j <= lim; j += i)
                mark[j] = true;
        }
    }

    vector<bool> isPrime(R - L + 1, true);
    for (u64 p : primes)
    {
        u64 start = max(p * p, (L + p - 1) / p * p);
        for (u64 j = start; j <= R; j += p)
            isPrime[j - L] = false;
    }

    if (L == 0)
    {
        isPrime[0] = false;
        if (R >= 1)
            isPrime[1] = false;
    }
    else if (L == 1)
    {
        isPrime[0] = false;
    }

    vector<u64> range_primes;
    for (u64 i = 0; i <= R - L; ++i)
    {
        if (isPrime[i])
            range_primes.push_back(L + i);
    }
    return range_primes;
}
\end{lstlisting}
\subsubsection{sum numberOfDivisors}
\begin{lstlisting}[language=C++]
/**
 *  THE ULTIMATE CP TEMPLATE BLUEPRINT: CUBIC ROOT DIVISOR SUM
 *
 * 1.  WHAT PROBLEM DOES THIS SOLVE?
 * - Keywords: "Sum of Divisors", "Sigma0(i)", "Sum of floor(N/i)".
 * - Classic Scenarios: You are asked to compute sum_{i=1}^{N} d(i), which is
 *   mathematically equivalent to sum_{i=1}^{N} floor(N / i), for a staggering N = 10^18.
 * - The Magic: The naive O(N) loop crashes. The standard Dirichlet hyperbola
 *   optimization takes O(sqrt(N)), which runs in ~1 second for 10^12 but fails at 10^18.
 *   This extremely advanced algorithm uses Stern-Brocot tree fraction approximation
 *   to step through the hyperbola in blocks, reducing the time to roughly O(N^(1/3))!
 *
 * 2.  HOW TO USE IT (THE BLACK BOX)
 * - Query:
 *       __uint128_t ans = sum_sigma0(N);
 *       // Remember to write a custom print function for __uint128_t if outputting directly.
 *
 * - Complexity:
 *       Time: ~O(N^(1/3)). Easily executes N = 10^18 in < 100 milliseconds.
 *       Space: O(log N) for the stack depth.
 *
 * 3.  HOW TO ADAPT IT (THE DIALS & SWITCHES)
 * - 128-bit Types: Be very careful if modifying the inner loop logic. The condition
 *   `uint128(x) * x * dy >= uint128(n) * dx` uses 128-bit integers
 *   because cubing boundaries near 10^18 guarantees overflow in standard `uint64_t`.
 */

using uint32 = unsigned int;
using uint64 = unsigned long long;
using uint128 = __uint128_t;

uint128 sum_sigma0(uint64 n)
{
    auto out = [n](uint64 x, uint32 y)
    { return x * y > n; };
    auto cut = [n](uint64 x, uint32 dx, uint32 dy)
    { return uint128(x) * x * dy >= uint128(n) * dx; };
    const uint64 sn = sqrtl(n);
    const uint64 cn = pow(n, 0.34);
    uint64 x = n / sn;
    uint32 y = n / x + 1;
    uint128 ret = 0;
    stack<pair<uint32, uint32>> st;
    st.emplace(1, 0);
    st.emplace(1, 1);
    while (true)
    {
        uint32 lx, ly;
        tie(lx, ly) = st.top();
        st.pop();
        while (out(x + lx, y - ly))
        {
            ret += x * ly + uint64(ly + 1) * (lx - 1) / 2;
            x += lx, y -= ly;
        }
        if (y <= cn)
            break;
        uint32 rx = lx, ry = ly;
        while (true)
        {
            tie(lx, ly) = st.top();
            if (out(x + lx, y - ly))
                break;
            rx = lx, ry = ly;
            st.pop();
        }
        while (true)
        {
            uint32 mx = lx + rx, my = ly + ry;
            if (out(x + mx, y - my))
            {
                st.emplace(lx = mx, ly = my);
            }
            else
            {
                if (cut(x + mx, lx, ly))
                    break;
                rx = mx, ry = my;
            }
        }
    }
    for (--y; y > 0; --y)
        ret += n / y;
    return ret * 2 - sn * sn;
}
\end{lstlisting}
\section{Combinatorics}
\subsection{Main}
\subsubsection{Basic}
\begin{lstlisting}[language=C++]
// ============================================================================
// MODULE 1: MODULAR ARITHMETIC & CORE COMBINATORICS
// ============================================================================
struct Combinatorics
{
    int N;
    vector<ll> fact, invFact;

    // Fast Modular Exponentiation: Computes (base^exp) % MOD in O(log exp)
    ll power(ll base, ll exp)
    {
        ll res = 1;
        base %= MOD;
        while (exp > 0)
        {
            if (exp % 2 == 1)
                res = (res * base) % MOD;
            base = (base * base) % MOD;
            exp /= 2;
        }
        return res;
    }

    // Modular Multiplicative Inverse: Computes (1 / n) % MOD in O(log MOD)
    // Uses Fermat's Little Theorem (Requires MOD to be prime)
    ll modInverse(ll n)
    {
        return power(n, MOD - 2);
    }

    // Initialize Factorials and Inverse Factorials in O(N)
    // CALL THIS ONCE IN MAIN BEFORE SOLVING TEST CASES!
    void init(int n)
    {
        N = n;
        fact.resize(N + 1);
        invFact.resize(N + 1);

        fact[0] = 1;
        invFact[0] = 1;
        for (int i = 1; i <= N; i++)
        {
            fact[i] = (fact[i - 1] * i) % MOD;
        }

        // Compute the inverse of the largest factorial, then work backwards
        invFact[N] = modInverse(fact[N]);
        for (int i = N - 1; i >= 1; i--)
        {
            invFact[i] = (invFact[i + 1] * (i + 1)) % MOD;
        }
    }

    // Choose: Number of ways to pick r items from n items. O(1)
    ll nCr(int n, int r)
    {
        if (r < 0 || r > n)
            return 0;
        return fact[n] * invFact[r] % MOD * invFact[n - r] % MOD;
    }

    // Permutation: Number of ways to arrange r items out of n items. O(1)
    ll nPr(int n, int r)
    {
        if (r < 0 || r > n)
            return 0;
        return fact[n] * invFact[n - r] % MOD;
    }

    // Stars and Bars: Ways to distribute n identical items into k distinct bins. O(1)
    // Example: Solutions to x1 + x2 + ... + xk = n (where xi >= 0)
    ll starsAndBars(int n, int k)
    {
        if (n == 0 && k == 0)
            return 1;
        return nCr(n + k - 1, k - 1);
    }

    // Stars and Bars (Positive): Distribute n identical items into k distinct bins,
    // such that every bin gets AT LEAST ONE item. O(1)
    // Example: Solutions to x1 + x2 + ... + xk = n (where xi >= 1)
    ll starsAndBarsPositive(int n, int k)
    {
        if (n < k)
            return 0;
        return nCr(n - 1, k - 1);
    }

    // Catalan Numbers: O(1)
    // Solves: Valid parenthesis expressions, binary trees, polygon triangulations.
    // Formula: C(n) = (1 / (n + 1)) * (2n C n)
    ll catalan(int n)
    {
        if (n < 0)
            return 0;
        ll res = nCr(2 * n, n);
        res = (res * modInverse(n + 1)) % MOD;
        return res;
    }
};

// ============================================================================
// MODULE 2: ADVANCED COUNTING (DERANGEMENTS, STIRLING, BELL)
// ============================================================================
struct AdvancedCombinatorics
{
    // Derangements: Number of permutations of N elements where NO element is in its original position.
    // O(N) precomputation. Formula: D(n) = (n-1) * (D(n-1) + D(n-2))
    vector<ll> derangement;

    void initDerangements(int n)
    {
        derangement.assign(n + 1, 0);
        if (n >= 0)
            derangement[0] = 1;
        if (n >= 1)
            derangement[1] = 0;
        for (int i = 2; i <= n; i++)
        {
            derangement[i] = ((i - 1) * (derangement[i - 1] + derangement[i - 2]) % MOD) % MOD;
        }
    }

    // Stirling Numbers of the First Kind: S1[n][k]
    // Solves: Number of permutations of 'n' elements that contain exactly 'k' disjoint cycles.
    // O(N^2) precomputation.
    vector<vector<ll>> S1;

    void initStirling1(int n)
    {
        S1.assign(n + 1, vector<ll>(n + 1, 0));
        S1[0][0] = 1;
        for (int i = 1; i <= n; i++)
        {
            for (int j = 1; j <= i; j++)
            {
                S1[i][j] = (S1[i - 1][j - 1] + (i - 1) * S1[i - 1][j]) % MOD;
            }
        }
    }

    // Stirling Numbers of the Second Kind: S2[n][k]
    // Solves: Number of ways to partition a set of 'n' DISTINCT elements into exactly 'k' non-empty indistinguishable subsets.
    // O(N^2) precomputation.
    vector<vector<ll>> S2;

    void initStirling2(int n)
    {
        S2.assign(n + 1, vector<ll>(n + 1, 0));
        S2[0][0] = 1;
        for (int i = 1; i <= n; i++)
        {
            for (int j = 1; j <= i; j++)
            {
                S2[i][j] = (S2[i - 1][j - 1] + j * S2[i - 1][j]) % MOD;
            }
        }
    }

    // Bell Numbers: Bell[n]
    // Solves: Total number of ways to partition a set of 'n' elements into ANY number of non-empty subsets.
    // O(N^2) precomputation (Sum of Stirling2(n, k) for k = 0 to n).
    vector<ll> bell;

    void initBell(int n)
    {
        if (S2.empty())
            initStirling2(n);
        bell.assign(n + 1, 0);
        for (int i = 0; i <= n; i++)
        {
            for (int j = 0; j <= i; j++)
            {
                bell[i] = (bell[i] + S2[i][j]) % MOD;
            }
        }
    }
};

// ============================================================================
// MODULE 3: LUCAS THEOREM (LARGE N, SMALL PRIME MODULO)
// ============================================================================
// Solves: nCr % P when N and R are massive (up to 10^18) but P is small (up to 10^5) and P is PRIME.
// Note: This requires a separate Combinatorics instance initialized up to P.
ll lucas_nCr(ll n, ll r, int p, Combinatorics &combo)
{
    if (r < 0 || r > n)
        return 0;
    if (r == 0)
        return 1;

    // Split n and r into base P representation
    ll ni = n % p;
    ll ri = r % p;

    if (ni < ri)
        return 0; // Optimization: P divides nCr

    // nCr = Product of (n_i C r_i) mod P
    return (combo.nCr(ni, ri) * lucas_nCr(n / p, r / p, p, combo)) % p;
}

\end{lstlisting}
\subsubsection{Pascal triangle}
\begin{lstlisting}[language=C++]
ll c[n][n]={};
c[0][0]=1;
for(int i=1;i<n;i++)
{
    c[i][0]=c[i][i]=1;
    for(int o=1;o<i;o++)
    {
        c[i][o]=c[i-1][o-1]+c[i-1][o];
    }
}

\end{lstlisting}
\subsection{PIE}
\subsubsection{multiples}
\begin{lstlisting}[language=C++]
/**
 *  THE ULTIMATE CP TEMPLATE BLUEPRINT: INCLUSION-EXCLUSION (PIE) ON MULTIPLES
 *
 * 1.  WHAT PROBLEM DOES THIS SOLVE?
 * - Keywords: "Divisible by any", "Coprime to array", "Multiples below X".
 * - Classic Scenarios: You are given a target number X, and an array of K integers.
 *   You need to find how many numbers from 1 to X are divisible by AT LEAST ONE
 *   of the integers in the array.
 * - The Magic: It systematically adds numbers divisible by 1 element, subtracts
 *   those divisible by 2 elements, adds for 3, etc. A naive bitmask loop goes out
 *   of bounds (LCM overflow) quickly. This DFS safely prunes branches if the LCM
 *   ever exceeds X, making it insanely fast and overflow-proof.
 *
 * 2.  HOW TO USE IT (THE BLACK BOX)
 * - Query:
 *       long long ans = count_divisible(X, A);
 *       // To find numbers COPRIME to the array: ans = X - count_divisible(X, A);
 *
 * 3.  HOW TO ADAPT IT (THE DIALS & SWITCHES)
 * - Array Cleaning: For maximum speed, it is highly recommended to sort the array `A`,
 *   remove duplicates, and remove any number that is a multiple of another number
 *   in the array before calling this function.
 */

using ll = long long;

ll pie_dfs(int idx, ll current_lcm, int elements_chosen, ll max_val, const vector<ll> &a)
{
    if (idx == a.size())
    {
        if (elements_chosen == 0)
            return 0;
        if (elements_chosen % 2 == 1)
            return max_val / current_lcm;
        return -(max_val / current_lcm);
    }

    ll ans = pie_dfs(idx + 1, current_lcm, elements_chosen, max_val, a);

    ll g = __gcd(current_lcm, a[idx]);
    if (current_lcm / g <= max_val / a[idx])
    {
        ans += pie_dfs(idx + 1, (current_lcm / g) * a[idx], elements_chosen + 1, max_val, a);
    }

    return ans;
}

ll count_divisible(ll max_val, const vector<ll> &a)
{
    return pie_dfs(0, 1, 0, max_val, a);
}
\end{lstlisting}
\subsection{Partition}
\subsubsection{Integer partition}
\begin{lstlisting}[language=C++]
/**
 *  THE ULTIMATE CP TEMPLATE BLUEPRINT: INTEGER PARTITION
 *
 * 1.  WHAT PROBLEM DOES THIS SOLVE?
 * - Keywords: "Ways to sum to N", "Unordered integer partitions".
 * - Classic Scenarios: Find the number of distinct ways to represent the integer N
 *   as a sum of positive integers (e.g., 4 = 4, 3+1, 2+2, 2+1+1, 1+1+1+1).
 * - The Magic: A standard coin-change DP takes O(N^2), which crashes for N = 10^5.
 *   Euler's Pentagonal Number Theorem computes the partition function P(N) in O(N sqrt(N))
 *   by subtracting and adding previous P() values at specifically spaced pentagonal
 *   intervals: P(k) = P(k-1) + P(k-2) - P(k-5) - P(k-7) + P(k-12) + ...
 *
 * 2.  HOW TO USE IT (THE BLACK BOX)
 * - Query:
 *       vector<long long> P = partition_function(N, MOD);
 *       // P[i] contains the number of partitions for integer i.
 *
 * - Complexity:
 *       Time: O(N * sqrt(N))
 *       Space: O(N)
 *
 * 3.  HOW TO ADAPT IT (THE DIALS & SWITCHES)
 * - Array Precomputation: It returns the entire array up to N. If the problem has
 *   multiple test cases querying different N's, run this function ONCE globally for
 *   the maximum possible N, then answer each query in O(1).
 */

using ll = long long;

vector<ll> partition_function(int n, ll mod)
{
    vector<ll> p(n + 1, 0);
    p[0] = 1;
    for (int i = 1; i <= n; i++)
    {
        for (int k = 1;; k++)
        {
            int p1 = k * (3 * k - 1) / 2;
            int p2 = k * (3 * k + 1) / 2;
            if (p1 > i && p2 > i)
                break;

            ll sign = (k % 2 == 1) ? 1 : -1;

            if (p1 <= i)
                p[i] = (p[i] + sign * p[i - p1]) % mod;
            if (p2 <= i)
                p[i] = (p[i] + sign * p[i - p2]) % mod;
        }
        p[i] = (p[i] + mod) % mod;
    }
    return p;
}
\end{lstlisting}
\subsection{nCr}
\subsubsection{nCr Mod any mod}
\begin{lstlisting}[language=C++]
/**
 *  THE ULTIMATE CP TEMPLATE BLUEPRINT: nCr MODULO ANY MOD
 *
 * 1.  WHAT PROBLEM DOES THIS SOLVE?
 * - Keywords: "nCr", "Composite Modulo", "No Modular Inverse".
 * - Classic Scenarios: Compute nCr(N, K) modulo M where M is NOT a prime number
 *   (e.g., M = 10^9, or some random composite).
 * - The Magic: Since M is composite, you cannot divide by factorials (no inverse).
 *   This algorithm factors M into prime powers p^k. It then counts the exact
 *   number of times p divides N!, K!, and (N-K)!, isolates the coprime parts,
 *   calculates the binomial coefficient modulo p^k, and finally uses the Chinese
 *   Remainder Theorem (CRT) to stitch them all back together into a single answer.
 *
 * 2.  HOW TO USE IT (THE BLACK BOX)
 * - Query:
 *       long long ans = nCr_arbitrary_mod(N, K, M);
 *
 * - Complexity:
 *       Time: O(sum(p^k) + log N). Extremely fast for M <= 10^6.
 *       Space: O(sum(p^k)) to store factorials modulo p^k.
 *
 * 3.  HOW TO ADAPT IT (THE DIALS & SWITCHES)
 * - Dependencies: This heavily relies on your existing `extgcd` and `crt` templates
 *   from your Modular Arithmetic section. Make sure they are included above this.
 */

using ll = long long;

ll custom_pow(ll base, ll exp, ll mod)
{
    ll res = 1;
    base %= mod;
    while (exp > 0)
    {
        if (exp % 2 == 1)
            res = (res * base) % mod;
        base = (base * base) % mod;
        exp /= 2;
    }
    return res;
}

ll extgcd(ll a, ll b, ll &x, ll &y);
pair<ll, ll> crt(const vector<ll> &a, const vector<ll> &m);

ll modInverse(ll a, ll m)
{
    ll x, y;
    extgcd(a, m, x, y);
    return (x % m + m) % m;
}

ll fact_mod_pk(ll n, ll p, ll pk, int &e)
{
    e = 0;
    if (n == 0)
        return 1;
    ll res = 1;
    for (ll i = 2; i <= pk; i++)
    {
        if (i % p != 0)
            res = (res * i) % pk;
    }
    res = custom_pow(res, n / pk, pk);
    for (ll i = 2; i <= n % pk; i++)
    {
        if (i % p != 0)
            res = (res * i) % pk;
    }
    int next_e;
    res = (res * fact_mod_pk(n / p, p, pk, next_e)) % pk;
    e += n / p + next_e;
    return res;
}

ll nCr_pk(ll n, ll r, ll p, ll pk)
{
    if (n < r || r < 0)
        return 0;
    int e1, e2, e3;
    ll num = fact_mod_pk(n, p, pk, e1);
    ll den1 = fact_mod_pk(r, p, pk, e2);
    ll den2 = fact_mod_pk(n - r, p, pk, e3);
    ll p_pow = custom_pow(p, e1 - e2 - e3, pk);
    if (p_pow == 0)
        return 0;
    ll ans = (num * modInverse((den1 * den2) % pk, pk)) % pk;
    return (ans * p_pow) % pk;
}

ll nCr_arbitrary_mod(ll n, ll r, ll m)
{
    if (n < r || r < 0)
        return 0;
    if (m == 1)
        return 0;
    vector<ll> rem, mod;
    ll temp = m;
    for (ll i = 2; i * i <= temp; i++)
    {
        if (temp % i == 0)
        {
            ll pk = 1;
            while (temp % i == 0)
            {
                pk *= i;
                temp /= i;
            }
            rem.push_back(nCr_pk(n, r, i, pk));
            mod.push_back(pk);
        }
    }
    if (temp > 1)
    {
        rem.push_back(nCr_pk(n, r, temp, temp));
        mod.push_back(temp);
    }
    return crt(rem, mod).first;
}
\end{lstlisting}
\subsubsection{nCr over congurence class}
\begin{lstlisting}[language=C++]
/**
 *  THE ULTIMATE CP TEMPLATE BLUEPRINT: nCr OVER CONGRUENCE CLASS
 *
 * 1.  WHAT PROBLEM DOES THIS SOLVE?
 * - Keywords: "Sum of nCr where r mod M = K".
 * - Classic Scenarios: You need to evaluate sum(nCr(N, i)) for all i where
 *   i % M == K. N can be up to 10^18, but M is small (M <= 100).
 * - The Magic: It maps the combinatorial states into a matrix. If you have an
 *   array tracking the sum mod M, multiplying it by a transition matrix simulates
 *   increasing N by 1. By exponentiating the matrix to the power of N in
 *   O(M^3 log N), you instantly compute the exact sum across the entire boundary.
 *
 * 2.  HOW TO USE IT (THE BLACK BOX)
 * - Query:
 *       long long ans = nCr_congruence(N, K, M, MOD);
 *
 * - Complexity:
 *       Time: O(M^3 * log N)
 *       Space: O(M^2) for the matrix.
 */

using ll = long long;

using Matrix = vector<vector<ll>>;

Matrix multiply_matrix(const Matrix &A, const Matrix &B, ll mod)
{
    int n = A.size();
    Matrix C(n, vector<ll>(n, 0));
    for (int i = 0; i < n; i++)
    {
        for (int k = 0; k < n; k++)
        {
            for (int j = 0; j < n; j++)
            {
                C[i][j] = (C[i][j] + A[i][k] * B[k][j]) % mod;
            }
        }
    }
    return C;
}

Matrix power_matrix(Matrix A, ll exp, ll mod)
{
    int n = A.size();
    Matrix res(n, vector<ll>(n, 0));
    for (int i = 0; i < n; i++)
        res[i][i] = 1;
    while (exp > 0)
    {
        if (exp % 2 == 1)
            res = multiply_matrix(res, A, mod);
        A = multiply_matrix(A, A, mod);
        exp /= 2;
    }
    return res;
}

ll nCr_congruence(ll n, int k, int m, ll mod)
{
    if (m == 1)
        return custom_pow(2, n, mod);
    Matrix T(m, vector<ll>(m, 0));
    for (int i = 0; i < m; i++)
    {
        T[i][i] = 1;
        T[i][(i + m - 1) % m] = 1;
    }
    T = power_matrix(T, n, mod);
    return T[k][0];
}
\end{lstlisting}
\section{Geometry}
\subsection{Ball 3D Class}
\begin{lstlisting}[language=C++]
/**
 *  THE ULTIMATE CP TEMPLATE BLUEPRINT: 3D CLOSEST PAIR OF POINTS (DIVIDE & CONQUER)
 *
 * 1.  WHAT PROBLEM DOES THIS SOLVE?
 * - Keywords: "Minimum Euclidean distance in 3D", "Closest points in space", "N up to 10^5".
 * - Classic Scenarios: You are given N points in a 3D coordinate system. Finding the closest 
 *   pair naively takes O(N^2), which will Time Limit Exceed (TLE) for N > 10^4.
 * - The Magic: Divide and Conquer! We sort the points by the X-axis and recursively divide the 
 *   space into halves. The tricky part in 3D is merging the halves (the "strip"). If we just 
 *   check all points in the strip naively, it can still degrade to O(N^2). 
 *   To fix this, the algorithm groups the strip into 1D "columns" along the Z-axis using the 
 *   current minimum distance `d` as the column width. This brilliant trick guarantees that each 
 *   point only checks a strict O(1) number of neighboring points. The overall complexity 
 *   beautifully drops to O(N log N).
 *
 * 2.  HOW TO USE IT (THE BLACK BOX)
 * - Initialization: Create a `BallManager` object. It automatically reads `N` from `std::cin`, 
 *   followed by the `X Y Z` coordinates of the `N` points.
 *       BallManager bm;
 *
 * - Execution: Call the main solver function to get the minimum distance.
 *       long double min_distance = bm.ClosestPair();
 *
 * - Complexity:
 *       Time: O(N log N)  The X and Y sorting takes O(N log N), and the recursive strip 
 *             processing takes O(N) per level.
 *       Space: O(N) to store the objects, vectors, and recursion stack.
 *
 * 3.  HOW TO ADAPT IT (THE DIALS & SWITCHES)
 * - Decoupling I/O: The constructor currently reads input directly. For better flexibility, 
 *   you can modify the constructor to accept a `vector<vector<ll>>` or `vector<Ball>` instead 
 *   of using `cin` inside the class.
 * - Precision Issues: The algorithm uses `long double` and `sqrtl`. While this is highly accurate, 
 *   some problems require exact integer squared distances. You can return the squared distance 
 *   at the very end by squaring `finalMinDist`, but do NOT remove `sqrtl` from `DistanceTo`, 
 *   as the `SearchStrip` bucketing logic mathematically relies on the true linear distance.
 */

#include <iostream>
#include <vector>
#include <cmath>
#include <algorithm>
#include <limits>

using namespace std;
#define ll long long
#define ld long double

class Ball
{
public:
    ll posX, posY, posZ; //Position values in 3D space
    ll columnNumber[3]; //ID number of the column which ball belongs to (left, middle, right)
    ll columnIndex[3]; //Position inside the column (left, middle, right)
    
    // Time Complexity: O(1)
    //            .
    ld DistanceTo(Ball* targetBall)
    {
        ll squareDistX = (posX - targetBall->posX) * (posX - targetBall->posX);
        ll squareDistY = (posY - targetBall->posY) * (posY - targetBall->posY);
        ll squareDistZ = (posZ - targetBall->posZ) * (posZ - targetBall->posZ);
        return sqrtl(squareDistX + squareDistY + squareDistZ);
    }
};

class BallManager
{
    ll ballCount = 0; //Total number of balls
    Ball* ballArray; //Dynamically allocated ball array which holds all balls inside
 
    // Time Complexity: O(1)
    static bool CompareX(Ball ball1, Ball ball2)
    {
        return (ball1.posX < ball2.posX);
    }
    
    // Time Complexity: O(1)
    static bool CompareY(Ball* ball1, Ball* ball2)
    {
        return (ball1->posY < ball2->posY);
    }
 
    // Time Complexity: O(N^2) 
    //  N    (balls.size()).      .
    ld NaiveMethod(std::vector<Ball*> balls)
    {
        int size = balls.size();
        ld minDist = std::numeric_limits<double>::max();
 
        //Take balls one by one
        for (int i = 0; i < size - 1; ++i)
        {
            //Check other remaining balls
            for (int j = i + 1; j < size; ++j)
            {
                ld dist = balls[i]->DistanceTo(balls[j]);
                if (dist < minDist)
                    minDist = dist;
            }
        }
        return minDist;
    }
    
    // Time Complexity: O(S) 
    //  S      (stripSize). 
    //    (Columns)    Z        
    //         (     ).
    ld SearchStrip(std::vector<Ball*> strip, ld stripDist)
    {
        //Save the strip distance value as minimum distance known
        ld minDist = stripDist;
 
        //If ball count in strip is less than 2, return strip distance
        int stripSize = strip.size();
        if (stripSize < 2)
            return minDist;
 
        //If ball count in strip is 2 or 3, find distance with naive method
        if (stripSize < 4)
        {
            ld newDist = NaiveMethod(strip);
            if (newDist < minDist)
                minDist = newDist;
            return minDist;
        }
 
        //If ball count is more than 3, start processing of creating columns
        //Find the smallest z value among balls -> O(S)
        int smallestZ = strip[0]->posZ;
        for (int i = 1; i < stripSize; i++)
            if (strip[i]->posZ < smallestZ)
                smallestZ = strip[i]->posZ;
 
        //Find the largest z value among balls -> O(S)
        int largestZ = strip[0]->posZ;
        for (int i = 1; i < stripSize; i++)
            if (strip[i]->posZ > largestZ)
                largestZ = strip[i]->posZ;
 
        //Find range between z values and column count
        int rangeZ = largestZ - smallestZ;
        int columnCount = (int)(rangeZ / stripDist) + 1;
 
        //Create an array of ball pointer vectors to put pointers of balls into proper columns
        std::vector<Ball*>* columns = new std::vector<Ball*>[columnCount];
 
        //Put all balls into proper columns (Phase 1) -> O(S)
        //Start picking balls one by one with the smallest Y first
        for (int i = 0; i < stripSize; i++)
        {
            //Find middle(main) column of ball, put ball into it, save index
            strip[i]->columnNumber[1] = (ll)((ld)(strip[i]->posZ - smallestZ) / stripDist);
            columns[strip[i]->columnNumber[1]].push_back(strip[i]);
            strip[i]->columnIndex[1] = columns[strip[i]->columnNumber[1]].size() - 1;
 
            //Find and save left and right columns of ball
            strip[i]->columnNumber[0] = strip[i]->columnNumber[1] - 1;
            strip[i]->columnNumber[2] = strip[i]->columnNumber[1] + 1;
 
            //There is no column at left of leftmost column, and right of rightmost column
            //Save indices in left and right columns
            if (strip[i]->columnNumber[1] != 0) //If not leftmost column
                strip[i]->columnIndex[0] = columns[strip[i]->columnNumber[0]].size() - 1;
 
            if (strip[i]->columnNumber[1] != columnCount - 1) //If not rightmost column
                strip[i]->columnIndex[2] = columns[strip[i]->columnNumber[2]].size() - 1;
        }
 
        //Find if such lesser distance exists between any pair of balls inside strip (Phase 2)
        // O(S) overall. The inner loops run a constant number of times O(1) per ball.
        //Start picking balls one by one with the smallest Y first
        for (int i = 0; i < stripSize; i++)
        {
            //For left, middle and right columns of the selected ball
            for (int j = 0; j < 3; j++)
            {
                //Skip invalid out of bounds columns
                if ((strip[i]->columnNumber[j] == -1) || (strip[i]->columnNumber[j] == columnCount))
                    continue;
 
                //Look through increasing Y value = O(1)
                for (int k = strip[i]->columnIndex[j] + 1; k < (ll)columns[strip[i]->columnNumber[j]].size(); k++)
                {
                    //If Y distance is already more than strip distance, no need to check remaning balls in column
                    if (columns[strip[i]->columnNumber[j]][k]->posY > strip[i]->posY + stripDist)
                        break;
 
                    //Check distance and if it is smaller than current minimum distance, save it as new minimum distance
                    ld newDist = strip[i]->DistanceTo(columns[strip[i]->columnNumber[j]][k]);
                    if (newDist < minDist)
                        minDist = newDist;
                }
            }
        }
 
        //Delete dynamically allocated columns array
        delete[] columns;
 
        //Return minimum distance found in strip or initially given strip distance
        return minDist;
    }
 
    // Time Complexity: T(N) = 2*T(N/2) + O(N) => O(N log N)
    //  N     (sortedY.size()).        O(N).
    ld RecursiveCall(int sortedXPos, std::vector<Ball*> sortedY)
    {
        //If size is small enough, find distance with naive method -> O(1)
        int size = sortedY.size();
        if (size < 4)
            return NaiveMethod(sortedY);
 
        //Get middle ball for X axis
        int midIndex = sortedXPos + size / 2;
        Ball* midBall = &ballArray[midIndex];
 
        //Initiliaze left and right subarrays
        std::vector<Ball*> leftVector;
        std::vector<Ball*> rightVector;
 
        //In case of there are other balls which share same X position as middle point
        bool putToRightArr = true;
 
        //Distribute balls into left and right vectors -> O(N)
        for (int i = 0; i < size; i++)
        {
            if (sortedY[i]->posX < midBall->posX)
                leftVector.push_back(sortedY[i]);
            else if (sortedY[i]->posX > midBall->posX)
                rightVector.push_back(sortedY[i]);
            else //if (sortedY[i].posX == midBall.posX)
            {
                if (putToRightArr)
                    rightVector.push_back(sortedY[i]);
                else
                    leftVector.push_back(sortedY[i]);
                putToRightArr = !putToRightArr;
            }
        }
 
        //Pass subarrays to same function recursively -> 2 * T(N/2)
        ld minLeftDist = RecursiveCall(sortedXPos, leftVector);
        ld minRightDist = RecursiveCall(sortedXPos + leftVector.size(), rightVector);
 
        //Initiliaze strip vector
        std::vector<Ball*> strip;
        ld stripDist = std::min(minLeftDist, minRightDist);
 
        //Put balls which are close enough to middle into strip -> O(N)
        for (int i = 0; i < size; i++)
        {
            ll distToMid = sortedY[i]->posX - midBall->posX;
            if (distToMid < 0)
                distToMid *= -1;
            if (distToMid < stripDist)
                strip.push_back(sortedY[i]);
        }
 
        //Search in strip and return the minimum distance -> O(N) (Since strip size <= N)
        return SearchStrip(strip, stripDist);
    }
 
public:
    // Time Complexity: O(N) 
    //  N    (ballCount)       .
    BallManager(){
        cin >> ballCount;
        ballArray = new Ball[ballCount];
        for (int i = 0; i < ballCount; ++i) {
            cin >> ballArray[i].posX;
            cin >> ballArray[i].posY;
            cin >> ballArray[i].posZ;
        }
    }
    
    // Time Complexity: O(1)
    ~BallManager() { delete[] ballArray;};
    
    // Time Complexity: O(N log N)
    //           .
    ld ClosestPair()
    {
        //Sort original ballArray by increasing X order = O(N log N)
        std::sort(ballArray, ballArray + ballCount, &BallManager::CompareX);
 
        //Initiliaze sortedY vector -> O(N)
        std::vector<Ball*> sortedY;
        for (int i = 0; i < ballCount; i++)
            sortedY.push_back(&ballArray[i]);
 
        //Sort sortedY by increasing Y order = O(N log N)
        std::sort(sortedY.begin(), sortedY.end(), &BallManager::CompareY);
 
        //Start recursive part of the program and get result -> O(N log N)
        ld finalMinDist = RecursiveCall(0, sortedY);
 
        //Return the final minimum distance
        return finalMinDist;
    }
};

\end{lstlisting}
\subsection{Li Chao Tree}
\begin{lstlisting}[language=C++]
#include <bits/stdc++.h>
using namespace std;

typedef long long ll;

const ll INF = 2e18; // Infinity for initializations

// A structure to represent the line y = mx + c
struct Line {
    ll m, c;
    // For Maximization, initialize empty lines with -INF. 
    // (If you need Minimization, change this to +INF).
    Line(ll m = 0, ll c = -INF) : m(m), c(c) {}
    
    ll eval(ll x) { 
        return m * x + c; 
    }
};

struct LiChaoTree {
    struct Node {
        Line line;
        int lc, rc; // Left child and Right child indices
        Node() : line(), lc(-1), rc(-1) {}
    };

    vector<Node> tree;
    ll min_x, max_x;

    // Initialize with your expected coordinate range (e.g., -1e9 to 1e9)
    LiChaoTree(ll min_x, ll max_x) : min_x(min_x), max_x(max_x) {
        tree.push_back(Node());
    }

    void add_line(Line nw, int u, ll l, ll r) {
        ll m = l + (r - l) / 2;
        
        // Change > to < if you want to solve a Minimization problem
        bool left_better = nw.eval(l) > tree[u].line.eval(l);
        bool mid_better = nw.eval(m) > tree[u].line.eval(m);

        // The line that is better at the midpoint \"owns\" this node
        if (mid_better) {
            swap(tree[u].line, nw);
        }

        if (r - l == 1) return; // Leaf node reached

        // If the old line was better at one of the boundaries, 
        // it means the lines intersect. Push the losing line down to the children!
        if (left_better != mid_better) {
            if (tree[u].lc == -1) {
                tree[u].lc = tree.size();
                tree.push_back(Node());
            }
            add_line(nw, tree[u].lc, l, m);
        } else {
            if (tree[u].rc == -1) {
                tree[u].rc = tree.size();
                tree.push_back(Node());
            }
            add_line(nw, tree[u].rc, m, r);
        }
    }

    void add_line(ll m, ll c) {
        add_line(Line(m, c), 0, min_x, max_x);
    }

    ll query(ll x, int u, ll l, ll r) {
        if (u == -1) return -INF; // Change to +INF for Minimization
        
        ll m = l + (r - l) / 2;
        ll cur = tree[u].line.eval(x);
        
        if (r - l == 1) return cur; // Leaf node reached
        
        // Change max() to min() for Minimization
        if (x < m) {
            return max(cur, query(x, tree[u].lc, l, m));
        } else {
            return max(cur, query(x, tree[u].rc, m, r));
        }
    }

    ll query(ll x) {
        return query(x, 0, min_x, max_x);
    }
};


How to use it:Initialize: Define the range of your $X$-coordinates (your $k$ queries).

// If your query values k range from -10^5 to 10^5:
LiChaoTree lct(-100005, 100005);

Add Lines: Loop through your pairs and insert the lines.// Insert line y = mx + c
lct.add_line(m, c);

Query: Find the maximum $y$ for a specific $x$.
ll max_y = lct.query(k);


add_line(m, c): Add a line to the set in $O(\log(\text{range}))$ time.
query(x): Find the maximum $y$ among all lines at coordinate $x$ in $O(\log(\text{range}))$ time.
\end{lstlisting}
\subsection{additonal geometry}
\begin{lstlisting}[language=C++]

// cosine rule:

c^2 = a^2 + b^2 - 2ab * cos(C) a^2 = b^2 + c^2 - 2bc * cos(A) b^2 = a^2 + c^2 - 2ac * cos(B)

cos(C) = (a^2 + b^2 - c^2) / (2ab) cos(A) = (b^2 + c^2 - a^2) / (2bc) cos(B) = (a^2 + c^2 - b^2) / (2ac)

C = cos^(-1) [(a^2 + b^2 - c^2) / (2ab)] A = cos^(-1) [(b^2 + c^2 - a^2) / (2bc)] B = cos^(-1) [(a^2 + c^2 - b^2) / (2ac)]


   // if I am given a set of points how to sort counter clock wise around the origin

// The arg() Function: This function automatically returns a mathematical value strictly within the
// interval $(-\pi, \pi]$, which perfectly aligns with the required counter-clockwise sweep.
    
bool comp(pt &a, pt &b) { 
    return arg(a) < arg(b); 
}

// cog
pt centerOfGravity(const vector<pt>& p) {
    int n = p.size();
    pt centroid = {0, 0};
    ld doubleArea = 0.0;

    for (int i = 0; i < n; i++) {
        int next = (i + 1) % n;
        ld c = cross(p[i], p[next]); 
        
        doubleArea += c;
        centroid += (p[i] + p[next]) * c; 
    }

    // Safety check to prevent division by zero for degenerate polygons
    if (fabs(doubleArea) < EPS) return centroid; 

    return centroid / (3.0 * doubleArea); 
}



typedef long double ld;
const ld eps = 1e-9;

struct P {
    ld x, y;

    // Constructor
    P(ld _x = 0, ld _y = 0) : x(_x), y(_y) {}

    // Overload subtraction to get a directional vector between two points
    P operator-(const P& other) const {
        return P(x - other.x, y - other.y);
    }

    // Overload ^ for 2D Cross Product
    ld operator^(const P& other) const {
        return x * other.y - y * other.x;
    }

    // Overload * for 2D Dot Product (Optional, but highly recommended)
    ld operator*(const P& other) const {
        return x * other.x + y * other.y;
    }

    // Overload < for sorting (sorts by X first, then by Y)
    // Used in Convex Hull and as a fallback in your polar sort
    bool operator<(const P& other) const {
        if (fabs(x - other.x) > eps) {
            return x < other.x;
        }
        return y < other.y;
    }

    // Overload == to check if two points are identical
    bool operator==(const P& other) const {
        return fabs(x - other.x) < eps && fabs(y - other.y) < eps;
    }
};

// Helper function used throughout your template to get a vector from A to B
inline P vec(P a, P b) {
    return b - a;
}




double get_full_angle(P a) {
    // atan2 takes (y, x), NOT (x, y)
    double ang = atan2(a.y, a.x);
    
    // If the angle is negative (quadrants 3 and 4), add 2*PI
    if (ang < 0) {
        ang += 2 * PI;
    }
    
    return ang;
}

// If you need the result in degrees from 0 to 360
double get_full_angle_degrees(P a) {
    double ang = atan2(a.y, a.x);
    if (ang < 0) {
        ang += 2 * PI;
    }
    return ang * (180.0 / PI);
}

// check if a,b,c are collinear
bool collinear(P a, P b, P c)  
{  
    return (vec(a, b) ^ vec(a, c)) == 0;  
}

// to sort points in counterclockwise order around a point
struct cmp  
{  
    P about;  
    cmp(P c)  
    {  
       about = c;  
    }  
    bool operator()(const P& a, const P& b) const  
    {  
       ld cr = vec(about, a) ^ vec(about, b);  
       if (fabs(cr) < eps)  
          return a < b;  
       return cr > 0;  
    }  
};  

void sortAntiClockWise(vector<P>& pnts)  
{  
    P mn(*min_element(all(pnts)));  
    sort(pnts.begin(), pnts.end(), cmp(mn));  
}  
inline bool pibb(P const& a, P const& b1, P const& b2)  
{  
    return a.x >= min(b1.x, b2.x) &&  
       a.x <= max(b1.x, b2.x) &&  
       a.y >= min(b1.y, b2.y) &&  
       a.y <= max(b1.y, b2.y);  
} 


// Returns the center of the circle passing through points A, B, and C
P getCircleCenter(P A, P B, P C) {  
    if (isCollinear(A, B, C)) {  
       collinear = true;  
       return {0, 0};   
    }  
    collinear = false;  
    ld D = 2 * (A.x * (B.y - C.y) + B.x * (C.y - A.y) + C.x * (A.y - B.y));  
    ld Ux = ((A.x * A.x + A.y * A.y) * (B.y - C.y) + (B.x * B.x + B.y * B.y) * (C.y - A.y) + (C.x * C.x + C.y * C.y) * (A.y - B.y)) / D;  
    ld Uy = ((A.x * A.x + A.y * A.y) * (C.x - B.x) + (B.x * B.x + B.y * B.y) * (A.x - C.x) + (C.x * C.x + C.y * C.y) * (B.x - A.x)) / D;  
    return {Ux, Uy};  
}



// compute the convex hull of a set of points using Andrew's monotone chain algorithm
// hull will contain the points of the convex hull in counterclockwise order

void convexHull(vector<P> p, vector<P>& hull)  
{  
  
    sort(all(p), [&](P& a, P& b)  
    {  
      if (a.x != b.x)  
         return a.x < b.x;  
      return a.y < b.y;  
    });  
    if (p.size() == 1)  
    {  
       hull.push_back(p[0]);  
       return;  
    }  
    for (int rep = 0; rep < 2; rep++)  
    {  
       int s = hull.size();  
       for (int i = 0; i < p.size(); i++)  
       {  
          while (hull.size() >= s + 2)  
          {  
             P p1 = hull.end()[-2];  
             P p2 = hull.end()[-1];  
             if ((vec(p1, p2) ^ vec(p1, p[i])) < -eps)  
                break;  
  
             hull.pop_back();  
          }  
          hull.push_back(p[i]);  
       }  
       reverse(all(p));  
       hull.pop_back();  
    }  
}


// Returns the centroid of a polygon defined by the points in vector p
// centroid is the center of mass of the polygon, assuming uniform density

P getCentroid(vector<P>& p)  
{  
    ld x, y;  
    long long tarea = 0;  
    x = 0;  
    y = 0;  
    for (int i = 1; i < p.size() - 2; i++)  
    {  
       long long area = (vec(p[0], p[i]) ^ vec(p[0], p[i + 1]));  
       x += area * ((p[0].x + p[i].x + p[i + 1].x) / 3.0);  
       y += area * ((p[0].y + p[i].y + p[i + 1].y) / 3.0);  
       tarea += area;  
    }  
    x /= tarea;  
    y /= tarea;  
    return { x, y };  
}


// some properties of regular polygons are listed below.

Properties of Regular polygons Some of the properties of regular polygons are listed below.

All the sides of a regular polygon are equal
All the interior angles are equal 
The perimeter of a regular polygon with n sides is equal to the n times of a side measure. 
The sum of all the interior angles of a simple n-gon or regular polygon = (n  2)  180 
The number of diagonals in a polygon with n sides = n(n  3)/2 
The number of triangles formed by joining the diagonals from one corner of a polygon = n  2 
The measure of each interior angle of n-sided regular polygon = [(n  2)  180]/n 
The measure of each exterior angle of an n-sided regular polygon = 360/n

area of reqular polygon = ((l^2)*n)/(4tan(PI/n)) l is the side length

n is the number of sides

\end{lstlisting}
\subsection{allGemometry}
\begin{lstlisting}[language=C++]
#include <iostream>
#include <complex>
#include <vector>
#include <cmath>
#include <algorithm>
#include <set>
#include <cassert>
#include <deque>
#include <numeric>

using namespace std;

// =============================================================================
//                              BASICS & TYPES
// =============================================================================

typedef long double ld;
typedef ld T;
typedef complex<T> pt;

const T EPS = 1e-9;
const long double PI = acos(-1.0);

#define x real()
#define y imag()

bool operator==(pt a, pt b) { return fabs(a.x - b.x) < EPS && fabs(a.y - b.y) < EPS; }
bool operator!=(pt a, pt b) { return !(a == b); }

int sgn(T val) {
    if (val > EPS) return 1;
    if (val < -EPS) return -1;
    return 0;
}

T sq(pt p) { return p.x * p.x + p.y * p.y; }
ld abs(pt p) { return sqrt(sq(p)); }

pt perp(pt p) { return {-p.y, p.x}; }

T dot(pt v, pt w) { return v.x * w.x + v.y * w.y; }
T cross(pt v, pt w) { return v.x * w.y - v.y * w.x; }

bool isPerp(pt v, pt w) { return fabs(dot(v, w)) < EPS; }

pt unitVector(pt p) {
    ld len = abs(p);
    if (len < EPS) return {0, 0};
    return p / len;
}

// =============================================================================
//                              TRANSFORMATIONS
// =============================================================================

pt translate(pt v, pt p) { return p + v; }

// Scale point p by a factor around center c
pt scale(pt c, T factor, pt p) {
    return c + (p - c) * factor;
}

// Rotate point p by angle a around center c
pt rot(pt p, pt c, ld a) {
    pt v = p - c;
    pt rotate = {cos(a), sin(a)};
    return c + rotate * v;
}

// If point p maps to fp and q maps to fq, find the image of r
pt linearTransfo(pt p, pt q, pt r, pt fp, pt fq) {
    return fp + (r - p) * (fq - fp) / (q - p);
}

// =============================================================================
//                              ANGLES
// =============================================================================

// (AB X AC) --> relative to AB: if(C left) return pos, else if (C right) return neg
T orient(pt a, pt b, pt c) { return cross(b - a, c - a); }

bool inAngle(pt a, pt b, pt c, pt p) {
    T abp = orient(a, b, p), acp = orient(a, c, p), abc = orient(a, b, c);
    if (abc < 0) swap(abp, acp);
    return (abp >= 0 && acp <= 0) ^ (abc < 0);
}

T angle(pt v, pt w) {
    return atan2(fabs(cross(v, w)), dot(v, w));
}

ld orientedAngle(pt a, pt b, pt c) {
    if (orient(a, b, c) >= 0)
        return angle(b - a, c - a);
    else
        return 2 * PI - angle(b - a, c - a);
}

ld angleTravelled(pt a, pt p, pt q) {
    ld ampli = angle(p - a, q - a);
    if (orient(a, p, q) > 0) return ampli;
    else return -ampli;
}

bool half(pt p) {
    return p.y > EPS || (fabs(p.y) <= EPS && p.x < -EPS);
}

// =============================================================================
//                              LINES
// =============================================================================

struct line {
    pt v; T c;

    // From direction vector v and offset c
    line(pt v, T c) : v(v), c(c) {}

    // From equation ax+by=c
    line(T a, T b, T _c) {
        v = {b, -a};
        c = _c;
    }

    // From points P and Q
    line(pt p, pt q) {
        v = q - p;
        c = cross(v, p);
    }

    T side(pt p) { return cross(v, p) - c; }
    double dist(pt p) { return abs(side(p)) / abs(v); }
    double sqDist(pt p) { return side(p) * side(p) / (T)sq(v); }
    line perpThrough(pt p) { return {p, p + perp(v)}; }
    bool cmpProj(pt p, pt q) { return dot(v, p) < dot(v, q); }
    line translate(pt t) { return {v, c + cross(v, t)}; }
    line shiftLeft(double dist) { return {v, c + dist * abs(v)}; }
    pt proj(pt p) { return p - perp(v) * side(p) / sq(v); }
    pt refl(pt p) { return p - perp(v) * (T)2.0 * side(p) / sq(v); }
};

bool inter(line l1, line l2, pt &out) {
    T d = cross(l1.v, l2.v);
    if (fabs(d) <= EPS) return false;
    out = (l2.v * l1.c - l1.v * l2.c) / d;
    return true;
}

line bisector(line l1, line l2, bool interior) {
    assert(cross(l1.v, l2.v) != 0); // Cannot be parallel
    T sign = interior ? 1 : -1;
    return {l2.v / (T)abs(l2.v) + l1.v / (T)abs(l1.v) * sign,
            l2.c / abs(l2.v) + l1.c / abs(l1.v) * sign};
}

// =============================================================================
//                              SEGMENTS
// =============================================================================

bool inDisk(pt a, pt b, pt p) {
    return dot(a - p, b - p) <= EPS;
}

bool onSegment(pt a, pt b, pt p) {
    return fabs(orient(a, b, p)) <= EPS && inDisk(a, b, p);
}

bool properInter(pt a, pt b, pt c, pt d, pt &out) {
    T oa = orient(c, d, a), ob = orient(c, d, b);
    T oc = orient(a, b, c), od = orient(a, b, d);
    if (sgn(oa) * sgn(ob) < 0 && sgn(oc) * sgn(od) < 0) {
        out = (a * ob - b * oa) / (ob - oa);
        return true;
    }
    return false;
}

set<pair<ld, ld>> inters(pt a, pt b, pt c, pt d) {
    set<pair<ld, ld>> s;
    pt out;
    if (a == c || a == d) s.insert({a.x, a.y});
    if (b == c || b == d) s.insert({b.x, b.y});
    if (s.size()) return s;

    if (properInter(a, b, c, d, out)) return {{out.x, out.y}};
    if (onSegment(c, d, a)) s.insert({a.x, a.y});
    if (onSegment(c, d, b)) s.insert({b.x, b.y});
    if (onSegment(a, b, c)) s.insert({c.x, c.y});
    if (onSegment(a, b, d)) s.insert({d.x, d.y});

    return s;
}

ld segPoint(pt a, pt b, pt p) {
    if (a != b) {
        line l(a, b);
        if (l.cmpProj(a, p) && l.cmpProj(p, b)) return l.dist(p);
    }
    return min(abs(p - a), abs(p - b));
}

ld rayPoint(pt a, pt b, pt p) {
    pt ab = b - a, ap = p - a;
    if (dot(ab, ap) < 0) return abs(ap);
    line l(a, b);
    return l.dist(p);
}

ld segSeg(pt a, pt b, pt c, pt d) {
    pt dummy;
    if (properInter(a, b, c, d, dummy)) return 0;
    return min({segPoint(a, b, c), segPoint(a, b, d),
                segPoint(c, d, a), segPoint(c, d, b)});
}

// =============================================================================
//                              CIRCLES
// =============================================================================

pair<pt, T> circumCircle(pt a, pt b, pt c) {
    b = b - a, c = c - a;
    assert(cross(b, c) != 0); // No circumcircle if aligned
    return {a + perp(b * sq(c) - c * sq(b)) / cross(b, c) / (T)2, 
            abs(perp(b * sq(c) - c * sq(b)) / cross(b, c) / (T)2)};
}

int circleLine(pt o, double r, line l, pair<pt, pt> &out) {
    double h2 = r * r - l.sqDist(o);
    if (h2 >= 0) {
        pt p = l.proj(o);
        pt h = l.v * (T)(sqrt(h2) / abs(l.v));
        out = {p - h, p + h};
    }
    return 1 + sgn(h2);
}

int circleCircle(pt o1, T r1, pt o2, T r2, pair<pt, pt> &out) {
    pt d = o2 - o1; T d2 = sq(d);
    if (d2 == 0) { assert(r1 != r2); return 0; }
    T pd = (d2 + r1 * r1 - r2 * r2) / 2;
    T h2 = r1 * r1 - pd * pd / d2;
    if (h2 >= 0) {
        pt p = o1 + d * pd / d2, h = perp(d) * sqrt(h2 / d2);
        out = {p - h, p + h};
    }
    return 1 + sgn(h2);
}

int tangents(pt o1, T r1, pt o2, T r2, bool inner, vector<pair<pt, pt>> &out) {
    if (inner) r2 = -r2;
    pt d = o2 - o1;
    T dr = r1 - r2, d2 = sq(d), h2 = d2 - dr * dr;
    if (d2 == 0 || h2 < 0) { assert(h2 != 0); return 0; }
    for (T sign : {-1, 1}) {
        pt v = (d * dr + perp(d) * sqrt(h2) * sign) / d2;
        out.push_back({o1 + v * r1, o2 + v * r2});
    }
    return 1 + (h2 > 0);
}

// =============================================================================
//                              BASIC POLYGONS
// =============================================================================

ld areaTriangle(pt a, pt b, pt c) { return abs(cross(b - a, c - a)) / 2.0; }

ld areaPolygon(vector<pt> p) {
    ld area = 0.0;
    for (int i = 0, n = p.size(); i < n; i++)
        area += cross(p[i], p[(i + 1) % n]);
    return abs(area) / 2.0;
}

bool isConvex(vector<pt> p) {
    bool hasPos = false, hasNeg = false;
    for (int i = 0, n = p.size(); i < n; i++) {
        int o = sgn(orient(p[i], p[(i + 1) % n], p[(i + 2) % n]));
        if (o > 0) hasPos = true;
        if (o < 0) hasNeg = true;
    }
    return !(hasPos && hasNeg);
}

bool above(pt a, pt p) { return p.y >= a.y; }

bool crossesRay(pt a, pt p, pt q) {
    return (above(a, q) - above(a, p)) * orient(a, p, q) > 0;
}

bool inPolygon(vector<pt> p, pt a, bool strict = true) {
    int numCrossings = 0;
    for (int i = 0, n = p.size(); i < n; i++) {
        if (onSegment(p[i], p[(i + 1) % n], a)) return !strict;
        numCrossings += crossesRay(a, p[i], p[(i + 1) % n]);
    }
    return numCrossings & 1;
}

bool onPolygonSide(vector<pt> p, pt a) {
    for (int i = 0, n = p.size(); i < n; i++) {
        if (onSegment(p[i], p[(i + 1) % n], a)) return true;
    }
    return false;
}

pt polygonCentroid(const vector<pt>& p) {
    int n = p.size();
    ld area = 0.0, cx = 0.0, cy = 0.0;
    for (int i = 0; i < n; i++) {
        int next = (i + 1) % n;
        ld cross_val = cross(p[i], p[next]);
        area += cross_val;
        cx += (p[i].x + p[next].x) * cross_val;
        cy += (p[i].y + p[next].y) * cross_val;
    }
    area /= 2.0;
    return {cx / (6.0 * area), cy / (6.0 * area)};
}

// =============================================================================
//                              PICK'S THEOREM
// =============================================================================

int boundaryPoints(const vector<pt>& p) {
    int b = 0, n = p.size();
    for (int i = 0; i < n; i++) {
        int next = (i + 1) % n;
        int dx = abs(p[i].x - p[next].x);
        int dy = abs(p[i].y - p[next].y);
        b += std::gcd(dx, dy);
    }
    return b;
}

int doubleAreaPolygon(const vector<pt>& p) {
    int doubleArea = 0, n = p.size();
    for (int i = 0; i < n; i++) {
        int next = (i + 1) % n;
        doubleArea += cross(p[i], p[next]);
    }
    return abs(doubleArea);
}

int interiorPoints(const vector<pt>& p) {
    return (doubleAreaPolygon(p) - boundaryPoints(p) + 2) / 2;
}

// =============================================================================
//                              ADVANCED POLYGONS
// =============================================================================

bool cw(pt a, pt b, pt c, bool include_collinear) {
    int o = sgn(orient(a, b, c));
    return o < 0 || (include_collinear && o == 0);
}

bool collinear(pt a, pt b, pt c) { return sgn(orient(a, b, c)) == 0; }

void convex_hull(vector<pt>& a, bool include_collinear = false) {
    if (a.empty()) return;
    pt p0 = *min_element(a.begin(), a.end(), [](pt a, pt b) {
        return make_pair(a.y, a.x) < make_pair(b.y, b.x);
    });
    sort(a.begin(), a.end(), [&p0](const pt& a, const pt& b) {
        int o = sgn(orient(p0, a, b));
        if (o == 0) return sq(p0 - a) < sq(p0 - b);
        return o < 0;
    });
    if (include_collinear) {
        int i = (int)a.size() - 1;
        while (i >= 0 && collinear(p0, a[i], a.back())) i--;
        reverse(a.begin() + i + 1, a.end());
    }
    vector<pt> st;
    for (int i = 0; i < (int)a.size(); i++) {
        while (st.size() > 1 && !cw(st[st.size() - 2], st.back(), a[i], include_collinear))
            st.pop_back();
        if (st.empty() || a[i] != st.back())
            st.push_back(a[i]);
    }
    if (!include_collinear && st.size() == 2 && st[0] == st[1])
        st.pop_back();
    a = st;
}

void reorder_polygon(vector<pt> & P) {
    size_t pos = 0;
    for (size_t i = 1; i < P.size(); i++) {
        if (P[i].y < P[pos].y || (P[i].y == P[pos].y && P[i].x < P[pos].x))
            pos = i;
    }
    rotate(P.begin(), P.begin() + pos, P.end());
}

// P and Q must be counter clockwise
vector<pt> minkowski(vector<pt> P, vector<pt> Q) {
    reorder_polygon(P);
    reorder_polygon(Q);
    P.push_back(P[0]); P.push_back(P[1]);
    Q.push_back(Q[0]); Q.push_back(Q[1]);
    
    vector<pt> result;
    size_t i = 0, j = 0;
    while (i < P.size() - 2 || j < Q.size() - 2) {
        result.push_back(P[i] + Q[j]);
        auto crs = cross(P[i + 1] - P[i], Q[j + 1] - Q[j]);
        if (crs >= 0 && i < P.size() - 2) ++i;
        if (crs <= 0 && j < Q.size() - 2) ++j;
    }
    return result;
}

double maximum_dist_from_polygon_to_polygon(vector<pt> &u, vector<pt> &v) {
    int n = (int)u.size(), m = (int)v.size();
    double ans = 0;
    if (n < 3 || m < 3) {
        for (int i = 0; i < n; i++)
            for (int j = 0; j < m; j++) 
                ans = max(ans, (double)sq(u[i] - v[j]));
        return sqrt(ans);
    }
    if (u[0].x > v[0].x) { swap(n, m); swap(u, v); }
    
    int i = 0, j = 0, step = n + m + 10;
    while (j + 1 < m && v[j].x < v[j + 1].x) j++;
    while (step--) {
        if (cross(u[(i + 1) % n] - u[i], v[(j + 1) % m] - v[j]) >= 0) j = (j + 1) % m;
        else i = (i + 1) % n;
        ans = max(ans, (double)sq(u[i] - v[j]));
    }
    return sqrt(ans);
}

// =============================================================================
//                              HALF-PLANE INTERSECTION
// =============================================================================

struct Halfplane {
    pt p, pq;
    long double angle;

    Halfplane() {}
    Halfplane(const pt& a, const pt& b) : p(a), pq(b - a) {
        angle = atan2l(pq.y, pq.x);
    }

    bool out(const pt& r) {
        return cross(pq, r - p) < -EPS;
    }

    bool operator < (const Halfplane& e) const { return angle < e.angle; }

    friend pt inter(const Halfplane& s, const Halfplane& t) {
        long double alpha = cross((t.p - s.p), t.pq) / cross(s.pq, t.pq);
        return s.p + (s.pq * alpha);
    }
};

vector<pt> hp_intersect(vector<Halfplane>& H) {
    const int inf = 1e9;
    pt box[4] = { pt(inf, inf), pt(-inf, inf), pt(-inf, -inf), pt(inf, -inf) };
    for (int i = 0; i < 4; i++) {
        Halfplane aux(box[i], box[(i + 1) % 4]);
        H.push_back(aux);
    }

    sort(H.begin(), H.end());
    deque<Halfplane> dq;
    int len = 0;

    for (int i = 0; i < int(H.size()); i++) {
        while (len > 1 && H[i].out(inter(dq[len - 1], dq[len - 2]))) {
            dq.pop_back(); --len;
        }
        while (len > 1 && H[i].out(inter(dq[0], dq[1]))) {
            dq.pop_front(); --len;
        }

        if (len > 0 && fabsl(cross(H[i].pq, dq[len - 1].pq)) < EPS) {
            if (dot(H[i].pq, dq[len - 1].pq) < 0.0) return vector<pt>();
            if (H[i].out(dq[len - 1].p)) {
                dq.pop_back(); --len;
            } else continue;
        }
        dq.push_back(H[i]); ++len;
    }

    while (len > 2 && dq[0].out(inter(dq[len - 1], dq[len - 2]))) {
        dq.pop_back(); --len;
    }
    while (len > 2 && dq[len - 1].out(inter(dq[0], dq[1]))) {
        dq.pop_front(); --len;
    }

    if (len < 3) return vector<pt>();

    vector<pt> ret(len);
    for (int i = 0; i + 1 < len; i++) {
        ret[i] = inter(dq[i], dq[i + 1]);
    }
    ret.back() = inter(dq[len - 1], dq[0]);
    return ret;
}

// =============================================================================
//                              O(LOG N) CONVEX INSIDE
// =============================================================================

bool inConvex(vector<pt> &p, pt a) {
    if (p.size() < 3) return false;
    int n = p.size(), l = 1, r = n - 2, mid, ans = 1;
    while (l <= r) {
        mid = (l + r) / 2;
        if (sgn(orient(p[0], p[mid], a)) > 0) {
            l = mid + 1;
            ans = mid;
        } else {
            r = mid - 1;
        }
    }
    return inPolygon({p[0], p[ans], p[ans + 1]}, a, false);
}

bool StrictInConvex(vector<pt> &p, pt a) {
    if (p.size() < 3) return false;
    int n = p.size(), l = 1, r = n - 2, mid, ans = 1;
    while (l <= r) {
        mid = (l + r) / 2;
        if (sgn(orient(p[0], p[mid], a)) > 0) {
            l = mid + 1;
            ans = mid;
        } else {
            r = mid - 1;
        }
    }
    bool f = true;
    if (ans == 1) f = !onSegment(p[1], p[2], a) && !onSegment(p[0], p[1], a);
    else if (ans == n - 2) f = !onSegment(p[n - 2], p[n - 1], a) && !onSegment(p[n - 1], p[0], a);

    return inPolygon({p[0], p[ans], p[ans + 1]}, a, false) && f;
}

// =============================================================================
//                              ANTIPODAL POINTS
// =============================================================================

vector<pair<int, int>> all_anti_podal(int n, vector<pt> &p) {
    vector<pair<int, int>> result;
    auto nx = [&](int i) { return (i + 1) % n; };
    auto pv = [&](int i) { return (i - 1 + n) % n; };
    vector<bool> vis(n, false);

    for (int p1 = 0, p2 = 0; p1 < n; ++p1) {
        pt base = p[nx(p1)] - p[p1];
        while (p2 == p1 || p2 == nx(p1) || sgn(cross(base, p[nx(p2)] - p[p2])) == sgn(cross(base, p[p2] - p[pv(p2)]))) {
            p2 = nx(p2);
        }
        if (vis[p1]) continue;
        vis[p1] = true;
        result.push_back({p1, p2});
        result.push_back({nx(p1), p2});

        if (sgn(cross(base, p[nx(p2)] - p[p2])) == 0) {
            result.push_back({p1, nx(p2)});
            result.push_back({nx(p1), nx(p2)});
            vis[p2] = true;
        }
    }
    return result;
}

\end{lstlisting}
\subsection{any two lines intersect}
\begin{lstlisting}[language=C++]
/ Detects if ANY two line segments intersect in a set of N segments
// Uses a simplified Bentley-Ottmann Sweep-line Algorithm in O(N log N) time

#include <bits/stdc++.h>
using namespace std;
 
#define int long long
#define endl '\n'
 
// 1. Struct for exact integer math (Zero precision issues!)
struct pt {
    int x, y;
};
 
// Standard 2D cross product
int cross(pt a, pt b, pt c) {
    return (b.x - a.x) * (c.y - a.y) - (b.y - a.y) * (c.x - a.x);
}
 
int sgn(int val) {
    return (val > 0) - (val < 0);
}
 
// 1D bounding box check
bool intersect1D(int a, int b, int c, int d) {
    if (a > b) swap(a, b);
    if (c > d) swap(c, d);
    return max(a, c) <= min(b, d);
}
 
// Standard integer exact segment intersection
// Checks bounding boxes first, then uses cross products to ensure segments straddle each other
bool check_intersect(pt a, pt b, pt c, pt d) {
    return intersect1D(a.x, b.x, c.x, d.x) &&
           intersect1D(a.y, b.y, c.y, d.y) &&
           sgn(cross(a, b, c)) * sgn(cross(a, b, d)) <= 0 &&
           sgn(cross(c, d, a)) * sgn(cross(c, d, b)) <= 0;
}
 
// 2. Floating point struct just for ordering the Sweep Line
struct pt_f {
    double x, y;
};
 
struct seg {
    pt_f p, q;
    int id;
    pt orig_p, orig_q; // Keep original points to do the exact integer check
    
    // Get the Y coordinate of the segment at a specific sweep line X
    double get_y(double x) const {
        if (abs(p.x - q.x) < 1e-9) return p.y;
        return p.y + (q.y - p.y) * (x - p.x) / (q.x - p.x);
    }
};
 
double current_x;
 
// Custom comparator for the active window (Sweep Line)
// Maintains the vertical order of segments at the current X coordinate
struct cmp {
    bool operator()(const seg& a, const seg& b) const {
        double ya = a.get_y(current_x);
        double yb = b.get_y(current_x);
        if (abs(ya - yb) > 1e-9) return ya < yb;
        return a.id < b.id; // Tie-breaker
    }
};
 
struct event {
    double x;
    int type; // 1 for left endpoint (insert), -1 for right endpoint (remove)
    seg s;
    bool operator<(const event& e) const {
        if (abs(x - e.x) > 1e-9) return x < e.x;
        return type > e.type; // Process left endpoints before right ones if X is tied
    }
};
 
void solve() {
    int n;
    if (!(cin >> n)) return;
 
    vector<event> events;
    vector<seg> segments(n);
 
    // Trick: Rotate coordinates slightly to avoid vertical line nightmares.
    // Vertical lines break the y=mx+b sweep line logic because multiple points share the same X.
    double angle = 0.001;
    double sinA = sin(angle), cosA = cos(angle);
 
    for (int i = 0; i < n; i++) {
        int x1, y1, x2, y2;
        cin >> x1 >> y1 >> x2 >> y2;
        
        segments[i].orig_p = {x1, y1};
        segments[i].orig_q = {x2, y2};
        segments[i].id = i + 1; // 1-based index for the output
 
        // Apply rotation to the floating-point representations
        pt_f p = {x1 * cosA - y1 * sinA, x1 * sinA + y1 * cosA};
        pt_f q = {x2 * cosA - y2 * sinA, x2 * sinA + y2 * cosA};
 
        // Always sweep left to right
        if (p.x > q.x) swap(p, q);
 
        segments[i].p = p;
        segments[i].q = q;
 
        // Create left and right events for the sweep line
        events.push_back({p.x, 1, segments[i]});
        events.push_back({q.x, -1, segments[i]});
    }
 
    sort(events.begin(), events.end());
 
    set<seg, cmp> active;
    vector<set<seg>::iterator> where(n + 1); // Keep track of elements in the set for quick deletion
 
    for (auto& ev : events) {
        current_x = ev.x;
        int id = ev.s.id;
 
        if (ev.type == 1) { 
            // 1. Insert segment into the active window
            auto nxt = active.insert(ev.s).first;
            where[id] = nxt;
            
            // 2. Check if it intersects with the segment immediately above or below it
            auto prv = nxt, nxt2 = nxt;
            if (prv != active.begin()) {
                prv--;
                if (check_intersect(ev.s.orig_p, ev.s.orig_q, prv->orig_p, prv->orig_q)) {
                    cout << "YES\n" << ev.s.id << " " << prv->id << "\n";
                    return;
                }
            }
            nxt2++;
            if (nxt2 != active.end()) {
                if (check_intersect(ev.s.orig_p, ev.s.orig_q, nxt2->orig_p, nxt2->orig_q)) {
                    cout << "YES\n" << ev.s.id << " " << nxt2->id << "\n";
                    return;
                }
            }
        } else {
            // 1. Find the segment in the active window
            auto it = where[id];
            
            // 2. Before removing it, check if its top and bottom neighbors now intersect
            auto prv = it, nxt = it;
            nxt++;
            if (prv != active.begin() && nxt != active.end()) {
                prv--;
                if (check_intersect(prv->orig_p, prv->orig_q, nxt->orig_p, nxt->orig_q)) {
                    cout << "YES\n" << prv->id << " " << nxt->id << "\n";
                    return;
                }
            }
            // 3. Remove the segment
            active.erase(it);
        }
    }
    
    // If the sweep line finishes with no hits
    cout << "NO\n";
}
\end{lstlisting}
\subsection{areaOfAnyShape}
\begin{lstlisting}[language=C++]
// Calculates the area of any simple polygon (convex or concave)
// Uses the Shoelace Formula (Surveyor's formula)

typedef long long ll;
typedef pair<ll, ll> pt;

// Computes signed area  2 (keeps operations in integer math to avoid precision loss)
ll twiceArea(const vector<pt> &p) {
    int n = p.size();
    ll area = 0;
    for (int i = 0; i < n; i++) {
        ll x1 = p[i].first, y1 = p[i].second;
        // Wrap around to the first point when i is at the last vertex
        ll x2 = p[(i + 1) % n].first, y2 = p[(i + 1) % n].second;  
        
        // Accumulate the 2D cross product of adjacent vertices
        area += (x1 * y2 - x2 * y1);
    }
    // Return absolute value so it works regardless of clockwise or counter-clockwise input
    return abs(area);
}

void solve() {
    int n;
    cin >> n;
    vector<pt> p(n);
    for (int i = 0; i < n; ++i)
        cin >> p[i].first >> p[i].second;

    ll twice = twiceArea(p);
    // Divide by 2.0 at the very end to get the actual floating-point area
    cout << fixed << setprecision(1) << (twice / 2.0) << '\n';
}
\end{lstlisting}
\subsection{geometryReferenceSheet}
\begin{lstlisting}[language=C++]
// =============================================================================
//                              BASICS & TYPES
// =============================================================================
int sgn(T val);                  // Returns -1 (negative), 1 (positive), or 0.
T sq(pt p);                      // Squared length of vector p.
ld abs(pt p);                    // Length of vector p.
pt perp(pt p);                   // Returns vector rotated 90 degrees CCW.
T dot(pt v, pt w);               // Dot product.
T cross(pt v, pt w);             // Cross product (v X w).
bool isPerp(pt v, pt w);         // True if dot product is 0.
pt unitVector(pt p);             // Vector of length 1 in the same direction.

// =============================================================================
//                              TRANSFORMATIONS
// =============================================================================
pt translate(pt v, pt p);        // Moves point p by vector v.
pt scale(pt c, T factor, pt p);  // Scales p away from center c by factor.
pt rot(pt p, pt c, ld a);        // Rotates p around c by angle a (in radians).
pt linearTransfo(pt p, pt q, pt r, pt fp, pt fq); // Maps r given that p->fp and q->fq.

// =============================================================================
//                              ANGLES
// =============================================================================
T orient(pt a, pt b, pt c);      // Relative to AB: > 0 (C is left), < 0 (C is right), 0 (collinear).
bool inAngle(pt a, pt b, pt c, pt p); // True if p is inside CCW angle BAC.
T angle(pt v, pt w);             // Shortest positive non-oriented angle [0, PI].
ld orientedAngle(pt a, pt b, pt c); // CCW angle BAC in range [0, 2*PI).
ld angleTravelled(pt a, pt p, pt q); // Signed angle travelled from P to Q around A.

// =============================================================================
//                              LINES (struct line)
// =============================================================================
line(pt v, T c);                 // From direction vector and offset.
line(T a, T b, T _c);            // From equation ax + by = c.
line(pt p, pt q);                // From two points P and Q.
T side(pt p);                    // > 0 (left of line), < 0 (right of line), 0 (on line).
double dist(pt p);               // Shortest distance from p to line.
line perpThrough(pt p);          // Returns perpendicular line passing through p.
line translate(pt t);            // Translates line by vector t.
line shiftLeft(double dist);     // Shifts line to the left by given distance.
pt proj(pt p);                   // Projection of point p onto the line.
pt refl(pt p);                   // Reflection of point p across the line.

bool inter(line l1, line l2, pt &out); // Line-line intersection. Returns false if parallel.
line bisector(line l1, line l2, bool interior); // Angle bisector of l1 and l2.

// =============================================================================
//                              SEGMENTS & RAYS
// =============================================================================
bool onSegment(pt a, pt b, pt p); // True if p lies on segment AB.
bool properInter(pt a, pt b, pt c, pt d, pt &out); // True if segments AB and CD intersect at exactly one interior point.
set<pair<ld, ld>> inters(pt a, pt b, pt c, pt d); // Returns all intersection points of segments AB and CD.
ld segPoint(pt a, pt b, pt p);    // Shortest distance from p to segment AB.
ld rayPoint(pt a, pt b, pt p);    // Shortest distance from p to ray AB.
ld segSeg(pt a, pt b, pt c, pt d); // Shortest distance between segment AB and segment CD.

// =============================================================================
//                              CIRCLES
// =============================================================================
pair<pt, T> circumCircle(pt a, pt b, pt c); // Returns {center, radius} of triangle ABC.
int circleLine(pt o, double r, line l, pair<pt, pt> &out); // Intersections of line and circle. Returns count (0, 1, or 2).
int circleCircle(pt o1, T r1, pt o2, T r2, pair<pt, pt> &out); // Intersections of 2 circles. Returns count.
int tangents(pt o1, T r1, pt o2, T r2, bool inner, vector<pair<pt, pt>> &out); // Tangent lines between 2 circles. 

// =============================================================================
//                              BASIC POLYGONS
// =============================================================================
ld areaPolygon(vector<pt> p);       // Area of any simple polygon.
bool isConvex(vector<pt> p);        // True if polygon is strictly convex.
bool inPolygon(vector<pt> p, pt a, bool strict = true); // O(N). True if a is inside. strict=true excludes boundaries.
pt polygonCentroid(const vector<pt>& p); // Center of gravity. Vertices must be ordered (CW/CCW).

// =============================================================================
//                              PICK'S THEOREM (Integer Coordinates)
// =============================================================================
int boundaryPoints(const vector<pt>& p);   // Number of integer points on boundary.
int doubleAreaPolygon(const vector<pt>& p); // 2 * Area (exact integer).
int interiorPoints(const vector<pt>& p);   // Number of strictly interior integer points.

// =============================================================================
//                              ADVANCED POLYGONS
// =============================================================================
void convex_hull(vector<pt>& a, bool include_collinear = false); // O(N log N). Modifies vector in-place.
vector<pt> minkowski(vector<pt> P, vector<pt> Q); // O(N+M). Minkowski sum. P and Q must be CCW convex.
double maximum_dist_from_polygon_to_polygon(vector<pt> &u, vector<pt> &v); // O(N+M). Max distance between two convex polygons.

// =============================================================================
//                              HALF-PLANE INTERSECTION
// =============================================================================
Halfplane(const pt& a, const pt& b); // Halfplane to the LEFT of directed line AB.
vector<pt> hp_intersect(vector<Halfplane>& H); // O(N log N). Returns CCW convex polygon of the intersection. Empty vector if no intersection.

// =============================================================================
//                              O(LOG N) CONVEX INSIDE
// =============================================================================
bool inConvex(vector<pt> &p, pt a); // O(log N). True if point inside convex polygon (including boundary).
bool StrictInConvex(vector<pt> &p, pt a); // O(log N). True if point strictly inside convex polygon.

// =============================================================================
//                              ANTIPODAL POINTS
// =============================================================================
vector<pair<int, int>> all_anti_podal(int n, vector<pt> &p); // O(N). Returns pairs of indices of antipodal points in a convex polygon.

\end{lstlisting}
\subsection{minimum distance calculation between points}
\begin{lstlisting}[language=C++]
// Finds the Closest Pair of Points in a 2D plane in O(N log N) time
// Uses a Sweep-line algorithm combined with a balanced BST (std::set)

#include <iostream>
#include <vector>
#include <cmath>
#include <algorithm>
#include <set>
#include <iomanip>

using namespace std;

// Struct to hold point data, including original index for the final answer
struct Point {
    double x, y;
    int id; 
};

// 1. Sort primarily by X-coordinate to process points from left to right
bool compareX(const Point& a, const Point& b) {
    if (a.x != b.x) return a.x < b.x;
    return a.y < b.y;
}

// 2. Custom comparator for the active window set to sort points by Y-coordinate
struct compareY {
    bool operator()(const Point& a, const Point& b) const {
        if (a.y != b.y) return a.y < b.y;
        return a.x < b.x; // Tie-breaker ensures points with same Y aren't treated as identical
    }
};

// Standard Euclidean distance
double get_dist(Point a, Point b) {
    return sqrt((a.x - b.x) * (a.x - b.x) + (a.y - b.y) * (a.y - b.y));
}

void solve() {
    int n;
    if (!(cin >> n)) return;
    
    vector<Point> pts(n);
    for (int i = 0; i < n; i++) {
        cin >> pts[i].x >> pts[i].y;
        pts[i].id = i; // Save original 0-based index
    }

    // Sort all points from left to right for the sweep-line
    sort(pts.begin(), pts.end(), compareX);

    // Active window storing points currently within 'min_d' distance of the sweep line
    set<Point, compareY> active_set;
    
    double min_d = 1e18; // Start with a massive distance
    int best_i = -1, best_j = -1;
    
    int left = 0; // Left boundary pointer of our sweeping window
    
    for (int i = 0; i < n; i++) {
        // Step A: Shrink the window
        // Remove points too far left to possibly form a closer pair with the current point
        while (left < i && pts[i].x - pts[left].x >= min_d) {
            active_set.erase(pts[left]);
            left++;
        }
        
        // Step B: Check Y-bounds
        // Only check points in the active set whose Y is within [current_y - min_d, current_y + min_d]
        Point search_point = {pts[i].x, pts[i].y - min_d, -1};
        auto it = active_set.lower_bound(search_point); 
        
        while (it != active_set.end() && it->y - pts[i].y < min_d) {
            double current_d = get_dist(pts[i], *it);
            
            // Step C: Update if a strictly closer pair is found
            if (current_d < min_d) {
                min_d = current_d;
                best_i = pts[i].id;
                best_j = it->id;
            }
            it++;
        }
        
        // Step D: Add the current point into the active window for future comparisons
        active_set.insert(pts[i]);
    }
    
    // Ensure the smaller index is printed first
    if (best_i > best_j) swap(best_i, best_j);
    
    cout << best_i << " " << best_j << " " << fixed << setprecision(6) << min_d << "\n";
}

int main() {
    ios_base::sync_with_stdio(false);
    cin.tie(NULL);
    solve();
    return 0;
}
\end{lstlisting}
\subsection{rectangle union}
\begin{lstlisting}[language=C++]
// To find the total area covered by a union of given rectangles
// Uses a Sweep-line algorithm paired with a Lazy Segment Tree

#include <bits/stdc++.h>
 
using namespace std;
 
typedef long long ll;
typedef long double ld;
 
#define endl '\n'
#define all(a)   a.begin(),a.end()
#define all_r(a)   a.rbegin(),a.rend()
#define int ll
 
const int oo = 1e9;
 
struct segtree{
    struct Node{
        int mn, freqmn, lazy, isLazy;
        Node(){
            mn = oo, freqmn = 0;
            isLazy = 0;
            lazy = 0;
        }
        void change(int x, int lx, int rx){
            mn += x;
            lazy += x;
            isLazy = 1;
        }
        Node(int x){
            mn = x;
            freqmn = 1; // Represents the length of this minimum value's segment
            isLazy = 0;
            lazy = 0;
        }
    };
 
    vector<Node> segData;
    int treeSize;
 
    void propagate(int nx, int lx, int rx){
        if(!segData[nx].isLazy || rx - lx == 1)
            return;
        int mid = lx + (rx - lx)/2;
        segData[2 * nx + 1].change(segData[nx].lazy, lx, mid);
        segData[2 * nx + 2].change(segData[nx].lazy, mid, rx);
        segData[nx].lazy = segData[nx].isLazy = 0;
    }
 
    segtree(int n){
        treeSize = 1;
        while(treeSize < n)
            treeSize *= 2;
        segData.assign(2 * treeSize, Node());
    }
 
    Node merge(Node lf, Node ri){
        Node ans = Node();
        ans.mn = min(lf.mn, ri.mn);
        // Combine frequencies of the minimums to track how much of the range is uncovered
        if(lf.mn == ans.mn) ans.freqmn += lf.freqmn;
        if(ri.mn == ans.mn) ans.freqmn += ri.freqmn;
        return ans;
    }
 
    void set(int l, int r, int val, int nx, int lx, int rx){
        propagate(nx, lx, rx);
        if(rx <= l || lx >= r) return;
        if(lx >= l && rx <= r){
            segData[nx].change(val, lx, rx);
            return;
        }
 
        int mid = lx + (rx - lx)/2;
 
        set(l, r, val, 2 * nx + 1, lx, mid);
        set(l, r, val, 2 * nx + 2, mid, rx);
        segData[nx] = merge(segData[2 * nx + 1], segData[2 * nx + 2]);
    }
 
    void set(int l, int r, int val){
        set(l, r, val, 0, 0, treeSize);
    }
 
    Node get(int l, int r, int nx, int lx, int rx){
        propagate(nx, lx, rx);
        if(l <= lx && r >= rx) return segData[nx];
        if(r <= lx || l >= rx) return Node();
 
        int mid = lx + (rx - lx)/2;
        return merge(get(l, r, 2 * nx + 1, lx, mid) ,get(l, r, 2 * nx + 2, mid, rx));
    }
 
    // Returns the total length of the active Y-intervals
    int get(int l, int r){
        Node ans = get(l, r, 0, 0, treeSize);
        // If the minimum is 0, subtract the frequency of 0s (uncovered space) from total range
        // If the minimum > 0, the entire range is completely covered by at least one rectangle
        return (r - l) - (ans.mn ? 0 : ans.freqmn);
    }
 
    void Build(vector<int> & v, int nx, int lx, int rx){
        if(rx - lx == 1){
            if(lx < v.size()){
                segData[nx]= Node(v[lx]);
            }
            return;
        }
 
        int mid = lx + (rx - lx)/2;
 
        Build(v, 2 * nx + 1, lx, mid);
        Build(v, 2 * nx + 2, mid, rx);
 
        segData[nx] = merge(segData[2 * nx + 1], segData[2 * nx + 2]);
    }
 
    void Build(vector<int> & v){
        Build(v, 0, 0, treeSize);
    }
};
 
// Represents a vertical line of a rectangle
struct event{
    int x, t, y1, y2; // x: coordinate, t: 1 for left edge / 0 for right edge
};
 
const int N = 2e6 + 1, sh = 1e6; // 'sh' offsets negative Y coordinates

void solve() {
    int n; cin >> n;
    vector<event> lines(2 * n);
    for (int i = 0; i < 2 * n; i += 2) {
        int x1, y1, x2, y2;
        cin >> x1 >> y1 >> x2 >> y2;
        // Map Y coordinates to positive values using shift 'sh'
        lines[i] = { x1, 1, y1 + sh, y2 + sh};
        lines[i + 1] = {x2, 0, y1 + sh, y2 + sh};
    }
 
    // Sort events from left to right. Process left edges before right edges if at the same X.
    sort(all(lines), [&](event & a, event & b)->bool{
        if (a.x != b.x) return a.x < b.x;
        return a.t < b.t;
    });
 
    int ans = 0;
 
    vector<int> b(N);
    segtree seg(N);
    seg.Build(b);
 
    int lastx = -1;
    
    for(auto& ev : lines){
        int x = ev.x;
        int t = ev.t;
        int y1 = ev.y1;
        int y2 = ev.y2;
        
        // Add (width * active_height) to total area
        ans += seg.get(0, N) * (x - lastx);
        
        // If it's a left edge, add 1 to the interval. If right edge, subtract 1.
        if(t) seg.set(y1, y2, 1);
        else seg.set(y1, y2, -1);
        
        lastx = x;
    }
 
    cout <<  ans << endl;
}
 
signed main() {
    ios_base::sync_with_stdio(0);cin.tie(0);cout.tie(0);
    cout << fixed << setprecision(9);
    int test = 1; //cin>>test;
    while (test--) solve();
    
    return 0;
}

\end{lstlisting}
\section{String}
\subsection{Aho Corasick}
\subsubsection{Aho corasick}
\begin{lstlisting}[language=C++]
/**
 *  THE ULTIMATE CP TEMPLATE BLUEPRINT: AHO-CORASICK
 *
 * 1.  WHAT PROBLEM DOES THIS SOLVE?
 * - Keywords: "Multiple patterns", "Dictionary matching", "String DP".
 * - Classic Scenarios: Find all occurrences of an array of strings inside a massive text,
 *   or count how many strings of length N don't contain any of the dictionary words.
 * - The Magic: Builds a Trie and connects nodes via "fail" links. If you mismatch,
 *   the fail link teleports you to the longest proper suffix that exists in the Trie.
 *
 * 2.  HOW TO USE IT (THE BLACK BOX)
 * - Initialization:
 *       AhoCorasick ac;
 *       for (string s : patterns) ac.insert(s);
 *       ac.build();
 * - Query:
 *       int total_matches = ac.search(text);
 *
 * 3.  HOW TO ADAPT IT (THE DIALS & SWITCHES)
 * - DP Transitions: The `next_state(u, c)` transitions are precalculated during `build()`.
 *   You can use `ac.trie[u].next[c]` to transition seamlessly in Dynamic Programming.
 * - Exit Link (dict_link): Fast-forwards to the next valid dictionary word ending to
 *   avoid O(N) worst-case walking up the failure tree.
 */

const int K = 26;

struct Vertex
{
    int next[K];
    bool leaf = false;
    int p = -1;
    char pch;
    int link = -1;
    int exit_link = -1;
    int matches = 0; // Count of dictionary words ending here

    Vertex(int p = -1, char ch = '$') : p(p), pch(ch)
    {
        fill(begin(next), end(next), -1);
    }
};

struct AhoCorasick
{
    vector<Vertex> t;

    AhoCorasick()
    {
        t.emplace_back();
    }

    void insert(const string &s)
    {
        int v = 0;
        for (char ch : s)
        {
            int c = ch - 'a';
            if (t[v].next[c] == -1)
            {
                t[v].next[c] = t.size();
                t.emplace_back(v, ch);
            }
            v = t[v].next[c];
        }
        t[v].leaf = true;
        t[v].matches++;
    }

    void build()
    {
        queue<int> q;
        q.push(0);
        t[0].link = 0;
        t[0].exit_link = 0;

        while (!q.empty())
        {
            int v = q.front();
            q.pop();

            for (int c = 0; c < K; c++)
            {
                int u = t[v].next[c];
                if (u != -1)
                {
                    t[u].link = (v == 0) ? 0 : t[t[v].link].next[c];
                    t[u].exit_link = t[t[u].link].leaf ? t[u].link : t[t[u].link].exit_link;
                    t[u].matches += t[t[u].link].matches;
                    q.push(u);
                }
                else
                {
                    t[v].next[c] = (v == 0) ? 0 : t[t[v].link].next[c];
                }
            }
        }
    }

    long long search(const string &text)
    {
        long long total = 0;
        int v = 0;
        for (char ch : text)
        {
            v = t[v].next[ch - 'a'];
            total += t[v].matches;
        }
        return total;
    }
    long long search_non_overlapping(const string &text)
    {
        long long total = 0;
        int v = 0;
        for (char ch : text)
        {
            v = t[v].next[ch - 'a'];

            if (t[v].matches > 0)
            {
                total++;
                v = 0;
            }
        }
        return total;
    }
};

\end{lstlisting}
\subsection{Hashing}
\subsubsection{2D string hashing}
\begin{lstlisting}[language=C++]
/**
 *  THE ULTIMATE CP TEMPLATE BLUEPRINT: 2D STRING HASHING
 *
 * 1.  WHAT PROBLEM DOES THIS SOLVE?
 * - Keywords: "Subgrid matching", "2D Pattern matching".
 * - Classic Scenarios: Find how many times an R x C pattern appears inside
 *   an N x M grid.
 * - The Magic: Extends the rolling hash to two dimensions using two different
 *   bases (one for rows, one for columns). Calculates 2D prefix sums of the hash.
 *
 * 2.  HOW TO USE IT (THE BLACK BOX)
 * - Initialization:
 *       Hash2D h(grid);
 * - Query:
 *       ll val = h.get(r1, c1, r2, c2); // 0-indexed, inclusive boundaries
 */

using ll = long long;

struct Hash2D
{
    const ll MOD = 1e9 + 7;
    const ll BASE_R = 313, BASE_C = 317;
    vector<vector<ll>> hash;
    vector<ll> pR, pC;

    Hash2D(const vector<string> &grid)
    {
        if (grid.empty())
            return;
        int n = grid.size();
        int m = grid[0].size();

        hash.assign(n + 1, vector<ll>(m + 1, 0));
        pR.assign(n + 1, 1);
        pC.assign(m + 1, 1);

        for (int i = 1; i <= n; i++)
            pR[i] = (pR[i - 1] * BASE_R) % MOD;
        for (int i = 1; i <= m; i++)
            pC[i] = (pC[i - 1] * BASE_C) % MOD;

        for (int i = 1; i <= n; i++)
        {
            for (int j = 1; j <= m; j++)
            {
                hash[i][j] = (hash[i - 1][j] * BASE_R + hash[i][j - 1] * BASE_C - hash[i - 1][j - 1] * BASE_R % MOD * BASE_C % MOD + grid[i - 1][j - 1]) % MOD;
                if (hash[i][j] < 0)
                    hash[i][j] += MOD;
            }
        }
    }

    ll get(int r1, int c1, int r2, int c2)
    {
        r1++;
        c1++;
        r2++;
        c2++;
        ll h = hash[r2][c2];
        h = (h - hash[r1 - 1][c2] * pR[r2 - r1 + 1]) % MOD;
        h = (h - hash[r2][c1 - 1] * pC[c2 - c1 + 1]) % MOD;
        h = (h + hash[r1 - 1][c1 - 1] * pR[r2 - r1 + 1] % MOD * pC[c2 - c1 + 1]) % MOD;
        if (h < 0)
            h += MOD;
        return h;
    }
};
\end{lstlisting}
\subsubsection{Hash(Hassan)}
\begin{lstlisting}[language=C++]
const int mod1 = 1e9 + 7;
const int mod2 = 1e9 + 9;
const int base1 = 313;
const int base2 = 317;

const int MAXN = 3e4 + 100;
int pw1[MAXN], pw2[MAXN];
bool isCalPower = false;

void pre() {
    if (isCalPower) return;

    pw1[0] = pw2[0] = 1;
    for (int i = 1; i < MAXN; i++) {
        pw1[i] = (pw1[i - 1] * base1) % mod1;
        pw2[i] = (pw2[i - 1] * base2) % mod2;
    }
    isCalPower = true;
}

struct Hash {
    int n;
    vector<pair<int, int>> pref;
    vector<pair<int, int>> suf; // Used only if palindrome checking is needed
 
    // Pass `compute_rev = true` if you need O(1) palindrome checks
    Hash(const string& s, bool compute_rev = false) {
        pre();
        n = s.size();
        pref.assign(n + 1, {0, 0});

        for (int i = 0; i < n; i++) {
            pref[i + 1].first = (pref[i].first * base1 + s[i]) % mod1;
            pref[i + 1].second = (pref[i].second * base2 + s[i]) % mod2;
        }

        if (compute_rev) {
            suf.assign(n + 1, {0, 0});
            for (int i = n - 1; i >= 0; i--) {
                suf[i].first = (suf[i + 1].first * base1 + s[i]) % mod1;
                suf[i].second = (suf[i + 1].second * base2 + s[i]) % mod2;
            }
        }
    }

    // Get hash of substring s[l...r] (0-indexed) in O(1)
    pair<int, int> get(int l, int r) {
        if (l > r) return {0, 0};
        int len = r - l + 1;
        
        int h1 = (pref[r+1].first - (pref[l].first * pw1[len]) % mod1) % mod1;
        if (h1 < 0) h1 += mod1;

        int h2 = (pref[r+1].second - (pref[l].second * pw2[len]) % mod2) % mod2;
        if (h2 < 0) h2 += mod2;

        return {h1, h2};
    }

    // Get reverse hash of substring s[l...r] in O(1)
    pair<int, int> get_rev(int l, int r) {
        if (l > r || suf.empty()) return {0, 0};
        int len = r - l + 1;
        int h1 = (suf[l].first - suf[r + 1].first * pw1[len]) % mod1;
        if (h1 < 0) h1 += mod1;

        int h2 = (suf[l].second - suf[r + 1].second * pw2[len]) % mod2;
        if (h2 < 0) h2 += mod2;

        return {h1, h2};
    }

    // Check if s[l...r] is a palindrome in O(1)
    bool is_pal(int l, int r) {
        return get(l, r) == get_rev(l, r);
    }
};


// --- USEFUL HASHING UTILITY FUNCTIONS ---

// Concatenate two hashes (h1 followed by h2 of length len2) in O(1)
pair<int, int> concat_hash(pair<int, int> h1, pair<int, int> h2, int len2) {
    int ans1 = (h1.first * pw1[len2] + h2.first) % mod1;
    int ans2 = (h1.second * pw2[len2] + h2.second) % mod2;
    return {ans1, ans2};
}

// Find Longest Common Prefix (LCP) length of s1[l1...r1] and s2[l2...r2] in O(log N)
int get_lcp(Hash& sh1, int l1, int r1, Hash& sh2, int l2, int r2) {
    int low = 1, high = min(r1 - l1 + 1, r2 - l2 + 1), ans = 0;
    while (low <= high) {
        int mid = low + (high - low) / 2;
        if (sh1.get(l1, l1 + mid - 1) == sh2.get(l2, l2 + mid - 1)) {
            ans = mid;
            low = mid + 1;
        } else {
            high = mid - 1;
        }
    }
    return ans;
}

// Compare s1[l1...r1] and s2[l2...r2] lexicographically in O(log N)
// Returns: -1 if s1 < s2, 0 if s1 == s2, 1 if s1 > s2
int compare_substrings(const string& s1, Hash& sh1, int l1, int r1, 
                       const string& s2, Hash& sh2, int l2, int r2) {
    int lcp_len = get_lcp(sh1, l1, r1, sh2, l2, r2);
    int len1 = r1 - l1 + 1;
    int len2 = r2 - l2 + 1;
    
    // If one is completely a prefix of another
    if (lcp_len == len1 && lcp_len == len2) return 0;
    if (lcp_len == len1) return -1;
    if (lcp_len == len2) return 1;
    
    // Compare the first differing character
    return s1[l1 + lcp_len] < s2[l2 + lcp_len] ? -1 : 1;
}

\end{lstlisting}
\subsubsection{Hash}
\begin{lstlisting}[language=C++]
struct Hash
{
    int base, mod1, mod2, inv1, inv2, n;
    vector<int> pw1, pw2, invpw1, invpw2, pre1, pre2;

    int power(int b, int p, int m)
    {
        int res{1};
        b %= m;
        while (p)
        {
            if (p & 1)
                res = res * b % m;
            b = b * b % m;
            p >>= 1;
        }
        return res;
    }

    Hash(const string &s, int b)
    {
        base = b;
        mod1 = 1e9 + 7;
        mod2 = 1e9 + 9;
        n = int(s.size());

        inv1 = power(base, mod1 - 2, mod1);
        inv2 = power(base, mod2 - 2, mod2);

        pw1.assign(n + 1, 1), pw2.assign(n + 1, 1);
        invpw1.assign(n + 1, 1), invpw2.assign(n + 1, 1);
        pre1.assign(n + 1, {}), pre2.assign(n + 1, {});

        for (int i = 0; i < n; i++)
        {
            pw1[i + 1] = pw1[i] * base % mod1;
            pw2[i + 1] = pw2[i] * base % mod2;
            invpw1[i + 1] = invpw1[i] * inv1 % mod1;
            invpw2[i + 1] = invpw2[i] * inv2 % mod2;
            pre1[i + 1] = (pre1[i] + pw1[i] * s[i]) % mod1;
            pre2[i + 1] = (pre2[i] + pw2[i] * s[i]) % mod2;
        }
    }

    pair<int, int> get(int l, int r)
    {
        return {(pre1[r + 1] - pre1[l] + mod1) * invpw1[l] % mod1, (pre2[r + 1] - pre2[l] + mod2) * invpw2[l] % mod2};
    }
};

\end{lstlisting}
\subsubsection{MultisetHashing}
\begin{lstlisting}[language=C++]
struct MultisetHashing
{
    int base, m, n;
    vector<int> pre;

    int power(int b, int p)
    {
        b %= m;
        int res{1};
        while (p)
        {
            if (p & 1)
                res = res * b % m;
            b = b * b % m;
            p >>= 1;
        }
        return res;
    }

    MultisetHashing(vector<int> &arr, int b)
    {
        n = int(arr.size());
        m = 1e9 + 7;
        base = b;
        pre.assign(n + 1, {});
        for (int i = 0; i < n; i++)
            pre[i + 1] = (pre[i] + power(base, arr[i])) % m;
    }

    int get(int l, int r)
    {
        // 0-based
        return (pre[r + 1] - pre[l] + m) % m;
    }
};

\end{lstlisting}
\subsubsection{XOR hashing}
\begin{lstlisting}[language=C++]
struct XorHashing
{
    int n;
    map<int, int> h;
    vector<int> pre;

    XorHashing(vector<int> &arr)
    {
        n = int(arr.size());
        pre.assign(n + 1, {});
        for (int i = 0; i < n; i++)
        {
            if (!h.count(arr[i]))
                h[arr[i]] = getRandom(1, 1e18);
            pre[i + 1] = pre[i] ^ h[arr[i]];
        }
    }

    int get(int l, int r)
    {
        // 0-based
        return pre[r + 1] ^ pre[l];
    }
};

\end{lstlisting}
\subsection{LCS}
\subsubsection{Bitset LCS}
\begin{lstlisting}[language=C++]
/**
 *  THE ULTIMATE CP TEMPLATE BLUEPRINT: BITSET LCS
 *
 * 1.  WHAT PROBLEM DOES THIS SOLVE?
 * - Keywords: "LCS of massive strings", "O(N^2) TLE".
 * - Classic Scenarios: Find the Longest Common Subsequence of two strings where
 *   N and M are up to 50,000.
 * - The Magic: C++ `std::bitset` does not support subtraction, which is required
 *   for the classic bitwise LCS formula. This custom block engine maps the DP states
 *   directly onto native 64-bit unsigned integers (uint64_t). It computes 64 DP
 *   cells simultaneously using carry-flag subtraction, slashing the O(N * M) time
 *   by a factor of 64.
 *
 * 2.  HOW TO USE IT (THE BLACK BOX)
 * - Query:
 *       int length = bitset_lcs(string_A, string_B);
 *
 * - Complexity:
 *       Time: O(N * M / 64)
 *       Space: O(M / 64)
 */

#include <cstdint>

int bitset_lcs(const string &a, const string &b)
{
    int n = a.length(), m = b.length();
    if (n == 0 || m == 0)
        return 0;

    int words = (m >> 6) + 1;
    vector<vector<uint64_t>> pos(256, vector<uint64_t>(words, 0));

    for (int i = 0; i < m; i++)
    {
        pos[b[i]][i >> 6] |= (1ULL << (i & 63));
    }

    vector<uint64_t> D(words, 0);

    for (int i = 0; i < n; i++)
    {
        uint64_t b_carry = 1;
        uint64_t shift_carry = 0;

        for (int j = 0; j < words; j++)
        {
            uint64_t P = D[j] | pos[a[i]][j];
            uint64_t d_shl = (D[j] << 1) | shift_carry;
            shift_carry = D[j] >> 63;

            uint64_t sub = d_shl + b_carry;
            uint64_t next_carry = (P < sub) || (sub < d_shl);
            uint64_t diff = P - sub;

            D[j] = P & (~P ^ diff);
            b_carry = next_carry;
        }
    }

    int lcs_length = 0;
    for (int j = 0; j < words; j++)
    {
        uint64_t val = D[j];
        while (val)
        {
            val &= (val - 1);
            lcs_length++;
        }
    }

    return lcs_length;
}
\end{lstlisting}
\subsection{Matching}
\subsubsection{FFT wildcard matching}
\begin{lstlisting}[language=C++]
/**
 *  THE ULTIMATE CP TEMPLATE BLUEPRINT: FFT WILDCARD MATCHING
 *
 * 1.  WHAT PROBLEM DOES THIS SOLVE?
 * - Keywords: "Wildcards in both text and pattern", "Massive strings".
 * - Classic Scenarios: Both strings have '?' wildcards. Bitsets will TLE if
 *   N = 200,000.
 * - The Magic: Uses the algebraic identity (T - P)^2 * T * P = 0. Maps wildcards
 *   to 0 and letters to 1-26. Expands the polynomial into three separate NTT
 *   convolutions. A result of 0 at index i means a perfect match.
 *
 * 2.  HOW TO USE IT (THE BLACK BOX)
 * - Query:
 *       vector<int> matches = fft_wildcard_match(text, pattern);
 *
 * 3.  HOW TO ADAPT IT
 * - Dependencies: This heavily relies on your `multiply(A, B)` function from
 *   your NTT template. Make sure it is included above this.
 */

using ll = long long;

vector<ll> multiply(vector<ll> const &a, vector<ll> const &b);

vector<int> fft_wildcard_match(const string &text, const string &pattern)
{
    int n = text.length(), m = pattern.length();
    if (m > n)
        return {};

    vector<ll> t1(n), t2(n), t3(n);
    vector<ll> p1(m), p2(m), p3(m);

    for (int i = 0; i < n; i++)
    {
        ll val = (text[i] == '?' ? 0 : text[i] - 'a' + 1);
        t1[i] = val;
        t2[i] = val * val;
        t3[i] = val * val * val;
    }

    for (int i = 0; i < m; i++)
    {
        ll val = (pattern[i] == '?' ? 0 : pattern[i] - 'a' + 1);
        p1[m - 1 - i] = val;
        p2[m - 1 - i] = val * val;
        p3[m - 1 - i] = val * val * val;
    }

    vector<ll> res1 = multiply(t3, p1);
    vector<ll> res2 = multiply(t2, p2);
    vector<ll> res3 = multiply(t1, p3);

    vector<int> matches;
    for (int i = m - 1; i < n; i++)
    {
        ll diff = res1[i] - 2 * res2[i] + res3[i];
        if (diff == 0)
        {
            matches.push_back(i - (m - 1));
        }
    }
    return matches;
}
\end{lstlisting}
\subsubsection{KMP}
\begin{lstlisting}[language=C++]
vector<int> KMP(string s, string p = "")
{
    if (!p.empty())
        s = p + "$" + s;
    int n = int(s.size());
    vector<int> LPS(n);
    int i{1}, j{};
    while (i < n)
    {
        if (s[i] == s[j])
        {
            j++;
            LPS[i] = j;
            i++;
        }
        else
        {
            if (j == 0)
            {
                LPS[i] = 0;
                i++;
            }
            else
                j = LPS[j - 1];
        }
    }
    return LPS;
}

int numberOfMatches(string &s, string &p)
{
    vector<int> pi = KMP(s, p);
    int ans{};
    for (auto &val : pi)
        ans += val == int(p.size());
    return ans;
}

vector<int> countPrefixMatching(string &s)
{
    int n = int(s.size());
    vector<int> pi = KMP(s);
    vector<int> frq(n);
    for (auto &x : pi)
        frq[x]++;
    for (int i = n - 1; i; i--)
        frq[pi[i - 1]] += frq[i];
    return frq;
}

vector<vector<int>> KMPAutomaton(string &s)
{
    vector<int> pi = KMP(s);
    int n = int(s.size());
    vector aut(n + 1, vector<int>(26));
    aut[0][s[0] - 'a'] = 1;
    for (int i = 1; i <= n; i++)
    {
        for (int j = 0; j < 26; j++)
        {
            if (i < n && s[i] - 'a' == j)
                aut[i][j] = i + 1;
            else
                aut[i][j] = aut[pi[i - 1]][j];
        }
    }
    return aut;
}

// compressed version of string ... number of repetitions

/*
// s = a^n  >> ans = n; if s = ababab   ans = 3
        int ans = 1;
        if (n % (n-pi[n-1]) == 0) {
            ans = n/(n-pi[n-1]);
        }
        cout <<ans <<'\n';

*/

\end{lstlisting}
\subsubsection{String matching using bitset}
\begin{lstlisting}[language=C++]
/**
 *  THE ULTIMATE CP TEMPLATE BLUEPRINT: BITSET STRING MATCHING
 *
 * 1.  WHAT PROBLEM DOES THIS SOLVE?
 * - Keywords: "Wildcards in text", "Small alphabet", "Multiple patterns".
 * - Classic Scenarios: Text has wildcards. Find all occurrences of pattern.
 * - The Magic: Represents occurrences of each character as a bitset. By shifting
 *   and ANDing the bitsets corresponding to the pattern's characters, you instantly
 *   identify valid shift offsets in O(N * M / 64) time.
 *
 * 2.  HOW TO USE IT (THE BLACK BOX)
 * - Query:
 *       bitset<MAX_N> matches = bitset_match(text, pattern);
 *       // If matches[i] is 1, the pattern exists starting at text[i].
 *
 * 3.  HOW TO ADAPT IT
 * - MAX_N must be a constant (e.g., 100005) because bitsets require compile-time sizes.
 */

const int MAX_N = 100005;

bitset<MAX_N> bitset_match(const string &text, const string &pattern)
{
    vector<bitset<MAX_N>> mask(26);
    for (int i = 0; i < text.length(); i++)
    {
        if (text[i] != '?')
        {
            mask[text[i] - 'a'].set(i);
        }
        else
        {
            for (int c = 0; c < 26; c++)
                mask[c].set(i);
        }
    }

    bitset<MAX_N> ans;
    ans.set();

    for (int i = 0; i < pattern.length(); i++)
    {
        if (pattern[i] != '?')
        {
            ans &= (mask[pattern[i] - 'a'] >> i);
        }
    }
    return ans;
}
\end{lstlisting}
\subsubsection{Z}
\begin{lstlisting}[language=C++]
struct ZAlgorithm
{
    vector<int> build(const string &s)
    {
        int n = s.size();
        vector<int> z(n, 0);
        int l = 0, r = 0;
        for (int i = 1; i < n; i++)
        {
            if (i < r)
                z[i] = min(r - i, z[i - l]);
            while (i + z[i] < n && s[z[i]] == s[i + z[i]])
                z[i]++;
            if (i + z[i] > r)
            {
                l = i;
                r = i + z[i];
            }
        }
        return z;
    }

    int numberOfMatches(const string &s, const string &p)
    {
        if (p.empty() || s.empty())
            return 0;
        string combined = p + "$" + s;
        vector<int> z = build(combined);
        int ans = 0;
        int p_len = p.size();
        for (int i = p_len + 1; i < z.size(); i++)
        {
            if (z[i] == p_len)
                ans++;
        }
        return ans;
    }

    vector<int> findOccurrences(const string &s, const string &p)
    {
        if (p.empty() || s.empty())
            return {};
        string combined = p + "$" + s;
        vector<int> z = build(combined);
        vector<int> occ;
        int p_len = p.size();
        for (int i = p_len + 1; i < z.size(); i++)
        {
            if (z[i] == p_len)
                occ.push_back(i - p_len - 1);
        }
        return occ;
    }
};

\end{lstlisting}
\subsection{Misc}
\subsubsection{Debruijn sequence}
\begin{lstlisting}[language=C++]
/**
 *  THE ULTIMATE CP TEMPLATE BLUEPRINT: DE BRUIJN SEQUENCE
 *
 * 1.  WHAT PROBLEM DOES THIS SOLVE?
 * - Keywords: "Contains all permutations", "Shortest string with all substrings".
 * - Classic Scenarios: Generate the shortest possible string that contains every
 *   possible sequence of length N using an alphabet of size K.
 * - The Magic: Mathematical guarantees state this string will always have length
 *   K^N + N - 1. This uses a recursive DFS (essentially finding an Eulerian path)
 *   to generate the lexicographically smallest De Bruijn sequence in optimal time.
 *
 * 2.  HOW TO USE IT (THE BLACK BOX)
 * - Query:
 *       string sequence = de_bruijn(K, N);
 * - Example:
 *       de_bruijn(2, 3) generates "00010111".
 *       (Note: To act as a true cyclic sequence, you conceptually wrap the ends.
 *       If you need it strictly linear, append the first N-1 characters to the end).
 *
 * - Complexity:
 *       Time: O(K^N)
 *       Space: O(K * N)
 */

string de_bruijn(int k, int n)
{
    if (k == 1)
        return string(n, '0');
    if (n == 1)
    {
        string res = "";
        for (int i = 0; i < k; i++)
            res += (char)('0' + i);
        return res;
    }

    vector<int> a(k * n, 0);
    string sequence;

    auto db = [&](auto &self, int t, int p) -> void
    {
        if (t > n)
        {
            if (n % p == 0)
            {
                for (int i = 1; i <= p; i++)
                {
                    sequence += (char)('0' + a[i]);
                }
            }
        }
        else
        {
            a[t] = a[t - p];
            self(self, t + 1, p);
            for (int j = a[t - p] + 1; j < k; j++)
            {
                a[t] = j;
                self(self, t + 1, t);
            }
        }
    };

    db(db, 1, 1);
    return sequence;
}
\end{lstlisting}
\subsubsection{Expression parser}
\begin{lstlisting}[language=C++]
/**
 *  THE ULTIMATE CP TEMPLATE BLUEPRINT: EXPRESSION PARSER
 *
 * 1.  WHAT PROBLEM DOES THIS SOLVE?
 * - Keywords: "Evaluate expression", "Calculator", "Infix to Postfix".
 * - Classic Scenarios: You are given a raw string like "3 + 5 * (2 - 8)" and
 *   need the exact mathematical result.
 * - The Magic: Uses Recursive Descent to strictly enforce PEMDAS (Parentheses,
 *   Exponents, Multiplication/Division, Addition/Subtraction) in pure O(N) time.
 *
 * 2.  HOW TO USE IT (THE BLACK BOX)
 * - Query:
 *       string expr = "10 + 2 * (6 - 3)";
 *       int pos = 0;
 *       long long result = parse_expr(expr, pos);
 *
 * 3.  HOW TO ADAPT IT (THE DIALS & SWITCHES)
 * - Division by Zero: Currently, it performs integer division. If a problem
 *   tests division by zero, add a check inside `parse_term`.
 * - Modulo: If the expression must be evaluated modulo M, wrap all additions,
 *   subtractions, and multiplications in the respective functions with `% MOD`.
 */

using ll = long long;

ll parse_expr(const string &s, int &i);
ll parse_term(const string &s, int &i);
ll parse_factor(const string &s, int &i);
ll parse_number(const string &s, int &i);

void skip_spaces(const string &s, int &i)
{
    while (i < s.length() && isspace(s[i]))
        i++;
}

ll parse_expr(const string &s, int &i)
{
    ll res = parse_term(s, i);
    while (i < s.length())
    {
        skip_spaces(s, i);
        if (i < s.length() && (s[i] == '+' || s[i] == '-'))
        {
            char op = s[i++];
            ll next_term = parse_term(s, i);
            if (op == '+')
                res += next_term;
            else
                res -= next_term;
        }
        else
        {
            break;
        }
    }
    return res;
}

ll parse_term(const string &s, int &i)
{
    ll res = parse_factor(s, i);
    while (i < s.length())
    {
        skip_spaces(s, i);
        if (i < s.length() && (s[i] == '*' || s[i] == '/'))
        {
            char op = s[i++];
            ll next_factor = parse_factor(s, i);
            if (op == '*')
                res *= next_factor;
            else
                res /= next_factor;
        }
        else
        {
            break;
        }
    }
    return res;
}

ll parse_factor(const string &s, int &i)
{
    skip_spaces(s, i);
    if (i < s.length() && s[i] == '(')
    {
        i++;
        ll res = parse_expr(s, i);
        skip_spaces(s, i);
        if (i < s.length() && s[i] == ')')
            i++;
        return res;
    }
    return parse_number(s, i);
}

ll parse_number(const string &s, int &i)
{
    skip_spaces(s, i);
    ll res = 0;
    while (i < s.length() && isdigit(s[i]))
    {
        res = res * 10 + (s[i++] - '0');
    }
    return res;
}
\end{lstlisting}
\subsubsection{Min cyclic shift}
\begin{lstlisting}[language=C++]
/**
 *  THE ULTIMATE CP TEMPLATE BLUEPRINT: DUVAL'S & MIN CYCLIC SHIFT
 *
 * 1.  WHAT PROBLEM DOES THIS SOLVE?
 * - Keywords: "Lexicographically smallest rotation", "Lyndon words".
 * - Classic Scenarios: You can rotate a string S any number of times. Find the
 *   rotation that comes first in the dictionary.
 * - The Magic: Duval's algorithm processes the string in strict O(N) time and
 *   O(1) auxiliary space, bypassing the need to build a heavy Suffix Automaton
 *   on the concatenated string `S + S`.
 *
 * 2.  HOW TO USE IT (THE BLACK BOX)
 * - Query:
 *       int best_index = min_cyclic_shift(S);
 *       string smallest_rotation = S.substr(best_index) + S.substr(0, best_index);
 *       vector<string> factorization = duval_factorization(S);
 *
 * - Complexity:
 *       Time: O(N)
 *       Space: O(1) for shift, O(N) for factorization return.
 */

int min_cyclic_shift(string s)
{
    s += s;
    int n = s.length();
    int i = 0, ans = 0;
    while (i < n / 2)
    {
        ans = i;
        int j = i + 1, k = i;
        while (j < n && s[k] <= s[j])
        {
            if (s[k] < s[j])
                k = i;
            else
                k++;
            j++;
        }
        while (i <= k)
            i += j - k;
    }
    return ans;
}

vector<string> duval_factorization(const string &s)
{
    int n = s.length();
    int i = 0;
    vector<string> factorization;
    while (i < n)
    {
        int j = i + 1, k = i;
        while (j < n && s[k] <= s[j])
        {
            if (s[k] < s[j])
                k = i;
            else
                k++;
            j++;
        }
        while (i <= k)
        {
            factorization.push_back(s.substr(i, j - k));
            i += j - k;
        }
    }
    return factorization;
}
\end{lstlisting}
\subsection{Palindromes}
\subsubsection{Manacher}
\begin{lstlisting}[language=C++]
struct Manacher
{
    vector<int> p;
    string s;

    Manacher(const string &str)
    {
        s = str;
        string t = "";
        for (char c : s)
        {
            t += "#";
            t += c;
        }
        t += "#";

        int n = t.size();
        string temp = "$" + t + "^";
        vector<int> p_temp(n + 2, 0);
        int l = 0, r = 1;

        for (int i = 1; i <= n; i++)
        {
            if (i <= r)
                p_temp[i] = min(r - i, p_temp[l + (r - i)]);
            while (temp[i - p_temp[i]] == temp[i + p_temp[i]])
                p_temp[i]++;
            if (i + p_temp[i] > r)
            {
                l = i - p_temp[i];
                r = i + p_temp[i];
            }
        }

        p.assign(n, 0);
        for (int i = 1; i <= n; i++)
        {
            p[i - 1] = p_temp[i] - 1;
        }
    }

    bool is_palindrome(int l, int r)
    {
        if (l > r)
            swap(l, r);
        int len = r - l + 1;
        int idx = l + r + 1;
        return p[idx] >= len;
    }

    string longest_palindrome()
    {
        int max_len = 0, center_idx = 0;
        for (int i = 0; i < p.size(); i++)
        {
            if (p[i] > max_len)
            {
                max_len = p[i];
                center_idx = i;
            }
        }
        int start = (center_idx - max_len) / 2;
        return s.substr(start, max_len);
    }
};

\end{lstlisting}
\subsubsection{Palindromic tree}
\begin{lstlisting}[language=C++]
/**
 *  THE ULTIMATE CP TEMPLATE BLUEPRINT: PALINDROMIC TREE (EERTREE)
 *
 * 1.  WHAT PROBLEM DOES THIS SOLVE?
 * - Keywords: "Distinct palindromes", "Occurrences of palindromes".
 * - Classic Scenarios: Find the number of distinct palindromic substrings, or
 *   find the palindrome P that maximizes length(P) * occurrences(P).
 * - The Magic: It builds a graph containing exactly two roots: one for even-length
 *   palindromes and one for odd-length. Every node represents a distinct palindrome,
 *   and fail links point to the longest proper palindromic suffix.
 *
 * 2.  HOW TO USE IT (THE BLACK BOX)
 * - Initialization:
 *       Eertree pt;
 *       for (char c : S) pt.add_char(c);
 *       pt.count_occurrences(); // Must call this after building!
 *
 * - Complexity:
 *       Time: Amortized O(N)
 *       Space: O(N * ALPHABET_SIZE)
 */

struct Eertree
{
    static const int K = 26;
    struct Node
    {
        int next[K];
        int len;
        int link;
        long long cnt;

        Node(int len = 0, int link = 0) : len(len), link(link), cnt(0)
        {
            fill(begin(next), end(next), 0);
        }
    };

    vector<Node> tree;
    string s;
    int last;

    Eertree()
    {
        tree.emplace_back(0, 1);  // Node 0: Even length root
        tree.emplace_back(-1, 1); // Node 1: Odd length root
        last = 0;
    }

    void add_char(char ch)
    {
        s += ch;
        int c = ch - 'a';
        int curr = last;
        int pos = s.length() - 1;

        while (pos - 1 - tree[curr].len < 0 || s[pos - 1 - tree[curr].len] != ch)
        {
            curr = tree[curr].link;
        }

        if (tree[curr].next[c])
        {
            last = tree[curr].next[c];
            tree[last].cnt++;
            return;
        }

        int now = tree.size();
        tree.emplace_back(tree[curr].len + 2);
        tree[curr].next[c] = now;

        if (tree[now].len == 1)
        {
            tree[now].link = 0;
        }
        else
        {
            int fail = tree[curr].link;
            while (pos - 1 - tree[fail].len < 0 || s[pos - 1 - tree[fail].len] != ch)
            {
                fail = tree[fail].link;
            }
            tree[now].link = tree[fail].next[c];
        }

        last = now;
        tree[last].cnt++;
    }

    void count_occurrences()
    {
        for (int i = tree.size() - 1; i > 1; i--)
        {
            tree[tree[i].link].cnt += tree[i].cnt;
        }
    }
};
\end{lstlisting}
\subsection{Suffix Structures}
\subsubsection{Suffix array}
\begin{lstlisting}[language=C++]
struct SuffixArray
{
    string s;
    int n;
    vector<int> p;
    vector<vector<int>> C;
    vector<int> lcp;
    bool is_cyclic;

    SuffixArray(string str, bool cyclic = false)
    {
        is_cyclic = cyclic;
        s = str;
        if (!is_cyclic)
            s += '$';
        n = s.size();
        const int alphapet = 256;
        p.assign(n, 0);
        vector<int> c(n), cnt(max(n, alphapet), 0);

        for (int i = 0; i < n; i++)
            cnt[s[i]]++;
        for (int i = 1; i < alphapet; i++)
            cnt[i] += cnt[i - 1];
        for (int i = 0; i < n; i++)
            p[--cnt[s[i]]] = i;

        c[p[0]] = 0;
        int classes = 1;
        for (int i = 1; i < n; i++)
        {
            if (s[p[i]] != s[p[i - 1]])
                classes++;
            c[p[i]] = classes - 1;
        }
        C.push_back(c);

        vector<int> pn(n), cn(n);
        for (int h = 0; (1 << h) < n; h++)
        {
            for (int i = 0; i < n; i++)
            {
                pn[i] = p[i] - (1 << h);
                if (pn[i] < 0)
                    pn[i] += n;
            }
            fill(cnt.begin(), cnt.begin() + classes, 0ll);
            for (int i = 0; i < n; i++)
                cnt[C.back()[pn[i]]]++;
            for (int i = 1; i < classes; i++)
                cnt[i] += cnt[i - 1];
            for (int i = n - 1; i >= 0; i--)
                p[--cnt[C.back()[pn[i]]]] = pn[i];

            cn[p[0]] = 0;
            classes = 1;
            for (int i = 1; i < n; i++)
            {
                pair<int, int> curr = {C.back()[p[i]], C.back()[(p[i] + (1 << h)) % n]};
                pair<int, int> prev = {C.back()[p[i - 1]], C.back()[(p[i - 1] + (1 << h)) % n]};
                if (curr != prev)
                    classes++;
                cn[p[i]] = classes - 1;
            }
            C.push_back(cn);
        }

        if (!is_cyclic)
        {
            p.erase(p.begin());
            int orig_n = n - 1;
            vector<int> rank(orig_n, 0);
            for (int i = 0; i < orig_n; i++)
                rank[p[i]] = i;
            int k = 0;
            lcp.assign(orig_n - 1, 0);
            for (int i = 0; i < orig_n; i++)
            {
                if (rank[i] == orig_n - 1)
                {
                    k = 0;
                    continue;
                }
                int j = p[rank[i] + 1];
                while (i + k < orig_n && j + k < orig_n && s[i + k] == s[j + k])
                    k++;
                lcp[rank[i]] = k;
                if (k)
                    k--;
            }
        }
    }

    string smallestCyclicShift()
    {
        return s.substr(p[0]) + s.substr(0, p[0]);
    }

    vector<int> subStringInString(string t)
    {
        int m = t.size();
        int orig_n = n - 1;
        int l = 0, r = orig_n - 1;
        int first = -1;
        while (l <= r)
        {
            int mid = l + (r - l) / 2;
            auto can = [&]()
            {
                for (int i = 0; i < m; i++)
                {
                    if (p[mid] + i >= orig_n)
                        return 1LL;
                    if (s[p[mid] + i] != t[i])
                    {
                        if (s[p[mid] + i] < t[i])
                            return 1LL;
                        if (s[p[mid] + i] > t[i])
                            return 2LL;
                    }
                }
                return 0LL;
            };
            int res = can();
            if (res == 0)
            {
                first = mid;
                r = mid - 1;
            }
            else if (res == 1)
            {
                l = mid + 1;
            }
            else
            {
                r = mid - 1;
            }
        }
        if (first == -1)
            return {};
        l = 0, r = orig_n - 1;
        int last = -1;
        while (l <= r)
        {
            int mid = l + (r - l) / 2;
            auto can = [&]()
            {
                for (int i = 0; i < m; i++)
                {
                    if (p[mid] + i >= orig_n)
                        return 1LL;
                    if (s[p[mid] + i] != t[i])
                    {
                        if (s[p[mid] + i] < t[i])
                            return 1LL;
                        if (s[p[mid] + i] > t[i])
                            return 2LL;
                    }
                }
                return 0LL;
            };
            int res = can();
            if (res == 0)
            {
                last = mid;
                l = mid + 1;
            }
            else if (res == 1)
            {
                l = mid + 1;
            }
            else
            {
                r = mid - 1;
            }
        }
        vector<int> ans;
        for (int i = first; i <= last; i++)
        {
            ans.push_back(p[i]);
        }
        sort(ans.begin(), ans.end());
        return ans;
    }

    int compareSubstrings(int i, int len1, int j, int len2)
    {
        int l = min(len1, len2);
        if (l > 0)
        {
            int k = __lg(l);
            pair<int, int> a = {C[k][i], C[k][i + l - (1 << k)]};
            pair<int, int> b = {C[k][j], C[k][j + l - (1 << k)]};
            if (a != b)
                return a < b ? -1 : 1;
        }
        if (len1 == len2)
            return 0;
        return len1 < len2 ? -1 : 1;
    }

    int numberOfDifferentSubStrings()
    {
        int orig_n = n - 1;
        return (orig_n * orig_n + orig_n) / 2 - accumulate(lcp.begin(), lcp.end(), 0ll);
    }

    string longestCommonSubstring(string t)
    {
        string orig_s = s;
        if (!is_cyclic)
            orig_s.pop_back();
        int len1 = orig_s.size();
        SuffixArray sa(orig_s + char(1) + t);
        int max_len = 0;
        int start_idx = -1;
        for (int i = 0; i < (int)sa.lcp.size(); i++)
        {
            if ((sa.p[i] < len1 && sa.p[i + 1] > len1) || (sa.p[i] > len1 && sa.p[i + 1] < len1))
            {
                if (sa.lcp[i] > max_len)
                {
                    max_len = sa.lcp[i];
                    start_idx = sa.p[i];
                }
            }
        }
        if (max_len == 0)
            return "";
        return sa.s.substr(start_idx, max_len);
    }

    string longestRepeatedSubstring()
    {
        int max_len = 0;
        int start_idx = -1;
        for (int i = 0; i < (int)lcp.size(); i++)
        {
            if (lcp[i] > max_len)
            {
                max_len = lcp[i];
                start_idx = p[i];
            }
        }
        if (max_len == 0)
            return "";
        return s.substr(start_idx, max_len);
    }
};

\end{lstlisting}
\subsubsection{Suffix automaton}
\begin{lstlisting}[language=C++]
struct SuffixAutomaton
{
    struct State
    {
        int len, link;
        int first_pos;
        bool is_clone;
        int cnt;
        long long paths;
        map<char, int> next;
        vector<int> inv_link;
    };

    vector<State> st;
    int sz, last;
    string s;

    SuffixAutomaton()
    {
        st.assign(2, State());
        st[0].len = 0;
        st[0].link = -1;
        st[0].cnt = 0;
        sz = 1;
        last = 0;
    }

    void extend(char c)
    {
        int cur = sz++;
        if (cur >= st.size())
            st.resize(st.size() * 2);
        st[cur].len = st[last].len + 1;
        st[cur].first_pos = st[cur].len - 1;
        st[cur].is_clone = false;
        st[cur].cnt = 1;

        s += c;

        int p = last;
        while (p != -1 && !st[p].next.count(c))
        {
            st[p].next[c] = cur;
            p = st[p].link;
        }
        if (p == -1)
        {
            st[cur].link = 0;
        }
        else
        {
            int q = st[p].next[c];
            if (st[p].len + 1 == st[q].len)
            {
                st[cur].link = q;
            }
            else
            {
                int clone = sz++;
                if (clone >= st.size())
                    st.resize(st.size() * 2);
                st[clone].len = st[p].len + 1;
                st[clone].next = st[q].next;
                st[clone].link = st[q].link;
                st[clone].first_pos = st[q].first_pos;
                st[clone].is_clone = true;
                st[clone].cnt = 0;

                while (p != -1 && st[p].next[c] == q)
                {
                    st[p].next[c] = clone;
                    p = st[p].link;
                }
                st[q].link = st[cur].link = clone;
            }
        }
        last = cur;
    }

    void build(const string &text)
    {
        for (char c : text)
            extend(c);
        precompute();
    }

    void precompute()
    {
        for (int v = 1; v < sz; v++)
        {
            st[st[v].link].inv_link.push_back(v);
        }
        vector<pair<int, int>> order(sz);
        for (int i = 0; i < sz; i++)
            order[i] = {st[i].len, i};
        sort(order.rbegin(), order.rend());
        for (int i = 0; i < sz - 1; i++)
        {
            int v = order[i].second;
            if (st[v].link != -1)
            {
                st[st[v].link].cnt += st[v].cnt;
            }
        }
    }

    bool checkOccurrence(string P)
    {
        int v = 0;
        for (char c : P)
        {
            if (!st[v].next.count(c))
                return false;
            v = st[v].next[c];
        }
        return true;
    }

    long long numberOfDifferentSubstrings()
    {
        long long tot = 0;
        for (int i = 1; i < sz; i++)
        {
            tot += st[i].len - st[st[i].link].len;
        }
        return tot;
    }

    long long totalLengthOfDifferentSubstrings()
    {
        long long tot = 0;
        for (int i = 1; i < sz; i++)
        {
            long long shortest = st[st[i].link].len + 1;
            long long longest = st[i].len;
            long long num_strings = longest - shortest + 1;
            long long cur = num_strings * (longest + shortest) / 2;
            tot += cur;
        }
        return tot;
    }

    void precomputePaths()
    {
        vector<pair<int, int>> order(sz);
        for (int i = 0; i < sz; i++)
            order[i] = {st[i].len, i};
        sort(order.rbegin(), order.rend());
        for (int i = 0; i < sz; i++)
        {
            int v = order[i].second;
            st[v].paths = st[v].cnt;
            for (auto &edge : st[v].next)
            {
                st[v].paths += st[edge.second].paths;
            }
        }
    }

    string kthSubstring(long long k)
    {
        precomputePaths();
        int v = 0;
        string ans = "";
        while (k > 0)
        {
            for (auto &edge : st[v].next)
            {
                int w = edge.second;
                if (st[w].paths >= k)
                {
                    ans += edge.first;
                    k -= st[w].cnt;
                    v = w;
                    break;
                }
                else
                {
                    k -= st[w].paths;
                }
            }
        }
        return ans;
    }

    static string smallestCyclicShift(string S)
    {
        SuffixAutomaton sa;
        sa.build(S + S);
        int v = 0;
        string ans = "";
        for (int i = 0; i < S.size(); i++)
        {
            auto edge = *sa.st[v].next.begin();
            ans += edge.first;
            v = edge.second;
        }
        return ans;
    }

    int numberOfOccurrences(string P)
    {
        int v = 0;
        for (char c : P)
        {
            if (!st[v].next.count(c))
                return 0;
            v = st[v].next[c];
        }
        return st[v].cnt;
    }

    int firstOccurrencePosition(string P)
    {
        int v = 0;
        for (char c : P)
        {
            if (!st[v].next.count(c))
                return -1;
            v = st[v].next[c];
        }
        return st[v].first_pos - P.length() + 1;
    }

    void dfsOccurrences(int v, int P_length, vector<int> &res)
    {
        if (!st[v].is_clone)
        {
            res.push_back(st[v].first_pos - P_length + 1);
        }
        for (int u : st[v].inv_link)
        {
            dfsOccurrences(u, P_length, res);
        }
    }

    vector<int> allOccurrencePositions(string P)
    {
        int v = 0;
        for (char c : P)
        {
            if (!st[v].next.count(c))
                return {};
            v = st[v].next[c];
        }
        vector<int> res;
        dfsOccurrences(v, P.length(), res);
        sort(res.begin(), res.end());
        return res;
    }

    string shortestNonAppearingString(char min_char = 'a', char max_char = 'z')
    {
        vector<int> d(sz, -1);
        auto dfs = [&](auto &self, int v) -> int
        {
            if (d[v] != -1)
                return d[v];
            d[v] = 1e9;
            for (char c = min_char; c <= max_char; c++)
            {
                if (!st[v].next.count(c))
                {
                    d[v] = 1;
                    break;
                }
                else
                {
                    d[v] = min(d[v], 1 + self(self, st[v].next[c]));
                }
            }
            return d[v];
        };
        dfs(dfs, 0);
        string ans = "";
        int v = 0;
        while (true)
        {
            for (char c = min_char; c <= max_char; c++)
            {
                if (!st[v].next.count(c))
                {
                    ans += c;
                    return ans;
                }
            }
            for (char c = min_char; c <= max_char; c++)
            {
                if (d[v] == 1 + d[st[v].next[c]])
                {
                    ans += c;
                    v = st[v].next[c];
                    break;
                }
            }
        }
    }

    string longestCommonSubstringTwo(string T)
    {
        int v = 0, l = 0, best = 0, bestpos = 0;
        for (int i = 0; i < T.size(); i++)
        {
            while (v && !st[v].next.count(T[i]))
            {
                v = st[v].link;
                l = st[v].len;
            }
            if (st[v].next.count(T[i]))
            {
                v = st[v].next[T[i]];
                l++;
            }
            if (l > best)
            {
                best = l;
                bestpos = i;
            }
        }
        return T.substr(bestpos - best + 1, best);
    }

    static string longestCommonSubstringMultiple(vector<string> strings)
    {
        if (strings.empty())
            return "";
        if (strings.size() == 1)
            return strings[0];

        int min_idx = 0;
        for (int i = 1; i < strings.size(); i++)
        {
            if (strings[i].size() < strings[min_idx].size())
            {
                min_idx = i;
            }
        }
        swap(strings[0], strings[min_idx]);

        SuffixAutomaton sa;
        sa.build(strings[0]);

        vector<int> order(sa.sz);
        iota(order.begin(), order.end(), 0);
        sort(order.begin(), order.end(), [&](int a, int b)
             { return sa.st[a].len > sa.st[b].len; });

        vector<int> min_max_l(sa.sz);
        for (int i = 0; i < sa.sz; i++)
        {
            min_max_l[i] = sa.st[i].len;
        }

        for (int i = 1; i < strings.size(); i++)
        {
            vector<int> max_l(sa.sz, 0);
            int v = 0, l = 0;
            for (char c : strings[i])
            {
                while (v > 0 && !sa.st[v].next.count(c))
                {
                    v = sa.st[v].link;
                    l = sa.st[v].len;
                }
                if (sa.st[v].next.count(c))
                {
                    v = sa.st[v].next[c];
                    l++;
                }
                max_l[v] = max(max_l[v], l);
            }

            for (int u : order)
            {
                min_max_l[u] = min(min_max_l[u], max_l[u]);
                int lnk = sa.st[u].link;
                if (lnk != -1)
                {
                    max_l[lnk] = max(max_l[lnk], min(max_l[u], sa.st[lnk].len));
                }
            }
        }

        int best_len = 0;
        int best_v = 0;
        for (int i = 1; i < sa.sz; i++)
        {
            if (min_max_l[i] > best_len)
            {
                best_len = min_max_l[i];
                best_v = i;
            }
        }
        if (best_len == 0)
            return "";
        return strings[0].substr(sa.st[best_v].first_pos - best_len + 1, best_len);
    }
};

\end{lstlisting}
\section{DP}
\subsection{DP Bitmask}
\subsubsection{Bitmask}
\begin{lstlisting}[language=C++]
int n;
vector<vector<int>> v;
int dp[22][1 << 22];

int sol(int man, int mask)
{
    if (man >= n)
        return 1;
    auto &ret = dp[man][mask];
    if (~ret)
        return ret;
    ret = 0;
    for (int i = 0; i < n; i++)
    {
        if (mask >> i & 1)
            continue;
        if (v[man][i])
        {
            ret += sol(man + 1, mask | 1 << i);
            ret %= MOD;
        }
    }
    return ret;
}

\end{lstlisting}
\subsubsection{SOS}
\begin{lstlisting}[language=C++]
vector<int> SOS(int n, vector<int> &a)
{
    vector<int> p = a;
    for (int dim = 0; dim < n; dim++)
    {
        for (int i = 0; i < 1 << n; i++)
        {
            if (i >> dim & 1)
                p[i] += p[i - (1 << dim)];
        }
    }
    return p;
}

\end{lstlisting}
\subsubsection{Submask}
\begin{lstlisting}[language=C++]
int dp[1<<16];
int cost[1<<16];
int memo (int mask) {
    if (mask == 0) return 0;
    auto &ret = dp[mask];
    if (~ret) return ret;
    for (int i = mask;i > 0;i = (i-1) & mask) {
        ret = max(ret,cost[i] + memo(mask ^ i));
    }
    return ret;
}
void hassan() {
    int n;cin >>n;
    int a[n][n];
    for (int i = 0; i < n; i++) {
        for (int j = 0; j < n; j++) {
            cin >> a[i][j];
        }
    }
    for (int mask = 0; mask < (1<<n); mask++) {
        for (int i = 0; i < n; i++) {
            if (mask & (1<<i)) {
                for (int j = 0; j < i; j++) {
                    if (mask & (1<<j)) {
                        cost[mask]+=a[i][j];
                    }
                }
            }
        }
      //  cout << mask << " >> " << cost[mask] << endl;
    }
    memset(dp, -1, sizeof(dp));
    cout << memo((1<<n)-1);

}

\end{lstlisting}
\subsubsection{subset convolution}
\begin{lstlisting}[language=C++]
/**
 *  THE ULTIMATE CP TEMPLATE BLUEPRINT: SUBSET CONVOLUTION
 *
 * 1.  WHAT PROBLEM DOES THIS SOLVE?
 * - Keywords: "Disjoint subset sum", "i OR j = k AND i AND j = 0".
 * - Classic Scenarios: Combining two DP states where the masks must perfectly
 *   interlock (share no common set bits) to form a target mask.
 * - The Magic: By adding an extra dimension for the `__builtin_popcount`, we
 *   guarantee that if two masks combine to form a mask of size K, their individual
 *   set bits must sum to exactly K, implicitly enforcing that they are disjoint.
 *   It runs SOS DP on each popcount layer.
 *
 * 2.  HOW TO USE IT (THE BLACK BOX)
 * - Query:
 *       vector<long long> C = subset_convolution(A, B, N, MOD);
 *
 * - Complexity:
 *       Time: O(N^2 * 2^N)
 *       Space: O(N * 2^N)
 */

using ll = long long;

vector<ll> subset_convolution(const vector<ll> &a, const vector<ll> &b, int n, ll mod)
{
    vector<vector<ll>> f_hat(n + 1, vector<ll>(1 << n, 0));
    vector<vector<ll>> g_hat(n + 1, vector<ll>(1 << n, 0));
    vector<vector<ll>> h_hat(n + 1, vector<ll>(1 << n, 0));

    for (int mask = 0; mask < (1 << n); mask++)
    {
        int pc = __builtin_popcount(mask);
        f_hat[pc][mask] = a[mask] % mod;
        g_hat[pc][mask] = b[mask] % mod;
    }

    for (int i = 0; i <= n; i++)
    {
        for (int j = 0; j < n; j++)
        {
            for (int mask = 0; mask < (1 << n); mask++)
            {
                if (mask & (1 << j))
                {
                    f_hat[i][mask] = (f_hat[i][mask] + f_hat[i][mask ^ (1 << j)]) % mod;
                    g_hat[i][mask] = (g_hat[i][mask] + g_hat[i][mask ^ (1 << j)]) % mod;
                }
            }
        }
    }

    for (int mask = 0; mask < (1 << n); mask++)
    {
        for (int i = 0; i <= n; i++)
        {
            ll sum = 0;
            for (int j = 0; j <= i; j++)
            {
                sum = (sum + f_hat[j][mask] * g_hat[i - j][mask]) % mod;
            }
            h_hat[i][mask] = sum;
        }
    }

    for (int i = 0; i <= n; i++)
    {
        for (int j = 0; j < n; j++)
        {
            for (int mask = 0; mask < (1 << n); mask++)
            {
                if (mask & (1 << j))
                {
                    h_hat[i][mask] = (h_hat[i][mask] - h_hat[i][mask ^ (1 << j)] + mod) % mod;
                }
            }
        }
    }

    vector<ll> c(1 << n);
    for (int mask = 0; mask < (1 << n); mask++)
    {
        c[mask] = h_hat[__builtin_popcount(mask)][mask];
    }
    return c;
}
\end{lstlisting}
\subsection{Digit Dp}
\subsubsection{main}
\begin{lstlisting}[language=C++]
string s;
int d, k;
int dp[20][20][2][2];

int sol(int idx, int cnt, int tight, int num)
{
    if (idx >= int(s.size()))
    {
        if (cnt == k)
            return 1;
        return 0;
    }
    auto &ret = dp[idx][cnt][tight][num];
    if (~ret)
        return ret;
    ret = 0;
    int limit = tight ? s[idx] - '0' : 9;
    for (int i = 0; i <= limit; i++)
    {
        int nn = num || i;
        int nt = tight && i == limit;
        int nc = cnt + (d == i && nn);
        ret += sol(idx + 1, nc, nt, nn);
    }
    return ret;
}

void solve()
{
    cin >> s >> d >> k;
    memset(dp, -1, sizeof dp);
    cout << sol(0, 0, 1, 0);
}

\end{lstlisting}
\subsection{Iterative}
\subsubsection{main}
\begin{lstlisting}[language=C++]
int answer(int n)
{
    if (n == 0 || n == 1)
        return n;
    return dp[n];
}

void solve()
{
    int n;
    cin >> n;
    for (int i = 2; i <= n; i++)
        dp[i] = answer(i - 1) + answer(i - 2);
    cout << answer(n);
}

\end{lstlisting}
\subsection{Knapsack}
\subsubsection{main}
\begin{lstlisting}[language=C++]
int n, W;
int dp[1001][31];
vector<pair<int, int>> a;

int sol(int i, int w)
{
    if (i >= n)
        return 0;
    auto &ret = dp[i][w];
    if (~ret)
        return ret;
    ret = sol(i + 1, w);
    if (w + a[i].second <= W)
        ret = max(ret, a[i].first + sol(i + 1, w + a[i].second));
    return ret;
}

\end{lstlisting}
\subsection{LCS}
\subsubsection{LCS}
\begin{lstlisting}[language=C++]
vector<vector<int>> dp;
int n1, n2;
string s1, s2;

int LCS(int i1, int i2)
{
    if (i1 >= n1 || i2 >= n2)
        return 0;
    auto &ret = dp[i1][i2];
    if (ret != -1)
        return ret;
    int choice1 = LCS(i1 + 1, i2);
    int choice2 = LCS(i1, i2 + 1);
    int choice3{};
    if (s1[i1] == s2[i2])
        choice3 = 1 + LCS(i1 + 1, i2 + 1);
    return ret = max({choice1, choice2, choice3});
}

\end{lstlisting}
\subsection{LIS}
\subsubsection{Basic 2D}
\begin{lstlisting}[language=C++]
int dp1[2010][2010];
int dp2[2010][2010];

int LDS(int i, int prev)
{
    if (i > n)
        return 0;
    auto &ret = dp1[i][prev];
    if (~ret)
        return ret;
    ret = LDS(i + 1, prev);
    if (prev == 0 || a[i] < a[prev])
        ret = max(ret, 1 + LDS(i + 1, i));
    return ret;
}

int LIS(int i, int prev)
{
    if (i > n)
        return 0;
    auto &ret = dp2[i][prev];
    if (~ret)
        return ret;
    ret = LIS(i + 1, prev);
    if (prev == 0 || a[i] > a[prev])
        ret = max(ret, 1 + LIS(i + 1, i));
    return ret;
}

\end{lstlisting}
\subsubsection{LIS1D}
\begin{lstlisting}[language=C++]
int LIS(int i)
{
    if (i > n)
        return 0;
    auto &ret = dp[i];
    if (dp[i] == -1)
        return ret;
    ret = 0;
    for (int j = i + 1; j <= n; j++)
    {
        if (v[j] > v[i])
            ret = max(ret, LIS(j));
    }
    ret++;
    return ret;
}

// int res{};
// for (int i = 1; i <= n; i++)
// {
//     int num = LIS(i);
//     res = max(num, res);
// }

\end{lstlisting}
\subsubsection{LIS BS}
\begin{lstlisting}[language=C++]
vector<int> LIS;
for (int i = 0; i < n; i++)
{
    auto it = lower_bound(LIS.begin(), LIS.end(), ans[i]);
    if (it == LIS.end())
        LIS.push_back(ans[i]);
    else
        *it = ans[i];
}

\end{lstlisting}
\subsubsection{LIS seg}
\begin{lstlisting}[language=C++]
struct Node
{
    int val;

    Node()
    {
        val = 0;
    }

    Node(int val)
        : val(val)
    {
    }

    void change(int x)
    {
        val = max(x, val);
    }
};

struct segmentTree
{
    int sz;
    vector<Node> seg;

    segmentTree(int n)
    {
        sz = 1;
        while (sz < n)
            sz *= 2;
        seg.assign(2 * sz, Node());
    }

    void init(vector<int> &arr, int node, int lx, int rx)
    {
        if (rx - lx == 1)
        {
            if (lx < int(arr.size()))
                seg[node] = Node(arr[lx]);
            return;
        }
        int mid = (lx + rx) / 2;
        init(arr, 2 * node + 1, lx, mid);
        init(arr, 2 * node + 2, mid, rx);
        seg[node] = merge(seg[2 * node + 1], seg[2 * node + 2]);
    }

    void init(vector<int> &arr)
    {
        init(arr, 0, 0, sz);
    }

    Node merge(Node &lf, Node &rt)
    {
        Node ans = Node();
        ans.val = max(lf.val, rt.val);
        return ans;
    }

    void set(int idx, int val, int node, int lx, int rx)
    {
        if (rx - lx == 1)
        {
            seg[node].change(val);
            return;
        }
        int mid = (lx + rx) / 2;
        if (idx < mid)
            set(idx, val, 2 * node + 1, lx, mid);
        else
            set(idx, val, 2 * node + 2, mid, rx);
        seg[node] = merge(seg[2 * node + 1], seg[2 * node + 2]);
    }

    void set(int idx, int val)
    {
        set(idx, val, 0, 0, sz);
    }

    Node get(int l, int r, int node, int lx, int rx)
    {
        if (lx >= l && rx <= r)
            return seg[node];
        if (rx <= l || lx >= r)
            return Node();
        int mid = (lx + rx) / 2;
        Node lf = get(l, r, 2 * node + 1, lx, mid);
        Node ri = get(l, r, 2 * node + 2, mid, rx);
        return merge(lf, ri);
    }

    int get(int l, int r)
    {
        return get(l, r, 0, 0, sz).val;
    }
};

int LIS(int n, vector<int> &p)
{
    segmentTree seg(n + 5);
    vector<int> dp(n + 1);
    int ans{};
    for (int i = 0; i < n; i++)
    {
        dp[i] = seg.get(0, p[i]) + 1;
        seg.set(p[i], dp[i]);
        ans = max(ans, dp[i]);
    }
    return ans;
}

\end{lstlisting}
\subsection{Lichao tree}
\subsubsection{lichao tree}
\begin{lstlisting}[language=C++]
/**
 *  THE ULTIMATE CP TEMPLATE BLUEPRINT: LI CHAO TREE (EXTENDED)
 *
 * 1.  WHAT PROBLEM DOES THIS SOLVE?
 * - Keywords: "Minimum of lines", "Linear cost over a range", "y = mx + c".
 * - Classic Scenarios: You need to minimize a linear DP transition cost, but
 *   certain transitions/lines are only valid within a specific interval [L, R].
 * - The Magic: A sparse segment tree where each node holds the line that
 *   provides the optimal value at the midpoint of its interval. It dynamically
 *   allocates memory (sparse) so it can comfortably handle x-coordinates up to 10^18.
 *
 * 2.  HOW TO USE IT (THE BLACK BOX)
 * - Initialization:
 *       LiChaoTree lct(-1e18, 1e18); // Pass the minimum and maximum possible X queries
 * - Operations:
 *       lct.add_line(M, C);                  // Adds an infinite line y = Mx + C
 *       lct.add_segment(M, C, L, R);         // Adds line y = Mx + C strictly for x in [L, R]
 *       long long min_y = lct.query(X);      // Gets the minimum y at point X
 *
 * 3.  HOW TO ADAPT IT (THE DIALS & SWITCHES)
 * - Maximization: This template is strictly built to find the MINIMUM value.
 *   If you need the MAXIMUM, insert your lines with inverted signs:
 *   `lct.add_line(-M, -C)` and then negate the result: `-lct.query(X)`.
 */

using ll = long long;
const ll INF = 4e18; // Safe upper bound to prevent overflow during additions

struct Line
{
    ll m, c;
    ll eval(ll x) const { return m * x + c; }
};

struct Node
{
    Line line;
    int lc = -1, rc = -1;
    Node() : line({0, INF}) {}
    Node(Line l) : line(l) {}
};

struct LiChaoTree
{
    vector<Node> tree;
    ll MIN_X, MAX_X;

    LiChaoTree(ll min_x, ll max_x) : MIN_X(min_x), MAX_X(max_x)
    {
        tree.emplace_back();
    }

    void add_line(Line nw, int u, ll l, ll r)
    {
        ll mid = l + (r - l) / 2;
        bool left = nw.eval(l) < tree[u].line.eval(l);
        bool middle = nw.eval(mid) < tree[u].line.eval(mid);

        if (middle)
            swap(tree[u].line, nw);
        if (l == r)
            return;

        if (left != middle)
        {
            if (tree[u].lc == -1)
            {
                tree[u].lc = tree.size();
                tree.emplace_back();
            }
            add_line(nw, tree[u].lc, l, mid);
        }
        else
        {
            if (tree[u].rc == -1)
            {
                tree[u].rc = tree.size();
                tree.emplace_back();
            }
            add_line(nw, tree[u].rc, mid + 1, r);
        }
    }

    void add_line(ll m, ll c)
    {
        add_line({m, c}, 0, MIN_X, MAX_X);
    }

    void add_segment(Line nw, ll ql, ll qr, int u, ll l, ll r)
    {
        if (ql <= l && r <= qr)
        {
            add_line(nw, u, l, r);
            return;
        }
        ll mid = l + (r - l) / 2;
        if (ql <= mid)
        {
            if (tree[u].lc == -1)
            {
                tree[u].lc = tree.size();
                tree.emplace_back();
            }
            add_segment(nw, ql, qr, tree[u].lc, l, mid);
        }
        if (qr > mid)
        {
            if (tree[u].rc == -1)
            {
                tree[u].rc = tree.size();
                tree.emplace_back();
            }
            add_segment(nw, ql, qr, tree[u].rc, mid + 1, r);
        }
    }

    void add_segment(ll m, ll c, ll ql, ll qr)
    {
        add_segment({m, c}, ql, qr, 0, MIN_X, MAX_X);
    }

    ll query(ll x, int u, ll l, ll r)
    {
        if (u == -1)
            return INF;
        ll res = tree[u].line.eval(x);
        if (l == r)
            return res;

        ll mid = l + (r - l) / 2;
        if (x <= mid)
            res = min(res, query(x, tree[u].lc, l, mid));
        else
            res = min(res, query(x, tree[u].rc, mid + 1, r));

        return res;
    }

    ll query(ll x)
    {
        return query(x, 0, MIN_X, MAX_X);
    }
};
\end{lstlisting}
\subsection{Misc}
\subsubsection{Broken profile}
\begin{lstlisting}[language=C++]
/**
 *  THE ULTIMATE CP TEMPLATE BLUEPRINT: BROKEN PROFILE DP (PLUG DP)
 *
 * 1.  WHAT PROBLEM DOES THIS SOLVE?
 * - Keywords: "Tile the grid", "Dominoes", "Hamiltonian Path on Grid".
 * - Classic Scenarios: Find the number of ways to completely cover an N x M grid
 *   using 1x2 dominoes where N is huge (10^5) and M is very small (<= 15).
 * - The Magic: It sweeps through the grid cell by cell (not row by row).
 *   It maintains a bitmask of length M representing the "profile" (the boundary
 *   of cells protruding into the current row). As we move to the next cell, we
 *   shift the mask and try placing dominoes horizontally or vertically.
 *
 * 2.  HOW TO USE IT (THE BLACK BOX)
 * - Query:
 *       long long ways = plug_dp_dominoes(N, M, MOD);
 *
 * - Complexity:
 *       Time: O(N * M * 2^M)
 *       Space: O(2^M) using rolling arrays.
 *
 * 3.  HOW TO ADAPT IT (THE DIALS & SWITCHES)
 * - Obstacles: If the grid has blocked cells, simply check if the current cell
 *   (i, j) is blocked. If it is, ensure the current bit in the mask is 0, shift it
 *   to 1 (meaning it's "filled"), and move to the next cell without placing dominoes.
 */

using ll = long long;

ll plug_dp_dominoes(int n, int m, ll mod)
{
    if ((n * m) % 2 != 0)
        return 0; // Odd number of cells cannot be tiled by 1x2

    // We only need the current and next state, so we use a rolling array of size 2
    vector<vector<ll>> dp(2, vector<ll>(1 << m, 0));

    int cur = 0, next = 1;
    dp[cur][0] = 1; // Base case: completely empty profile

    for (int i = 0; i < n; i++)
    {
        for (int j = 0; j < m; j++)
        {
            fill(dp[next].begin(), dp[next].end(), 0);

            for (int mask = 0; mask < (1 << m); mask++)
            {
                if (dp[cur][mask] == 0)
                    continue;

                // Case 1: The current cell is already filled by a vertical domino
                // from the previous row. We must leave it alone. The bit shifts out.
                if (mask & (1 << (m - 1)))
                {
                    int next_mask = (mask ^ (1 << (m - 1))) << 1;
                    dp[next][next_mask] = (dp[next][next_mask] + dp[cur][mask]) % mod;
                }
                // Case 2: The cell is empty. We must place a domino.
                else
                {
                    // Option A: Place a vertical domino (protrudes into the next row).
                    // We set the 0-th bit of the shifted mask to 1.
                    int next_mask_vert = (mask << 1) | 1;
                    dp[next][next_mask_vert] = (dp[next][next_mask_vert] + dp[cur][mask]) % mod;

                    // Option B: Place a horizontal domino.
                    // Only valid if we aren't at the right edge AND the next cell is empty.
                    if (j + 1 < m && !(mask & 1))
                    {
                        // The horizontal domino fills the current and next cell in THIS row.
                        // We shift the mask by 1, but we don't protrude into the next row,
                        // so the 0-th bit remains 0. However, the cell to our right is now filled,
                        // which we must account for in our transition logic by skipping it or
                        // setting its bit. In this cell-by-cell state machine, placing a horizontal
                        // domino effectively sets the bit for the *next* cell in the current row.
                        int next_mask_horiz = (mask << 1) | 2;
                        dp[next][next_mask_horiz] = (dp[next][next_mask_horiz] + dp[cur][mask]) % mod;
                    }
                }
            }
            swap(cur, next);
        }
    }

    // The answer is the number of ways to have a completely flat profile (mask == 0) at the end.
    return dp[cur][0];
}
\end{lstlisting}
\subsection{More Tricks}
\subsubsection{Slope trick}
\begin{lstlisting}[language=C++]
/**
 *  THE ULTIMATE CP TEMPLATE BLUEPRINT: SLOPE TRICK
 *
 * 1.  WHAT PROBLEM DOES THIS SOLVE?
 * - Keywords: "Minimum operations to make array non-decreasing".
 * - Classic Scenarios: You can increment or decrement elements of an array.
 *   Find the minimum total cost to make the array sorted (non-decreasing).
 * - The Magic: A 2D DP takes O(N * MAX_VAL) which TLEs. By observing that the
 *   DP cost function is convex and piecewise linear, we can just maintain the
 *   points where the slope changes using a max-heap.
 *
 * 2.  HOW TO USE IT (THE BLACK BOX)
 * - Query:
 *       long long min_cost = make_non_decreasing(A);
 *
 * 3.  HOW TO ADAPT IT (THE DIALS & SWITCHES)
 * - Strictly Increasing: If the problem requires strictly increasing (A[i] < A[i+1]),
 *   subtract `i` from every element initially: `A[i] -= i`, then run this exact
 *   function. It maps the strictly increasing problem directly to a non-decreasing one.
 *
 * - Complexity:
 *       Time: O(N log N)
 *       Space: O(N)
 */

using ll = long long;

ll make_non_decreasing(const vector<ll> &a)
{
    priority_queue<ll> pq;
    ll total_cost = 0;

    for (int i = 0; i < a.size(); i++)
    {
        pq.push(a[i]);

        // If the max element so far is strictly greater than the current element,
        // it means we have a slope violation. We "flatten" it to match A[i].
        if (pq.top() > a[i])
        {
            total_cost += pq.top() - a[i];
            pq.pop();
            pq.push(a[i]);
        }
    }

    return total_cost;
}
\end{lstlisting}
\subsection{Optimization}
\subsubsection{1D1D}
\begin{lstlisting}[language=C++]
/**
 *  THE ULTIMATE CP TEMPLATE BLUEPRINT: 1D1D OPTIMIZATION (LINE CONTAINER)
 *
 * 1.  WHAT PROBLEM DOES THIS SOLVE?
 * - Keywords: "Convex Hull Trick", "y = mx + c", "Li Chao Tree alternative".
 * - Classic Scenarios: dp[i] = min_{j < i} (dp[j] + M[j] * X[i] + C[j]).
 * - The Magic: Maintains the upper/lower envelope of a set of lines. It supports
 *   adding lines dynamically and querying the extreme value at any X in O(log N).
 *   This specific implementation (KACTL's LineContainer) handles lines added with
 *   NON-MONOTONIC slopes, making it universally safe.
 *
 * 2.  HOW TO USE IT (THE BLACK BOX)
 * - Initialization:
 *       LineContainer cht;
 * - Operations:
 *       cht.add(m, c);      // Adds line y = mx + c
 *       ll max_y = cht.query(x);
 *
 * 3.  HOW TO ADAPT IT (THE DIALS & SWITCHES)
 * - Minimization: This natively computes the MAXIMUM value at x. If you need the
 *   MINIMUM value, insert lines with inverted signs: `cht.add(-m, -c)` and negate
 *   the queried result: `-cht.query(x)`.
 */

using ll = long long;

struct Line
{
    mutable ll k, m, p;
    bool operator<(const Line &o) const { return k < o.k; }
    bool operator<(ll x) const { return p < x; }
};

struct LineContainer : multiset<Line, less<>>
{
    static const ll inf = LLONG_MAX;
    ll div(ll a, ll b)
    {
        return a / b - ((a ^ b) < 0 && a % b);
    }
    bool isect(iterator x, iterator y)
    {
        if (y == end())
            return x->p = inf, 0;
        if (x->k == y->k)
            x->p = x->m > y->m ? inf : -inf;
        else
            x->p = div(y->m - x->m, x->k - y->k);
        return x->p >= y->p;
    }
    void add(ll k, ll m)
    {
        auto z = insert({k, m, 0}), y = z++, x = y;
        while (isect(y, z))
            z = erase(z);
        if (x != begin() && isect(--x, y))
            isect(x, y = erase(y));
        while ((y = x) != begin() && (--x)->p >= y->p)
            isect(x, erase(y));
    }
    ll query(ll x)
    {
        assert(!empty());
        auto l = *lower_bound(x);
        return l.k * x + l.m;
    }
};
\end{lstlisting}
\subsubsection{Aliens}
\begin{lstlisting}[language=C++]
#include <bits/stdc++.h>

using namespace std;

#define El_Hefnawys                   \
    ios_base::sync_with_stdio(false); \
    cin.tie(nullptr);                 \
    cout.tie(nullptr)

typedef long long ll;
#define int ll

const int OO = 1e15;

// choose k subarrays from the array and trying to maximize the sum
// dp[i][0/1]

// The DP function returns a pair:
// 1. The maximum sum (including the penalties applied)
// 2. The NEGATIVE number of subarrays used (to break ties by preferring fewer subarrays)
pair<int, int> check(int penalty, int n, const vector<int> &a)
{
    // dp[i][0] = max sum up to index i, currently NOT in a subarray
    // dp[i][1] = max sum up to index i, currently INSIDE a subarray
    vector<vector<pair<int, int>>> dp(n, vector<pair<int, int>>(2));

    dp[0][0] = {0, 0};
    dp[0][1] = {a[0] - penalty, -1}; // Starting an array costs 'penalty'

    for (int i = 1; i < n; i++)
    {
        // Option 0: We are not in a subarray.
        // We either were already not in one, or we just closed the previous one.
        dp[i][0] = max(dp[i - 1][0], dp[i - 1][1]);

        // Option 1: We are inside a subarray.
        // Transition A: Start a brand NEW subarray today (costs penalty, adds -1 to count)
        pair<int, int> start_new = make_pair(
            dp[i - 1][0].first - penalty + a[i],
            dp[i - 1][0].second - 1);

        // Transition B: Extend the EXISTING subarray (no penalty, count doesn't change)
        pair<int, int> extend_existing = make_pair(
            dp[i - 1][1].first + a[i],
            dp[i - 1][1].second);

        dp[i][1] = max(start_new, extend_existing);
    }

    return max(dp[n - 1][0], dp[n - 1][1]);
}

void solve()
{
    int n, k;
    if (!(cin >> n >> k))
        return;

    vector<int> a(n);
    for (int i = 0; i < n; i++)
    {
        cin >> a[i];
    }

    // The penalty bounds. Can be negative depending on the problem!
    int l = -1e14, r = 1e14;
    int best_ans = 0;

    while (l <= r)
    {
        int mid_penalty = l + (r - l) / 2;

        pair<int, int> res = check(mid_penalty, n, a);
        int arrays_used = -res.second; // Flip back to positive

        // If we used at least K arrays, this penalty is valid to consider.
        // We might have tied and used exactly K, or more.
        if (arrays_used >= k)
        {
            // Restore the true sum by adding back the penalty for EXACTLY K arrays.
            // Even if the DP used more than K due to ties (0-sum arrays),
            // the mathematical projection at the hull tangent guarantees this formula is exact!
            best_ans = res.first + (k * mid_penalty);

            // Try to force it to use FEWER arrays by INCREASING the penalty
            l = mid_penalty + 1;
        }
        else
        {
            // We used too few arrays. The penalty is too harsh. Decrease it.
            r = mid_penalty - 1;
        }
    }

    cout << best_ans << "\n";
}

signed main()
{
    El_Hefnawys;
    int t{1};
    // cin >> t;
    while (t--)
    {
        solve();
    }
    return 0;
}

\end{lstlisting}
\subsubsection{DC}
\begin{lstlisting}[language=C++]
/**
 *  THE ULTIMATE CP TEMPLATE BLUEPRINT: DIVIDE AND CONQUER OPTIMIZATION
 *
 * 1.  WHAT PROBLEM DOES THIS SOLVE?
 * - Keywords: "Partition array into K contiguous segments", "Minimize maximum cost".
 * - Classic Scenarios: dp[i][j] = min_{k < j} (dp[i-1][k] + cost(k, j)).
 * - The Magic: If the optimal split point `opt(i, j)` satisfies the condition
 *   opt(i, j) <= opt(i, j+1), we don't need to check all `k`. We recursively
 *   compute the middle state and restrict the search bounds for the left and
 *   right halves, dropping O(K * N^2) to O(K * N log N).
 *
 * 2.  HOW TO USE IT (THE BLACK BOX)
 * - Setup: You need two 1D arrays: `dp_prev` (results from i-1 segments)
 *   and `dp_cur` (results for i segments).
 * - Query:
 *       for (int i = 2; i <= K; i++) {
 *           compute_dc(1, N, 1, N, dp_prev, dp_cur);
 *           dp_prev = dp_cur;
 *       }
 */

using ll = long long;
const ll INF = 1e18;

// Replace this with your actual O(1) cost function (usually via prefix sums)
ll cost(int l, int r)
{
    if (l > r)
        return 0;
    return 0; // return actual cost
}

void compute_dc(int l, int r, int opt_l, int opt_r, const vector<ll> &dp_prev, vector<ll> &dp_cur)
{
    if (l > r)
        return;
    int mid = l + (r - l) / 2;
    pair<ll, int> best = {INF, -1};

    for (int k = opt_l; k <= min(mid, opt_r); k++)
    {
        ll current_cost = dp_prev[k - 1] + cost(k, mid);
        if (current_cost < best.first)
        {
            best = {current_cost, k};
        }
    }

    dp_cur[mid] = best.first;
    int opt = best.second;

    compute_dc(l, mid - 1, opt_l, opt, dp_prev, dp_cur);
    compute_dc(mid + 1, r, opt, opt_r, dp_prev, dp_cur);
}
\end{lstlisting}
\subsubsection{Knuth}
\begin{lstlisting}[language=C++]
/**
 *  THE ULTIMATE CP TEMPLATE BLUEPRINT: DIVIDE AND CONQUER OPTIMIZATION
 *
 * 1.  WHAT PROBLEM DOES THIS SOLVE?
 * - Keywords: "Partition array into K contiguous segments", "Minimize maximum cost".
 * - Classic Scenarios: dp[i][j] = min_{k < j} (dp[i-1][k] + cost(k, j)).
 * - The Magic: If the optimal split point `opt(i, j)` satisfies the condition
 *   opt(i, j) <= opt(i, j+1), we don't need to check all `k`. We recursively
 *   compute the middle state and restrict the search bounds for the left and
 *   right halves, dropping O(K * N^2) to O(K * N log N).
 *
 * 2.  HOW TO USE IT (THE BLACK BOX)
 * - Setup: You need two 1D arrays: `dp_prev` (results from i-1 segments)
 *   and `dp_cur` (results for i segments).
 * - Query:
 *       for (int i = 2; i <= K; i++) {
 *           compute_dc(1, N, 1, N, dp_prev, dp_cur);
 *           dp_prev = dp_cur;
 *       }
 */

using ll = long long;
const ll INF = 1e18;

// Replace this with your actual O(1) cost function (usually via prefix sums)
ll cost(int l, int r)
{
    if (l > r)
        return 0;
    return 0; // return actual cost
}

void compute_dc(int l, int r, int opt_l, int opt_r, const vector<ll> &dp_prev, vector<ll> &dp_cur)
{
    if (l > r)
        return;
    int mid = l + (r - l) / 2;
    pair<ll, int> best = {INF, -1};

    for (int k = opt_l; k <= min(mid, opt_r); k++)
    {
        ll current_cost = dp_prev[k - 1] + cost(k, mid);
        if (current_cost < best.first)
        {
            best = {current_cost, k};
        }
    }

    dp_cur[mid] = best.first;
    int opt = best.second;

    compute_dc(l, mid - 1, opt_l, opt, dp_prev, dp_cur);
    compute_dc(mid + 1, r, opt, opt_r, dp_prev, dp_cur);
}
\end{lstlisting}
\subsubsection{Using data structures}
\begin{lstlisting}[language=C++]
struct Node
{
    int val;

    Node()
    {
        val = 0;
    }

    Node(int val)
        : val(val)
    {
    }

    void change(int x)
    {
        val = max(x, val);
    }
};

struct segmentTree
{
    int sz;
    vector<Node> seg;

    segmentTree(int n)
    {
        sz = 1;
        while (sz < n)
            sz *= 2;
        seg.assign(2 * sz, Node());
    }

    void init(vector<int> &arr, int node, int lx, int rx)
    {
        if (rx - lx == 1)
        {
            if (lx < int(arr.size()))
                seg[node] = Node(arr[lx]);
            return;
        }
        int mid = (lx + rx) / 2;
        init(arr, 2 * node + 1, lx, mid);
        init(arr, 2 * node + 2, mid, rx);
        seg[node] = merge(seg[2 * node + 1], seg[2 * node + 2]);
    }

    void init(vector<int> &arr)
    {
        init(arr, 0, 0, sz);
    }

    Node merge(Node &lf, Node &rt)
    {
        Node ans = Node();
        ans.val = max(lf.val, rt.val);
        return ans;
    }

    void set(int idx, int val, int node, int lx, int rx)
    {
        if (rx - lx == 1)
        {
            seg[node].change(val);
            return;
        }
        int mid = (lx + rx) / 2;
        if (idx < mid)
            set(idx, val, 2 * node + 1, lx, mid);
        else
            set(idx, val, 2 * node + 2, mid, rx);
        seg[node] = merge(seg[2 * node + 1], seg[2 * node + 2]);
    }

    void set(int idx, int val)
    {
        set(idx, val, 0, 0, sz);
    }

    Node get(int l, int r, int node, int lx, int rx)
    {
        if (lx >= l && rx <= r)
            return seg[node];
        if (rx <= l || lx >= r)
            return Node();
        int mid = (lx + rx) / 2;
        Node lf = get(l, r, 2 * node + 1, lx, mid);
        Node ri = get(l, r, 2 * node + 2, mid, rx);
        return merge(lf, ri);
    }

    int get(int l, int r)
    {
        return get(l, r, 0, 0, sz).val;
    }
};

int LIS(int n, vector<int> &p)
{
    segmentTree seg(n + 5);
    vector<int> dp(n + 1);
    int ans{};
    for (int i = 0; i < n; i++)
    {
        dp[i] = seg.get(0, p[i]) + 1;
        seg.set(p[i], dp[i]);
        ans = max(ans, dp[i]);
    }
    return ans;
}

\end{lstlisting}
\section{Misc}
\subsection{Big Integer}
\subsubsection{Big integers}
\begin{lstlisting}[language=C++]
/**
 *  THE ULTIMATE CP TEMPLATE BLUEPRINT: COMPRESSED BIGINT
 *
 * 1.  WHAT PROBLEM DOES THIS SOLVE?
 * - Keywords: "No modulo given", "Massive addition/multiplication", "1000 digits".
 * - Classic Scenarios: DP state sums or combinatorics where the answer easily
 *   exceeds 2^64, but no modulo is provided.
 * - The Magic: C++ has no native BigInt. Using strings character-by-character
 *   is O(N^2) with a massive constant factor. This struct bundles 9 decimal digits
 *   into a single integer, processing data 9x faster and using significantly less memory.
 *
 * 2.  HOW TO USE IT (THE BLACK BOX)
 * - Initialization:
 *       BigInt a = "12345678901234567890";
 *       BigInt b = 987654321;
 * - Operations:
 *       BigInt c = a + b;
 *       c.print();
 */

const int BASE = 1e9;

struct BigInt
{
    vector<int> a;

    BigInt() {}
    BigInt(long long v)
    {
        while (v > 0)
        {
            a.push_back(v % BASE);
            v /= BASE;
        }
    }
    BigInt(string s)
    {
        if (s.empty())
            return;
        for (int i = s.length(); i > 0; i -= 9)
        {
            if (i < 9)
                a.push_back(stoll(s.substr(0, i)));
            else
                a.push_back(stoll(s.substr(i - 9, 9)));
        }
        trim();
    }

    void trim()
    {
        while (a.size() > 1 && a.back() == 0)
            a.pop_back();
    }

    BigInt operator+(const BigInt &b) const
    {
        BigInt res;
        int carry = 0;
        for (int i = 0; i < max(a.size(), b.a.size()) || carry; i++)
        {
            long long cur = carry;
            if (i < a.size())
                cur += a[i];
            if (i < b.a.size())
                cur += b.a[i];
            res.a.push_back(cur % BASE);
            carry = cur / BASE;
        }
        return res;
    }

    BigInt operator*(const BigInt &b) const
    {
        BigInt res;
        if (a.empty() || b.a.empty())
            return res;
        res.a.assign(a.size() + b.a.size(), 0);
        for (int i = 0; i < a.size(); i++)
        {
            long long carry = 0;
            for (int j = 0; j < b.a.size() || carry > 0; j++)
            {
                long long cur = res.a[i + j] + a[i] * 1LL * (j < b.a.size() ? b.a[j] : 0) + carry;
                res.a[i + j] = cur % BASE;
                carry = cur / BASE;
            }
        }
        res.trim();
        return res;
    }

    bool operator<(const BigInt &b) const
    {
        if (a.size() != b.a.size())
            return a.size() < b.a.size();
        for (int i = a.size() - 1; i >= 0; i--)
        {
            if (a[i] != b.a[i])
                return a[i] < b.a[i];
        }
        return false;
    }

    void print() const
    {
        if (a.empty())
        {
            cout << 0 << "\n";
            return;
        }
        cout << a.back();
        for (int i = (int)a.size() - 2; i >= 0; i--)
        {
            cout << setfill('0') << setw(9) << a[i];
        }
        cout << "\n";
    }
};
\end{lstlisting}
\subsection{Chess}
\subsubsection{Knight's tour}
\begin{lstlisting}[language=C++]
/**
 *  THE ULTIMATE CP TEMPLATE BLUEPRINT: WARNSDORFF'S HEURISTIC
 *
 * 1.  WHAT PROBLEM DOES THIS SOLVE?
 * - Keywords: "Visit every square", "Knight's Tour", "Hamiltonian path on grid".
 * - Classic Scenarios: Output a valid sequence of knight moves to visit all N x N
 *   cells exactly once without repeating.
 * - The Magic: Backtracking has a massive branching factor and TLEs immediately.
 *   Warnsdorff's Rule states: always prioritize the next valid move that has the
 *   LEAST number of subsequent valid moves. This local greedy choice finds the
 *   global Hamiltonian path almost instantly.
 *
 * 2.  HOW TO USE IT (THE BLACK BOX)
 * - Query:
 *       vector<pair<int, int>> tour;
 *       bool success = knight_tour(START_X, START_Y, BOARD_SIZE, tour);
 */

const int dx[] = {1, 1, 2, 2, -1, -1, -2, -2};
const int dy[] = {2, -2, 1, -1, 2, -2, 1, -1};

int get_degree(int x, int y, int n, const vector<vector<int>> &board)
{
    int count = 0;
    for (int i = 0; i < 8; i++)
    {
        int nx = x + dx[i], ny = y + dy[i];
        if (nx >= 0 && ny >= 0 && nx < n && ny < n && board[nx][ny] == -1)
        {
            count++;
        }
    }
    return count;
}

bool solve_kt(int x, int y, int move_i, int n, vector<vector<int>> &board, vector<pair<int, int>> &tour)
{
    board[x][y] = move_i;
    tour.push_back({x, y});

    if (move_i == n * n - 1)
        return true;

    vector<pair<int, int>> next_moves;
    for (int i = 0; i < 8; i++)
    {
        int nx = x + dx[i], ny = y + dy[i];
        if (nx >= 0 && ny >= 0 && nx < n && ny < n && board[nx][ny] == -1)
        {
            next_moves.push_back({get_degree(nx, ny, n, board), i});
        }
    }

    // Warnsdorff's Heuristic: Sort by degree ascending
    sort(next_moves.begin(), next_moves.end());

    for (auto move : next_moves)
    {
        int i = move.second;
        int nx = x + dx[i], ny = y + dy[i];
        if (solve_kt(nx, ny, move_i + 1, n, board, tour))
        {
            return true;
        }
    }

    // Backtrack
    board[x][y] = -1;
    tour.pop_back();
    return false;
}

bool knight_tour(int start_x, int start_y, int n, vector<pair<int, int>> &tour)
{
    vector<vector<int>> board(n, vector<int>(n, -1));
    return solve_kt(start_x, start_y, 0, n, board, tour);
}
\end{lstlisting}
\subsection{Matroids}
\subsubsection{Matroid intersection}
\begin{lstlisting}[language=C++]
/**
 *  THE ULTIMATE CP TEMPLATE BLUEPRINT: MATROID INTERSECTION
 *
 * 1.  WHAT PROBLEM DOES THIS SOLVE?
 * - Keywords: "Spanning tree with unique colors", "Maximize intersection of rules".
 * - Classic Scenarios: Find the maximum number of elements from a set that satisfy
 *   TWO independent structural rules simultaneously (e.g., forms a forest AND has
 *   at most K edges from any specific set).
 * - The Magic: It builds a bipartite graph between the currently picked elements (IN)
 *   and unpicked elements (OUT). It searches for an augmenting path that swaps elements
 *   while maintaining independence in BOTH matroids.
 *
 * 2.  HOW TO USE IT (THE BLACK BOX)
 * - Setup: You must define what "independence" means for your two rules by writing
 *   the `oracle1_can_add` and `oracle2_can_add` logic.
 * - Query:
 *       vector<int> independent_set = matroid_intersection(N);
 *
 * 3.  ENGINEERING NOTE
 * - The robust, modular design here perfectly mirrors the isolation of concerns
 *   required in high-tier backend systems engineering. The core augmenting engine
 *   never cares about the underlying rules, only the true/false oracle responses.
 */

vector<int> matroid_intersection(int n)
{
    vector<bool> in_set(n, false);

    // Abstract augmenting path loop
    while (true)
    {
        // Build the bipartite exchange graph
        vector<int> X1, X2; // X1: valid starts, X2: valid ends
        vector<vector<int>> adj(n);

        // --- ORACLE 1 (e.g., Graphic Matroid) ---
        // For every out_node 'v', if {in_set U v} is independent, add 'v' to X1.
        // If dependent, find which in_node 'u' forms the circuit and add edge v -> u.

        // --- ORACLE 2 (e.g., Colorful Matroid) ---
        // For every out_node 'v', if {in_set U v} is independent, add 'v' to X2.
        // If dependent, find which in_node 'u' forms the circuit and add edge u -> v.

        // Example placeholders:
        /*
        for (int i = 0; i < n; i++) {
            if (!in_set[i]) {
                if (oracle1_can_add(i)) X1.push_back(i);
                if (oracle2_can_add(i)) X2.push_back(i);
            }
        }
        */

        // BFS to find the shortest augmenting path from X1 to X2
        queue<int> q;
        vector<int> dist(n, 1e9), parent(n, -1);

        for (int x : X1)
        {
            dist[x] = 0;
            q.push(x);
        }

        int end_node = -1;
        while (!q.empty())
        {
            int u = q.front();
            q.pop();

            // Check if we hit the target set
            if (find(X2.begin(), X2.end(), u) != X2.end())
            {
                end_node = u;
                break;
            }

            for (int v : adj[u])
            {
                if (dist[v] > dist[u] + 1)
                {
                    dist[v] = dist[u] + 1;
                    parent[v] = u;
                    q.push(v);
                }
            }
        }

        // If no augmenting path exists, we have reached the maximum independent set
        if (end_node == -1)
            break;

        // Augment the independent set by flipping states along the path
        int curr = end_node;
        while (curr != -1)
        {
            in_set[curr] = !in_set[curr];
            curr = parent[curr];
        }
    }

    vector<int> result;
    for (int i = 0; i < n; i++)
    {
        if (in_set[i])
            result.push_back(i);
    }
    return result;
}
\end{lstlisting}
\subsection{Search}
\subsubsection{Ternary search}
\begin{lstlisting}[language=C++]
#include <bits/stdc++.h>
using namespace std;

// The function we want to minimize (e.g., total cost, distance, time)
// For this example, let's pretend it's a parabola: f(x) = (x - 5)^2 + 10
// The minimum should clearly be at x = 5.
double f_double(double x)
{
    return (x - 5.0) * (x - 5.0) + 10.0;
}

long long f_int(long long x)
{
    return (x - 5) * (x - 5) + 10;
}

// ---------------------------------------------------------
// TEMPLATE 1: Continuous Ternary Search (Floating Point)
// ---------------------------------------------------------
double ternary_search_double(double l, double r)
{
    // THE CP TRICK: Never use `while (r - l > eps)`. Floating point precision
    // errors can cause an infinite loop and TLE your submission.
    // Instead, looping exactly 200 times guarantees immense precision
    // (the range shrinks by 2/3^200, which is basically 0).
    for (int i = 0; i < 200; i++)
    {
        double m1 = l + (r - l) / 3.0;
        double m2 = r - (r - l) / 3.0;

        if (f_double(m1) < f_double(m2))
        {
            r = m2; // We are minimizing, so throw away the higher right side
        }
        else
        {
            l = m1; // Throw away the higher left side
        }
    }
    return l; // l and r are essentially identical now
}

// ---------------------------------------------------------
// TEMPLATE 2: Discrete Ternary Search (Integers)
// ---------------------------------------------------------
long long ternary_search_int(long long l, long long r)
{
    // We stop when the search space shrinks to 3 elements or less.
    // If we try to split a range of size 2 into thirds with integers,
    // m1 and m2 will evaluate to the same number and cause an infinite loop!
    while (r - l > 2)
    {
        long long m1 = l + (r - l) / 3;
        long long m2 = r - (r - l) / 3;

        if (f_int(m1) < f_int(m2))
        {
            r = m2;
        }
        else
        {
            l = m1;
        }
    }

    // Now we have a tiny range (l, l+1, r). Just brute force the answer!
    long long best_x = l;
    long long min_val = f_int(l);

    for (long long i = l + 1; i <= r; i++)
    {
        if (f_int(i) < min_val)
        {
            min_val = f_int(i);
            best_x = i;
        }
    }

    return best_x;
}

void solve()
{
    cout << "Continuous Minimum at X = " << ternary_search_double(-1e9, 1e9) << "\n";
    cout << "Discrete Minimum at X = " << ternary_search_int(-1000000, 1000000) << "\n";
}

int main()
{
    ios_base::sync_with_stdio(false);
    cin.tie(NULL);
    solve();
    return 0;
}

\end{lstlisting}
\subsubsection{fraction binary search}
\begin{lstlisting}[language=C++]
/**
 *  THE ULTIMATE CP TEMPLATE BLUEPRINT: FRACTION BINARY SEARCH
 *
 * 1.  WHAT PROBLEM DOES THIS SOLVE?
 * - Keywords: "Best rational approximation", "Fraction with bounded denominator".
 * - Classic Scenarios: You are given a target floating point value and must find
 *   a fraction P/Q exactly equal to it (or closest to it) such that Q <= MAX_D.
 * - The Magic: The Stern-Brocot tree generates all positive rational numbers in
 *   simplest form. Given bounds a/b and c/d, the mediant (a+c)/(b+d) is strictly
 *   between them. We binary search down this tree.
 *
 * 2.  HOW TO USE IT (THE BLACK BOX)
 * - Query:
 *       pair<ll, ll> best_fraction = fraction_search(target, MAX_D);
 *
 * - Complexity:
 *       Time: O(log(MAX_D))
 */

using ll = long long;

// Returns {numerator, denominator}
pair<ll, ll> fraction_search(double target, ll max_d)
{
    ll a = 0, b = 1; // Lower bound: 0/1
    ll c = 1, d = 0; // Upper bound: 1/0 (Infinity)

    ll best_num = 0, best_den = 1;
    double min_diff = 1e18;

    while (b + d <= max_d)
    {
        ll mid_num = a + c;
        ll mid_den = b + d;

        double mid_val = (double)mid_num / mid_den;
        double diff = abs(mid_val - target);

        if (diff < min_diff)
        {
            min_diff = diff;
            best_num = mid_num;
            best_den = mid_den;
        }

        if (mid_val == target)
        {
            return {mid_num, mid_den};
        }
        else if (mid_val < target)
        {
            a = mid_num;
            b = mid_den;
        }
        else
        {
            c = mid_num;
            d = mid_den;
        }
    }

    return {best_num, best_den};
}
\end{lstlisting}
\subsubsection{parallel binary search}
\begin{lstlisting}[language=C++]
/**
 *  THE ULTIMATE CP TEMPLATE BLUEPRINT: PARALLEL BINARY SEARCH
 *
 * 1.  WHAT PROBLEM DOES THIS SOLVE?
 * - Keywords: "Earliest time all components connect", "Minimum weight to satisfy Q queries".
 * - Classic Scenarios: You are given an empty graph and an array of edges added over time.
 *   You have Q queries asking: "At what time do nodes U and V finally connect?"
 * - The Magic: It runs a binary search for ALL queries simultaneously. It sweeps
 *   through the timeline 1 to T exactly log(T) times. At each step, it checks if
 *   the current timeline state satisfies the query's midpoint.
 *
 * 2.  HOW TO USE IT (THE BLACK BOX)
 * - Array Setup:
 *       L[i] = 1, R[i] = T for all queries i.
 *       ans[i] = -1.
 * - Loop:
 *       Run the `while` loop until no queries need updating.
 *
 * 3.  HOW TO ADAPT IT (THE DIALS & SWITCHES)
 * - State Machine: The `DSU` and `apply_event` logic inside the timeline sweep MUST
 *   be replaced with whatever data structure the specific problem demands
 *   (e.g., Fenwick Tree for range sums, Segment Tree, etc.).
 */

struct Query
{
    int u, v; // Replace with problem-specific query parameters
};

void parallel_binary_search(int num_events, const vector<Query> &queries)
{
    int q = queries.size();
    vector<int> L(q, 1), R(q, num_events), ans(q, -1);
    bool changed = true;

    while (changed)
    {
        changed = false;
        // check_at[t] stores a list of query indices that need to be evaluated at time t
        vector<vector<int>> check_at(num_events + 1);

        for (int i = 0; i < q; i++)
        {
            if (L[i] <= R[i])
            {
                int mid = L[i] + (R[i] - L[i]) / 2;
                check_at[mid].push_back(i);
                changed = true;
            }
        }

        if (!changed)
            break;

        // --- RESET DATA STRUCTURE HERE ---
        // e.g., DSU dsu(N);
        // ---------------------------------

        for (int t = 1; t <= num_events; t++)
        {
            // --- APPLY EVENT 't' TO DATA STRUCTURE ---
            // e.g., dsu.unite(edges[t].u, edges[t].v);
            // -----------------------------------------

            for (int query_idx : check_at[t])
            {
                // --- EVALUATE QUERY CONDITION ---
                // e.g., bool ok = dsu.same(queries[query_idx].u, queries[query_idx].v);
                bool ok = true; // Replace with actual condition
                // --------------------------------

                if (ok)
                {
                    ans[query_idx] = t;
                    R[query_idx] = t - 1; // Try to find an earlier time
                }
                else
                {
                    L[query_idx] = t + 1; // Need more time/events
                }
            }
        }
    }

    // 'ans' now contains the exact minimum time for every query in O((N+Q) log T)
}
\end{lstlisting}
\subsection{Techniques}
\subsubsection{MANHATTAN CHEBYSHEV}
\begin{lstlisting}[language=C++]
/**
 *  THE ULTIMATE CP TEMPLATE BLUEPRINT: MANHATTAN <-> CHEBYSHEV
 *
 * 1.  WHAT PROBLEM DOES THIS SOLVE?
 * - Keywords: "Rotate coordinate system", "Simplify absolute value distance".
 * - Classic Scenarios: A problem asks you to find the number of points within
 *   Manhattan distance D from (X, Y). This is an "octagonal" area, which is
 *   extremely hard to query.
 * - The Magic:
 *   Transformation: (x, y) -> (x+y, x-y)
 *   Manhattan distance in the original plane is equivalent to Chebyshev distance
 *   in the transformed plane.
 *   Chebyshev distance: max(|x1 - x2|, |y1 - y2|) creates a *square* area, which
 *   is trivially solvable using standard 2D Segment Trees or Fenwick Trees.
 *
 * 2.  HOW TO USE IT
 * - Transform:
 *       new_x = x + y;
 *       new_y = x - y;
 */
\end{lstlisting}
\subsubsection{Meet in the middle}
\begin{lstlisting}[language=C++]
#include <bits/stdc++.h>
using namespace std;

// This template solves the classic "Find if any subset sums to Target"
// Complexity: O(2^(N/2) * N)

void generate_sums(const vector<long long> &arr, vector<long long> &sums)
{
    int n = arr.size();
    // Use bitmasking to generate all 2^n subsets
    for (int i = 0; i < (1 << n); i++)
    {
        long long current_sum = 0;
        for (int j = 0; j < n; j++)
        {
            if ((i >> j) & 1)
            {
                current_sum += arr[j];
            }
        }
        sums.push_back(current_sum);
    }
}

void solve()
{
    int n;
    long long target;
    cin >> n >> target;

    vector<long long> arr(n);
    for (int i = 0; i < n; i++)
        cin >> arr[i];

    // 1. Split into two halves
    vector<long long> left_part, right_part;
    for (int i = 0; i < n; i++)
    {
        if (i < n / 2)
            left_part.push_back(arr[i]);
        else
            right_part.push_back(arr[i]);
    }

    // 2. Generate all subset sums for both halves
    vector<long long> left_sums, right_sums;
    generate_sums(left_part, left_sums);
    generate_sums(right_part, right_sums);

    // 3. Sort one half for binary searching
    sort(left_sums.begin(), left_sums.end());

    // 4. Search for the complement in the other half
    bool found = false;
    for (long long s : right_sums)
    {
        long long needed = target - s;
        // Check if 'needed' exists in left_sums
        if (binary_search(left_sums.begin(), left_sums.end(), needed))
        {
            found = true;
            break;
        }
    }

    if (found)
        cout << "YES\n";
    else
        cout << "NO\n";
}

int main()
{
    ios_base::sync_with_stdio(false);
    cin.tie(NULL);
    solve();
    return 0;
}

\end{lstlisting}
\subsubsection{Sparse table lsa trick}
\begin{lstlisting}[language=C++]
/**
 *  THE ULTIMATE CP TEMPLATE BLUEPRINT: SPARSE TABLE (LSA TRICK)
 *
 * 1.  WHAT PROBLEM DOES THIS SOLVE?
 * - Keywords: "Range Maximum Query", "Immutable range query".
 * - Classic Scenarios: You have an array that never changes, and you need to
 *   query the maximum/minimum in any range [L, R] in O(1).
 * - The Magic: Precomputes logs for all ranges of length 2^k. Since `max` is
 *   idempotent (max(a, a) = a), we can answer queries by taking the max of two
 *   overlapping 2^k blocks that cover the entire [L, R].
 */

const int K = 20; // Enough for 10^6 elements (2^20 > 10^6)
int st[K][100005];
int lg[100005];

void build_sparse_table(const vector<int> &a)
{
    int n = a.size();
    lg[1] = 0;
    for (int i = 2; i <= n; i++)
        lg[i] = lg[i / 2] + 1;
    for (int i = 0; i < n; i++)
        st[0][i] = a[i];

    for (int j = 1; j < K; j++)
    {
        for (int i = 0; i + (1 << j) <= n; i++)
        {
            st[j][i] = max(st[j - 1][i], st[j - 1][i + (1 << (j - 1))]);
        }
    }
}

int query(int L, int R)
{
    int j = lg[R - L + 1];
    return max(st[j][L], st[j][R - (1 << j) + 1]);
}
\end{lstlisting}
\subsection{min plus convolution.cpp}
\subsubsection{shawk}
\begin{lstlisting}[language=C++]
/**
 *  THE ULTIMATE CP TEMPLATE BLUEPRINT: SMAWK ALGORITHM
 *
 * 1.  WHAT PROBLEM DOES THIS SOLVE?
 * - Keywords: "Totally monotone matrix", "Min-plus convolution", "O(N) optimization".
 * - Classic Scenarios: You need to find the minimum value in every row of an N x M
 *   matrix where the matrix satisfies the Monge property (or is totally monotone).
 * - The Magic: Bypasses standard O(N M) brute force and even O(N log M) Divide &
 *   Conquer. It recursively filters out columns that can never contain a row minimum,
 *   solving the entire problem in strictly O(N + M) time.
 *
 * 2.  HOW TO USE IT (THE BLACK BOX)
 * - Query:
 *       vector<int> min_indices = smawk(R, C, cost_function);
 * - Result:
 *       Returns an array `ans` of size R, where `ans[i]` is the column index
 *       containing the minimum value for row i.
 */

using ll = long long;

template <class F>
vector<int> smawk(int r, int c, F cost)
{
    auto solve = [&](auto &self, const vector<int> &rows, const vector<int> &cols) -> vector<int>
    {
        if (rows.empty())
            return {};

        // Reduce: filter columns
        vector<int> st;
        for (int c : cols)
        {
            while (!st.empty())
            {
                if (cost(rows[st.size() - 1], st.back()) > cost(rows[st.size() - 1], c))
                {
                    st.pop_back();
                }
                else if (st.size() < rows.size())
                {
                    break;
                }
                else
                {
                    goto next_col;
                }
            }
            st.push_back(c);
        next_col:;
        }

        vector<int> odd_rows;
        for (int i = 1; i < rows.size(); i += 2)
            odd_rows.push_back(rows[i]);

        // Solve recursively for odd-indexed rows
        vector<int> odd_ans = self(self, odd_rows, st);
        vector<int> ans(rows.size());
        for (int i = 0; i < odd_ans.size(); i++)
            ans[2 * i + 1] = odd_ans[i];

        // Interpolate for even-indexed rows
        int c_idx = 0;
        for (int i = 0; i < rows.size(); i += 2)
        {
            int end_c = (i + 1 == rows.size()) ? st.back() : ans[i + 1];
            int best_c = st[c_idx];
            ll best_val = cost(rows[i], best_c);

            while (st[c_idx] != end_c)
            {
                c_idx++;
                ll val = cost(rows[i], st[c_idx]);
                if (val < best_val)
                {
                    best_val = val;
                    best_c = st[c_idx];
                }
            }
            ans[i] = best_c;
        }
        return ans;
    };

    vector<int> rows(r), cols(c);
    iota(rows.begin(), rows.end(), 0);
    iota(cols.begin(), cols.end(), 0);
    return solve(solve, rows, cols);
}
\end{lstlisting}
\subsection{misc}
\subsubsection{Dynamic bitset}
\begin{lstlisting}[language=C++]
/**
 *  THE ULTIMATE CP TEMPLATE BLUEPRINT: DYNAMIC BITSET
 *
 * 1.  WHAT PROBLEM DOES THIS SOLVE?
 * - Keywords: "Bitset DP with variable size", "Memory Limit Exceeded".
 * - Classic Scenarios: You need the O(N / 64) speedup of std::bitset, but the
 *   required size of the bitset is determined by the input (e.g., sum of array),
 *   meaning you cannot use std::bitset<MAX_SIZE> without hitting Memory Limits.
 * - The Magic: Wraps a std::vector<uint64_t> and manually implements the
 *   left-shift (<<) and bitwise OR (|) operators with hardware-level carry logic.
 *
 * 2.  HOW TO USE IT (THE BLACK BOX)
 * - Initialization:
 *       DynamicBitset dp(total_sum + 1);
 *       dp.set(0);
 * - Transitions:
 *       dp |= (dp << item_weight);
 */

#include <cstdint>

struct DynamicBitset
{
    vector<uint64_t> blocks;
    int bits;

    DynamicBitset(int n = 0)
    {
        bits = n;
        blocks.assign((n + 63) / 64, 0);
    }

    void set(int i)
    {
        if (i < bits)
            blocks[i / 64] |= (1ULL << (i % 64));
    }

    bool test(int i) const
    {
        if (i >= bits)
            return false;
        return (blocks[i / 64] >> (i % 64)) & 1;
    }

    // Crucial for Knapsack / Reachability DP
    DynamicBitset operator<<(int shift) const
    {
        DynamicBitset res(bits + shift);
        int block_shift = shift / 64;
        int bit_shift = shift % 64;
        uint64_t carry = 0;

        for (int i = 0; i < blocks.size(); i++)
        {
            res.blocks[i + block_shift] = (blocks[i] << bit_shift) | carry;
            if (bit_shift > 0)
            {
                carry = blocks[i] >> (64 - bit_shift);
            }
            else
            {
                carry = 0;
            }
        }
        if (carry && (blocks.size() + block_shift < res.blocks.size()))
        {
            res.blocks[blocks.size() + block_shift] = carry;
        }
        // Restore original logical size constraints if necessary
        res.bits = bits + shift;
        return res;
    }

    void operator|=(const DynamicBitset &other)
    {
        for (int i = 0; i < min(blocks.size(), other.blocks.size()); i++)
        {
            blocks[i] |= other.blocks[i];
        }
    }
};
\end{lstlisting}
\subsubsection{LGV lemma}
\begin{lstlisting}[language=C++]
/**
 *  THE ULTIMATE CP TEMPLATE BLUEPRINT: LGV LEMMA (NON-INTERSECTING PATHS)
 *
 * 1.  WHAT PROBLEM DOES THIS SOLVE?
 * - Keywords: "Non-intersecting paths", "Multiple robots on a grid".
 * - Classic Scenarios: You have K robots starting at points A_1...A_K and they
 *   must reach B_1...B_K on a DAG (like a grid with obstacles). They cannot share
 *   any vertices. How many valid routing configurations exist?
 * - The Magic: You only need to calculate the standard number of paths from
 *   every A_i to every B_j. You construct a K x K matrix of these path counts.
 *   The determinant of this matrix magically cancels out all intersecting path
 *   configurations via parity inversions.
 *
 * 2.  HOW TO USE IT (THE BLACK BOX)
 * - Dependencies: You MUST have your `det_prime` function (Template 371) included
 *   above this from your Linear Algebra section.
 * - Logic: Replace `calc_ways(start, end)` with a standard O(V+E) DP or nCr math.
 */

using ll = long long;

// Forward declaration of your existing determinant template
ll det_prime(vector<vector<ll>> a, ll mod);

// Replace this with your specific problem's path counting logic
ll calc_ways(pair<int, int> start, pair<int, int> end, ll mod)
{
    // Example: Grid paths without obstacles using Combinatorics: C(dx + dy, dx)
    // Example: Grid paths with obstacles using O(N^2) DP
    return 0;
}

ll lgv_lemma(const vector<pair<int, int>> &starts, const vector<pair<int, int>> &ends, ll mod)
{
    int k = starts.size();
    if (k != ends.size())
        return 0;

    vector<vector<ll>> M(k, vector<ll>(k, 0));

    // Build the path matrix
    for (int i = 0; i < k; i++)
    {
        for (int j = 0; j < k; j++)
        {
            M[i][j] = calc_ways(starts[i], ends[j], mod);
        }
    }

    // The answer is literally just the determinant of the path matrix
    return det_prime(M, mod);
}
\end{lstlisting}
\section{Helpers}
\subsection{2d prefixSum}
\begin{lstlisting}[language=C++]
#include <bits/stdc++.h>
using namespace std;
#define ll long long


// Compute 2D prefix sum
void computePrefixSum(const vector<vector<int>>& a, vector<vector<long long>>& prefix) {
    int n = a.size(), m = a[0].size();
    prefix.assign(n + 1, vector<long long>(m + 1, 0));
    for (int i = 1; i <= n; i++) {
        for (int j = 1; j <= m; j++) {
            prefix[i][j] = a[i - 1][j - 1]
                         + prefix[i - 1][j]
                         + prefix[i][j - 1]
                         - prefix[i - 1][j - 1];
        }
    }
}

// Query sum of submatrix from (r1, c1) to (r2, c2), inclusive (0-based indexing)
long long getSum(const vector<vector<long long>>& prefix, int r1, int c1, int r2, int c2) {
    r1++; c1++; r2++; c2++; // convert to 1-based for prefix array
    return prefix[r2][c2]
         - prefix[r1 - 1][c2]
         - prefix[r2][c1 - 1]
         + prefix[r1 - 1][c1 - 1];
}


void solve() {
    int n , m , k; cin >> n >> m >> k;
    vector<vector<int>> mat(n, vector<int>(m));
    for (int i = 0; i < n; i++) {
        for (int j = 0; j < m; j++) {
            cin >> mat[i][j];
        }
    }


    vector<vector<ll>>pre;
    computePrefixSum(mat, pre);



}


int main() {
    ios::sync_with_stdio(false);
    cin.tie(nullptr);
    int t = 1;
    cin >> t;
    while (t--)solve();
    return 0;
}
\end{lstlisting}
\subsection{Adhocs}
\begin{lstlisting}[language=C++]
void PratialSum()
{
    int n; cin >> n;
    vector<int>arr(n + 1);
    vector<ll> pratial(n + 1, 0);
    
    for (int i = 1; i <= n; i++) cin >> arr[i];
    
    int l, r, value; cin >> l >> r >> value;   // sure l < r

    pratial[l] += value;
    pratial[r + 1] -= value;

    for (int i = 1; i <= n; i++)
        pratial[i] += pratial[i - 1];

    for (int i = 1; i <= n; i++)
        pratial[i] += arr[i];

    for (int i = 1; i <= n; i++)
        cout << pratial[i] << "   ";

}
// Function to get the sum of elements in the subarray defined by (x1, y1) to (x2, y2)
long long getSum(const vector<vector<ll>>& Prefix, int x1, int y1, int x2, int y2) {
    return Prefix[x2][y2] - Prefix[x1 - 1][y2] - Prefix[x2][y1 - 1] + Prefix[x1 - 1][y1 - 1];
}
void PrefixSum_2D() {
    int n, m, q;
    cin >> n >> m >> q;
    vector<vector<ll>> arr(n + 1, vector<ll>(m + 1, 0));
    vector<vector<ll>> Prefix(n + 1, vector<ll>(m + 1, 0));
    for (int i = 1; i <= n; i++) {
        for (int j = 1; j <= m; j++) {
            cin >> arr[i][j];
        }
    }
    for (int i = 1; i <= n; i++) {
        for (int j = 1; j <= m; j++) {
            Prefix[i][j] = arr[i][j]
                + Prefix[i - 1][j]
                + Prefix[i][j - 1]
                - Prefix[i - 1][j - 1];
        }
    }
    while (q--) {
        int x1, y1, x2, y2;
        cin >> x1 >> y1 >> x2 >> y2; // sure x1 < x2 and y1 < y2
        long long result = getSum(Prefix, x1, y1, x2, y2);
        cout << result << endl;
    }
}

void PratialSum_2D()
{
    int n, m, q;
    cin >> n >> m >> q;
    vector<vector<ll>> arr(n + 1, vector<ll>(m + 1, 0));
    vector<vector<ll>> Pratial(n + 2, vector<ll>(m + 2, 0));
    while (q--) {
        int x1, y1, x2, y2 , val;
        cin >> x1 >> y1 >> x2 >> y2 >> val;
        Pratial[x1][y1] += val;
        Pratial[x1][y2 + 1] -= val;
        Pratial[x2 + 1][y1] -= val;
        Pratial[x2+1][y2+1] += val;
    }
    for (int i = 1; i <= n; i++) {
        for (int j = 1; j <= m; j++) {
            Pratial[i][j] = Pratial[i][j]
                + Pratial[i - 1][j]
                + Pratial[i][j - 1]
                - Pratial[i - 1][j - 1];
        }
    }

    for (int i = 1; i <= n; i++) {
        for (int j = 1; j <= m; j++) {
            cout << Pratial[i][j] << " ";

        }
        cout << '\n';
    }
     
}

\end{lstlisting}
\subsection{BitMask\&Complete Search}
\begin{lstlisting}[language=C++]

vector <ll> getRepresention(ll num, ll base) // convert any number to (any)base 
{
    vector<ll>ans;
    while (num)
    {
        ans.push_back(num % base);
        num /= base;
    }
    reverse(ans.begin(), ans.end());
    return ans;
}
 
// 1. Check if the i-th bit is set (1) or not (0)
bool checkBit(long long mask, int i) {
    return (mask & (1LL << i)) != 0;
}

// 2. Set the i-th bit to 1 (leave other bits unchanged)
long long setBit(long long mask, int i) {
    return mask | (1LL << i);
}

// 3. Clear the i-th bit to 0 (leave other bits unchanged)
long long clearBit(long long mask, int i) {
    return mask & ~(1LL << i);
}

// 4. Toggle the i-th bit (flip 1 to 0, or 0 to 1)
long long toggleBit(long long mask, int i) {
    return mask ^ (1LL << i);
}

// 5. Get the value of the lowest set bit (e.g., 1010 -> 0010)
// Useful for Fenwick trees (Binary Indexed Trees)
long long lowestSetBit(long long mask) {
    return mask & -mask; 
}

// 6. Check if a number is a perfect power of 2
// A power of 2 has exactly one bit set (e.g., 8 is 1000 in binary)
bool isPowerOfTwo(long long n) {
    return n > 0 && (n & (n - 1)) == 0;
}

// 7. Strip the lowest set bit (change the rightmost 1 to a 0)
long long stripLowestSetBit(long long mask) {
    return mask & (mask - 1);
}
int LastBitValue(ll n )
{
    return n & (~(n - 1));
}

// 1. Count the number of set bits (1s) in a 32-bit integer
int count = __builtin_popcount(mask);

// 2. Count the number of set bits (1s) in a 64-bit integer (long long)
int countLL = __builtin_popcountll(mask);

// 3. Count the number of leading zeros
int leadingZeros = __builtin_clz(mask);

// 4. Count the number of trailing zeros (useful to find the index of the first set bit)
int trailingZeros = __builtin_ctz(mask);

 ////////// Ranges ////////////////

void prefix_Xor()
{
    int n; cin >> n; vector<ll>v(n + 1);
    for (int i = 1; i <= n; i++) cin >> v[i];
    vector<ll>pre_Xor(n + 1);
    for (int i = 1; i <= n; i++)
        pre_Xor[i] = pre_Xor[i - 1] ^ v[i];   
    int q; cin >> q; while (q--) {
        int l, r; cin >> l >> r;
        cout << (pre_Xor[r] ^ pre_Xor[l - 1]) << '\n';
    }
}
void prefix_OR()
{
    int n; cin >> n; vector<ll>v(n + 1);
    vector<vector<ll>>prefix_bit(64, vector<ll>(n + 1));
    for (int i = 1; i <= n; i++) cin >> v[i];

    for (int bit = 0; bit < 64; bit++)
    {
        for (int i = 1; i <= n; i++)
        {
            prefix_bit[bit][i] = prefix_bit[bit][i - 1] + ((v[i] >> bit) & 1);
        }
    }
    int q; cin >> q; while (q--) {
        int l, r; cin >> l >> r;
        ll OR = 0;
        for (int bit = 0; bit < 64; bit++)
        {
            int on_bits = prefix_bit[bit][r] - prefix_bit[bit][l - 1]; 
            if (on_bits > 0)
            {
                OR += (1LL << bit);
            }
        }
        cout << OR << '\n';
    }
}
void prefix_AND()
{
    int n; cin >> n; vector<ll>v(n + 1);
    vector<vector<ll>>prefix_bit(64, vector<ll>(n + 1));
    for (int i = 1; i <= n; i++) cin >> v[i];

    for (int bit = 0; bit < 64; bit++)
    {
        for (int i = 1; i <= n; i++)
        {
            prefix_bit[bit][i] = prefix_bit[bit][i - 1] + ((v[i] >> bit) & 1);
        }
    }
    int q; cin >> q; while (q--) {
        int l, r; cin >> l >> r;
        ll AND = 0;
        for (int bit = 0; bit < 64; bit++)
        {
            int on_bits = prefix_bit[bit][r] - prefix_bit[bit][l - 1];
            if (on_bits == r-l+1)
            {
                AND += (1LL << bit);
            }
        }
        cout << AND << '\n';
    }
}


void CompleteSearch()
{
    int n; cin >> n;
    for (int mask = 0; mask < (1 << n); mask++)
    {
        for (int i = 0; i < n; i++)
        {
            if ((mask >> i) &1)
            //if (getBit(mask,i))
                cout << 1;
            else
                cout << 0;
        }
        cout << '\n';
    }
}

void iterateSubsets(int mask) {
    // Start with the full mask, and keep stripping away bits
    for (int sub = mask; sub > 0; sub = (sub - 1) & mask) {
        // 'sub' is a valid subset of 'mask'
        std::cout << sub << "\n";
    }
    // Note: The loop above misses the empty subset (0). 
    // If you need to process 0 as well, you can do it after the loop.
}


void BitSet() {
    bitset<32> b = 15;

    b[3] = 0;
    cout << b.to_ullong() << '\n'; // 7
    cout << b << '\n';  // 00000000000000000000000000000111
    cout << b.size() << "\n"; //32 bit
    cout << b.test(2) << '\n'; // 1 >> true
    cout << b.test(3) << '\n'; // 0 >> false
    cout << b.count() << '\n'; // Count number of 1  == 3
    cout << b.flip() << '\n';  // 11111111111111111111111111111000
    cout << b.set() << '\n';   //11111111111111111111111111111111
    cout << b.reset() << '\n'; //00000000000000000000000000000000
    b[4] = 1;
    cout << b.to_string() << '\n'; //00000000000000000000000000010000
}

\end{lstlisting}
\subsection{Bitmasks(AMR)}
\begin{lstlisting}[language=C++]
//
// Created by AMR on 4/16/2025.
//

#include <bits/stdc++.h>
using namespace std;
typedef long long ll;
#define El_Hefnawys ios_base::sync_with_stdio(false); cin.tie(nullptr); cout.tie(nullptr)

const int MOD = 1e9 + 7;



// A^B = (A|B) - (A&B)
// A^B = (A+B) - 2*(A&B)
// A^B = 2*(A|B) - (A+B)

bool is_power_2(int x) {
    if (x == 0)
        return 0;
    return x&(x-1) == 0;
}

int MSB_value(int n) {
    if (n == 0)
        return 0;
    int index = 31 - __builtin_clz(n);
    return 1 << index;
}

int MSB_index(int n) {
    if (n == 0)
        return -1;
    return 31 - __builtin_clz(n);
}

// you can use __builtin_ctz(n) to get the LSB also directly
int LSB_value(int n) {
    if (n == 0)
        return 0;
    return n&(-n);
}

int LSB_index(int n) {
    if (n == 0)
        return -1;
    int value = n&(-n);
    return 31 - __builtin_clz(value);
}

bool check_bit(ll n, ll i) {
    return n&(1ll<<i);
}

// Another way
bool check_bit_2(ll n, ll i) {
    return 1ll&(n>>i);
}

ll toggle_bit(ll n, ll i) {
    return n^(1ll<<i);
}

ll set_bit(ll n, ll i) {
    return n | (1ll<<i);
}

ll clear_bit(ll n, ll i) {
    return n & (~(1ll<<i));
}

bool check_pow_2(ll n) {
    return __builtin_popcountll(n) == 1; // function to count the number of 1s in a number
    // O(1) function
}

ll num_modulo_pow_2(ll n, ll k) {
    return n & ((1ll << k) - 1);
}

int main() {
    El_Hefnawys;
    cout << num_modulo_pow_2(25,4);
    return 0;
}

//////////////////////////

//
// Created by AMR on 4/17/2025.
//

#include <bits/stdc++.h>
using namespace std;
typedef long long ll;
#define El_Hefnawys                   \
    ios_base::sync_with_stdio(false); \
    cin.tie(nullptr);                 \
    cout.tie(nullptr)

const int MOD = 1e9 + 7;

int Xor_nums_from_1_to_n(int n)
{
    if (!(n % 4))
        return n;
    if (n % 4 == 1)
        return 1;
    if (n % 4 == 2)
        return n + 1;
    return 0;
}
\end{lstlisting}
\subsection{Compression}
\begin{lstlisting}[language=C++]
vector<int> compress(vector<int> &v)
{
    vector<int> help = v;
    sort(help.begin(), help.end());
    help.erase(unique(help.begin(), help.end()), help.end());
    // 1-based
    // to get original value -> help[compressed value - 1];
    for (auto &val : v)
        val = distance(help.begin(), lower_bound(help.begin(), help.end(), val)) + 1;
    return help;
}

\end{lstlisting}
\subsection{Get Words}
\begin{lstlisting}[language=C++]
vector<string> get_words(string str)
{
    vector<string> ret;
    string word;
    stringstream ss(str);
    while (ss >> word) {
        ret.push_back(word);
    }
    return ret;
}

\end{lstlisting}
\subsection{Hashing segment tree}
\begin{lstlisting}[language=C++]
/**
 *  THE ULTIMATE CP TEMPLATE BLUEPRINT: DYNAMIC STRING HASHING (SEGMENT TREE)
 *
 * 1.  WHAT PROBLEM DOES THIS SOLVE?
 * - Keywords: "Dynamic Palindrome Queries", "Substring Equality with Point Updates", "Double Hashing".
 * - Classic Scenarios: You have a string and you need to answer queries like "Is the substring [L, R] 
 *   a palindrome?" OR "Does substring [L1, R1] match [L2, R2]?". Crucially, the string is NOT static; 
 *   you will receive point updates (changing the character at index X).
 * - The Magic: "Polynomial Rolling Hash inside a Segment Tree". A rolling hash is just a sum of terms 
 *   (char * base^i). Segment trees are perfect for point updates and range sums! By maintaining the hash 
 *   in a Segment Tree, updates take O(log N). To query a range [L, R], we sum the hashes in that range 
 *   and "normalize" it by multiplying by the modular inverse of base^(L-1).
 *
 * 2.  HOW TO USE IT (THE BLACK BOX)
 * - Initialization (CRITICAL): You MUST call `init()` once at the start of `main()` to precompute 
 *   the powers and modular inverses.
 * - 1-Based Indexing ONLY: Your string indices and queries MUST be 1-based (from 1 to N).
 *       HashingSegmentTree forward_tree(N), backward_tree(N);
 * - Building: Add characters one by one. For the backward tree, insert at `N - idx + 1`.
 *       forward_tree.update(i, s[i]);
 *       backward_tree.update(n - i + 1, s[i]);
 * - Execution: 
 *       To update a char: forward_tree.update(idx, new_c); backward_tree.update(n - idx + 1, new_c);
 *       To check palindrome: isPalindrome(forward_tree, backward_tree, L, R, n);
 *
 * - Complexity:
 *       Time: Initialization O(N), Updates O(log N), Queries O(log N).
 *       Space: O(N) for the precomputed arrays and Segment Tree nodes.
 *
 * 3.  HOW TO ADAPT IT (THE DIALS & SWITCHES)
 * -  0-Based Indexing Crash Warning: The `query` function uses `inv1[l - 1]`. If you pass `l = 0`, 
 *   it will access index `-1` and cause Undefined Behavior. ALWAYS use 1-based indexing for queries.
 * - Hack Prevention: This template uses Double Hashing with primes 1e9+7 and 2e9+11. It is virtually 
 *   impossible to break (anti-hash tests will fail). If you face a very strict Time Limit (TLE), you 
 *   can convert it to Single Hashing (remove mod2, base2, pw2, inv2) and use a randomized base instead.
 * - `#define int long long`: This template heavily relies on this macro. If you ever remove it, you 
 *   MUST change all hash values, modulos, and products to `long long` to prevent integer overflow 
 *   during `(A * B) % MOD`.
 */

#define int long long
const int N = 1e5 + 5, mod1 = 1e9 + 7, mod2 = 2e9 + 11;
ll base1 = 31, base2 = 37, pw1[N + 1], pw2[N + 1], inv1[N + 1], inv2[N + 1];

ll powmod(ll a, ll b, ll m)
{
    ll ans = 1;
    while (b > 0)
    {
        if (b & 1)
        {
            ans = (ans * a) % m;
        }
        a = (a * a) % m;
        b >>= 1;
    }
    return ans;
}

void init()
{
    pw1[0] = pw2[0] = inv1[0] = inv2[0] = 1;
    int temp1 = powmod(base1, mod1 - 2, mod1);
    int temp2 = powmod(base2, mod2 - 2, mod2);
    for (int i = 1; i < N; i++)
    {
        pw1[i] = (base1 * pw1[i - 1]) % mod1;
        pw2[i] = (base2 * pw2[i - 1]) % mod2;
        inv1[i] = (inv1[i - 1] * temp1) % mod1;
        inv2[i] = (inv2[i - 1] * temp2) % mod2;
    }
}

struct HashingSegmentTree
{
private:
    vector<pair<int, int>> seg;
    int sz;

    pair<int, int> merge(pair<int, int> l, pair<int, int> r)
    {
        pair<int, int> ret = l;
        ret.first = (ret.first + r.first) % mod1;
        ret.second = (ret.second + r.second) % mod2;
        return ret;
    }

    void update(int l, int r, int node, int idx, int ch)
    {
        if (l == r)
        {
            seg[node] = {(ch * pw1[idx]) % mod1, (ch * pw2[idx]) % mod2};
            return;
        }
        int mid = l + r >> 1;
        if (idx <= mid)
            update(l, mid, 2 * node + 1, idx, ch);
        else
            update(mid + 1, r, 2 * node + 2, idx, ch);
        seg[node] = merge(seg[2 * node + 1], seg[2 * node + 2]);
    }

    pair<int, int> query(int l, int r, int node, int lx, int rx)
    {
        if (l >= lx && r <= rx)
        {
            return seg[node];
        }
        if (l > rx || r < lx)
            return {0, 0};
        int mid = l + r >> 1;
        pair<int, int> lft = query(l, mid, 2 * node + 1, lx, rx);
        pair<int, int> rgt = query(mid + 1, r, 2 * node + 2, lx, rx);
        return merge(lft, rgt);
    }

public:
    HashingSegmentTree(int n)
    {
        sz = 1;
        while (sz <= n)
            sz *= 2;
        seg = vector<pair<int, int>>(sz * 2);
    }

    void update(int idx, char ch)
    {
        update(0, sz - 1, 0, idx, ch - 'a' + 1);
    }

    pair<int, int> query(int l, int r)
    {
        pair<int, int> ret = query(0, sz - 1, 0, l, r);
        ret.first = (ret.first * inv1[l - 1]) % mod1;
        ret.second = (ret.second * inv2[l - 1]) % mod2;
        return ret;
    }
};

bool isPalindrome(HashingSegmentTree &a, HashingSegmentTree &b, int &l, int &r, int &n)
{
    return (a.query(l, r) == b.query(n - r + 1, n - l + 1));
}

\end{lstlisting}
\subsection{Hashing treap}
\begin{lstlisting}[language=C++]
/**
 *  THE ULTIMATE CP TEMPLATE BLUEPRINT: IMPLICIT TREAP WITH STRING HASHING
 *
 * 1.  WHAT PROBLEM DOES THIS SOLVE?
 * - Keywords: "Dynamic String", "Insert/Delete characters", "Shift indices", "Palindrome queries".
 * - Classic Scenarios: A Segment Tree can update a character at index X, but it CANNOT insert 
 *   a new character (shifting everything right) or delete a substring (shifting everything left). 
 *   If a problem asks you to insert/delete characters AND query substrings/palindromes, 
 *   standard Segment Trees or Fenwick trees will fail.
 * - The Magic: "The Implicit Treap". It acts as a dynamic array where every node represents 
 *   a character. By using the size of the left subtree as the "implicit index", we can Split 
 *   and Merge the string in O(log N) time. We augment each node to store the Rolling Hash of 
 *   its subtree. When we merge nodes, we mathematically combine their hashes using precomputed 
 *   powers of our base.
 *
 * 2.  HOW TO USE IT (THE BLACK BOX)
 * - Initialization: You MUST call `init()` in `main()` to precompute the `pw` array.
 *       Treap a(s), b(reversed_s);
 * - Insertions: Insert character `c` at 1-based position `pos`.
 *       a.addChar(pos, c);
 * - Deletions: Delete the entire substring from 1-based index `L` to `R`.
 *       a.deleteRange(L, R);
 * - Queries: Get the rolling hash of the substring [L, R].
 *       int h = a.getHash(L, R);
 *
 * - Complexity:
 *       Time: O(log N) for Insert, Delete, and Query!
 *       Space: O(N) for the Treap nodes.
 *
 * 3.  HOW TO ADAPT IT (THE DIALS & SWITCHES)
 * - The Palindrome Reverse Math: When tracking a reversed string dynamically, indices flip. 
 *   If the current string length is `N`, the equivalent of index `idx` in the reversed string 
 *   is `N - idx + 1`. Note that `N` MUST be manually updated (e.g., `n++` or `n -= length`) 
 *   exactly as shown in `main()` whenever you insert or delete!
 * - Memory Limit Exceeded (MLE) Warning: In `deleteRange`, the code drops the `b[0]` component 
 *   (the deleted segment) without freeing the memory. In competitive programming, this is 
 *   usually fine. But if you face MLE on huge test cases, you should write a recursive `destroy()` 
 *   function to `delete` the extracted nodes.
 * - Single vs Double Hashing: This template uses Single Hashing (`mod = 1e9+7`). If you get 
 *   Wrong Answer (WA) due to malicious anti-hash tests, you can upgrade `h` to a `pair<int,int>` 
 *   and apply double hashing just like in the Segment Tree template.
 */
#include <bits/stdc++.h>
using namespace std;
#define ll long long
#define int long long
const int N = 1e6 + 5, mod = 1e9 + 7;
ll base = 31, pw[N + 1];

void init()
{
    pw[0] = 1;
    for (int i = 1; i < N; i++)
    {
        pw[i] = (base * pw[i - 1]) % mod;
    }
}

std::mt19937_64 rnd(std::chrono::system_clock::now().time_since_epoch().count());

struct Treap
{
    struct Node
    {
        int val, h = 0;
        int pri = rnd(), sz = 1;
        array<Node *, 2> c = {0, 0};

        Node()
        {
        }

        Node(int k)
        {
            val = h = k;
        }
    };

    Node *root = 0;
    int getSize(Node *t)
    {
        return t ? t->sz : 0;
    }

    int getHash(Node *t)
    {
        return t ? t->h : 0;
    }

    Treap(string &s)
    {
        for (int i = 0; i < s.size(); i++)
        {
            root = merge(root, new Node(s[i] - 'a' + 1));
        }
    }

    Node *fix(Node *t)
    {
        t->sz = getSize(t->c[0]) + getSize(t->c[1]) + 1;
        t->h = (((getHash(t->c[0]) * base + t->val) % mod) * pw[getSize(t->c[1])] + getHash(t->c[1])) % mod;
        return t;
    }

    array<Node *, 2> split(Node *t, int k)
    {
        if (!t)
            return {0, 0};
        if (getSize(t->c[0]) >= k)
        {
            auto ret = split(t->c[0], k);
            t->c[0] = ret[1];
            return {ret[0], fix(t)};
        }
        else
        {
            auto ret = split(t->c[1], k - getSize(t->c[0]) - 1);
            t->c[1] = ret[0];
            return {fix(t), ret[1]};
        }
    }

    Node *merge(Node *u, Node *v)
    {
        if (!u || !v)
            return u ? u : v;
        if (u->pri > v->pri)
        {
            u->c[1] = merge(u->c[1], v);
            return fix(u);
        }
        else
        {
            v->c[0] = merge(u, v->c[0]);
            return fix(v);
        }
    }

    void deleteRange(int l, int r)
    {
        auto a = split(root, l - 1);
        auto b = split(a[1], r - l + 1);
        root = merge(a[0], b[1]);
    }

    int getHash(int l, int r)
    {
        auto a = split(root, l - 1);
        auto b = split(a[1], r - l + 1);
        int ret = b[0]->h;
        root = merge(a[0], merge(b[0], b[1]));
        return ret;
    }

    void addChar(int pos, char c)
    {
        auto a = split(root, pos - 1);
        root = merge(merge(a[0], new Node(c - 'a' + 1)), a[1]);
    }

    void print(Node *t)
    {
        if (!t)
            return;
        print(t->c[0]);
        cout << t->h << ' ';
        print(t->c[1]);
    }
};

int rev(int i, int n)
{
    return n - i + 1;
}

signed main()
{
    ios_base::sync_with_stdio(0);
    cin.tie(0);
    cout.tie(0);
    init();
    int n, q;
    cin >> n >> q;
    string s;
    cin >> s;
    string f = s;
    reverse(f.begin(), f.end());
    Treap a(s), b(f);
    while (q--)
    {
        int op;
        cin >> op;
        if (op == 1)
        {
            int l, r;
            cin >> l >> r;
            int l2 = rev(l, n), r2 = rev(r, n);
            n -= (r - l + 1);
            a.deleteRange(l, r);
            b.deleteRange(r2, l2);
        }
        else if (op == 2)
        {
            n++;
            int pos;
            char c;
            cin >> c >> pos;
            a.addChar(pos, c);
            b.addChar(rev(pos, n), c);
        }
        else
        {
            int l, r;
            cin >> l >> r;
            int l2 = rev(l, n), r2 = rev(r, n);
            if (a.getHash(l, r) == b.getHash(r2, l2))
            {
                cout << "yes\n";
            }
            else
            {
                cout << "no\n";
            }
        }
    }
    return 0;
}

\end{lstlisting}
\subsection{INT 128}
\begin{lstlisting}[language=C++]
__int128 read() {
    __int128 x = 0, f = 1;
    char ch = getchar();
    while (ch < '0' || ch > '9') {
        if (ch == '-') f = -1;
        ch = getchar();
    }
    while (ch >= '0' && ch <= '9') {
        x = x * 10 + ch - '0';
        ch = getchar();
    }
    return x * f;
}
void print(__int128 x) {
    if (x < 0) {
        putchar('-');
        x = -x;
    }
    if (x > 9) print(x / 10);
    putchar(x % 10 + '0');
}
bool cmp(__int128 x, __int128 y) { return x > y; }


\end{lstlisting}
\subsection{MEX}
\begin{lstlisting}[language=C++]
int MEX(vector<int> &v)
{
    int n = int(v.size());
    vector<int> frq(n);
    for (int i = 0; i < n; i++)
    {
        if (v[i] < n)
            frq[v[i]]++;
    }
    for (int i = 0; i < n; i++)
    {
        if (frq[i] == 0)
            return i;
    }
    return n;
}

\end{lstlisting}
\subsection{Ordered set}
\begin{lstlisting}[language=C++]
#include <ext/pb_ds/assoc_container.hpp>
#include <ext/pb_ds/tree_policy.hpp>
using namespace __gnu_pbds;

typedef tree<int, null_type, less<>, rb_tree_tag, tree_order_statistics_node_update> ordered_set;

void myerase(ordered_set &t, int v)
{
    int rank = t.order_of_key(v);
    auto it = t.find_by_order(rank);
    t.erase(it);
}
/*
 *  t.order_of_key(v);           // Returns the NUMBER of elements strictly less than v (the rank)
 *  *t.find_by_order(k);         // Returns the value at index k (0-indexed)
 *  myerase(t, v);               // Safely erases exactly ONE instance of v
*/

\end{lstlisting}
\subsection{Power}
\begin{lstlisting}[language=C++]
ll power(ll base, ll pow)
{
    base %= MOD;
    ll res{1};
    while (pow)
    {
        if (pow & 1)
            res = res * base % MOD;
        base = base * base % MOD;
        pow >>= 1;
    }
    return res;
}

\end{lstlisting}
\subsection{Query}
\begin{lstlisting}[language=C++]
struct Query
{
    int x, a, id;

    bool operator<(const Query &other) const
    {
        return a < other.a;
    }
};

\end{lstlisting}
\subsection{Solve Equations}
\begin{lstlisting}[language=C++]
pair<ll,ll> equ(ll a,ll b,ll c)
{
    ll x1 = (-b + sqrt(b * b - 4 * a * c)) / (2 *a);
    ll x2 = (-b - sqrt(b * b - 4 * a * c)) / (2 * a);
    return {x1,x2};
}

\end{lstlisting}
\subsection{Split mix}
\begin{lstlisting}[language=C++]
struct custom_hash
{
    static uint64_t splitmix64(uint64_t x)
    {
        x += 0x9e3779b97f4a7c15;
        x = (x ^ (x >> 30)) * 0xbf58476d1ce4e5b9;
        x = (x ^ (x >> 27)) * 0x94d049bb133111eb;
        return x ^ (x >> 31);
    }

    size_t operator()(uint64_t x) const
    {
        static const uint64_t FIXED_RANDOM = chrono::steady_clock::now().time_since_epoch().count();
        return splitmix64(x + FIXED_RANDOM);
    }
};

\end{lstlisting}
\subsection{debug}
\begin{lstlisting}[language=C++]
#define LOCAL ;

template <typename A, typename B>
string to_string(pair<A, B> p);

template <typename A, typename B, typename C>
string to_string(tuple<A, B, C> p);

template <typename A, typename B, typename C, typename D>
string to_string(tuple<A, B, C, D> p);

string to_string(const string &s)
{
    return '"' + s + '"';
}

string to_string(const char *s)
{
    return to_string((string)s);
}

string to_string(bool b)
{
    return (b ? "true" : "false");
}

string to_string(vector<bool> v)
{
    bool first = true;
    string res = "{";
    for (auto &&i : v)
    {
        if (!first)
        {
            res += ", ";
        }
        first = false;
        res += to_string(i);
    }
    res += "}";
    return res;
}

template <size_t N>
string to_string(bitset<N> v)
{
    string res = "";
    for (size_t i = 0; i < N; i++)
    {
        res += static_cast<char>('0' + v[i]);
    }
    return res;
}

template <typename A>
string to_string(A v)
{
    bool first = true;
    string res = "{";
    for (const auto &x : v)
    {
        if (!first)
        {
            res += ", ";
        }
        first = false;
        res += to_string(x);
    }
    res += "}";
    return res;
}

template <typename A, typename B>
string to_string(pair<A, B> p)
{
    return "(" + to_string(p.first) + ", " + to_string(p.second) + ")";
}

template <typename A, typename B, typename C>
string to_string(tuple<A, B, C> p)
{
    return "(" + to_string(get<0>(p)) + ", " + to_string(get<1>(p)) + ", " + to_string(get<2>(p)) + ")";
}

template <typename A, typename B, typename C, typename D>
string to_string(tuple<A, B, C, D> p)
{
    return "(" + to_string(get<0>(p)) + ", " + to_string(get<1>(p)) + ", " + to_string(get<2>(p)) + ", " +
           to_string(get<3>(p)) + ")";
}

template <typename T>
string to_string(queue<T> q)
{
    string res = "{";
    bool first = true;
    while (!q.empty())
    {
        if (!first)
            res += ", ";
        first = false;
        res += to_string(q.front());
        q.pop();
    }
    res += "}";
    return res;
}

template <typename T>
string to_string(stack<T> s)
{
    string res = "{";
    bool first = true;
    while (!s.empty())
    {
        if (!first)
            res += ", ";
        first = false;
        res += to_string(s.top());
        s.pop();
    }
    res += "}";
    return res;
}

template <typename T, typename Container, typename Compare>
string to_string(std::priority_queue<T, Container, Compare> pq)
{
    string res = "{";
    bool first = true;
    while (!pq.empty())
    {
        if (!first)
            res += ", ";
        first = false;
        res += to_string(pq.top());
        pq.pop();
    }
    res += "}";
    return res;
}

void debug_out() { cerr << endl; }

template <typename Head, typename... Tail>
void debug_out(Head H, Tail... T)
{
    cerr << " " << to_string(H);
    debug_out(T...);
}

#ifdef LOCAL
#define debug(...) cerr << "[" << #__VA_ARGS__ << "]:", debug_out(__VA_ARGS__)
#else
#define debug(...) 42
#endif

\end{lstlisting}
\subsection{dist}
\begin{lstlisting}[language=C++]
double dist(double x1, double y1, double x2, double y2)
{
    return hypot(x2 - x1, y2 - y1);
}

\end{lstlisting}
\subsection{div}
\begin{lstlisting}[language=C++]
vector<int> div(int n)
{
    vector<int> divs;
    for (int i = 1; i * i <= n; i++)
    {
        if (n % i == 0)
        {
            divs.push_back(i);
            if (i * i != n)
                divs.push_back(n / i);
        }
    }
    sort(divs.begin(), divs.end());
    return divs;
}

\end{lstlisting}
\subsection{double compare}
\begin{lstlisting}[language=C++]
int dcomp(double a, double b)
{
    return fabs(a - b) < EPS ? 0 : a > b ? 1
                                         : -1;
}

\end{lstlisting}
\subsection{edge}
\begin{lstlisting}[language=C++]
struct edge
{
    ll from, to, w;

    edge(ll from, ll to, ll w) : from(from), to(to), w(w)
    {
    }

    bool operator<(const edge &e) const
    {
        return w < e.w;
    }

    bool operator==(const edge &e)
    {
        return from == e.from && to == e.to && w == e.w;
    }
};

\end{lstlisting}
\subsection{find last occurance}
\begin{lstlisting}[language=C++]
find last occurance for char to specific index 
vector<vector<int>> nxt;
vector<int> last;
string s;
int n,k;
void pre() {
    nxt.assign(n+5, vector<int>(k+5, 0));
    last.assign(k+5, n);

    for (int i = n-1; i >= 0; i--) {
        for (int j = 0; j < k; j++) {
            nxt[i][j] = last[j];
        }
        last[s[i] - 'a'] = i;
    }

}

\end{lstlisting}
\subsection{inv mod}
\begin{lstlisting}[language=C++]
ll power(ll base, ll pow)
{
    base %= MOD;
    ll res{1};
    while (pow)
    {
        if (pow & 1)
            res = res * base % MOD;
        base = base * base % MOD;
        pow >>= 1;
    }
    return res;
}

// MOD must be prime number
ll modInv(ll x)
{
    return power(x, MOD - 2);
}

\end{lstlisting}
\subsection{is prime}
\begin{lstlisting}[language=C++]
bool is_prime(int x)
{
    if (x <= 1)
        return false;
    if (x % 2 == 0)
        return x == 2;
    for (int i = 3; i * i <= x; i += 2)
    {
        if (x % i == 0)
            return false;
    }
    return true;
}

\end{lstlisting}
\subsection{kadanes}
\begin{lstlisting}[language=C++]
int curr{a[0]};
int glob{a[0]};
for (int i = 1; i < n; i++)
{
    curr = max(a[i], curr + a[i]);
    glob = max(curr, glob);
}

//------------------------------------------------------------------//

int mx{}, curr{};
pair<int, int> p{0, 0};
pair<int, int> best{};
for (int i = 0;i < n - 1;i++)
{
    curr += a[i];
    if (curr < 0)
    {
        curr = 0;
        p.first = i + 1;
    }
    p.second = i;
    if (curr > mx)
    {
        best = p;
        mx = curr;
    }
    else if (mx == curr)
    {
        if (best.second - best.first < p.second - p.first)
            best = p;
    }
}

//-------------------------------------------------//

int mx = -OO;

for (int left = 0;left < n;left++)
{
    vector<int> temp(n, 0);

    for (int right = left; right < n; right++)
    {
        for (int i = 0; i < n; i++)
        {
            temp[i] += a[i][right];
        }

        int current_max = temp[0];
        int current_sum = temp[0];

        for (int i = 1; i < n; i++)
        {
            current_sum = max(temp[i], current_sum + temp[i]);
            current_max = max(current_max, current_sum);
        }

        mx = max(mx, current_max);
    }
}

\end{lstlisting}
\subsection{nCr}
\begin{lstlisting}[language=C++]
int nCr(int n, int r)
{
    if (r < 0 || n < 0 || r > n)
        return 0;
    int ret{1};
    for (int i = 1; i <= min(r, n - r); i++)
    {
        ret *= n - i + 1;
        ret /= i;
    }
    return ret;
}

int nPr(int n, int r)
{
    if (r < 0 || n < 0 || r > n)
        return 0;
    int ret{1};
    for (int i = 0; i < r; i++)
        ret = mul(ret, n - i);
    return ret;
}

\end{lstlisting}
\subsection{next greater}
\begin{lstlisting}[language=C++]
vector<int> nextGreater(n, n);
for (int i = 0;i < n;i++)
{
    while (!s.empty() && v[i] > v[s.top()])
    {
        nextGreater[s.top()] = i;
        s.pop();
    }
    s.push(i);
}

\end{lstlisting}
\subsection{permutation hashing}
\begin{lstlisting}[language=C++]
#include <bits/stdc++.h>
using namespace std;
#define ll long long
const int N = 2e5 + 5;
std::mt19937_64 rng(std::chrono::system_clock::now().time_since_epoch().count());
ll rand(ll l, ll r)
{
    return uniform_int_distribution<ll>(l, r)(rng);
}

struct node
{
    ll a, b, c, d, e;

    node() : a(0), b(0), c(0), d(0), e(0)
    {
    }

    node(ll x, ll xx, ll xxx, ll xxxx, ll xxxxx) : a(x), b(xx), c(xxx), d(xxxx), e(xxxxx)
    {
    }

    node operator+(const node &x)
    {
        node w;
        w.a = a + x.a, w.b = b + x.b, w.c = c + x.c, w.d = d + x.d, w.e = e + x.e;
        return w;
    }

    bool operator==(const node &x)
    {
        return a == x.a && b == x.b && c == x.c && d == x.d && e == x.e;
    }

    void operator+=(const node &x)
    {
        a += x.a, b += x.b, c += x.c, d += x.d, e += x.e;
    }

    node operator=(const node &x)
    {
        a = x.a, b = x.b, c = x.c, d = x.d, e = x.e;
        return node(a, b, c, d, e);
    }
};

node f[N], per[N];

void init()
{
    for (int i = 1; i < N; i++)
        f[i] = per[i] = node(rand(1, 1e9), rand(1, 1e9), rand(1, 1e9), rand(1, 1e9),
                             rand(1, 1e9));
    for (int i = 2; i < N; i++)
        per[i] += per[i - 1];
}

signed main()
{
    ios_base::sync_with_stdio(0);
    cin.tie(0);
    cout.tie(0);
    init();
    int t;
    cin >> t;
    while (t--)
    {
        int n;
        cin >> n;
        ll s[n + 1] = {};
        for (int i = 1; i <= n; i++)
            cin >> s[i];
        node pre[n + 5], suf[n + 5];
        for (int i = 1; i <= n; i++)
            pre[i] = suf[i] = f[s[i]];
        for (int i = 2; i <= n; i++)
            pre[i] += pre[i - 1];
        for (int i = n - 1; i >= 1; i--)
            suf[i] += suf[i + 1];
    }
    return 0;
}

\end{lstlisting}
\subsection{random}
\begin{lstlisting}[language=C++]
int getRandom(int min, int max)
{
    static std::mt19937 gen(chrono::steady_clock::now().time_since_epoch().count());
    std::uniform_int_distribution<int> distrib(min, max);
    return distrib(gen);
}

\end{lstlisting}
\subsection{sweepLineAndOfflineQueries}
\begin{lstlisting}[language=C++]
// Counts pairs in subarrays where one element divides the other
// Uses Offline Query Processing and a Segment Tree in O((N + M) log N) time

#include <bits/stdc++.h>
using namespace std;

// Standard Point-Update, Range-Sum Segment Tree (1-based indexing)
struct SegmentTree {
    int n;
    vector<int> tree;
    
    SegmentTree(int n) {
        this->n = n;
        tree.assign(4 * n + 1, 0);
    }
    
    void add(int node, int start, int end, int idx, int val) {
        if (start == end) {
            tree[node] += val;
            return;
        }
        int mid = start + (end - start) / 2;
        if (idx <= mid) {
            add(2 * node, start, mid, idx, val);
        } else {
            add(2 * node + 1, mid + 1, end, idx, val);
        }
        tree[node] = tree[2 * node] + tree[2 * node + 1];
    }
    
    void add(int idx, int val) {
        add(1, 1, n, idx, val);
    }
    
    int query(int node, int start, int end, int l, int r) {
        if (r < start || end < l) return 0;
        
        if (l <= start && end <= r) return tree[node];
        
        int mid = start + (end - start) / 2;
        int left_sum = query(2 * node, start, mid, l, r);
        int right_sum = query(2 * node + 1, mid + 1, end, l, r);
        
        return left_sum + right_sum;
    }
    
    int query(int l, int r) {
        return query(1, 1, n, l, r);
    }
};

struct Query {
    int l, r, id;
};

int main() {
    // Fast I/O
    ios_base::sync_with_stdio(false);
    cin.tie(NULL);

    int n, m;
    if (!(cin >> n >> m)) return 0;

    vector<int> p(n + 1);
    vector<int> pos(n + 1);
    
    for (int i = 1; i <= n; i++) {
        cin >> p[i];
        pos[p[i]] = i; // Store the 1-based index where each value is located
    }

    // pre-calculate all valid pairs (i, j) where i divides j
    // pairs_ending_at[R] will store the left index 'L' for a valid pair ending at index 'R'
    vector<vector<int>> pairs_ending_at(n + 1);
    for (int i = 1; i <= n; i++) {
        // Harmonic series loop: steps through all multiples of i
        // Runs in O(N log N) total time across all iterations
        for (int j = i; j <= n; j += i) {
            int u = pos[i];
            int v = pos[j];
            
            int L = min(u, v);
            int R = max(u, v);
            pairs_ending_at[R].push_back(L);
        }
    }

    vector<Query> queries(m);
    for (int i = 0; i < m; i++) {
        cin >> queries[i].l >> queries[i].r;
        queries[i].id = i;
    }

    // Offline Query Trick: Sort queries by their right endpoint in ascending order
    sort(queries.begin(), queries.end(), [](const Query& a, const Query& b) {
        return a.r < b.r;
    });

    vector<int> ans(m);
    SegmentTree segTree(n);

    int q_idx = 0;
    
    // Sweep through the array from left to right (simulating expanding the subarray)
    for (int r = 1; r <= n; r++) {
        // "Activate" all pairs whose rightmost element is exactly at the current 'r'
        for (int l : pairs_ending_at[r]) {
            segTree.add(l, 1);
        }

        // Answer all queries that ask for a subarray ending at the current 'r'
        while (q_idx < m && queries[q_idx].r == r) {
            int l = queries[q_idx].l;
            // Since all active pairs in the tree currently have their right endpoint <= r,
            // we just need to count how many of them have their left endpoint >= l.
            ans[queries[q_idx].id] = segTree.query(l, r);
            q_idx++;
        }
    }

    for (int i = 0; i < m; i++) {
        cout << ans[i] << "\n";
    }

    return 0;
}
\end{lstlisting}
\subsection{valid}
\begin{lstlisting}[language=C++]
bool valid(int i, int j)
{
    if (i >= r || i < 0 || j >= c || j < 0)
        return false;
    return true;
}

\end{lstlisting}
\end{document}
